\providecommand{\toplevelprefix}{../..}  %
\documentclass[../../book-main_es.tex]{subfiles}

\begin{document}

\chapter{Aprendizaje de representaciones en datos y tareas del mundo real}
\label{ch:applications}

\begin{quote}
    \hfill    ``{\em La mejor teoría se inspira en la práctica y la
    mejor práctica se inspira en la teoría}.''

    $~$ \hfill --- Donald Knuth
\end{quote}
\vspace{5mm}

\section{Introducción}
En los capítulos anteriores se ha presentado una introducción sistemática a
los problemas matemáticos, los marcos computacionales y los algoritmos prácticos
relacionados con el aprendizaje de distribuciones de baja dimensión
a partir de datos de alta dimensión. Aunque la mayoría de las justificaciones teóricas
de estos métodos se han establecido para modelos idealizados de distribuciones de datos,
como (mezclas de) subespacios y/o gaussianas, los
principios e ideas que subyacen a estos métodos computacionales son,
no obstante, potentes y generales, pensados de hecho para ser
aplicables a conjuntos de datos y tareas del mundo real.

Para ayudar a los lectores a comprender mejor el contenido de este libro y aprender 
cómo aplicar lo aprendido hasta ahora a datos del mundo real, a continuación ofrecemos algunas demostraciones y ejemplos de
varias aplicaciones representativas. Cada aplicación propone una
solución a una tarea del mundo real con un conjunto de datos del mundo real (como
datos visuales, datos de movimiento y datos de texto), utilizando los métodos que hemos
presentado en este libro. Los resultados presentados en este capítulo
tienen como objetivo cumplir los dos propósitos siguientes:
\begin{itemize}
    \item en primer lugar, aportar detalles experimentales adicionales y
        evidencia empírica que validen los métodos presentados
        anteriormente en el libro y demostrar su importante
        potencial en contextos del mundo real;
    \item en segundo lugar, familiarizar al lector con ciertos
        métodos y tareas empíricos modernos del aprendizaje profundo que no están
        bien documentados fuera de los entornos de investigación o de los códigos de producción.
\end{itemize}
Sin embargo, en nuestra sincera opinión, las soluciones y los resultados que se presentan
aquí están pensados únicamente para comprobar que la metodología funciona. Por
lo tanto, existe un amplio margen de mejora futura, tanto en lo que respecta a la comprensión teórica como a la ingeniería práctica,
para seguir avanzando en el
rendimiento de vanguardia. A lo largo de este capítulo, y especialmente en la última Sección \ref{sec:scalable},
presentaremos algunas mejoras inmediatas y
evidentes. Dejaremos para futuras investigaciones en el \Cref{ch:future} final el análisis de ideas potencialmente mucho más transformers sobre el aprendizaje de representaciones y los importantes problemas aún sin resolver en materia de inteligencia en general.
\subsection{Estructura del Capítulo}
\label{sec:application_outline}
En lo que respecta a la inteligencia en la naturaleza, podría decirse que los tres tipos de datos más
importantes con los que nuestro cerebro construye nuestra memoria y
nuestro conocimiento son {\em los datos visuales}, {\em los movimientos de nuestro cuerpo} y los {\em lenguajes naturales}.
Por lo tanto, en este capítulo mostraremos cómo aplicar los principios y métodos presentados en este libro para obtener 
buenas representaciones de estos datos del mundo real y explicaremos cómo dichas representaciones facilitan tareas prácticas importantes relacionadas con ellos.

En primer lugar, desde \Cref{sec:contrastive_learning} hasta
la Sección \ref{sec:CUPID}, comenzaremos con los datos visuales y mostraremos cómo
aprender las distribuciones de imágenes 2D a objetos 3D, ya que la memoria visual
nos sirve (a nosotros y a los animales) como modelo de bajo nivel pero fundamental que almacena
el conocimiento del mundo físico y nos ayuda
a percibir y predecir en un nuevo entorno. A continuación, en la Sección
\ref{sec:egoallo}, mostraremos cómo aprender la distribución de los movimientos corporales,
ya que es fundamental para nosotros (y para los animales) actuar o interactuar con rapidez y
precisión en los entornos. Por último, en la Sección
\ref{sec:clm_text}, mostraremos cómo modelar y aprender la
distribución de los lenguajes naturales.

Evidentemente, en comparación con los lenguajes naturales, nuestra cobertura de
los datos físicos, visuales y de movimiento será mucho más amplia.
Aunque en los últimos años el sector de la inteligencia artificial ha realizado importantes inversiones en el modelado de 
los lenguajes naturales,\footnote{Incluidos sus derivados, como los lenguajes de programación, las deducciones lógicas y 
las demostraciones matemáticas.} la gente ha empezado a darse cuenta de la importancia de modelar el conocimiento sobre el 
mundo físico y nuestras acciones en él. En los últimos años, ha habido un rápido aumento del interés y un cambio de foco 
hacia el desarrollo de lo que se conoce como «modelo del mundo», «inteligencia espacial», «inteligencia física» o 
«inteligencia incorporada», ya que son elementos cruciales para los sistemas o agentes de IA---por ejemplo, un robot 
inteligente---que necesitan interactuar con un mundo físico abierto. Como hemos aprendido en el Capítulo \ref{ch:intro}, 
ese es precisamente el objetivo y el alcance del programa de {\em  Cibernética} iniciado por Norbert Wiener en la década de 1940.

En el plano técnico, desde la \Cref{sec:contrastive_learning} hasta la
\Cref{sec:image_classification}, mostraremos cómo utilizar los principios
presentados en este libro para aprender buenas representaciones de datos de imágenes
en entornos {\em no supervisados, débilmente supervisados} y {\em supervisados},
respectivamente.\footnote{El objetivo en los tres casos es aprender una buena representación 
de la distribución de los datos. Solo se diferencian en la cantidad de información adicional 
que se proporciona sobre la distribución y que puede facilitar el aprendizaje.} Esto servirá tanto como
introducción al procesamiento de datos de imágenes, a las técnicas de aumento de datos y a la popular 
arquitectura profunda Transformer, como paso previo a la implementación posterior de la arquitectura 
CRATE. Estos ejemplos también ofrecen la primera demostración de las drásticas simplificaciones que 
ya se pueden llevar a cabo utilizando los principios presentados en el libro. Continuaremos
con las modificaciones de la \textit{arquitectura de Transformer} en la \Cref{sec:image_classification} 
para imágenes y en la \Cref{sec:clm_text} para textos, respectivamente, que demuestran las capacidades 
de las arquitecturas transparentes white-box simplificadas, incluyendo CRATE y sus variantes, para 
la \textit{codificación} en los ámbitos de la imagen y el texto. También demostraremos las capacidades 
de las arquitecturas simplificadas para la \textit{auto-codificación} en la \Cref{sec:image_completion}. %

En la práctica, la razón principal por la que queremos aprender una buena
representación de una distribución de datos es para poder tomar muestras de la
distribución (como una prior) con el fin de regenerar, estimar o predecir el
estado del mundo, condicionado a la observación actual o a la información dada. 
Demostramos las capacidades generativas de las representaciones aprendidas
\begin{itemize}
\item para imágenes 2D en la \Cref{sec:auto-encode-img-gen} y la
\Cref{sec:cond_gen};
\item para objetos 3D en la
\Cref{sec:Michelangelo} y la \Cref{sec:shape-and-pose};
\item para los movimientos humanos en la \Cref{sec:egoallo};
\item y para los lenguajes naturales en la \Cref{sec:clm_text}.
\end{itemize} 

\subsection{Configuración técnica general}\label{sec:experiment_setup}

En capítulos anteriores, hemos mencionado diferentes configuraciones en las que utilizamos
técnicas de aprendizaje de representaciones para procesar datos reales a gran escala. En
este capítulo, describiremos dichas configuraciones con gran detalle. El
objetivo de esta sección es que el lector sea capaz de
reproducir cualquier experimento que se trate en esta sección (o, de hecho, en
todo el libro) utilizando únicamente la descripción que te daremos en el libro, los
principios introducidos en capítulos anteriores y ampliados en este
capítulo, así como los hiperparámetros extraídos de diversos artículos cuyos
resultados se analizan en este capítulo. Con este fin, describiremos
\textit{con precisión} todos los procedimientos utilizando un lenguaje detallado,
pseudocódigo o notación matemática que pueda implementarse directamente
en código. Siempre que sea posible, analizaremos cómo las implementaciones concretas
se relacionan con los principios presentados anteriormente en el libro.

\begin{figure}
    \includegraphics[width=\textwidth]{\toplevelprefix/chapters/chapter8/figs/encoder_pipeline.pdf}
    \caption{\small\textbf{Un diagrama del proceso del codificador.} Los datos
        \(\vX \in \cD\) se introducen en la función de incrustación
        \(f_{\theta}^{\emb}\) para obtener una secuencia en \((\R^{d})^{*}\). La
        incrustación pasa por una red troncal \(f_{\theta}^{\mathrm{bb}}\)
        para obtener características \(\vZ_{\theta}(\vX)\) para cada token. Podemos
        extraer una característica agregada \(\vz_{\theta}(\vX)\) utilizando el
        mapa de extracción \(f_{\theta}^{\mathrm{ext}}\). Por último, para utilizar la
        característica agregada en tareas posteriores, podemos emplear el
    cabezal específico para la tarea \(h_{\theta}\).}
    \label{fig:overall_encoder_pipeline}
\end{figure}

\begin{figure}
    \includegraphics[width=\textwidth]{\toplevelprefix/chapters/chapter8/figs/autoencoder_pipeline.pdf}
    \caption{\small\textbf{Un diagrama del proceso del auto-codificador.}
        Los datos \(\vX \in \cD\) se introducen en la función de incrustación
        \(f_{\theta}^{\emb}\) para obtener una secuencia en \((\R^{d})^{*}\). La
        incrustación pasa por una red troncal del codificador
        \(f_{\theta}^{\mathrm{bb}}\) para obtener características
        \(\vZ_{\theta}(\vX)\) para cada token. Para decodificar
        \(\vZ_{\theta}(\vX)\), lo pasamos a través de una red troncal del decodificador
        \(g_{\eta}^{\mathrm{bb}}\). Para mapear la salida de la red troncal del decodificador
        de vuelta al espacio de datos \(\cD\), utilizamos una capa de \textit{desincrustación}
        \(g_{\eta}^{\mathrm{unemb}}\), obteniendo en definitiva una reconstrucción
        \(\hat{\vX}_{\theta, \eta}(\vX)\) (aquí estilizada como una
    reconstrucción pixelada de la entrada).}
    \label{fig:overall_autoencoder_pipeline}
\end{figure}

%\paragraph{Configuración técnica.}
Definamos el conjunto de datos posibles como \(\cD\) (que, en última instancia,
    será el conjunto de imágenes \(\cI\), por ejemplo, o el conjunto de textos
\(\cT\)), así como el conjunto de secuencias finitas de tokens en \(\R^{d}\)
(es decir, el conjunto de matrices con \(d\) filas) como \((\R^{d})^{*} \doteq
\bigcup_{T = 1}^{\infty}\R^{d \times T}\). Para analizar la
gran variedad de aplicaciones que presentamos en este capítulo, recordamos primero
que en el resto del libro tratamos dos tipos diferentes de arquitecturas de modelos.
\begin{itemize}
    \item Una arquitectura de \textit{codificador}, parametrizada por
        \(\theta\), que se compone de varias componentes:
        \begin{itemize}
            \item Una \textit{incrustación} \(f_{\theta}^{\emb} \colon
                \cD \to (\R^{d})^{*}\), que convierte los datos de entrada
                \(\cD\) en una serie de \textit{tokens} que se
                mapean o se \textit{incrusten} en
                un espacio de \(d\) dimensiones. \textit{En el resto del
                    capítulo, a menudo identificaremos los tokens y las incrustaciones
                entre sí.}
            \item Una \textit{red troncal de codificación}
                \(f_{\theta}^{\backbone} \colon (\R^{d})^{*} \to
                (\R^{d})^{*}\), que procesa la serie de
                incrustaciones mediante una operación de secuencia a secuencia.
                Esta red troncal se implementa mediante las
                arquitecturas de red discutidas en los capítulos anteriores, pero
                daremos una descripción más formal a medida que avancemos.
            \item Un \textit{extractor de características agregadas}
                \(f_{\theta}^{\ext} \colon (\R^{d})^{*} \to \R^{d}\),
                que extrae una representación agregada de
                toda la secuencia. Esto se utiliza para definir una única
                característica para toda la muestra de datos. 
            \item Una \textit{cabeza específica para la tarea} \(h_{\theta} \colon
                \R^{d} \to \R^{m}\), que extrae una
                salida de \(m\) dimensiones para la predicción.
        \end{itemize}
        También definimos \(f_{\theta} \doteq f_{\theta}^{\backbone}
        \circ f_{\theta}^{\emb} \colon \cD \to (\R^{d})^{*}\). Dada
        una entrada \(\vX\), escribimos \(\vZ_{\theta}(\vX) \doteq
        f_{\theta}(\vX)\) y \(\vz_{\theta}(\vX) \doteq
        f_{\theta}^{\ext}(\vZ_{\theta}(\vX))\). El proceso general
        se muestra en la \Cref{fig:overall_encoder_pipeline}.
    \item Una arquitectura de \textit{auto-codificador}, que se compone de
        varias componentes:
        \begin{itemize}
            \item Una \textit{incrustación} \(f_{\theta}^{\emb} \colon
                \cD \to (\R^{d})^{*}\), que convierte los datos de entrada
                \(\cD\) en una serie de \textit{tokens} que se
                incrustan en un espacio de \(d\) dimensiones.
            \item Una \textit{red troncal codificadora}
                \(f_{\theta}^{\backbone} \colon (\R^{d})^{*} \to
                (\R^{d})^{*}\), que procesa la serie de
                incrustaciones mediante una operación de secuencia a secuencia.
            \item Una \textit{red troncal decodificadora} \(g_{\eta}^{\backbone}
                \colon (\R^{d})^{*} \to (\R^{d})^{*}\), que
                conceptualmente revierte la operación de la red troncal codificadora.
            \item Un \textit{desincrustación} \(g_{\eta}^{\unemb} \colon
                (\R^{d})^{*} \to \cD\), que actúa como inversa de la
                incrustación.
        \end{itemize}
        También definimos \(f_{\theta} \doteq f_{\theta}^{\backbone}
        \circ f_{\theta}^{\emb} \colon \cD \to (\R^{d})^{*}\) y
        \(g_{\eta} \doteq g_{\eta}^{\unemb} \circ
        g_{\eta}^{\backbone} \colon (\R^{d})^{*} \to \cD\). Dada una
        entrada \(\vX\), escribimos \(\vZ_{\theta}(\vX) \doteq
        f_{\theta}(\vX)\) y \(\hat{\vX}_{\theta, \eta}(\vX) \doteq
        g_{\eta}(\vZ_{\theta}(\vX))\). El proceso general se
        muestra en la \Cref{fig:overall_autoencoder_pipeline}.
\end{itemize}

Utilizaremos esta notación en numerosas ocasiones a lo largo de este capítulo, así que
no dude en volver a consultarla si algo no le queda claro.
Esta descomposición de nuestras redes también se ajusta en gran medida a la mayoría de las implementaciones de código
y usted puede comenzar sus proyectos de programación definiendo
estas redes.

\section{Aprendizaje no supervisado de
representaciones de imágenes}\label{sec:contrastive_learning}

El aprendizaje de representaciones fieles y de alta calidad de las distribuciones de datos de imágenes
es un problema fundamental en la inteligencia artificial, tal y como
hemos analizado en profundidad en el Capítulo \ref{ch:lossy-compression}.
En la literatura se han propuesto muchos enfoques populares
para esta tarea, muchos de los cuales pueden parecer que no utilizan las técnicas y
los principios descritos en este manuscrito. Uno de estos enfoques se denomina
\textit{aprendizaje contrastivo}, llamado así porque el objetivo del aprendizaje 
consiste (en términos generales) en garantizar que las características de
los datos «similares» sean similares y que las características de los datos «disimilares» estén
muy alejadas entre sí. Las soluciones de aprendizaje contrastivo suelen ser 
enfoques muy sofisticados y diseñados en forma empírica. En esta sección,
presentaremos uno de estos enfoques populares, conocido como DINO \citep{caron2021emerging}. Además, demostraremos cómo utilizar los principios descritos en este libro para simplificar drásticamente sus decisiones de diseño y, al mismo tiempo, mejorar las representaciones aprendidas. Esto se conoce como \textbf{SimDINO} \cite{wu2025simplifying}:
\\
\centerline{\url{https://robinwu218.github.io/SimDINO/}.}\\

\subsection{Datos}\label{sub:contrastive_learning_data}
Los datos que utilizaremos para analizar y simplificar la metodología DINO
son todos datos de imágenes bidimensionales (2D).

Para el \textit{entrenamiento},
utilizaremos los conjuntos de datos ImageNet-1K e ImageNet-21K. Cada muestra del conjunto 
de datos es una imagen RGB, de resolución variable, junto a una etiqueta que indica el objeto 
o la escena que contiene la imagen (es decir, la \textit{clase} de la imagen). El conjunto de 
datos ImageNet-1K contiene 1,28 millones de imágenes de entrenamiento y 50 000 imágenes de 
validación, divididas en 1 000 clases. El conjunto de datos ImageNet-21K contiene 14,2 millones 
de imágenes de entrenamiento y 21 800 clases, pero las clases no son disjuntas (es decir, 
algunas clases son subconjuntos de otras). Dado que estamos realizando un aprendizaje 
auto-supervisado, las etiquetas no se utilizarán durante el entrenamiento, sino solo durante 
la evaluación. Con el fin de aumentar la solidez de nuestro modelo, a menudo aplicamos 
\textit{pequeñas transformaciones mediante aumento de los datos} a cada imagen durante el procesamiento, 
tales como volteos, la adición de un pequeño ruido aleatorio o recortes aleatorios de 
gran tamaño; no incluimos esto en nuestra notación, ya que cada modificación de una 
imagen natural es en sí misma (casi) una imagen natural de nuestro conjunto de datos.

Para la \textit{evaluación}, utilizaremos una amplia gama de conjuntos de datos. De ellos, 
el más común es CIFAR-10. CIFAR-10 es un conjunto de datos compuesto por 60 000 imágenes 
naturales en formato RGB de 32 × 32, divididas en 10 clases, con un
conjunto de entrenamiento predeterminado de 50 000 muestras y un conjunto de validación 
de 10 000 muestras. El objetivo de utilizar CIFAR-10 es garantizar que los modelos que se 
entrenan con una distribución de imágenes (ImageNet) puedan generalizarse a otra distribución 
de imágenes (CIFAR-10). También hacemos referencia a otros
conjuntos de datos similares, como CIFAR-100 (independiente de CIFAR-10), Oxford Flowers y Oxford Pets. 
En la \Cref{fig:in1k_cifar10_examples} se muestran ejemplos de los datos de ImageNet-1K y CIFAR-10.

\begin{figure}
    \centering

    \hfill
    \begin{subfigure}{0.3\textwidth}
        \centering
        \includegraphics[width=\textwidth]{\toplevelprefix/chapters/chapter8/figs/imagenet.pdf}
        \caption{\small Muestras de ImageNet-1K.}
    \end{subfigure}
    \hfill
    \begin{subfigure}{0.26\textwidth}
        \centering
        \includegraphics[width=\textwidth]{\toplevelprefix/chapters/chapter8/figs/cifar10.pdf}
        \caption{\small Muestras de CIFAR10.}
    \end{subfigure}
    \hfill
    \phantom{}

    \caption{\small\textbf{Imágenes de ImageNet-1K \textit{(izquierda)} 
        y CIFAR-10 \textit{(derecha)}.} Obsérvese que las imágenes de CIFAR-10
        tienen, en general, una resolución mucho menor, lo que reduce la
    complejidad del aprendizaje de esa distribución.}
    \label{fig:in1k_cifar10_examples}
\end{figure}

En un plano un poco más formal, nuestros datos \(\vX\) serán imágenes;
definamos \(\cI\) como el conjunto de todas las imágenes. Dado que una imagen es una matriz rectangular
de píxeles y cada píxel tiene un color dado por RGB, CMYK u
otro formato de color, decimos que una imagen es un elemento de \(\R^{c
\times h \times w}\) --- aquí \(c\) es el número de canales (es decir,
\(3\) para RGB y \(4\) para CMYK), \(h\) es la altura de la imagen y
\(w\) es el ancho de la imagen. Por consiguiente, el conjunto de todas las imágenes \(\cI
\doteq \bigcup_{c, h, w = 1}^{\infty}\R^{c \times h \times w}\) es
el conjunto de todos los datos posibles de este tipo. Una vez más, utilizaremos esta notación en repetidas ocasiones.

\subsection{La tarea y el objetivo simplificados de DINO}
\label{sub:contrastive_learning_objective}

Nuestra tarea consiste en aprender una buena representación de los datos. El aprendizaje contrastivo,
en términos generales, lo consigue definiendo qué propiedades de la
imagen de entrada queremos que reflejen las características, generando muestras
que compartan esas propiedades pero varíen en otras, y estableciendo una función de pérdida
que favorezca que las características de las imágenes con propiedades comunes sean
similares y que las de las imágenes con propiedades diferentes sean distintas. La
solución óptima natural a este problema de aprendizaje consiste en que las
características aprendidas conserven las propiedades deseadas de la entrada.
Sin embargo, existen muchas complicaciones prácticas y empíricas que
surgen durante el entrenamiento de modelos contrastivos.

En el caso de DINO, los autores proponen utilizar una metodología que
genera un único vector de características para toda la imagen y pretende que
dicho vector contenga información «global» (es decir, a nivel de imagen).
En consecuencia, la función de pérdida favorecerá que las imágenes con información global
similar tengan características similares y que las imágenes con información global
diferente tengan características diferentes.

Esto parece intuitivo, pero, como se mencionó anteriormente, hay varias
consideraciones empíricas, incluso a la hora de definir la función de pérdida. En primer lugar,
¿cómo debemos promover las similitudes y las diferencias? La
respuesta de DINO \citep{caron2021emerging} es\footnote{En
opinión del autor, inexplicablemente...} convertir las características de salida en
«logits» correspondientes a alguna distribución de probabilidad y calcular
su entropía cruzada. Más concretamente, sea \(\Delta_{m} \doteq \{\vx
        \in \R^{m} \colon x_{i} \geq 0\ \forall i \in [m], \sum_ {i =
1}^{m}x_{i} = 1\}\) sea el espacio de vectores de probabilidad en \(\R^{m}\)
y definamos la función \(d_{\CE} \colon \R^{m} \times \R^{m} \to \R\) por
\begin{equation}\label{eq:cross_entropy_difference}
    d_{\CE}(\vp, \vq) \doteq \CE(\vp, \vq), \quad \forall \vp, \vq
    \in \Delta_{m}
\end{equation}
donde \(\CE \colon \Delta_{m} \times \Delta_{m} \to \R\) es la
entropía cruzada, definida como
\begin{equation}\label{eq:expts_def_ce}
    \CE(\vp, \vq) \doteq -\sum_{i = 1}^{m}p_{i}\log q_{i}, \quad
    \forall \vp = (p_{1}, \dots, p_{m}), \vq = (q_{1}, \dots, q_{m})
    \in \Delta_{m}.
\end{equation}
Antes de continuar con nuestro análisis, intentemos hacernos una idea de
esta función de distancia. En concreto, tenemos que:
\begin{align}
    \CE(\vp, \vq)
    &= -\sum_{i = 1}^{m}p_{i}\log q_{i} = \sum_{i = 1}^{m}p_{i}\log(p_{i}/q_{i})
    - \sum_{i = 1}^{m}p_{i}\log p_{i} \\ 
    &= \KL(\vp\mmid \vq) + H(\vp)
\end{align}
donde \(\KL \colon \Delta_{m} \times \Delta_{m} \to \R\) es la divergencia de KL,
definida como
\begin{equation}
    \KL(\vp\mmid \vq) \doteq \sum_{i = 1}^{m}p_{i}\log(p_{i}/q_{i}),
\end{equation}
y \(H \colon \Delta_{m} \to \R\) es la entropía de una variable aleatoria. Obsérvese
que \(\KL(\vp\mid \vq)\) se minimiza si y solo si \(\vp = \vq\).
Así que minimizar \(d_{\CE}\) hace dos cosas: hace que \(\vp = \vq\),
y hace que \(\vp\) y \(\vq\) tengan entropía mínima (es decir, vectores
    con \(1\) en una sola componente y \(0\) en el resto --- denominados
\textit{vectores «one-hot»}). En general, la meta de este objetivo no es
\textit{solo} hacer coincidir \(\vp\) y \(\vq\), sino también darles una forma
determinada para que tengan baja entropía. Tengamos esto en cuenta cuando
analicemos la formulación.

La siguiente pregunta es: ¿cómo podemos obtener muestras con información global similar?
La respuesta que ofrece DINO (al igual que casi todo el
aprendizaje contrastivo) es el \textit{aumento de datos}: a partir de cada
muestra, se generan varias muestras correlacionadas que comparten las propiedades
deseadas. En el caso de DINO, utilizamos diferentes recortes o
\textit{vistas} de la imagen de entrada. Recordemos que modelamos una imagen como
un elemento del conjunto \(\cI\). En esta notación, una vista es una función
\(v \colon \cI \to \cI\). En el caso de DINO, la vista es un
\textit{recorte redimensionado aleatoriamente}: toma un recorte rectangular
elegido al azar de la imagen (que ocupa un porcentaje fijo \(p_{v} \in [0, 1]\)
del área total de la imagen), le cambia el tamaño proporcionalmente para que
el \textit{lado más corto} tenga \(S_{\rsz}\) píxeles de largo y luego le cambia el tamaño
a una forma fija \((C, S_{v}, S_{v})\), donde \(S_{v} \geq 1\) es
el tamaño de la vista y \(C\) es el número de canales de la imagen original.

\begin{figure}
    \centering
    \includegraphics[width=0.7\textwidth]{\toplevelprefix/chapters/chapter8/figs/global_local_views.pdf}
    \caption{\textbf{Vistas locales y globales en DINO.} Las vistas locales y
        las vistas globales toman un recorte rectangular de la imagen de entrada y
        lo redimensionan a una forma cuadrada, que luego se introduce en la red
    para su procesamiento.}
    \label{fig:dino_local_global_views}
\end{figure}

Hay dos tipos de vistas que queremos utilizar, tal y como se muestra en la
\Cref{fig:dino_local_global_views}:
\begin{itemize}
    \item \textit{vistas globales}, que son recortes redimensionados aleatoriamente con
        el parámetro de porcentaje de área \(p_{\glo} \in [0, 1]\) y la forma de salida
        \((C, S_{\glo}, S_{\glo})\);
    \item \textit{vistas locales}, que son recortes redimensionados aleatoriamente con
        el parámetro de porcentaje de área \(p_{\loc} \in [0, 1]\) y la forma de salida
        \((C, S_{\loc}, S_{\loc})\). Aquí \(p_{\loc} <
        p_{\glo}\) y \(S_{\loc} < S_{\glo}\).
\end{itemize}

DINO pretende que las características agregadas \(\vz_{\theta}(\vX_{v})
\doteq (f_{\theta}^{\ext} \circ f_{\theta})(\vX_{v})\) de todas las vistas
\(\vX_{v} \doteq v(\vX)\) de una imagen de entrada \(\vX\) sean consistentes
entre sí. DINO lo hace utilizando un \textit{``DINO
head''}\footnote{Tenga en cuenta que \(h_{\vW, \vmu}\) es la cabeza específica de la tarea,
    la cual, en \Cref{sec:experiment_setup}, se parametriza únicamente mediante
    \(\theta\), en lugar de mediante parámetros específicos; sin embargo, dado que utilizamos
    dos invocaciones de \(h\) con valores distintos del segundo
parámetro, mantenemos la notación especificada.} \(h_{\vW, \vmu}\),
parametrizada por una matriz \(\vW \in \R^{s \times d}\) y un vector
\(\vmu \in \R^{s}\), para extraer un vector de probabilidad \(\vp_{\theta,
\vW, \vmu}(\vX_{v}) \doteq h_{\vW, \vmu}(\vz_{\theta}(\vX_{v}))\)
a partir de la característica agregada \(\vz_{\theta}(\vX_{v})\), utilizando la
siguiente receta simple:
\begin{equation}
    h_{\vW, \vmu}(\vz) \doteq \softmax([\vW\vz - \vmu]/\tau), \qquad
    \forall \vz \in \R^{d},
\end{equation}
donde la función \(\softmax \colon \R^{s} \to \Delta_{s}\) se define por
\begin{equation}
    \softmax\rp{\mat{x_{1} \\ \vdots \\ x_{s}}} \doteq
    \frac{1}{\sum_{i = 1}^{s}e^{x_{i}}}\mat{e^{x_{1}} \\ \vdots \\ e^{x_{s}}}
\end{equation}
y \(\tau > 0\) es el parámetro de temperatura que controla la
entropía de la distribución resultante del softmax.

En particular, DINO minimiza la diferencia entre el vector de probabilidad
\(\vp_{\theta, \vW, \vmu}(\vX_{g}) \doteq h_{\vW, 
\vmu}(\vz_{\theta}(\vX_{g}))\) para cada vista global \(\vX_{g} \doteq
v_{g}(\vX)\) y el vector de probabilidad \(\vp_{\theta, \vW}(\vX_{c})
\doteq h_{\vW, \vzero_{m}}(\vz_{\theta}(\vX_{c}))\) para cada vista
\(\vX_{c} \doteq v_{c}(\vX)\). En este caso, \(v_{c}\) puede ser
una vista local o una vista global. Analizaremos la implementación de
\(f_{\theta}\) y \(f_{\theta}^{\ext}\) en breve en
\Cref{sub:contrastive_learning_architecture}. En general, DINO resuelve el problema
\begin{eqnarray}
    &&\min_{\theta, \vW, \vmu}\cL_{\dino}(\theta, \vW, \vmu) \qquad
    \text{where} \nonumber \\
    &&\cL_{\dino}(\theta, \vW, \vmu) \doteq \Ex[d_{\CE}(\vp_{\theta,
    \vW, \vmu}(\vX_{g}), \vp_{\theta, \vW}(\vX_{c}))],\label{eq:dino_loss}
\end{eqnarray}
donde la esperanza se calcula sobre los datos \(\vX\), las vistas globales \(v_{g}\)
y otras vistas \(v_{c}\).

En este caso concreto, sin embargo, si intentas implementar
\eqref{eq:dino_loss} y optimizarla en una red real, es muy
probable que te encuentres con un problema: tras ejecutar unas cuantas
iteraciones del algoritmo de aprendizaje, la correspondencia de características
\(f_{\theta}^{\ext} \circ f_{\theta}\) ¡se \textit{convertirá en una
función constante}! Esto sin duda optimiza la pérdida anterior, ya que
minimiza la distancia entre las características de diferentes vistas de la
misma imagen. Pero obviamente no queremos aprender esta solución.

De hecho, evitar el colapso es una consideración muy común en
el aprendizaje contrastivo. Entonces, ¿cómo lo hacemos en este caso? La solución
propuesta por DINO, una vez más, está diseñada en forma empírica y ajusta cuidadosamente la
optimización del parámetro \(\vmu\) (que se actualiza utilizando todas
las muestras del lote) y un hiperparámetro de «temperatura» \(\tau\)
que forma parte de la implementación de \(h_{\vW, \vmu}\) y
se analiza en \Cref{sub:contrastive_learning_architecture}. Si se dispone de
un conjunto específico de hiperparámetros que funcionan bien, esto es, sin duda,
suficiente para garantizar que la representación no colapse. Sin embargo, fuera
de esta configuración específica, es difícil lograr que los modelos
converjan durante el entrenamiento, y este resulta muy inestable.

Para remediar esta situación, analicemos algunas simplificaciones de la
formulación. En primer lugar, en lugar de calcular un vector de probabilidad utilizando una
transformación aprendida \(h_{\vW, \vmu}\) de las características agregadas
\(\vz_{\theta}\), podemos \textit{utilizar directamente la representación agregada},
ignorando la cabeza específica de la tarea (o, lo que es lo mismo,
asignándole la aplicación de identidad). Pero ahora necesitamos una forma de comparar
los vectores directamente. Partiendo de nuestra hipótesis del
\Cref{ch:deep-networks}, según la cual las buenas representaciones deben tener
geometría euclidiana (subespacial), una medida mucho más natural de la
diferencia es la \textit{distancia cuadrada \(\ell^{2}\)}
\(d_{\ell^{2}} \colon \R^{d} \times \R^{d} \to \R\), definida como
\begin{equation}\label{eq:cosine_similarity}
    d_{\ell^{2}}(\vx, \vy) \doteq \frac{1}{2}\norm{\vx -
    \vy}_{2}^{2}, \qquad  \forall \vx, \vy \in \R^{d}.
\end{equation}
Esta medida basada en la distancia es aún más eficiente de calcular que la
medida de entropía cruzada. Por lo tanto, \(d_{\ell^{2}}\) sustituye a
\(d_{\CE}\) en nuestra simplificación.

Anteriormente, el colapso se evitaba utilizando trucos para actualizar \(\vmu\) y
\(\tau\). En nuestra simplificación, si comparamos las características dentro
del espacio de representación en lugar de convertirlas en probabilidades,
no contamos con ninguno de estos parámetros, por lo que debemos considerar una
forma diferente de evitar el colapso. Para ello, volvemos a los
fundamentos. La idea fundamental para evitar el colapso consiste en que 
diferentes muestras produzcan características distintas. En
otras palabras, aspiramos a que la \textit{matriz de covarianza} de las características tenga
una \textit{magnitud} significativa. Pero a partir de los
\Cref{ch:lossy-compression,ch:deep-networks}, ya tenemos una cantidad
que mide el tamaño de la matriz de covarianza. Concretamente, utilizamos la
sencilla \textit{tasa de codificación gaussiana}
\(R\) (a nivel de población) para garantizar que las características de las vistas globales de las diferentes
imágenes, que contienen información global distinta, estén bien separadas
y no colapsen (por lo tanto, se expandan). La función de pérdida modificada 
\(\cL_{\simdino}\) queda así:
\begin{align}\label{eq:simdino_loss}
    \cL_{\simdino}(\theta) &\doteq \Ex[d_{\ell^{2}}(\vz_{\theta}(\vX_{g}),
    \vz_{\theta}(\vX_{c}))] \\ 
    &\qquad - \frac{\gamma}{2}\log\det\rp{\vI +
    \frac{d}{\eps^{2}}\Cov(\vz_{\theta}(\vX_{g}))},
\end{align}
donde \(\eps > 0\) es un valor fijo y las expectativas correspondientes se calculan, como
antes, sobre los datos \(\vX\), la vista global \(v_{g}\) y otra
vista (local o global) \(v_{c}\). La pérdida en \eqref{eq:simdino_loss}
es esencialmente la misma pérdida introducida en la Sección \ref{sec:SimDINO}
del Capítulo \ref{ch:lossy-compression} para el DINO simplificado
(``SimDINO'' \cite{wu2025simplifying}). Como veremos, cuando se implementa correctamente, funciona
al menos tan bien como el DINO original.

\subsection{Arquitectura: Vision 
Transformer}\label{sub:contrastive_learning_architecture}

Para la arquitectura, utilizamos un transformer visual estándar. A continuación se explica
cómo funciona formalmente dicha arquitectura en el contexto de los datos de imagen.
Recordemos de la \Cref{sec:experiment_setup} que una arquitectura de codificador consta de cuatro
componentes: una incrustación, una red
troncal, un extractor de características y una cabeza específica para la tarea. A continuación analizaremos
estas cuatro partes.

\begin{figure}
    \centering
    \includegraphics[width=0.8\textwidth]{\toplevelprefix/chapters/chapter8/figs/patchify.pdf}
    \caption{\small\textbf{Un ejemplo de una imagen dividida en \(5
        \times 5\) parches cuadrados, que se colocan en orden de trama.}
        Cada parche tiene el mismo tamaño, y la cuadrícula de parches tiene
        la forma \((N_{H}, N_{W}) = (5, 5)\). La cuadrícula de parches se
    despliega entonces en una secuencia de longitud \(5 \times 5 = 25\) en orden de trama.}
    \label{fig:patchify_rasterize}
\end{figure}

\begin{figure}
    \centering
    \includegraphics[width=\textwidth]{\toplevelprefix/chapters/chapter8/figs/transformer_embedding.pdf}
    \caption{\small\textbf{El proceso de incrustación del transformer.} Dada
        una secuencia de parches desplegados en orden de trama \(\vX^{\patch}\),
        cada parche edsenrollado se proyecta linealmente en el espacio de características,
        y se le asigna una codificación posicional (aditiva) y un
        token adicional conocido como token de clase. La salida es la
        característica de entrada de la primera capa \(\vZ_{\theta}^{1}(\vX) =
    f_{\theta}^{\emb}(\vX)\).}
    \label{fig:transformer_embedding}
\end{figure}

\paragraph{Incrustación.} Dados unos datos de imagen \(\vX \in \cI\), los incrustamos
como una secuencia de tokens en \(\R^{d}\) utilizando la aplicación
\(f_{\theta}^{\emb}\), de la siguiente manera. Los dos primeros pasos se muestran
en la \Cref{fig:patchify_rasterize} y los dos últimos se muestran en la
\Cref{fig:transformer_embedding}.
\begin{enumerate}
    \item En primer lugar, convertimos los datos de la imagen \(\vX\) en una secuencia de
        parches de forma \((C, P_{H}, P_{W})\), donde \(P_{H}\) y
        \(P_{W}\) son las dimensiones de los parches. Suponemos que \(P_{H}\)
        y \(P_{W}\) dividen uniformemente la altura y el ancho de \(\vX\),
        respectivamente (en la notación de
            \Cref{sub:contrastive_learning_objective} suponemos que
            \(P_{H}\) y \(P_{W}\) dividen uniformemente \(S_{\loc}\) y
        \(S_{\glo}\)). Supongamos que la cuadrícula de parches resultante tiene
        \(N_{H}\) filas y \(N_{W}\) columnas.
    \item Desenrollamos cada parche en un vector de longitud \(D \doteq
        CP_{H}P_{W}\). Hay \(N \doteq N_{H}N_{W}\) vectores de parches,
        los cuales colocamos en «orden de trama» (arriba a la izquierda \(\to\)
        arriba a la derecha \(\to\) abajo a la izquierda \(\to\) abajo a la derecha) en una
        matriz \(\vX^{\patch} \in \R^{D \times N}\), donde
        \(\vX^{\patch} \doteq f^{\patch}(\vX)\). Obsérvese que \(D\)
        depende únicamente del tamaño del parche y del número de canales. Dado que
        esta última cantidad suele ser constante entre las muestras del
        mismo conjunto de datos, \(D\) es igual para todas las imágenes del
        conjunto de datos, mientras que \(N\) varía según el tamaño de las imágenes.
    \item A continuación, realizamos la siguiente operación sobre \(\vX^{\patch}
        \in \R^{D \times N}\) para proyectarla en \(\R^{d \times n}\)
        donde \(n \doteq N + 1\):
        \begin{equation}
            \vX^{\patch} \mapsto [\vz_{\cls}^{1}, \vW^{\emb}\vX] + \vE^{\pos}.
        \end{equation}
        Aquí tenemos tres parámetros entrenables: \(\vW^{\emb}\),
        \(\vz_{\cls}^{1}\) y \(\vE^{\pos}\), cuya finalidad es la siguiente:
        \begin{itemize}
            \item \(\vW^{\emb} \in \R^{d \times D}\) es una matriz
                que proyecta cada vector de parche en una característica de token.
            \item \(\vz_{\cls}^{1} \in \R^{d}\) es lo que se conoce como
                \textit{token de clase} o \textit{token de registro}. El
                token de clase contiene, de manera heurística, información global de
                todos los datos y se utiliza para tareas posteriores. En
                el marco de las redes profundas compresivas descrito en el
                \Cref{ch:deep-networks}, esperamos que el token de clase
                se proyecte en los mismos subespacios que los
                tokens destacados o semánticamente relevantes durante el
                desarrollo de la pasada hacia adelante.
            \item \(\vE^{\pos} \in \R^{d \times N}\) es lo que se conoce como
                \textit{codificación posicional}, que distingue
                entre los tokens de los distintos parches. Es decir,
                las características de los tokens deben tener información posicional, de modo que
                el mapa global \(f^{\pre}\) no sea invariante ante
                las permutaciones de los parches y \(\vE^{\pos}\)
                inserta esta información posicional.
                \begin{itemize}
                    \item En este caso de DINO, en el que el transformer
                        recibe imágenes de distintos tamaños, aprendemos una
                        codificación posicional para el tamaño más grande
                        recibido durante el entrenamiento e interpolamos para
                        obtener las codificaciones posicionales de las entradas de menor tamaño.
                \end{itemize}
        \end{itemize}
\end{enumerate}
Así pues, al final tenemos
\begin{equation}\label{eq:definition_of_embedding_module}
    f_{\theta}^{\emb}(\vX) \doteq \mat{\vz_{\cls}^{1},
    \vW^{\emb}f^{\patch}(\vX) + \vE^{\pos}}.
\end{equation}
Todos los parámetros \(\vz_{\cls}^{1}, \vW^{\emb}, \vE^{\pos}\) están
incluidos en el conjunto de parámetros \(\theta\).

\begin{figure}
    \centering
    \includegraphics[width=\textwidth]{\toplevelprefix/chapters/chapter8/figs/transformer_backbone.pdf}
    \caption{\small\textbf{Una capa \(f_{\theta}^{\ell}\) de la
        red troncal del transformer.} Las características de entrada pasan por
        bloques de normalización de capa, auto-atención multicanal y perceptrón multicapa
    en secuencia para formar las características de salida de la capa.}
    \label{fig:transformer_backbone}
\end{figure}

\paragraph{Red troncal.} Dada una secuencia de incrustaciones
\(\vZ_{\theta}^{1}(\vX) \doteq f_{\theta}^{\emb}(\vX) \in
(\R^{d})^{*}\), la procesamos utilizando el mapa de la red troncal
\(f_{\theta}^{\backbone}\) de la siguiente manera y tal como se muestra en la
\Cref{fig:transformer_backbone}. La función
\(f_{\theta}^{\backbone}\) se compone de \(L\) \textit{capas}
\(f_{\theta}^{\ell}\), es decir,
\begin{equation}
    f_{\theta}^{\backbone} = f_{\theta}^{L} \circ \cdots \circ f_{\theta}^{1}.
\end{equation}
La capa \(f_{\theta}^{\ell}\) tiene la siguiente implementación:
\begin{align}\label{eq:vit-res-block}
    \vZ_{\theta}^{\ell + 1/2}(\vX)
    &= \vZ_{\theta}^{\ell}(\vX) +
    \MHSA_{\theta}^{\ell}(\LN_{\theta}^{1, \ell}(\vZ_{\theta}^{\ell}(\vX))) \\
    \vZ_{\theta}^{\ell + 1}(\vX)
    &= \vZ_{\theta}^{\ell + 1/2}(\vX) +
    \MLP_{\theta}^{\ell}(\LN_{\theta}^{2, \ell}(\vZ_{\theta}^{\ell + 1/2}(\vX)))
\end{align}
y \(f_{\theta}^{\ell}\) se define de tal manera que
\(f_{\theta}^{\ell}(\vZ_{\theta}^{\ell}(\vX)) \doteq
\vZ_{\theta}^{\ell + 1}(\vX)\). Aquí hemos utilizado algunos operadores,
como \(\MHSA_{\theta}^{\ell}, \MLP_{\theta}^{\ell}\) y
\(\LN_{\theta}^{i, \ell}\), que se definen de la siguiente manera:
\begin{itemize}
    \item El operador \(\MHSA_{\theta}^{\ell}\) es
        una auto-atención multicabezal, predecesora de la auto-atención
        multicabezal en subespacios (véase \Cref{ch:deep-networks}). La
        formulación es la siguiente:
        \begin{align}\label{eq:mhsa}
            &\scalebox{0.95}{\(\displaystyle\MHSA_{\theta}^{\ell}(\vZ) 
            \doteq \vU_{\out}^{\ell}\mat{\SA([\vU_{\query}^{1,
                    \ell}]^{\top}\vZ, [\vU_{\attnkey}^{1, \ell}]^{\top}\vZ,
                [\vU_{\val}^{1, \ell}]^{\top}\vZ) \\ \vdots
                \\ \SA([\vU_{\query}^{K, \ell}]^{\top}\vZ,
                    [\vU_{\attnkey}^{K, \ell}]^{\top}\vZ, [\vU_{\val}^{K,
            \ell}]^{\top}\vZ)} + \vb_{\out}^{\ell}\vone_{n}^{\top},\)} \\
            \label{eq:self_attention}
            &\text{donde} \qquad \SA(\vQ, \vK, \vV)
            \doteq \vV
            \underbrace{\softmax\rp{\frac{\vK^{\top}\vQ}{\sqrt{p}}}}_{\doteq
            \vA(\vQ, \vK)}
        \end{align}
        donde \(p\) es un entero positivo, \(\vU_{\query}^{k, \ell},
            \vU_{\attnkey}^{k, \ell}, \vU_{\val}^{k, \ell} \in \R^{d
        \times p}\), \(\vU_{\out}^{\ell} \in \R^{d \times Kp}\), y
        \(\vb_{\out}^{\ell} \in \R^{d}\) son parámetros entrenables
        contenidos en el conjunto de parámetros \(\theta\), y la
        función \(\softmax\) se define por columnas como
        \begin{align}
            \softmax(\vM)
            &\doteq \mat{\softmax(\vm_{1}) & \cdots & \softmax(\vm_{p})}, \\
            \forall \vM
            &= \mat{\vm_{1}, \dots, \vm_{p}} \in \R^{n \times p}.
        \end{align}
        En la práctica, las dimensiones suelen elegirse de tal manera que \(Kp
        = d\). Los términos
        \begin{eqnarray}
            \vA_{\theta}^{k, \ell}(\vZ) &\doteq&
            \vA([\vU_{\query}^{k, \ell}]^{\top}\vZ,
            [\vU_{\attnkey}^{k, \ell}]^{\top}\vZ), \nonumber \\
            \SA_{\theta}^{k, \ell}(\vZ) &\doteq&
            \SA([\vU_{\query}^{k, \ell}]^{\top}\vZ,
                [\vU_{\attnkey}^{k, \ell}]^{\top}\vZ, [\vU_{\val}^{k,
            \ell}]^{\top}\vZ)
            \label{eq:attention_map}
        \end{eqnarray}
        también se conocen como el \textit{\(k\)-ésimo mapa de atención} y
        la \textit{\(k\)-ésima salida del cabezal de atención} en la capa \(\ell\),
        respectivamente. Además, la operación \(\SA(\vQ, \vK,
        \vV)\) se puede calcular de manera extremadamente eficiente utilizando
        software especializado como FlashAttention
        \citep{shah2025flashattention}.
    \item El \(\MLP_{\theta}^{\ell}\) es un perceptrón de dos capas, una
        no linealidad habitual en las redes profundas y tiene la forma
        \begin{equation}
            \MLP_{\theta}^{\ell}(\vZ) \doteq
            \vW_{\down}^{\ell}\ReLU(\vW_{\up}^{\ell}\vZ +
            \vb_{\up}^{\ell}\vone_{n}^{\top}) +
            \vb_{\down}^{\ell}\vone_{n}^{\top}
        \end{equation}
        donde \(\vW_{\up}^{\ell} \in \R^{q \times d},
            \vW_{\down}^{\ell} \in \R^{d \times q}, \vb_{\up}^{\ell} \in
        \R^{q}, \vb_{\down}^{\ell} \in \R^{d}\) son parámetros entrenables
        que también se encuentran en el conjunto de parámetros \(\theta\),
        y \(\ReLU\) es la no linealidad ReLU elemento por elemento, es decir,
        \(\ReLU(\vM)_{ij} = \max\{M_{ij}, 0\}\).
    \item Cada norma de capa \(\LN_{\theta}^{i, \ell}\) para \(i \in \{1,
        2\}\) es una normalización estándar, que se aplica columna por columna
        a cada característica de token de manera independiente:
        \begin{equation}
            \LN_{\theta}^{i, \ell}(\vZ) = \LN_{\theta}^{i,
            \ell}(\mat{\vz_{1}, \dots, \vz_{n}}) =
            \mat{\LN_{\theta}^{i, \ell}(\vz_{1}), \dots,
            \LN_{\theta}^{i, \ell}(\vz_{n})}
        \end{equation}
        y tiene la forma
        \begin{align}
            \LN_{\theta}^{i, \ell}(\vz) 
            &= \frac{\vz
            - \operatorname{mean}(\vz)\vone_{d}}{\norm{\vz
            - \operatorname{mean}(\vz)\vone_{d}}_{2}} \hada
            \valpha^{i, \ell} + \vbeta^{i, \ell} \\ 
            \text{donde}
            \qquad \operatorname{mean}(\vz) &= \frac{1}{d}\vone_{d}^{\top}\vz
        \end{align}
        donde \(\hada\) denota la multiplicación elemento por elemento, y
        \(\valpha^{i, \ell}, \vbeta^{i, \ell} \in \R^{d}\) son
        parámetros entrenables contenidos en el conjunto de parámetros
        \(\theta\). El operador de norma de capa sirve como una especie de
        normalización de cada token, donde la escala de cada token
        posteriormente es aprendible y compartida entre todos los tokens.
\end{itemize}

El transformer es una de las arquitecturas de redes neuronales
más populares de la historia y se utiliza en aplicaciones de casi todos los campos
del aprendizaje profundo.

\paragraph{Extractor de características.} Utilizamos un paso de posprocesamiento
\(f_{\theta}^{\ext}\) que extrae la \textit{característica del token de clase
}, la cual (recordemos) es la característica destinada a contener información agregada
sobre la imagen de entrada y le aplica una MLP y
una normalización. Es decir, tenemos
\begin{align}
    \vz_{\theta}(\vX) &\doteq f_{\theta}^{\ext}(\vZ_{\theta}(\vX)) =
    f_{\theta}^{\ext}([\vz_{\theta}^{1}(\vX), \dots,
    \vz_{\theta}^{n}(\vX)]) \\
    &\doteq
    \frac{\MLP_{\theta}^{\ext}(\vz_{\theta}^{1}(\vX))}{\norm{\MLP_{\theta}^{\ext}(\vz_{\theta}^{1}(\vX))}_{2}}.
\end{align}

\paragraph{Cabezal específico para cada tarea («DINO»).} Para DINO, utilizamos el
cabezal específico para la tarea \(h_{\vW, \vmu}\). Para SimDINO, no utilizamos
ningún cabezal específico, tal como se ha descrito anteriormente.

\subsection{Estrategia de optimización}\label{sub:contrastive_learning_optimization}

\begin{figure}
    \centering
    \includegraphics[width=\textwidth]{\toplevelprefix/chapters/chapter8/figs/dino_pipeline.pdf}
    \caption{\small \textbf{El canal de procesamiento DINO.} Las características del alumno y
        las del maestro se calculan para cada entrada. El objetivo
        intenta alinear las características del alumno con las del maestro
        proyectando ambos conjuntos de características en un simplex de probabilidad de alta dimensión
        y calculando una pérdida de entropía cruzada. Cabe destacar que,
        debido al «stop-grad», el gradiente solo se calcula
    con respecto a las salidas de los parámetros del alumno.}
    \label{fig:dino_pipeline}
\end{figure}

\paragraph{Optimización para DINO.} Contamos con una función de pérdida y una
arquitectura; a continuación describiremos la estrategia de optimización. La
estrategia de optimización para DINO utiliza \textit{dos conjuntos de pesos para
la misma arquitectura}: los pesos del \textit{alumno}
\(\theta_{\student}\) y los pesos del \textit{maestro}
\(\theta_{\teacher}\). Estos corresponden a dos redes neuronales diferentes,
denominadas red maestra y red de alumno, con la
misma arquitectura. La red maestra codifica todas las vistas globales,
mientras que la red de alumno codifica todas las demás vistas. El objetivo de
la función de pérdida consiste en destilar las salidas de la red maestra hacia el modelo de alumno.
Es decir, entrenamos con la función de pérdida \(\cL_{\dino{}-\student\teacher}\):
\begin{equation}\label{eq:dino_loss_teacherstudent}
    \cL_{\dino{}-\student\teacher}(\theta_{\student},
    \theta_{\teacher}, \vW_{\student}, \vW_{\teacher}, \vmu) \doteq
    \Ex[d_{\CE}(\vp_{\theta_{\teacher}, \vW_{\teacher},
    \vmu}(\vX_{g}), \vp_{\theta_{\student}, \vW_{\student}}(\vX_{c}))].
\end{equation}
Ahora podemos describir detalladamente el flujo de trabajo general de DINO, que se muestra en la
\Cref{fig:dino_pipeline}.

Aunque es fácil razonar sobre \eqref{eq:dino_loss_teacherstudent},
en la práctica resulta imposible implementar algoritmos de optimización
como el descenso por gradiente con una pérdida dada por
\(\cL_{\dino{}-\student\teacher}\). Esto se debe a que las expectativas
en la pérdida son imposibles de evaluar y mucho menos de calcular
su gradiente. En este escenario común, aproximamos la
expectativa mediante muestras finitas. Es decir, en cada paso de tiempo \(k\) hacemos lo siguiente:
\begin{itemize}
    \item Seleccionamos una submuestra \(B\) de puntos de datos de nuestro conjunto de datos
        \(\{\vX_{1}^{(k)}, \dots, \vX_{B}^{(k)}\} \subset \cI\).
    \item Para cada punto de datos \(\vX_{b}^{(k)}\), muestreamos \(M_{\glo}\)
        vistas globales \(v_{b, g}^{(k), i}\) y \(M_{\loc}\) vistas
        locales \(v_{b, \ell}^{(k), i}\). Se aplican las vistas a
        \(\vX_{b}^{(k)}\) para obtener \(\vX_{b, g}^{(k), i} \doteq
        v_{b, g}^{(k), i}(\vX_{b}^{(k)})\) y \(\vX_{b, \ell}^{(k),
        i} \doteq v_{b, \ell}^{(k), i}(\vX_{b}^{(k)})\).
    \item Para cada vista \textit{local} \(\vX_{b, \ell}^{(k), i}\),
        se calculan las siguientes magnitudes:
        \begin{align}
            &\vz_{\theta_{\student}}(\vX_{b, \ell}^{(k), i}) \doteq
            (f_{\theta_{\student}}^{\ext} \circ
            f_{\theta_{\student}})(\vX_{b, \ell}^{(k), i}), \\
            &\vp_{\theta_{\student}, \vW_{\student}}(\vX_{b,
            \ell}^{(k), i}) \doteq h_{\vW_{\student},
            \vzero_{m}}(\vz_{\theta_{\student}}(\vX_{b, \ell}^{(k), i}(\theta)))
        \end{align}
        y para cada vista \textit{global} \(\vX_{b, g}^{(k), i}\),
        se calculan las siguientes magnitudes (abusando de la notación):
        \begin{align}
            &\vz_{\theta_{\student}}(\vX_{b, g}^{(k), i}) \doteq
            (f_{\theta_{\student}}^{\ext} \circ
            f_{\theta_{\student}})(\vX_{b, g}^{(k), i}), \\
            &\vp_{\theta_{\student}, \vW_{\student}}(\vX_{b, g}^{(k),
            i}) \doteq h_{\vW_{\student},
            \vzero_{m}}(\vz_{\theta_{\student}}(\vX_{b, g}^{(k), i})), \\
            &\vz_{\theta_{\teacher}}(\vX_{b, g}^{(k), i}) \doteq
            (f_{\theta_{\teacher}}^{\ext} \circ
            f_{\theta_{\teacher}})(\vX_{b, g}^{(k), i}), \\
            &\vp_{\theta_{\teacher}, \vW_{\teacher}, \vmu}(\vX_{b,
            g}^{(k), i}) \doteq h_{\vW_{\teacher},
            \vmu}(\vZ_{\theta_{\teacher}}(\vX_{b, g}^{(k), i})).
        \end{align}
    \item Se calcula la \textit{pérdida aproximada} (sustituta)
        \(\hat{\cL}_{\dino-\student\teacher}^{(k)}\), definida así:
        \begin{align}\label{eq:dino_loss_teacherstudent_empirical}
            &\hat{\cL}_{\dino{}-\student\teacher}^{(k)}(\theta_{\student},
            \theta_{\teacher}, \vW_{\student}, \vW_{\teacher}, \vmu) \doteq
            \frac{1}{BM_{\glo}(M_{\glo} + M_{\loc} - 1)}\sum_{b =
            1}^{B}\sum_{i = 1}^{M_{\glo}}\\
            &\Bigg[\sum_{j =
                1}^{M_{\loc}}d_{\CE}(\vp_{\theta_{\teacher},
                    \vW_{\teacher}, \vmu}(\vX_{b, g}^{(k), i}),
                    \vp_{\theta_{\student}, \vW_{\student}}(\vX_{b,
                \ell}^{(k), j})) \\ 
                &\qquad \qquad \qquad \qquad + \sum_{\substack{j = 1 \\ j \neq
                i}}^{M_{\glo}}d_{\CE}(\vp_{\theta_{\teacher},
                    \vW_{\teacher}, \vmu}(\vX_{b, g}^{(k), i}),
                    \vp_{\theta_{\student}, \vW_{\student}}(\vX_{b, g}^{(k),
            j}))\Bigg]
        \end{align}
        así como sus gradientes con respecto a
        \(\theta_{\student}\) y \(\vW_{\student}\), que deben
        calcularse bajo el supuesto de que \(\theta_{\teacher}\),
        \(\vW_{\teacher}\) y \(\vmu\) son constantes --- es decir,
        que están \textit{desacoplados del grafo computacional}
        y no dependen de \(\theta_{\student}\) y \(\vW_{\student}\).
    \item Actualizar los parámetros del alumno \(\theta_{\student}\) y
        \(\vW_{\student}\) mediante un algoritmo de optimización iterativo basado en el gradiente,
        y actualizar \(\theta_{\teacher}\),
        \(\vW_{\teacher}\) y \(\vmu\) mediante medias móviles exponenciales
        con parámetros de decaimiento \(\nu^{(k)}\), \(\nu^{(k)}\),
        y \(\rho^{(k)}\) respectivamente, es decir,
        \begin{align}
            &(\theta_{\student}^{(k + 1)}, \vW_{\student}^{(k + 1)})
            = \textsc{OptUpdate}^{(k)}(\theta_{\student}^{(k)},
                \vW_{\student}^{(k)}; \nabla_{(\theta_{\student},
            \vW_{\student})}\hat{\cL}_{\dino-\student\teacher}^{(k)}) \\
            &\theta_{\teacher}^{(k + 1)}
            = \nu^{(k)}\theta_{\teacher}^{(k)} + (1 -
            \nu^{(k)})\theta_{\student}^{(k + 1)} \\
            &\vW_{\teacher}^{(k + 1)}
            = \nu^{(k)}\vW_{\teacher}^{(k)} + (1 -
            \nu^{(k)})\vW_{\student}^{(k + 1)} \\
            &\vmu^{(k + 1)}
            = \rho^{(k)}\vmu^{(k)} + (1 -
            \rho^{(k)})\cdot\frac{1}{BM_{\glo}}\sum_{b =
            1}^{B}\sum_{i =
            1}^{M_{\glo}}\vW^{(k)}\vz_{\theta_{\teacher}}(\vX_{b, g}^{(k), i}),
        \end{align}
        Por ejemplo, si el algoritmo de optimización elegido fuera
        el descenso de gradiente estocástico, tendríamos la actualización
        \(\theta_{\student}^{(k + 1)} \doteq \theta_{\student}^{(k)}
            -
        \delta^{(k)}\nabla_{\theta_{\student}}\hat{\cL}_{\dino{}-\student\teacher}^{(k)}\),
        y así sucesivamente.
\end{itemize}
Nótese que el procedimiento de optimización es bastante irregular: aunque
los cuatro parámetros se actualizan en cada iteración, solo dos de ellos se
actualizan directamente a partir de un método basado en gradientes. Los otros dos se
actualizan a partir de medias móviles exponenciales y, de hecho, se tratan como
constantes al calcular cualquier gradiente. Tras el entrenamiento, descartamos
los pesos de los alumnos y utilizamos los pesos de los maestros para nuestra red entrenada
\(f\), ya que se ha demostrado empíricamente que esta media móvil exponencial
estabiliza el modelo resultante (este concepto es conocido como
«promediado de Polyak» o «promedio iterativo»).

La forma en que \(\nu\) y \(\rho\) varían a lo largo de la trayectoria de optimización
(es decir, las funciones \(k \mapsto \nu^{(k)}\) y \(k
\mapsto \rho^{(k)}\)) son hiperparámetros o decisiones de diseño, siendo
\(\nu^{(1)} < 1\) y \(\lim_{k \to \infty}\nu^{(k)} = 1\) por lo general,
y lo mismo para \(\rho\). El hiperparámetro de temperatura \(\tau\),
utilizado en la cabeza DINO \(h_{\vW, \vmu}\), también cambia a lo largo de la
trayectoria de optimización (aunque esta dependencia no se indica explícitamente).

El uso de la pérdida sustituta («empírica») transforma nuestro problema de optimización
inmanejable---como es el caso de la optimización de la pérdida en
\eqref{eq:dino_loss_teacherstudent}---en un problema de optimización estocástica
manejable, que se utiliza para entrenar prácticamente todos los modelos de aprendizaje profundo
del mundo. Esta conversión resulta extremadamente natural
una vez que se han visto algunos ejemplos y esperamos ofrecer estos
ejemplos a lo largo del capítulo.

\begin{figure}
    \centering
    \includegraphics[width=0.7\textwidth]{\toplevelprefix/chapters/chapter8/figs/simdino_pipeline.pdf}
    \caption{\small\textbf{El canal de procesamiento de SimDINO.} Aquí, a diferencia
        del proceso de DINO que se muestra en \Cref{fig:dino_pipeline}, la pérdida se
        calcula directamente sobre las características sin necesidad de
        manipulaciones adicionales. Esto elimina una cantidad de parámetros equivalente a varias matrices grandes
        y simplifica el proceso, al tiempo que hace que
    el entrenamiento sea más estable.}\label{fig:simdino_pipeline}
\end{figure}

\paragraph{Optimización para SimDINO.} El objetivo simplificado de DINO a nivel de población
es muy similar en espíritu, pero considerablemente más sencillo de ejecutar, es decir,
\begin{align}\label{eq:simdino_loss_teacherstudent}
    \cL_{\simdino-\student\teacher}(\theta_{\student}, \theta_{\teacher}) &\doteq
    \Ex\rs{d_{\ell^{2}}(\vz_{\theta_{\teacher}}(\vX_{g}),
    \vz_{\theta_{\student}}(\vX_{c}))} \\
    &\qquad \qquad - \frac{\gamma}{2}\log\det\rp{\vI
+ \frac{d}{\eps^{2}}\Cov(\vz_{\theta_{\student}}(\vX_{g})))}.
\end{align}
Por lo tanto, tal y como se detalla en la \Cref{fig:simdino_pipeline}, el canal de procesamiento de SimDINO
es considerablemente más sencillo que el de DINO. Podemos utilizar una
versión más sencilla del proceso de entrenamiento de DINO para optimizar SimDINO. En
cada paso temporal \(k\), hacemos lo siguiente:
\begin{itemize}
\item Tomar una submuestra \(B\) de puntos de datos de nuestro conjunto de datos \(\{\vX_{1}^{(k)},
    \dots, \vX_{B}^{(k)}\} \subset \cI\).
\item Para cada punto de datos \(\vX_{b}^{(k)}\), mostrar \(M_{\glo}\)
    vistas globales \(v_{b, g}^{(k), i}\) y \(M_{\loc}\) vistas locales
    \(v_{b, \ell}^{(k), i}\). Aplicar las vistas a \(\vX_{b}^{(k)}\) para
    obtener \(\vX_{b, g}^{(k), i} \doteq v_{b, g}^{(k),
    i}(\vX_{b}^{(k)})\) y \(\vX_{b, \ell}^{(k), i} \doteq v_{b,
    \ell}^{(k), i}(\vX_{b}^{(k)})\).
\item Para cada vista \textit{local} \(\vX_{b, \ell}^{(k), i}\), calcular
    \(\vz_{\theta_{\student}}(\vX_{b, \ell}^{(k), i}) \doteq
        (f_{\theta_{\student}}^{\ext} \circ
    f_{\theta_{\student}})(\vX_{b, \ell}^{(k), i})\). Para cada
    vista \textit{global} \(\vX_{b, g}^{(k), i}\), calcular
    \(\vz_{\theta_{\student}}(\vX_{b, g}^{(k), i}) \doteq
        (f_{\theta_{\student}}^{\ext} \circ
    f_{\theta_{\student}})(\vX_{b, g}^{(k), i})\) y
    \(\vz_{\theta_{\teacher}}(\vX_{b, g}^{(k), i}) \doteq
        (f_{\theta_{\teacher}}^{\ext} \circ
    f_{\theta_{\teacher}})(\vX_{b, g}^{(k), i})\).
\item Calcular la \textit{pérdida sustitutiva o aproximada}
    \(\hat{\cL}_{\simdino-\student\teacher}^{(k)}\), definida de la siguiente manera:
    \begin{align}\label{eq:simdino_loss_teacherstudent_empirical}
        &\hat{\cL}_{\simdino{}-\student\teacher}^{(k)}(\theta_{\student},
        \theta_{\teacher})  \\
        \doteq &
        \scalebox{0.975}{\(\displaystyle\frac{1}{BM_{\glo}(M_{\glo} + M_{\loc} - 1)}\sum_{b =
        1}^{B}\sum_{i = 1}^{M_{\glo}}\Bigg[\sum_{j =
            1}^{M_{\loc}}d_{\ell^{2}}(\vz_{\theta_{\teacher}}(\vX_{b,
                g}^{(k), i}), \vz_{\theta_{\student}}(\vX_{b, \ell}^{(k),
            j}))\)} \\
            & + \sum_{j =
            1}^{M_{\glo}}d_{\ell^{2}}(\vz_{\theta_{\teacher}}(\vX_{b,
                g}^{(k), i}), \vz_{\theta_{\student}}(\vX_{b, g}^{(k),
        j}))\Bigg] \\
        & - \frac{\gamma}{M_{\glo}}\sum_{i =
        1}^{M_{\glo}}R_{\eps}([\vz_{\theta_{\student}}(\vX_{1,
                g}^{(k), i}), \dots, \vz_{\theta_{\student}}(\vX_{B, g}^{(k),
        i})])
    \end{align}
    donde \(R_{\eps}\) es la tasa de codificación gaussiana estimada mediante
    muestras finitas, tal como se describe en el \Cref{ch:deep-networks}. El
    gradiente de \(\hat{\cL}_{\simdino-\student\teacher}^{(k)}\) con
    respecto a \(\theta_{\student}\) debe (de nuevo) calcularse,
    bajo la suposición de que \(\theta_{\teacher}\) es constante.
\item Actualizar los parámetros del alumno \(\theta_{\student}\) mediante un
    algoritmo de optimización iterativo basado en descensos de gradiente y actualizar
    \(\theta_{\teacher}\) mediante una media móvil exponencial con
    el parámetro de decaimiento \(\nu^{(k)}\), es decir,
    \begin{align}
        \theta_{\student}^{(k + 1)}
        &= \textsc{OptUpdate}^{(k)}(\theta_{\student}^{(k)};
        \nabla_{\theta_{\student}}\hat{\cL}_{\simdino-\student\teacher}^{(k)})
        \\
        \theta_{\teacher}^{(k + 1)}
        &= \nu^{(k)}\theta_{\teacher}^{(k)} + (1 -
        \nu^{(k)})\theta_{\student}^{(k + 1)}.
    \end{align}
\end{itemize}
Una vez más, reiteramos que el gradiente solo se calcula
respecto a \(\theta_{\student}\), considerando \(\theta_{\teacher}\) como una
constante. En este caso, cabe señalar que, si bien la elección de \(\nu\) sigue siendo una
decisión de diseño, los hiperparámetros \(\rho\) y \(\tau\) se eliminan.

\subsection{Metodología de evaluación}\label{sub:contrastive_learning_evals}
Hay varias formas de evaluar un modelo de transformer entrenado. En
esta sección destacamos dos. Definamos la vista de \textit{recorte central}
\(v_{\cc} \colon \cI \to \cI\), que es un \textit{recorte redimensionado
determinístico}:
\begin{itemize}
\item ajusta el tamaño de la imagen de modo que el lado más corto tenga un tamaño de
    \(S_{\rsz}\) (similar a los recortes con tamaño aleatorio y el parámetro de porcentaje de área
    \(1\));
\item luego toma el recorte \textit{central} \(S_{\cc} \times S_{\cc}\);
\end{itemize}
de modo que la forma final es \((C, S_{\cc}, S_{\cc})\). Obsérvese que la
vista \(v_{\cc}\) es completamente determinista dada una entrada. Para una
entrada \(\vX\), escribimos \(\vX_{\cc} \doteq v_{\cc}(\vX)\). Aquí
\(S_{\cc} \leq S_{\rsz}\).

\paragraph{Sondeo lineal.}

La primera forma de evaluar un modelo de codificador
\(\vX \mapsto \vz_{\theta}(\vX)\), y la más independiente de la arquitectura, es emplear el \textit{sondeo lineal} (linear probing). El sondeo lineal consiste, en pocas palabras, en aplicar una regresión logística
sobre las características agregadas calculadas por el codificador. Esto
nos indica cuánta información semántica existe en las representaciones,
así como la facilidad con la que se puede extraer dicha información. (Es decir: ¿
hasta qué punto las características de las imágenes con diferentes significados se encuentran
en diferentes subespacios del espacio de características?)

De manera más formal, supongamos que queremos evaluar la calidad
y la fidelidad de las características del codificador en datos de imagen-etiqueta
\((\vX, \vy)\), donde hay \(N_{\cls}\) clases y \(\vy \in
\{0, 1\}^{N_{\cls}}\) es una «codificación one-hot» (es decir, ceros en todas las
posiciones excepto un \(1\) en la posición \(i\) si \(\vX\) está en la
clase \(i\)). Una forma de hacerlo es resolver el problema de regresión logística
\begin{equation}\label{eq:linear_probing}
\min_{\vW \in \R^{N_{\cls} \times d}}\Ex[\CE(\vy, \vW \vz_{\theta}(\vX_{\cc}))].
\end{equation}
En la práctica, si disponemos de datos etiquetados \(\{(\vX_{b}, \vy_{b})\}_{b
= 1}^{B}\), podemos resolver el problema de regresión logística \textit{empírica}
(similar a \eqref{eq:dino_loss_teacherstudent}
vs.~\eqref{eq:dino_loss_teacherstudent_empirical}) dado por
\begin{equation}\label{eq:linear_probing_empirical}
\min_{\vW \in \R^{N_{\cls} \times d}}\frac{1}{B}\sum_{b =
1}^{B}\CE(\vy_{b}, \vW \vz_{\theta}(\vX_{b, \cc})).
\end{equation}
Este problema es un problema de optimización convexa en \(\vW\), por lo que
puede resolverse de manera eficiente mediante el descenso de gradiente (estocástico) o una
gran variedad de otros algoritmos. Esta sonda lineal, junto con el
codificador, puede utilizarse como clasificador y podemos evaluar la
precisión de la clasificación. La práctica habitual consiste en entrenar el modelo
primero con un conjunto de datos de gran tamaño (como ImageNet-1K), luego entrenar la
sonda lineal con un conjunto de datos (como el conjunto de datos de entrenamiento de CIFAR-10) y
evaluarla con un tercer conjunto de datos (de «validación») que se extrae de la
misma distribución que el segundo (como el conjunto de datos de evaluación
de CIFAR-10).

\paragraph{Vecinos \(k\)-más cercanos.} También podemos evaluar el
rendimiento de las características en tareas de clasificación \textit{sin
necesidad de entrenar explícitamente un clasificador} utilizando el algoritmo de los vecinos \(k\)-más cercanos
para obtener una etiqueta predicha promedio. Es decir, dado un
conjunto de datos \(\{\vz_{b}\}_{b = 1}^{B} \subseteq \R^{d}\), definimos los
vecinos \(k\)-más cercanos de otro punto \(\vz \in \R^{d}\) como
\(\operatorname{NN}_{k}(\vz, \{\vz_{b}\}_{b = 1}^{B})\). Utilizando esta
notación, podemos calcular la etiqueta predicha \(\hat{\vy}_{\theta}(\vX
\mid \{(\vX_{b}, \vy_{b})\}_{b = 1}^{B})\) como
\begin{eqnarray}
&&\hat{\vy}_{\theta}(\vX \mid \{(\vX_{b}, \vy_{b})\}_{b = 1}^{B}) =
\vone(i^{\star}) \quad \text{donde} \nonumber \\
&& i^{\star} \doteq \argmax_{i \in [Q]}\sum_{b =
1}^{B}\vy_{b}\indvar[\vz_{\theta}(\vX_{\cc, b}) \in
\operatorname{NN}_{k}(\vz_{\theta}(\vX_{\cc}))].
\end{eqnarray}
Aquí, \(\vone(i) \in \Delta_{N_{\cls}}\) es (mediante un abuso de notación,
véanse las variables indicadoras) el vector de probabilidad «one-hot» con valores en
\(i\), es decir, \(1\) en la \(i\)-ésima coordenada y \(0\) en todas las demás.
Es decir, este procedimiento toma la etiqueta más común entre los \(k\)
puntos más cercanos en el espacio de características. La precisión de la clasificación de los \(k\) vecinos más cercanos
es simplemente la precisión de esta etiqueta predicha, es decir,
\begin{equation}
\Ex_{\vX, \vy}[\indvar(\hat{\vy}_{\theta}(\vX \mid \{(\vX_{b},
\vy_{b})\}_{b = 1}^{B}) = \vy)]
\end{equation}
o, más comúnmente, su versión empírica correspondiente, en la que \((\vX,
\vy)\) recorre un conjunto de datos finito (que no son las muestras existentes
\((\vX_{b}, \vy_{b})\) que se utilizan para los \(k\) vecinos).

\paragraph{Fiabilidad de los mapas de atención.}

Otra forma de evaluar el rendimiento de las representaciones, en el caso de un
codificador basado en transformers, consiste en examinar la fidelidad de los
mapas de atención \(\vA^{L, k} \in \R^{n \times n}\), tal y como se definen en la
\Cref{eq:attention_map}, en la última capa \(L\), y que se obtienen mediante el
siguiente canal de procesamiento:
\begin{equation}
\scalebox{0.925}{\(\displaystyle\vX \mapsto \cdots \mapsto \vZ^{L - 1} = [\underbrace{\vz_{1}^{L -
    1}}_{\text{class token}}, \underbrace{\vz_{2}^{L - 1} \dots,
\vz_{n}^{L - 1}}_{\text{tokens de parche}}] \mapsto \vA^{k, L} =
\mat{\vA_{1, 1}^{k, L} & \vA_{1, 2:}^{k, L} \\ \vA_{2:, 1}^{k, L} &
\vA_{2:, 2:}^{k, L}}.\)}
\end{equation}
En concreto, analizamos qué revelan los mapas de atención de una entrada determinada
sobre los objetos salientes de la imagen de entrada, es decir, qué
partes de la imagen proporcionan la información más relevante a nivel global para
el token de clase. Una forma concreta de hacerlo consiste en examinar la
componente del mapa de atención en el que el token de clase se extrae como
consulta y se elimina de la matriz de valores, es decir, \(\vA_{2:, 1}^{k,
L} \in \R^{1 \times (n - 1)} = \R^{1 \times N}\) o su transpuesta
\(\va^{k, L} = (\vA_{2:, 1}^{k, L})^{\top} \in \R^{N}\). Obsérvese que
este vector \(\va^{k, L}\), al que denominamos «vector de
saliencia en el \(k\)-ésimo cabezal de atención de la capa \(L\)», tiene un valor
para cada parche, \(1, \dots, N\) y utilizamos este valor para describir
la relevancia de cada parche respecto a la información global. En
particular, con fines de visualización, creamos una nueva imagen en la que cada
parche se sustituye por su valor correspondiente en el vector de saliencia,
lo que permite apreciar la contribución de cada parche; a esta imagen la llamamos
``\textit{mapa de saliencia} en la \(k\)-ésima cabeza de atención de la capa
\(L\)''. Para visualizar la relevancia total de cada parche respecto a la
información global en todas las cabezas, podemos calcular el promedio del vector de saliencia,
es decir, \(\tilde{\va}^{L} \doteq \frac{1}{K}\sum_{k =
1}^{K}\va^{k, L}\) y expandirlo en el \textit{mapa de saliencia promedio}.
Los mapas de saliencia promedio deberían resaltar las partes relevantes de la
imagen de entrada.

\paragraph{Detección y segmentación de objetos.}

Podemos evaluar cómo las representaciones captan las propiedades más detalladas
(es decir, más pequeñas o más precisas) de la entrada utilizándolas
para la \textit{segmentación semántica}. A grandes rasgos, esto significa que
utilizamos las características para construir cajas delimitadoras para todos los objetos de la
entrada. Hay varias formas de hacerlo, y varias formas de evaluar
las cajas delimitadoras resultantes en comparación con la verdad de referencia (ground truth). Cada
combinación de métodos se corresponde con una métrica de segmentación
concreta. No las describimos formalmente aquí, ya que no
aportan información especialmente relevante, pero el artículo sobre DINO \citep{caron2021emerging}
y el artículo sobre DINOv2 \citep{oquab2023dinov2} contienen referencias a todas
las métricas que se utilizan en la práctica.

\subsection{Configuración experimental y resultados}
\label{sub:contrastive_learning_experiment_results}

Dado que SimDINO se basa directamente en DINO, comparamos los parámetros óptimos
de DINO, tal y como se indican en su artículo original
\citep{caron2021emerging}, con los mismos parámetros aplicados a SimDINO \cite{wu2025simplifying}
para garantizar una comparación imparcial.

\paragraph{Función objetivo.} Utilizamos \(10\) vistas locales (es decir,
\(M_{\loc} = 10\)) con una resolución de \(96 \times 96\) (es decir, \(S_{\loc} =
96\)) y \(2\) vistas globales (es decir, \(M_{\glo} = 2\)) con una resolución
\(224 \times 224\) (es decir, \(S_{\glo} = 224\)) para todos los experimentos.
Las porciones correspondientes de las imágenes originales recortadas para las vistas locales
y globales son \(p_{\loc} \in [\frac{1}{20}, \frac{3}{10}]\)
y \(p_{\glo} \in [\frac{3}{10}, 1]\) (elegidas al azar por vista).
El tamaño de borde más pequeño dentro de los recortes redimensionados es \(S_{\rsz} = 256\),
y el tamaño de borde de la vista del recorte central (evaluación) es \(S_{\cc} = 224\).
Todos estos ajustes se aplican tanto a DINO como a SimDINO.

\paragraph{Arquitectura del Modelo.} Para todas las entradas, establecemos el tamaño del parche
en \(16 \times 16\) (es decir, \(P_{H} = P_{W} = 16\)). Utilizamos los
modelos pequeño, básico y grande de la arquitectura ViT \citep{dosovitskiy2020image}
como red de incrustación y red troncal. El extractor de características es
una red neuronal multicapas (MLP) de tres capas con un tamaño de capas ocultas de \(2048\) y una dimensión de salida
de \(256\), seguida de una normalización \(\ell^{2}\), tal y como
se especifica en \Cref{sub:contrastive_learning_architecture}. Para las arquitecturas DINO
(es decir, no las arquitecturas SimDINO), el cabezal DINO
\(\vW\) es una matriz en \(\R^{65536 \times 256}\), y el parámetro
\(\vmu\) es un vector en \(\R^{65536}\).

\paragraph{Conjuntos de datos y optimización.} Para el preentrenamiento, tanto nuestra
reproducción de DINO como SimDINO utilizan el conjunto de datos ImageNet-1K en todos los
métodos. Utilizamos AdamW \citep{Loshchilov2017DecoupledWD} como
optimizador, lo cual es una elección muy habitual. Seguimos las siguientes
recomendaciones de hiperparámetros:
\begin{itemize}
\item El tamaño del lote es \(B = 1024\).
\item La tasa de aprendizaje (para AdamW y el modelo alumno) tiene
    un valor «base» de \(2 \times 10^{-3}\). En las primeras \(10\) épocas,
    la tasa de aprendizaje aumenta linealmente de \(0\) al valor base
    (es decir, en la época \(i\) la tasa de aprendizaje es \((i/10) \cdot 2
    \times 10^{-3}\), para \(1 \leq i \leq 10\)). A continuación, durante las siguientes
    \(90\) épocas, la tasa de aprendizaje disminuye siguiendo el denominado
    \textit{programa coseno} hasta volver a \(0\). La definición de un
    programa coseno se encuentra en muchos sitios, incluida
    \href{https://pytorch.org/docs/stable/generated/torch.optim.lr_scheduler.CosineAnnealingLR.html}{la documentación de PyTorch} 
    y se utiliza habitualmente en el entrenamiento de modelos de visión profunda.
\item La disminución de la ponderación (la W de AdamW) sigue un programa coseno
    que va de 0,04 a \(0,4\) a lo largo del entrenamiento.
\item La tasa EMA \(\nu\) sigue un programa coseno desde \(0,996\)
    hasta \(1,0\) a lo largo del entrenamiento. En el caso concreto de DINO, la tasa EMA de centrado
    \(\rho\) se fija en \(0,9\).
\item En el caso concreto de DINO, la temperatura del maestro
    \(\tau_{\teacher}\) se fija en \(0,1\), mientras que la temperatura del alumno
    \(\tau_{\student}\) aumenta linealmente de \(0,04\)
    a \(0,07\) durante las primeras \(30\) épocas y se fija en
    \(0,07\) a partir de entonces.
\end{itemize}
Utilizamos algunas técnicas de aumento de datos (que, en esencia, conservan la información),
como el volteo, la variación de color, el desenfoque gaussiano y la solarización, para
cada imagen observada durante el entrenamiento, antes de obtener las vistas local y global.
Los hiperparámetros exactos que rigen estas técnicas no se enumeran aquí,
pero se mencionan en el artículo sobre DINO \citep{caron2021emerging}.

En el caso del sondeo lineal (linear probing), la sonda lineal suele entrenarse utilizando el
optimizador AdamW con una tasa de aprendizaje de \(2 \times 10^{-4}\), una
decadencia de pesos de \(0,01\) y un tamaño de lote de \(512\), aunque estos
parámetros suelen modificarse según cada caso concreto para minimizar la pérdida.

\begin{table}
\caption{\small\textbf{Rendimiento en clasificación} en datos de prueba de validación
    para DINO y SimDINO, utilizando tanto la precisión de los \(k\) vecinos más cercanos
    (\(k = 20\)) como el sondeo lineal. Con el mismo número de iteraciones
    (\(100\)), SimDINO es claramente mejor en términos de rendimiento y es
    más estable (el entrenamiento de DINO ejecutado en la red base ViT-L con la
        configuración proporcionada tiene una optimización muy inestable y obtiene una pérdida NaN
    en poco tiempo). También comparamos con otros métodos destacados, concretamente
SwAV y MoCov3, en los que se basa DINO.}

\centering
\begin{tabular}{@{}lcccc@{}} %
    \toprule
    Método & Modelo & Épocas & 20-NN & Sondeo lineal
    \\
    \midrule
    DINO & ViT-B & 100 & 72.9 & 76.3 \\
    SimDINO & ViT-B & 100 & \bf 74.9 & \bf 77.3 \\
    DINO & ViT-L & 100 & -- & -- \\
    SimDINO & ViT-L & 100 & \bf 75.6 & \bf 77.4 \\
    \midrule
    \color{gray} SwAV & \color{gray} ViT-S & \color{gray} 800 &
    \color{gray} 66.3 & \color{gray} 73.5 \\
    \color{gray} MoCov3 & \color{gray} ViT-B & \color{gray} 300  &
    \color{gray} -- & \color{gray} 76.7 \\
    \bottomrule
\end{tabular}
\label{tab:dino_imagenet_linear_probing}
\end{table}

\begin{figure}
\centering
\includegraphics[width=\textwidth]{\toplevelprefix/chapters/chapter8/figs/dino_attention_maps.png}
\caption{\small\textbf{Una comparación cualitativa de los mapas de saliencia}
    generados por DINO \textit{(fila del medio)} y por SimDINO
    \textit{(fila inferior)}. Para cada imagen, calculamos y mostramos el
    mapa de saliencia promedio en la última capa \(L\). Los mapas de saliencia son
    similares en todos los modelos, lo que significa que todos los modelos convergen en una
    concepción similar de qué objetos son importantes. Tenga en cuenta que, aunque
    \(X_{\evaluation}\) es una imagen cuadrada, se ha interpolado de nuevo a
una forma rectangular para crear esta visualización.}
\label{fig:dino_attention_maps_saliency}
\end{figure}

\begin{table}
\caption{\small\textbf{Rendimiento en segmentación} de los modelos DINO
    y SimDINO preentrenados en COCO val2017 \citep{lin2014microsoft}, un
    conjunto de datos de segmentación que contiene metadatos sobre la ubicación de los objetos. No
    entrenamos en COCO, sino que simplemente utilizamos la incrustación y la
    estructura preentrenadas, y las cajas delimitadoras se extraen de las características mediante
    un método llamado MaskCut \citep{wang2023cut}. No obstante, SimDINO
    supera a DINO en la detección y segmentación de objetos en una comparación
    imparcial, e incluso supera a DINO con un tamaño de parche más pequeño (longitud lateral
    de \(8\) en lugar de \(16\)). Se sabe que los tamaños de parche más pequeños
    mejoran el rendimiento, especialmente en tareas de detección y segmentación,
por lo que este resultado resulta bastante sorprendente y alentador.}
\centering
\begin{tabular}{@{}llcccccccc@{}}
    \toprule
    &  & \multicolumn{3}{c}{Detection $\uparrow$} &
    \multicolumn{3}{c}{Segmentation $\uparrow$} \\
    Método & Modelo & AP$_{50}$  & AP$_{75}$ & AP & AP$_{50}$ &
    AP$_{75}$ & AP  \\
    \midrule
    SimDINO &ViT-L/16 &\bf 5.4 &1.9 &2.4 &4.5 &1.4 &1.9 \\
    SimDINO &ViT-B/16 &5.2 & \bf 2.0 & \bf 2.5 & \bf4.7 & \bf 1.5 & \bf 2.0 \\
    DINO &ViT-B/16 &3.9 &1.5 &1.8 &3.1 &1.0 &1.4 \\
    \midrule
    \color{gray} DINO & \color{gray} ViT-B/8 & \color{gray}5.1 &
    \color{gray}2.3 & \color{gray}2.5 & \color{gray}4.1 &
    \color{gray}1.3 & \color{gray}1.8 \\
    \bottomrule
\end{tabular}
\label{tab:dino_segmentation}
\end{table}

\paragraph{Resultados de la evaluación.} En cuanto al rendimiento en la clasificación posterior,
obtenemos los resultados que se muestran en
\Cref{tab:dino_imagenet_linear_probing}. Observamos que el
rendimiento de SimDINO es mucho mayor que el de DINO en una comparación
equitativa. Además, es mucho más estable: los ajustes predeterminados de
DINO no permiten entrenar un modelo ViT-L(arge). Por otro lado,
\Cref{fig:dino_attention_maps_saliency} muestra visualizaciones de los
mapas de saliencia promedio en DINO y en nuestro DINO simplificado, lo que permite observar que
los mapas de saliencia son bastante similares entre los distintos modelos, lo que indica que
los modelos aprenden características que son, como mínimo, igual de eficaces a la hora de captar
detalles de gran precisión. Los resultados de segmentación y detección de objetos
en \Cref{tab:dino_segmentation} confirman esta afirmación
cuantitativamente, donde las características de SimDINO muestran una mejora sustancial
respecto a las de DINO.

\section{Representación de imágenes con supervisión débil mediante vinculación de texto}\label{sec:CLIP}
%\yima{Robin, please provide a vignette about CLIP or its simplified version. Sam will give a general description of concepts and methods for joint image/text embedding and related work in Chapter 7. }

Otro enfoque influyente del aprendizaje contrastivo se aleja de las comparaciones 
puramente visuales y, en su lugar, aprovecha la correspondencia natural entre las 
imágenes y el lenguaje. En lugar de definir la similitud únicamente a través de 
múltiples vistas de la misma imagen, esta línea de investigación aprovecha el hecho 
de que muchas imágenes del mundo real van acompañadas de descripciones textuales 
que proporcionan una supervisión semántica débil. Un ejemplo destacado de este 
paradigma es CLIP (Contrastive Language–Image Pretraining) \citep{Radford2021-ir}. 
En CLIP, las imágenes y sus correspondientes leyendas se tratan como pares 
positivos, mientras que los pares de imagen y texto que no coinciden se tratan como 
pares negativos. El objetivo del aprendizaje consiste en que la representación de 
una imagen se acerque a la representación de su leyenda asociada y se aleje de las 
leyendas de otras imágenes. Al basar las representaciones visuales en el lenguaje 
natural, CLIP aprende características semánticamente ricas que captan conceptos de 
alto nivel que van más allá de la mera apariencia.

\subsection{Datos}

Los datos que utilizaremos para analizar la metodología CLIP son {\em pares de texto e imagen}, 
en lugar de solo imágenes. El trabajo original de CLIP construyó un gran corpus a escala web 
de aproximadamente 400 millones de pares de texto e imagen recopilados de fuentes disponibles 
públicamente en Internet. Cada muestra consta de una imagen y una cadena de lenguaje natural 
asociada (por ejemplo, una leyenda, un título o una breve descripción) que coocurre con la 
imagen en la web, lo que proporciona una forma débil y ruidosa de supervisión. Para fomentar 
una amplia cobertura de conceptos visuales, el trabajo de CLIP describe la creación de este 
conjunto de datos mediante la búsqueda de pares cuyo texto coincida con una de las 
aproximadamente 500 000 consultas, y equilibrando los resultados de manera aproximada al 
incluir hasta 20 000 pares por consulta. Este conjunto de datos no se ha publicado, por lo 
que la mayoría de las reproducciones y ampliaciones de CLIP realizadas por la comunidad se 
basan, en su lugar, en conjuntos de datos abiertos a escala web de pares de texto e imagen. 
La más común de ellas es la familia LAION~\citep{schuhmann2022laion5bopenlargescaledataset}: 
LAION-400M ofrece 400 millones de pares filtrados con CLIP, y LAION-5B amplía este procedimiento 
a 5.850 millones de pares.

Dado que los pares de datos sin procesar extraídos de la web suelen presentar altos niveles de ruido 
(por ejemplo, imágenes que no coinciden con el texto asociado, texto alternativo genérico generado automáticamente, 
spam, duplicados o contenido peligroso), los conjuntos de datos al estilo CLIP suelen crearse 
mediante un proceso de posprocesamiento que filtra y limpia los datos candidatos antes del 
entrenamiento. Entre los métodos habituales se incluyen filtros básicos de calidad (por ejemplo, 
descartar pares con texto extremadamente breve o imágenes muy pequeñas o dañadas), la 
identificación del idioma para conservar las leyendas en el idioma de destino, la 
eliminación de duplicados y la eliminación de contenido potencialmente dañino.

\subsection{La tarea y el objetivo del CLIP}
Nuestro objetivo es aprender una buena representación de los datos. Sin embargo, en el marco 
de CLIP, el concepto de «similitud» no se deduce a partir de dos vistas ampliadas de la misma 
imagen, sino a partir de la supervisión semántica débil que proporcionan las descripciones en 
lenguaje natural. En concreto, nos gustaría que la representación de una imagen se pareciera 
a la representación de su leyenda asociada y se diferenciara de las leyendas de otras imágenes. 
Del mismo modo, si dos imágenes tienen descripciones textuales similares, sus características 
visuales aprendidas deberían ser similares, mientras que las imágenes descritas por textos no 
relacionados deberían estar separadas en el espacio latente.

Formalizamos el concepto de «similitud» entre representaciones utilizando la \textit{similitud 
coseno}. Concretamente, sea \(f_{\theta}\) un codificador visual y \(g_{\plainphi}\) un 
codificador de texto (cuya arquitectura se detallará en la siguiente sección), ambos mapeando 
a un espacio latente común \(\vZ \in \mathbb{R}^{d}\). Dada una imagen \(\vX \in \cI\) y una 
cadena de texto \(\vT \in \cT\), formamos incrustaciones normalizadas
\begin{equation}
    \vz^{I}_{\theta}(\vX) \doteq \frac{f_{\theta}(\vX)}{\|f_{\theta}(\vX)\|_2},
    \qquad
    \vz^{T}_{\plainphi}(\vT) \doteq \frac{g_{\plainphi}(\vT)}{\|g_{\plainphi}(\vT)\|_2}.
\end{equation}
Entonces, la similitud coseno entre las representaciones se define como:
\begin{equation}
    s(\vX,\vT)
    =
    \left\langle \vz^{I}_{\theta}(\vX), \vz^{T}_{\plainphi}(\vT)\right\rangle
    \in [-1,1].
\end{equation}
Por lo tanto, nuestro objetivo es maximizar la similitud coseno entre los pares coincidentes y 
minimizarla entre los pares no coincidentes. Para lograrlo, podemos utilizar una pérdida de 
entropía cruzada simétrica simple, tal como se explica en \Cref{sub:text-cond}. Concretamente, 
dado un mini-lote de \(n\) pares de texto-imagen \((\vX_i, \vT_i)_{i=1}^{n}\), podemos definir 
la pérdida como:
\begin{equation}\label{eq:practical_clip_loss}
    \scalebox{0.925}{\(\displaystyle\cL_{\text{CLIP}} = -\frac{1}{n}\left(
    \sum_{i=1}^{n}
    \frac{\exp(s(\vX_i,\vT_i)/\tau)}{\sum_{k=1}^{n}\exp(s(\vX_i,\vT_k)/\tau)}
    +
    \sum_{i=1}^{n}
    \frac{\exp(s(\vT_i,\vX_i)/\tau)}{\sum_{k=1}^{n}\exp(s(\vT_i,\vX_k)/\tau)}
    \right),\)}
\end{equation}
donde \(\tau > 0\) es un parámetro de temperatura que controla la brusquedad de la función softmax.

\subsection{Arquitectura: Vision Tower}\label{sub:clip_vision_tower}
El objetivo de CLIP es aprender representaciones tanto de imágenes como de texto en el mismo espacio 
latente. Por consiguiente, su arquitectura entrena un codificador visual 
\(f_{\theta}: \cI \to \mathbb{R}^{d} \) y un codificador de texto 
\(g_{\theta}: \cT \to \mathbb{R}^{d} \), lo que también se conoce como la «arquitectura de dos torres».

En cuanto a la arquitectura, se pueden utilizar la mayoría de las redes troncales de visión: por ejemplo, 
las redes convolucionales como ResNet~\citep{He2016-lc}, así como el Vision Transformer~\citep{dosovitskiy2020image} 
(descrito en detalle en \Cref{sub:contrastive_learning_architecture}) son opciones válidas para CLIP. En la práctica, 
el Vision Transformer se utiliza con mayor frecuencia debido a su mejor rendimiento y versatilidad. Al igual que en DINO, 
tomamos la característica del token de clase \(\vz_{\cls}\) como representación de cada imagen: \(f_{\theta}(\vX) = \vz_{\cls}\). 
Esto se debe a que el token de clase está diseñado para agregar información global sobre toda la imagen y se 
utiliza habitualmente en tareas de predicción posteriores.

% \paragraph{Feature Extractor.}
% Following the backbone, CLIP requires an additional feature extractor \(f^{\text{ext}}_{\theta}: \mathbb{R}^{t \times d^\prime} \to \mathbb{R}^{1 \times d^\prime} \), where \(t\) is the sequence length of the backbone's output and \(d^\prime\) is its hidden dimension. There are two mainstream designs of CLIP’s \(f^{\text{ext}}_{\theta}\), differing only in how they aggregate the backbone’s sequence output into a single \(d\)-dimensional image representation. The first is a \emph{mean extractor}: assume the backbone produces a sequence of latent features \([\vz_1,\dots,\vz_t] \in \mathbb{R}^{t\times d'}\), we take the pooled feature to be the average of token embeddings,
% \[
% \vz^{ext} =
% f^{\text{ext}}_{\theta}([\vz_1,\dots,\vz_t]) = \frac{1}{t}\sum_{i=1}^{t}\vz_i \in \mathbb{R}^{d'}.
% \]
% The second is a \emph{class-token extractor}, as is standard for Vision Transformers: assume the backbone output includes a designated class token along other latent features \([\vz_{\cls}, \vz_1,\dots,\vz_t] \in \mathbb{R}^{(t+1)\times d'}\), then we have
% \[
% \vz^{ext} = \vz_{\cls}.
% \]

% \paragraph{Task-specific Head.}
% The feature extractor is typically followed by a projection head \(f^{\text{emb}}_{\theta} : \mathbb{R}^{d'} \to \mathbb{R}^{d}\) (often a learned linear map) to project \( \vz^{ext} \) into the \(d\)-dimension latent space shared of by the final image and text representation
% \[
% \vz^{ext} = 
% \vW^{proj}\vz^{pool}
% \qquad
% \vW^{proj} \in \mathbb{R}^{d^\prime \times d}\,
% \]
% which is the final output of \(f^{I}_{\theta}\) and will later be used in the objective.

\subsection{Arquitectura: Text Tower}
La torre textual CLIP \(g_{\plainphi}\) suele adoptar una arquitectura basada en Transformer similar 
al Vision Transformer descrito anteriormente. La principal diferencia radica en la capa de incrustación, 
que se utiliza para convertir la secuencia de texto en una secuencia de representaciones latentes en el 
espacio latente \(\vZ \in \mathbb{R}^{d}\).

\paragraph{Incrustación.} Dada una secuencia de texto \(\vT \in \cT\) (por ejemplo, una
oración o un documento), la representamos como una secuencia de tokens en
\(\R^d\) mediante la aplicación \(g_{\plainphi}^{\text{emb}}\), de la siguiente manera.
\begin{enumerate}
    \item En primer lugar, tokenizamos \(\vT\) en una secuencia de símbolos discretos
        (subpalabras/caracteres) procedentes de un vocabulario \(\cV\) de
        tamaño \(|\cV|\). Concretamente, un tokenizador
        \(t^{\tok} : \cT \to \cV^{N}\) genera una secuencia de tokens de longitud \(N\)
        \footnote{En la práctica, a menudo se amplía esta
        secuencia con tokens especiales (por ejemplo, \(\langle \text{bos}\rangle\) y
        \(\langle \text{eos}\rangle\), que se utilizan para indicar el inicio y el final 
        de la secuencia) y se pueden rellenar secuencias más cortas, pero mantenemos la notación
        genérica y las tratamos como elementos de \(\cV\).} \((\tau_{1},\dots,\tau_{N})\), donde
        \(\tau_i \in \cV\).
    \item Convertimos los tokens en índices enteros mediante un mapa de vocabulario
        \(\cV\to\{1,\dots,|\cV|\}\), lo que da como resultado una secuencia de índices
        \((s_{1},\dots,s_{N})\). Colocamos estos índices en
        una matriz \(\vS \in \{0,1\}^{|\cV|\times N}\), donde la
        \(i\)-ésima columna es el vector «one-hot» \(\ve_{s_i}\).
    \item A continuación, realizamos la siguiente operación sobre \(\vS\) para proyectarla en \(\R^{d \times (N+1)}\):
        \begin{equation}
            \vS \mapsto [\vz_{\cls}^{1}, \vW^{\tok}\vS] + \vE^{\pos}.
        \end{equation}
        Este es el resultado final de \(g_{\plainphi}^{\text{emb}}\). Aquí tenemos tres parámetros entrenables: \(\vW^{\tok}\),
        \(\vz_{\cls}^{1}\) y \(\vE^{\pos}\), cuya función es la siguiente:
        \begin{itemize}
            \item \(\vW^{\tok} \in \R^{d \times |\cV|}\) es una matriz
                cuya \(j\)-ésima columna es el vector de incrustación aprendido
                para el \(j\)-ésimo elemento del vocabulario. Por lo tanto, \(\vW^{\tok}\vS
                \in \R^{d\times N}\) asigna cada token discreto a una
                característica continua en \(\R^{d}\).
            \item \(\vz_{\cls}^{1} \in \R^{d}\) es un token especial
                de \textit{clasificación} (o \textit{resumen}).
                Desde un punto de vista heurístico, su función es agregar información global
                sobre toda la secuencia y se utiliza habitualmente
                para tareas de predicción posteriores.
            \item \(\vE^{\pos} \in \R^{d \times n}\) es una
                \textit{codificación posicional} que distingue
                los tokens en diferentes posiciones. Dado que la auto-atención
                es invariante ante permutaciones en la dimensión de los tokens,
                \(\vE^{\pos}\) inyecta información de orden de modo que la
                función global no sea invariante ante permutaciones de los
                tokens.
        \end{itemize}
\end{enumerate}
Todos los parámetros \(\vz_{\cls}^{1}, \vW^{\tok}, \vE^{\pos}\) están
incluidos en el conjunto de parámetros \(\theta\). El vocabulario \(\cV\) y el tokenizador \(t^{\tok}\) son fijos.

Otras componentes de la torre textual, como la columna vertebral y el cabezal específico para cada 
tarea, siguen opciones de diseño similares a las de la torre visual y se omiten por motivos de 
brevedad. También cabe señalar que trabajos posteriores han demostrado una flexibilidad considerable 
en la elección del codificador de texto de CLIP: las representaciones de texto fijas contenidas en 
grandes modelos de lenguaje, e incluso las características simples de «bolsa de palabras», también 
pueden resultar eficaces. Analizamos estas variantes en la \Cref{sub:follow-up}.

% \paragraph{Backbone.} Now, our text input (a string) has been transformed into a sequence of latent representations
% \([\vz_{\cls}, \vz_{1}, \ldots, \vz_{n}]\) in the latent space \(\R^{d}\). Therefore, its subsequent processing follows the same procedure as in \cref{sub:contrastive_learning_architecture} for image latent representations.

% \paragraph{Feature Extractor.} The text transformer adopts a class-token extractor similar to the one in \cref{sub:clip_vision_tower}.

% \paragraph{Task-specific Head.} The text transformer adopts a linear head similar to the one in \cref{sub:clip_vision_tower}.

\subsection{Estrategia de optimización}
Entrenamos CLIP utilizando un procedimiento estándar de optimización estocástica de extremo a extremo
(por ejemplo, SGD con momentum, AdamW). A continuación presentamos un proceso genérico.

En cada paso temporal \(k\), se realizan las siguientes operaciones:
\begin{itemize}
    \item \textbf{Datos de submuestra en pares.} Se toma una miniserie de \(n\)
    ejemplos en pares
    \[
        \{(\vX_i^{(k)}, \vT_i^{(k)})\}_{i=1}^{n} \subseteq \cX \times \cT,
    \]
    donde \(\vX_i^{(k)}\) es una imagen y \(\vT_i^{(k)}\) es su
    texto asociado.

    \item \textbf{Calcular representaciones latentes.} Para cada par
    \((\vX_i^{(k)}, \vT_i^{(k)})\), calcular incrustaciones normalizadas utilizando los codificadores de visión y texto \(f_{\theta}\) y \(g_{\phi}\):
    \begin{equation}
        z_{\theta}^{I}(\vX_i^{(k)}) \doteq \frac{f_{\theta}(\vX_i^{(k)})}{\|f_{\theta}(\vX_i^{(k)})\|_2},
        \qquad
        z_{\plainphi}^{T}(\vT_i^{(k)}) \doteq \frac{g_{\plainphi}(\vT_i^{(k)})}{\|g_{\plainphi}(\vT_i^{(k)})\|_2}.
    \end{equation}

    \item \textbf{Calcular la pérdida CLIP.} Calcular la pérdida CLIP tal y como se define en \eqref{eq:practical_clip_loss}:
    \begin{equation}
        \scalebox{0.875}{\(\displaystyle\hat{\cL}_{\text{CLIP}}^{(k)}
        \doteq
        -\frac{1}{n}\left(    
        \sum_{i=1}^{n}
        \frac{\exp(s(\vX_i^{(k)},\vT_i^{(k)})/\tau)}{\sum_{k=1}^{n}\exp(s(\vX_i^{(k)},\vT_k^{(k)})/\tau)}
        +
        \sum_{i=1}^{n}
        \frac{\exp(s(\vT_i^{(k)},\vX_i^{(k)})/\tau)}{\sum_{k=1}^{n}\exp(s(\vT_i^{(k)},\vX_k^{(k)})/\tau)}
        \right).\)}
    \end{equation}

    \item \textbf{Calcular los gradientes.} Calcular los gradientes de la
    pérdida estocástica con respecto a ambos conjuntos de parámetros:
    \begin{equation}
        \nabla_{\theta}\hat{\cL}_{\text{CLIP}}^{(k)}(\theta,\plainphi),
        \qquad
        \nabla_{\plainphi}\hat{\cL}_{\text{CLIP}}^{(k)}(\theta,\plainphi).
    \end{equation}

    \item \textbf{Actualizar parámetros.} Aplicar un paso de un
    algoritmo de optimización a \((\theta,\plainphi)\), lo que da como resultado la
    iteración
    \begin{equation}
        (\theta^{(k+1)}, \plainphi^{(k+1)})
        \doteq
        \textsc{OptUpdate}^{(k)}\!\left(
        (\theta^{(k)}, \plainphi^{(k)});
        \nabla_{\theta}\hat{\cL}_{\text{CLIP}}^{(k)},
        \nabla_{\plainphi}\hat{\cL}_{\text{CLIP}}^{(k)}
        \right).
    \end{equation}
\end{itemize}


\subsection{Metodología de evaluación}\label{sub:clip_eval}
Dado que CLIP es independiente de la arquitectura, se pueden aplicar la mayoría de los 
protocolos de evaluación estándar para codificadores de visión y lenguaje. Por ejemplo, 
se puede realizar un sondeo lineal (linear probing) entrenando un módulo de clasificación ligero sobre 
el codificador visual para evaluar la calidad de las representaciones aprendidas en 
tareas de clasificación de imágenes. Además, se puede llevar a cabo una clasificación 
sin entrenamiento previo («zero-shot») y sin necesidad de entrenamiento adicional gracias a la arquitectura de doble 
torre de CLIP.

Recordemos que, en una configuración estándar de clasificación de imágenes de \(n\) clases, se nos proporciona un conjunto de datos de imágenes
\(\{\vX_i\}_{i=1}^{N}\subseteq\cI\), cada una etiquetada con una de las \(n\) etiquetas de clase pertenecientes al conjunto
\(\cL \doteq \{\ell_1,\ldots,\ell_n\}\) (por ejemplo, \(\ell_j \in \{\text{gato},\text{perro},\ldots\}\)).
Nuestro objetivo es construir un proceso de predicción \(P\), aprovechando un modelo preentrenado, de tal manera que
\(P:\cI\to\{1,\ldots,n\}\) asigne a cada imagen \(\vX_i\) un índice de clase predicho
\(\hat{y}_i \doteq P(\vX_i)\). Escribiendo la etiqueta de clase de referencia como su índice \(y_i \in \{1,\ldots,n\}\),
la precisión de clasificación (top-1) de \(P\) en este conjunto de datos se calcula entonces como
\[
\frac{1}{N}\sum_{i=1}^{N}\boldsymbol{1}\{\hat{y}_i = y_i\}.
\]

En CLIP, podemos utilizar el codificador de texto \(g_{\plainphi}\) para construir un diccionario 
de texto \(\vD\in\R^{n\times d}\), cuya \(j\)-ésima fila es la representación latente normalizada 
de la \(j\)-ésima etiqueta de clase.
Concretamente, para cada clase \(\ell_j \in \cL\), se construye una descripción textual 
(prompt) \(\vT_j\) (p. ej., «una foto de un \(\ell_j\)») y se calcula su incrustación normalizada
\(\vz_j^T \doteq \frac{g_{\plainphi}(\vT_j)}{\|g_{\plainphi}(\vT_j)\|_2} \in \R^{d}\). Al apilarlas por filas se obtiene
\[
\vD \doteq 
\mat{(\vz_1^T)^{\top}\\ \vdots \\ (\vz_n^T)^{\top}}
\in \R^{n\times d},
\]

Dada una imagen \(\vX_i \in \cI\), se calcula su representación de imagen mediante el codificador visual y se la normaliza como
\[
\vz_i^I \doteq \frac{f_{\theta}(\vX_i)}{\|f_{\theta}(\vX_i)\|_2} \in \R^{d}.
\]
We can then compute a similarity vector
\[
\vd_i \doteq \vD\,\vz_i^I \in \R^{n},
\]
cuya \(j\)-ésima entrada \(\vd_i[j]\) es la similitud coseno entre la \(i\)-ésima imagen y la \(j\)-ésima etiqueta de clase
en el espacio latente compartido. Según el objetivo de entrenamiento de CLIP, se fomenta que las imágenes sean más similares a sus
textos correspondientes que a los que no coinciden; por lo tanto, una mayor similitud coseno indica una alineación más estrecha entre
\(\vX_i\) y la clase \(\ell_j\). Por lo tanto, predecimos el índice de clase mediante
\[
\hat{y}_i \doteq \arg\max_{j \in \{1,\ldots,n\}} \vd_i[j].
\]
Por último, podemos calcular la precisión de la clasificación tal y como se ha mencionado anteriormente.

Este protocolo de decisión «zero-shot» no se limita a la clasificación, sino que también se aplica en forma natural a las tareas de \emph{recuperación},
en las que el objetivo es devolver los elementos más relevantes de una gran base de datos a partir de una consulta (por ejemplo, recuperar las imágenes más
relevantes para una consulta de texto, o los textos más relevantes para una consulta de imagen). Por ejemplo, en la recuperación de texto a imagen,
en lugar de crear un diccionario a partir de indicaciones de nombres de clases, creamos un \emph{diccionario de imágenes} codificando
y normalizando todas las imágenes de la base de datos \(\{\vX_m\}_{m=1}^{M}\), y apilando sus incrustaciones en
\(\vD\in\R^{M\times d}\). Dado un texto de consulta \(\vT\), calculamos su incrustación normalizada
\(\vz^T \doteq g_{\plainphi}(\vT)/\|g_{\plainphi}(\vT)\|_2\), generamos puntuaciones de similitud
\(\vd \doteq \vD\vz^T\in\R^{M}\), y luego clasificamos las imágenes según estas puntuaciones (de manera equivalente, devolvemos los \(K\)
índices principales).
% More broadly, this simple embedding-and-dot-product pipeline is arguably the most distinctive and
% practically important feature of CLIP.



\subsection{Configuración experimental y resultados}

A continuación, presentamos los resultados del artículo original sobre CLIP, en el que se compara 
CLIP con una red ResNet supervisada de referencia en varios bancos de pruebas clásicos de 
clasificación. En todos los benchmarks, CLIP se evalúa utilizando el protocolo de clasificación 
zero-shot descrito en \Cref{sub:clip_eval}. Los detalles completos de la configuración experimental 
de ResNet se proporcionan en~\citep{He2016-lc}; aquí, presentamos brevemente los conjuntos de datos 
de evaluación y resumimos los ajustes de entrenamiento y evaluación de CLIP relevantes para estas comparaciones.

\paragraph{Datos para la evaluación.}
\begin{itemize}
    \item \textbf{CIFAR-10~\citep{krizhevsky2009learning}.}
    CIFAR-10 es un conjunto de datos de referencia para la clasificación de imágenes a pequeña escala muy utilizado, que consta de 60 000 imágenes RGB de tamaño \(32\times 32\) distribuidas en 10 clases (6 000 imágenes por clase).
    Normalmente se divide en 50 000 imágenes de entrenamiento y 10 000 imágenes de prueba.

    \item \textbf{CIFAR-100~\citep{krizhevsky2009learning}.}
    CIFAR-100 tiene el mismo tamaño total y la misma resolución de imagen que CIFAR-10 (60 000 imágenes RGB de \(32\times 32\) con una división de 50 000/10 000 para entrenamiento/prueba), pero es más desafiante porque contiene 100 clases detalladas (600 imágenes por clase).
    Las 100 clases se agrupan además en 20 superclases.

    \item \textbf{PASCAL VOC 2007 (VOC2007)~\citep{everingham2010pascal}.}
    PASCAL VOC 2007 es un conjunto clásico de benchmarks en visión artificial para el reconocimiento de objetos, que admite tareas como la clasificación de imágenes, la detección de objetos y la segmentación.
    Se centra en 20 categorías de objetos (por ejemplo, persona, automóvil, perro).

    \item \textbf{ImageNet-1K~\citep{imagenet_cvpr09}.}
    ImageNet-1K es un subconjunto de ILSVRC 2012 que contiene 1{,}000 clases de objetos.
    Este subconjunto incluye 1 281 167 imágenes de entrenamiento, 50 000 imágenes de validación 
    y 100 000 imágenes de prueba, y ha desempeñado un papel central en el reconocimiento visual 
    a gran escala y el preentrenamiento.
\end{itemize}

\paragraph{Configuración del modelo.}
En los modelos CLIP que aquí se presentan, se utiliza Vision Transformer como codificador visual 
\(f_{\theta}\) y se implementa como una de las tres variantes de ViT: ViT-B/32, ViT-B/16 o ViT-L/14. 
Aquí, \(\text{B}\) («Base») y \(\text{L}\) («Large») especifican el ancho y la profundidad del 
transformer (por ejemplo, ViT-B suele utilizar 12 bloques de transformer con una dimensión de 
incrustación \(d=768\) y 12 cabezales de atención, mientras que ViT-L utiliza 24 bloques con 
\(d=1024\) y 16 cabezales), mientras que los sufijos \(/32\), \(/16\) y \(/14\) denotan el 
\emph{tamaño del parche} en píxeles utilizado para particionar la imagen de entrada en parches. Los tamaños 
de parche más pequeños (por ejemplo, 16 o 14) generan más tokens por imagen a una resolución fija, 
lo que aumenta la carga computacional, pero a menudo mejora la precisión gracias a una granularidad 
espacial más fina.

\paragraph{Configuración de la optimización.}
Todos los modelos CLIP basados en ViT se entrenan durante 32 épocas utilizando el optimizador Adam 
con decaimiento de pesos desacoplado (es decir, regularización al estilo AdamW), en el que el 
decaimiento de pesos se aplica a todos los parámetros excepto a las ganancias y los sesgos. La tasa 
de aprendizaje se reduce siguiendo un programa de coseno. El parámetro de temperatura aprendible 
\(\tau\) (utilizado para escalar los logits contrastivos) se inicializa en el valor equivalente a 
\(0,07\) y se limita para evitar que los logits se escalen más de 100, lo cual se considera 
necesario para la estabilidad del entrenamiento. Por último, el entrenamiento utiliza un tamaño 
de mini-lote muy grande, de \(32 768\).

\paragraph{Resultados.} La \Cref{tab:clip_vs_resnet} presenta la precisión top-1 en aprendizaje 
sin entrenamiento previo (zero-shot) de tres variantes de CLIP ViT (B/32, B/16, L/14) y las compara con las 
redes de referencia supervisadas ResNet-50 y ResNet-101.
En los cuatro conjuntos de datos, CLIP supera sistemáticamente a las redes ResNet supervisadas.
Ampliar aún más la estructura central de CLIP genera mejoras adicionales, con ViT-L/14 alcanzando 
\(98.0\) en CIFAR-10, \(87.5\) en CIFAR-100, \(89.6\) en VOC2007 y \(83.9\) en ImageNet.
No obstante, esta tabla debe interpretarse como una comparación entre dos \emph{sistemas} y no 
como una prueba controlada de eliminación de la red troncal: CLIP se beneficia del
preentrenamiento a gran escala de pares de imagen-texto y se evalúa mediante prompts de tipo «zero-shot», 
mientras que los modelos de referencia de ResNet se basan en un preentrenamiento supervisado y 
requieren datos etiquetados para ajustar una prueba en cada banco de pruebas.
Por lo tanto, los márgenes consistentes demuestran principalmente la solidez de las 
representaciones transferibles de CLIP en el marco del protocolo «zero-shot», en lugar 
de establecer una superioridad de ViT sobre ResNet basada exclusivamente en la 
arquitectura o en la correspondencia de datos.

\begin{table}[t]
\centering
\begin{tabular}{@{}llcccc@{}}
\toprule
Red troncal & Variante & CIFAR-10 & CIFAR-100 & VOC2007 & ImageNet \\
\midrule
CLIP-ViT & B/32 & 95.1 & 80.5 & 87.7 & 76.1 \\
CLIP-ViT & B/16 & 96.2 & 83.1 & 89.2 & 80.2 \\
CLIP-ViT & L/14 & 98.0 & 87.5 & 89.6 & 83.9 \\
\midrule
ResNet & 50  & 91.8 & 74.5 & 83.8 & 74.3 \\
ResNet & 101 & 93.0 & 77.2 & 84.4 & 75.8 \\
\bottomrule
\end{tabular}
\caption{Rendimiento en conjuntos de datos de tareas posteriores para diversas estructuras CLIP y modelos de referencia ResNet.}
\label{tab:clip_vs_resnet}
\end{table}

\subsection{Una extensión simplificada: LIFT}\label{sub:follow-up}

Aunque es eficaz, CLIP presenta varias limitaciones. En primer lugar, entrenar conjuntamente 
desde cero tanto un codificador de texto como un codificador visual supone un gran esfuerzo 
computacional, ya que a menudo requiere lotes extremadamente grandes y conjuntos de datos masivos 
para lograr una alineación sólida y una transferencia eficaz en las etapas posteriores. En segundo 
lugar, los modelos entrenados de esta manera pueden tener dificultades para comprender la 
composición---es decir, captar el orden de las palabras en el texto, la disposición espacial 
en las imágenes y las relaciones entre objetos y atributos o entre objetos. Esto se debe, en 
parte, a que la supervisión contrastiva de tipo «recuperación» puede favorecer soluciones 
simplistas que restan importancia a las características composicionales más detalladas. LIFT 
(Language–Image alignment with a Fixed Text encoder)~\citep{yang2025languageimagealignmentfixedtext} 
revisa una suposición fundamental en la que se basan estos flujos de trabajo: que una alineación 
óptima requiere el entrenamiento conjunto de extremo a extremo de ambos codificadores:\\
\centerline{\url{https://jingfeng0705.github.io/LIFT/}.}\\
En cambio, LIFT aprovecha el hecho de que los modelos de lenguaje a gran escala (LLM) modernos 
ya generan incrustaciones de texto muy informativas. Concretamente, LIFT establece un codificador 
de texto preentrenado de gran calidad (por ejemplo, uno derivado de un LLM o ajustado sobre él), 
calcula las incrustaciones de texto sin conexión y entrena únicamente el codificador visual 
para que se alinee con estos objetivos fijos.

\paragraph{Arquitectura.}
LIFT conserva la misma estructura de doble torre de CLIP en cuanto a un codificador visual y un 
codificador de texto que se proyectan en un espacio latente compartido, pero difiere 
fundamentalmente en que la torre textual es \emph{fija} y puede ser implementada por cualquier 
codificador de texto potente basado en un LLM \(g_{\plainphi}:\ cT\to\R^{d}\) para producir 
incrustaciones de texto, donde \(\plainphi\) no se actualiza durante el entrenamiento. En nuestra 
implementación, adoptamos el codificador de texto NV-Embed-V2, para el cual la dimensión de 
incrustación es \(d=4096\). Dado que la torre textual \(g_{\theta}\) es fija, podemos precalcular 
todas las incrustaciones de texto \(\{g_{\theta}(\vT)\}\) en forma offline y, durante el 
entrenamiento, solo calcular los gradientes y actualizar el codificador visual 
\(f_{\theta}\) en forma online. La torre visual \(f_{\theta}:\cI\to\R^{d}\) sigue la misma 
estructura que en \Cref{sub:clip_vision_tower}, con la dimensión de salida del cabezal de 
proyección establecida en \(d\) para que coincida con la dimensión del espacio de incrustación 
de texto fijo. Aparte de fijar \(g_{\plainphi}\) y ajustar la dimensión de proyección, las 
componentes restantes---incluida la pérdida contrastiva y el procedimiento de optimización---siguen 
los mismos principios de diseño que CLIP.

\paragraph{Metodología de evaluación.} 
La comprensión composicional es una limitación conocida de CLIP. Evaluamos LIFT y CLIP en 
siete tareas de SugarCrepe~\citep{hsieh2023sugarcrepe}, en las que cada imagen \(\vX\) se 
empareja con una etiqueta \emph {positiva} (correcta) \(\vT_{\text{pos}}\) y una etiqueta \emph{negativa} 
\(\vT_{\text{neg}}\) construida añadiendo, sustituyendo o intercambiando un objeto, atributo 
o relación en \(\vT_{\text{pos}}\). Véase la Figura \ref{fig:SugarCrepe-examples} para algunos ejemplos.
\[
\vz_i^{I} \doteq \frac{f_{\theta}(\vX_i)}{\|f_{\theta}(\vX_i)\|_2},\qquad
\vz_{i,\text{pos}}^{T} \doteq \frac{g_{\plainphi}(\vT_{i,\text{pos}})}{\|g_{\plainphi}(\vT_{i,\text{pos}})\|_2},\qquad
\vz_{i,\text{neg}}^{T} \doteq \frac{g_{\plainphi}(\vT_{i,\text{neg}})}{\|g_{\plainphi}(\vT_{i,\text{neg}})\|_2}.
\]
A continuación, se calculan las puntuaciones de similitud mediante similitudes coseno:
\[
s_{i,\text{pos}} \doteq \langle \vz_i^{I}, \vz_{i,\text{pos}}^{T} \rangle,
\qquad
s_{i,\text{neg}} \doteq \langle \vz_i^{I}, \vz_{i,\text{neg}}^{T} \rangle.
\]
Se considera que el modelo es correcto en la muestra \(i\) si
\(s_{i,\text{pos}} > s_{i,\text{neg}}\) y la precisión compositiva general es
\[
\frac{1}{N}\sum_{i=1}^{N}\boldsymbol{1}\{s_{i,\text{pos}} > s_{i,\text{neg}}\}.
\]

\begin{figure}
\centering
\includegraphics[width=\textwidth]{\toplevelprefix/chapters/chapter8/figs/negative_caption_demo.pdf}
\caption{Las leyendas originales (arriba) y sus correspondientes negativos (abajo) de dos tareas de SugarCrepe: \texttt{reemplazar relación} (izquierda) y \texttt{intercambiar atributo} (derecha).}
\label{fig:SugarCrepe-examples}
\end{figure}

\paragraph{Resultados experimentales.}
Como se muestra en la \Cref{tab:Compositionality}, cuando se entrena con DataComp-1B, LIFT supera 
a CLIP en las siete tareas con una mejora media de la precisión del \textbf{6,8\%}. LIFT logra 
mejoras significativas en las tareas de \texttt{añadir atributo}, \texttt{reemplazar atributo} 
y \texttt{reemplazar relación}. Estas mejoras constituyen una clara evidencia de que los 
codificadores de texto basados en LLM \(g_{\plainphi}\) capturan representaciones de texto más 
informativas que permiten modelar con mayor precisión las asociaciones entre objetos y atributos, 
así como las relaciones entre objetos. Por otro lado, podemos observar que LIFT muestra una precisión 
relativamente baja en las tareas \texttt{intercambiar objeto} y \texttt{intercambiar atributo} en comparación 
con otras tareas de SugarCrepe. Atribuimos esta limitación al objetivo del aprendizaje contrastivo, 
que se centra principalmente en alinear estadísticas de orden inferior. Para hacer frente a este reto, 
es necesario explorar medidas más refinadas de la teoría de la información para la alineación entre 
lenguaje e imagen, lo que constituye una línea de trabajo clave para el futuro.

\begin{table}[t]
\centering
\resizebox{\textwidth}{!}{%
\begin{tabular}{@{}lllccccccc@{}}
\toprule
& & & \multicolumn{2}{c}{Añadir} & \multicolumn{3}{c}{Reemplazar} & \multicolumn{2}{c}{Swap} \\
Método & C.Datos & Muestra vista & Obj & Atr & Obj & Atr & Rel & Obj & Atr \\
\midrule
OpenCLIP & DataComp & 1.28B & 82.3 & 73.7 & 91.7 & 79.4 & 61.2 & 59.6 & 56.9 \\
LIFT     & DataComp & 1.28B & \textbf{89.0} & \textbf{86.1} & \textbf{93.2} & \textbf{86.0} & \textbf{70.6} & \textbf{64.1} & \textbf{63.4} \\
\bottomrule
\end{tabular}%
}
\caption{El rendimiento de LIFT y CLIP en siete tareas de SugarCrepe. Utilizamos una 
implementación de código abierto de CLIP, OpenCLIP~\citep{ilharco_gabriel_2021_5143773}, 
y entrenamos tanto OpenCLIP como LIFT en el conjunto de datos de código abierto de pares 
de texto e imagen DataComp~\citep{gadre2023datacompsearchgenerationmultimodal}. 
Presentamos la precisión para cada tarea, con los mejores resultados en negrita.}
\label{tab:Compositionality}
\end{table}


\section{Representación de imágenes supervisada mediante clasificación}\label{sec:image_classification}

En las secciones anteriores, hemos mostrado cómo simplificar algunos objetivos de aprendizaje excesivamente
complejos —y a menudo diseñados de manera heurística— basándonos en
los principios de compresión presentados en este libro. Sin embargo,
los objetivos asociados a algunas de las tareas de aprendizaje más populares
ya son bastante sencillos. En estos casos, es difícil
simplificar aún más el objetivo. Por lo tanto, en esta y en futuras secciones, nos
centraremos en formas fundamentadas de modificar las \textit{arquitecturas de redes profundas} 
para una variedad de tareas.

Comencemos primero con la que podría considerarse la tarea más clásica del aprendizaje automático:
\textit{la clasificación de imágenes}, que a menudo se utiliza como
tarea estándar para evaluar algoritmos de reconocimiento de patrones o arquitecturas de redes profundas.
A partir de nuestro análisis de las arquitecturas white-box (transparentes)
en \Cref{ch:deep-networks}, solo necesitamos una tarea semánticamente significativa
para aprender buenas representaciones con arquitecturas white-box.
Validaremos esta idea en esta sección.

En primer lugar, el conjunto de datos se mantiene prácticamente igual que en la
\Cref{sub:contrastive_learning_data}. Tanto los datos de entrenamiento como los de prueba
consisten en imágenes etiquetadas, es decir, pares de imagen-etiqueta \((\vX, \vy) \in
\R^{C \times H \times W} \times \{0, 1\}^{N_{\cls}}\). Seguimos aplicando
varias técnicas de aumento de datos (por ejemplo, volteos, desenfoque gaussiano,
solarización, etc.) a cada muestra de cada nuevo lote.

\subsection{Tarea y objetivo} \label{sub:image_classification_objective}

A diferencia de antes, nuestra tarea no consiste solo en aprender una buena representación de
los datos, sino también en construir simultáneamente un clasificador. Formalmente,
disponemos de pares de datos etiquetados \((\vX, \vy)\), donde \(\vy \in \{0,
1\}^{N_{\cls}}\) es un vector «one-hot» que indica la clase a la que pertenece
\(\vX\). Consideramos una vista determinista de recorte central
\(v_{\cc}\) de los datos de entrada \(\vX\) (véase
\Cref{sub:contrastive_learning_objective}). Queremos entrenar conjuntamente
una función de mapeo \((f_{\theta}, f_{\theta}^{\ext})\) y un
\textit{cabezal de clasificación} \(h_{\theta}\), definidos de la siguiente manera:
\begin{equation}
h_{\theta}(\vz) \doteq \softmax(\vW^{\head}\vz + \vb^{\head}), \qquad
\forall \vz \in \R^{d}
\end{equation}
donde \((\vW^{\head}, \vb^{\head}) \in \R^{N_{\cls} \times d} \times
\R^{N_{\cls}}\) son parámetros entrenables en el conjunto de parámetros
\(\theta\), de tal manera que la aplicación \(\vX_{\cc} \mapsto
\vp_{\theta}(\vX_{\cc}) \doteq h_{\theta}(\vz_{\theta}(\vX_{\cc}))\)
predice una etiqueta suavizada para la vista \(\vX_{\cc} = v_{\cc}(\vX)\)
de la entrada \(\vX\). El problema de aprendizaje busca minimizar la
distancia entre \(\vp_{\theta}\) y \(\vy\) medida mediante entropía cruzada:
\begin{equation}\label{eq:classification_ce_loss}
\min_{\theta}\bc{\cL_{\CE}(\theta) \doteq \Ex[\CE(\vy,
\vp_{\theta}(\vX_{\cc}))]}.
\end{equation}

\subsection{La arquitectura CRATE}\label{sub:image_classification_architecture}

La arquitectura que utilizamos es la arquitectura CRATE, descrita con
cierto detalle en el \Cref{ch:deep-networks}. La configuración general es análoga
a la del transformer estándar de la 
\Cref{sub:contrastive_learning_architecture}, con algunos cambios.
Si bien el paso de incrustación es el mismo que en DINO y SimDINO en la
\Cref{sub:contrastive_learning_architecture}, el paso de extracción de características
es el mismo que en SimDINO en la
\Cref{sub:contrastive_learning_architecture}, ya que solo extrae la
característica asociada al token de clase, y el cabezal de clasificación
se describe en la \Cref{sub:image_classification_objective}, la
arquitectura de la red troncal es diferente. Cada capa toma la forma
\begin{align}\label{eq:CARTE updates}
\vZ_{\theta}^{\ell + 1/2}(\vX)
&= \vZ_{\theta}^{\ell}(\vX) + \MSSA_{\theta}^{\ell}(\LN_{\theta}^{1,
\ell}(\vZ_{\theta}^{\ell}(\vX))), \\
\vZ_{\theta}^{\ell + 1}(\vX)
&= \ISTA_{\theta}^{\ell}(\LN_{\theta}^{2, \ell}(\vZ_{\theta}^{\ell +
1/2}(\vX))),
\end{align}
donde los bloques \(\MSSA_{\theta}^{\ell}\) e \(\ISTA_{\theta}^{\ell}\)
son los descritos en el \Cref{ch:deep-networks}, a saber:
\begin{itemize}
\item El operador \(\MSSA\) es una auto-atención de subespacio de múltiples cabezas,
    definida de la siguiente manera:
    \begin{equation}
        \MSSA_{\theta}^{\ell}(\vZ) \doteq
        \vU_{\out}^{\ell}\mat{\SA([\vU^{1, \ell}]^{\top}\vZ, [\vU^{1,
            \ell}]^{\top}\vZ, [\vU^{1, \ell}]^{\top}\vZ)\\ \vdots
            \\ \SA([\vU^{K, \ell}]^{\top}\vZ, [\vU^{K, \ell}]^{\top}\vZ,
        [\vU^{1, \ell}]^{\top}\vZ)} + \vb_{\out}^{\ell}\vone_{n}^{\top}
    \end{equation}
    donde \(\vU^{k, \ell} \in \R^{d \times p}\), \(\vU_{\out}^{\ell}
    \in \R^{d \times Kp}\), y \(\vb_{\out}^{\ell} \in \R^{d}\) son
    parámetros entrenables que pertenecen al conjunto de parámetros \(\theta\),
    y (recordemos) el operador de auto-atención \(\SA\) se define en
    \eqref{eq:self_attention}.
\item El operador \(\ISTA\) es el
    operador del algoritmo de umbralización por contracción iterativa, definido de la siguiente manera:
    \begin{equation}
        \ISTA_{\theta}^{\ell}(\vZ) \doteq \ReLU(\vZ - \beta
            (\vD^{\ell})^{\top}(\vD^{\ell}\vZ - \vZ) + \beta\lambda
        \vone_{d}\vone_{n}^{\top}),
    \end{equation}
    llamado así porque la función \(\vX \mapsto \ReLU(\vX - \beta
            \vD^{\top}(\vD\vX - \vZ) + \beta  \lambda
    \vone_{d}\vone_{n}^{\top})\) constituye un paso del algoritmo
    para encontrar una representación dispersa no negativa elemento por elemento
    para \(\vZ\) con respecto al diccionario completo
    \(\vD\) (véase \Cref{sec:dictionary_learning}).
\end{itemize}

A esta arquitectura la llamamos CRATE y en la \Cref{fig:crate_backbone} se muestra
una capa de su estructura central. Los modelos CRATE, más allá de su interpretabilidad,
se caracterizan por su alto rendimiento y su eficiencia paramétrica.

\subsection{Optimización} \label{sub:image_classification_optimization}

Nuestro clasificador se entrena mediante un sencillo procedimiento de optimización estocástica de extremo a extremo,
en el que se toma una submuestra de los datos y las vistas, se calcula
la pérdida promedio y su gradiente sobre estas muestras y se utiliza un
algoritmo de optimización para actualizar los parámetros. En cada paso temporal \(k\), se hace lo siguiente:
\begin{itemize}
\item Se seleccionan \(B\) muestras etiquetadas diferentes: \(\{(\vX_{b}^{(k)},
    \vy_{b}^{(k)})\}_{b = 1}^{B} \subseteq \cI \times \{0, 1\}^{N_{\cls}}\).
\item Para cada muestra etiquetada \((\vX_{b}^{(k)}, \vy_{b}^{(k)})\),
    se calcula la vista central recortada \(v_{b, \cc}^{(k)}\) y se la aplíca
    a \(\vX_{b}^{(k)}\) para obtener \(\vX_{b, \cc}^{(k)} \doteq v_{b,
    \cc}^{(k)}(\vX_{b}^{(k)})\).
\item Se calculan las predicciones \(\vp_{\theta}(\vX_{b, \cc}^{(k)})
        \doteq (h_{\theta} \circ f_{\theta}^{\ext} \circ
    f_{\theta})(\vX_{b, \cc}^{(k)})\).
\item Se calcula la pérdida estocástica sustitutiva
    \begin{equation}
        \hat{\cL}_{\CE}^{(k)}(\theta) \doteq \frac{1}{B}\sum_{b =
        1}^{B}\CE(\vy_{b}^{(k)}, \vp_{\theta}(\vX_{b, \cc}^{(k)})).
    \end{equation}
\item Se calcula un paso de un algoritmo de optimización sobre \(\theta\),
    lo que da lugar a la siguiente iteración:
    \begin{equation}
        \theta^{(k + 1)} \doteq
        \textsc{OptUpdate}^{(k)}(\theta^{(k)};
        \nabla_{\theta}\hat{\cL}_{\CE}^{(k)}).
    \end{equation}
\end{itemize}

\subsection{Metodología de evaluación} \label{sub:image_classification_evals}

Se utiliza el mismo procedimiento de evaluación que en la
\Cref{sub:contrastive_learning_evals}. En resumen, para todas las
evaluaciones (así como para el entrenamiento) se utiliza una vista de recorte central
\(v_{\cc}\) que remodela la imagen de entrada y se extrae un recorte central 
de tamaño \((C, S_{\cc}, S_{\cc})\), donde \(C\) es el número de
canales en la imagen de entrada. Con esta vista se pueden realizar sondeos lineales, visualización de mapas de atención
y benchmarks de detección y segmentación, dada la
salida de esta vista.

\subsection{Configuración experimental y
resultados}\label{sub:image_classification_experiments}

Dado que CRATE se basa directamente en el modelo Transformer, se comparan los
parámetros óptimos para ViT, tal y como se indican en
\cite{dosovitskiy2020image,touvron2020training}, con los mismos
parámetros aplicados a CRATE, a fin de realizar una comparación imparcial.

\paragraph{Arquitectura del modelo.} El recorte central redimensiona toda la
imagen de modo que el lado más corto tenga un tamaño de \(256\) (es decir, \(S_{\rsz} =
256\)) antes de realizar un recorte central de tamaño \(224 \times 224\) (es decir,
\(S_{\cc} = 224\)), tanto en la evaluación como en el entrenamiento. Se toma un tamaño de parche
de \(16\) (es decir, \(P_{H} = P_{W} = 16\)). Se utilizan los modelos tiny, small,
base y large de la arquitectura ViT \cite{dosovitskiy2020image}
como red de incrustación y red troncal, sustituyendo las componentes MHSA y
MLP por MSSA e ISTA, respectivamente, utilizando el mismo número
de cabezas y la misma dimensión de cabeza en el caso de MSSA, lo que
reduce drásticamente el número de parámetros de entrenamiento. Para CRATE,
se establece \((\beta, \lambda) = (1, 0.1)\).

\paragraph{Conjuntos de datos y optimización.} Para el preentrenamiento, se utiliza el
conjunto de datos ImageNet-1K. Se emplea el optimizador LION
\citep{chen2024symbolic} para preentrenar tanto nuestra réplica de ViT como
CRATE. Se establece la tasa de aprendizaje base en \(2.4 \times
10^{-4}\), la decaimiento de pesos en \(0.5\) y el tamaño del lote en \(B =
2048\). Nuestro programa de tasas de aprendizaje aumenta la tasa de aprendizaje
en forma lineal hasta la tasa de aprendizaje base durante las primeras \(5\) épocas, y
disminuye hasta \(0\) siguiendo un programa coseno durante las siguientes \(145\)
épocas (entrenando todos los modelos durante \(150\) épocas cada uno). Para el
preentrenamiento, se aplica un régimen habitual de aumentos de datos (volteretas,
desenfoques gaussianos, solarización, etc.) a los datos de imagen y también se añade
un pequeño ruido a las etiquetas (esto se denomina \textit{suavizado de etiquetas}
\citep{muller2019does}).

Para el aprendizaje por sonda lineal, se utilizan varios conjuntos de datos de evaluación, tales como
CIFAR10, Oxford-Flowers y Oxford-IIT-Pets. Se emplea el optimizador AdamW
para entrenar la sonda lineal, con una tasa de aprendizaje de \(5 \times
10^{-5}\), una decaimiento de pesos de \(0.01\) y un tamaño de lote de \(B = 256\). También
se aplican las aumentaciones de datos mencionadas anteriormente a los datos de imagen.

\begin{table}
\caption{\small \textbf{Precisión de clasificación mediante sondaje lineal de
    CRATE y ViT} en diversos conjuntos de datos con diferentes tamaños de modelo cuando
    la red base se ha entrenado previamente para la clasificación en ImageNet-1K. Se observa
    que, con la misma configuración de modelo, CRATE ofrece un rendimiento de clasificación comparable
    con un diseño más sencillo, más fundamentado y más
eficiente en cuanto a parámetros.}
\centering
\resizebox{1\columnwidth}{!}
{
    \begin{tabular}{@{}lcccc|cc@{}}
        \toprule
        \textbf{Modelo} & CRATE-T  &  CRATE-S & CRATE-B & CRATE-L & {
        \color{gray} ViT-T} &  { \color{gray}ViT-S } \\
        \midrule
        \midrule
        \# parameters & 6.09M & 13.12M & 22.80M & 77.64M & {
        \color{gray} 5.72M} & { \color{gray} 22.05M} \\
        \midrule
        ImageNet-1K & 66.7 & 69.2 & 70.8 & 71.3 & { \color{gray}
        71.5} & { \color{gray} 72.4} \\
        ImageNet-1K ReaL & 74.0 & 76.0 & 76.5 & 77.4 & { \color{gray}
        78.3 } & { \color{gray} 78.4} \\
        CIFAR10 & 95.5 & 96.0 & 96.8 & 97.2 & { \color{gray} 96.6} &
        { \color{gray} 97.2} \\
        CIFAR100 & 78.9 & 81.0 & 82.7 & 83.6 & { \color{gray} 81.8} &
        { \color{gray} 83.2}\\
        Oxford Flowers-102 & 84.6 & 87.1 & 88.7 & 88.3 & {
        \color{gray} 85.1} & { \color{gray} 88.5}\\
        Oxford-IIIT-Pets & 81.4 & 84.9 & 85.3 & 87.4 & { \color{gray}
        88.5} & { \color{gray} 88.6} \\
        \bottomrule
    \end{tabular}
}
\label{tab:crate_classification_linear_probing}
\end{table}

\begin{figure}[t]
\centering
\includegraphics[width=\textwidth]{\toplevelprefix/chapters/chapter8/figs/crate_semantic_heads.pdf}
\caption{\small\textbf{Mapas de saliencia interpretables en CRATE} con
    un tamaño de parche de \(8\). Cuando se presentan imágenes con propiedades similares (quizás
    pero no necesariamente de la misma clase), los mapas de saliencia
    correspondientes a diferentes cabezales de atención en la última capa
    resaltan cada uno una propiedad específica. Se puede observar que el mapa de
    saliencia promedio (no incluido) resalta todos los objetos relevantes de
    la imagen, lo que demuestra que utiliza todos los detalles más sutiles de la imagen de entrada
    para la clasificación. Este es el \textit{primer} sistema de aprendizaje automático
    que hace esto, según el conocimiento de los autores, y mucho menos en forma automática
sin haber sido entrenado con datos de segmentación.}
\label{fig:crate_semantic_heads}
\end{figure}

\begin{table}
\caption{\small \textbf{Detección de objetos y segmentación de alta resolución
    mediante MaskCut en COCO {val2017}~\citep{lin2014microsoft}}.
    En este caso, todos los modelos se entrenan con un tamaño de parche de \(8\) en lugar de \(16\).
    CRATE obtiene, sin lugar a dudas, mejores resultados que ViT en las métricas de detección y
segmentación cuando ambos se entrenan mediante clasificación supervisada.}
\centering
\begin{tabular}{@{}lcccccccc@{}}
    \toprule
    &  \multicolumn{3}{c}{Detección (\(\uparrow\))} &
    \multicolumn{3}{c}{Segmentación (\(\uparrow\))} \\
    Modelo & AP$_{50}$ & AP$_{75}$ & AP & AP$_{50}$ & AP$_{75}$ & AP  \\
    \midrule
    CRATE-S/8 & \textbf{2.9} & \textbf{1.0} & 1.1 & 1.8 & \textbf{0.7} & 0.8 \\
    CRATE-B/8 & \textbf{2.9} & \textbf{1.0} & \textbf{1.3} &
    \textbf{2.2} & \textbf{0.7} & \textbf{1.0} \\
    ViT-S/8 & 0.1& 0.1 & 0.0 & 0.0 & 0.0 & 0.0 \\
    ViT-B/8 & 0.8 & 0.2 & 0.4 & 0.7 & 0.5 & 0.4 \\
    \bottomrule
\end{tabular}
\label{tab:crate_detection_segmentation}
\end{table}

\paragraph{Resultados experimentales.}

La \Cref{tab:crate_classification_linear_probing} demuestra que
los modelos CRATE alcanzan un rendimiento equivalente o superior al de la popular
arquitectura Vision Transformer (ViT) con un número similar de parámetros, al
menos en lo que respecta a la separabilidad lineal de sus características
con respecto a las diferentes clases. En cuanto a la fidelidad de los mapas de atención,
\Cref{fig:crate_semantic_heads} muestra un resultado verdaderamente extraordinario:
sin necesidad de entrenar con datos de segmentación o detección de objetos,
\textit{no solo} los mapas de saliencia capturan eficazmente
todas las partes relevantes de la imagen de entrada, sino que
\textit{se auto-organizan} para que cada uno corresponda a un conjunto discreto de 
conceptos, ¡incluso entre muestras y clases! Este es el primer sistema
que logra esto, según el conocimiento de los autores, y puede hacerlo sin
utilizar ningún dato adicional, salvo los datos de clasificación de imágenes.
La \Cref{tab:crate_detection_segmentation} confirma estas observaciones cualitativas
de manera cuantitativa, mostrando una mejora significativa con respecto a los ViT
entrenados en la misma configuración de clasificación supervisada.

\section{Representación de imágenes mediante auto-codificación enmascarada}\label{sec:image_completion}

En esta sección, se muestra cómo aplicar la arquitectura CRATE al
problema del \textit{completamiento de imágenes}, también conocido como \textit{auto-codificación enmascarada} 
(MAE), que puede considerarse una generalización del
problema del completamiento de matrices de rango bajo tratado en el \Cref{ch:classic-models}.
El auto-codificador enmascarado, desde su introducción en el contexto del aprendizaje profundo
por \cite{he2022masked}, ha sido un método básico y sencillo
de aprendizaje de representaciones auto-supervisado, cuyo objetivo es dotar
a cada característica de parche dentro de \(\vZ_{\theta}\) de información agregada
así como de información sobre sus vecinos, de modo que tanto la característica
del parche como las características agregadas sean fuentes ricas de información para
la muestra completa.

Los datos que se utilizan para esta tarea son los mismos que los conjuntos de datos de imágenes
analizados en la \Cref{sub:contrastive_learning_data}. Como de costumbre, se siguen
aplicando aumentos de datos a cada muestra de cada nuevo lote.

\subsection{La tarea MAE y su objetivo}\label{sub:image_completion_objective}

Como su nombre indica, la auto-codificación enmascarada implica una vista \(v_{m}\)
que, dada una entrada, realiza un recorte aleatorio con cambio de tamaño (véase
\Cref{sub:contrastive_learning_objective}) para convertir la imagen de entrada
en una imagen cuadrada de tamaño \((C, S_{\mask}, S_{\mask})\), y luego
\textit{enmascara} (es decir, pone a cero) un porcentaje fijo \(p_{\mask}
\in [0, 1]\) de los píxeles de la entrada. Por motivos de eficiencia
\footnote {La implementación original de MAE por
\cite{he2022masked} incrusta la imagen completa, \textit{elimina} los
tokens que se enmascararían, pasa el conjunto de tokens resultante a través
del codificador, vuelve a añadir los tokens de marcador de posición aprendidos en los espacios enmascarados
y vuelve a añadir la codificación posicional adecuada, y pasa el
conjunto de tokens resultante a través del decodificador para obtener la predicción de auto-codificación.
Esto resulta más eficiente, ya que el codificador tiene que procesar menos tokens,
pero, en esencia, es igual que el método descrito
en el texto, y el rendimiento de los modelos resultantes en la tarea de auto-codificación enmascarada
y en las evaluaciones posteriores es muy similar.}, el
enmascaramiento se realiza por parches; es decir, tras incrustar la imagen completa,
el porcentaje \(p_{\mask}\) de los parches se establece en cero.  El
objetivo del MAE es entrenar un codificador \(f_{\theta} \colon \cI \to
(\R^{d})^{*}\) y un decodificador \(g_{\eta} \colon (\R^{d})^{*} \to \cI\)
capaces de reconstruir una entrada a partir de su enmascarado; es decir, si se establece que
\(\hat{\vX}_{\theta, \eta} \doteq g_{\eta} \circ f_{\theta}\), se tiene
\begin{equation}
\min_{\theta, \eta}\bc{\cL_{\mathrm{MAE}}(\theta, \eta) \doteq
\Ex\norm{\hat{\vX}_{\theta, \eta}(\vX_{m}) - \vX}_{F}^{2}}
\end{equation}
En esencia, esto significa que las características \(\vZ_{\theta}(\vX_{m})\) de
la vista \(\vX_{m} \doteq v_{m}(\vX)\) deben contener información sobre
los \textit{parches enmascarados}, así como sobre los parches no enmascarados. Desde la
perspectiva de los modelos transparentes white-box basados en compresión descritos en
\Cref{ch:deep-networks}, si un auto-codificador white-box  \((f_{\theta},
g_{\eta})\) tiene éxito en esta tarea, significa que los
subespacios y diccionarios aprendidos realizan una codificación \textit{redundante} de con compresión
los datos, de tal manera que puede reconstruir las partes faltantes de los datos a partir de
otras partes codificadas de los datos. Esto significa que la información sobre
cada parche se almacena en otros parches. Por lo tanto, cada característica de un parche
debería contener tanto información sobre el parche como información sobre
las estadísticas de toda la imagen. Así pues, una vez más, se espera que las
representaciones contengan información semánticamente relevante tanto a nivel local como global,
y, por lo tanto, las representaciones de diferentes
parches con información local y global similar deberían estar relacionadas
(es decir, en el mismo subespacio o codificadas juntas por un diccionario).

\subsection{Arquitectura MAE basada en CRATE}\label{sub:image_completion_architecture}

\begin{figure}
\centering
\includegraphics[width=\textwidth]{\toplevelprefix/chapters/chapter8/figs/crate_ae_backbone.pdf}
\caption{\small\textbf{Una capa del codificador y del decodificador en una
    estructura central de auto-codificador CRATE.} Tanto las capas del codificador como las del decodificador pasan
    sus entradas a través de una auto-atención de subespacio de múltiples cabezas y un
    paso de aprendizaje de diccionario o codificación de diccionario. Tenga en cuenta que las
    capas del codificador y del decodificador están diseñadas simétricamente; el objetivo conceptual
    de cada capa del decodificador es invertir una capa del codificador, por lo que esta
    simetría es en gran medida intencionada (véase, por ejemplo, el
\Cref{ch:consistent-representations}).}
\end{figure}

Se utilizan un codificador y un decodificador CRATE, tal y como se muestra en la
\Cref{fig:transformer_backbone}, aunque, por supuesto, es posible
utilizar un codificador y un decodificador Transformer estándar. A continuación se ofrecen más detalles.

\paragraph{El codificador.} El codificador es el mismo que el codificador CRATE
de la \Cref{sub:image_classification_architecture}, con la salvedad de que
no hay ningún extractor de características \(f_{\theta}^{\ext}\). Sin embargo, tanto
la incrustación \(f_{\theta}^{\emb}\) como la red troncal
\(f_{\theta}^{\backbone}\) son las mismas.

\paragraph{La red troncal del decodificador.} La red troncal del decodificador es el decodificador CRATE
descrito en el \Cref{ch:consistent-representations}. Para completar la información,
lo describimos a continuación. Dada una secuencia de características \(\vZ_{\theta}(\vX) \doteq
f_{\theta}(\vX) \in (\R^{d})^{*}\), podemos procesarla utilizando la
red troncal del decodificador \(g_{\eta}^{\backbone}\) de la siguiente manera. La función
\(g_{\eta}^{\backbone}\) se compone de \(L\) capas \(g_{\eta}^{\ell}\), es decir,
\begin{equation}
g_{\eta}^{\backbone} = g_{\eta}^{L} \circ \cdots \circ g_{\eta}^{1}.
\end{equation}
La capa \(g_{\eta}^{\ell}\) se implementa de la siguiente manera.
En primer lugar, se define \(\tilde{\vZ}_{\theta, \eta}^{1}(\vX) \doteq
\vZ_{\theta}(\vX)\). Entonces, se obtiene
\begin{align}
\tilde{\vZ}_{\theta, \eta}^{\ell + 1/2}(\vX)
&= [\tilde{\vD}^{\ell}]^{\top}\LN_{\eta}^{1,
\ell}(\tilde{\vZ}_{\theta, \eta}^{\ell}(\vX)) \\
\tilde{\vZ}_{\theta, \eta}^{\ell + 1}(\vX)
&= \tilde{\vZ}_{\theta, \eta}^{\ell + 1/2}(\vX) -
\MSSA_{\eta}^{\ell}(\LN_{\eta}^{2, \ell}(\tilde{\vZ}_{\theta,
\eta}^{\ell + 1/2}))
\end{align}
y \(g_{\eta}^{\ell}\) se define de tal manera que
\(g_{\eta}^{\ell}(\tilde{\vZ}_{\theta, \eta}^{\ell}) \doteq
\tilde{\vZ}_{\theta, \eta}^{\ell + 1}(\vX)\). Aquí, el concepto relevante
es que \(g_{\eta}^{\ell}\) debe aprender una inversa aproximada
de \(f_{\theta}^{L + 1 - \ell}\), como discretizaciones de un
proceso de difusión hacia adelante y hacia atrás en el tiempo, respectivamente. En
particular, \(\tilde{\vD}^{\ell}\) debería aproximarse a \(\vD^{L + 1 -
\ell}\), y de manera similar, los parámetros de \(\MSSA_{\eta}^{\ell}\) deberían
ser similares a los parámetros de \(\MSSA_{\theta}^{L + 1 - \ell}\).
La salida es \(\tilde{\vZ}_{\theta, \eta} \doteq
\tilde{\vZ}_{\theta, \eta}^{L + 1}\).

\paragraph{El módulo de desincrustación.} Para transformar
\(\tilde{\vZ}_{\theta, \eta}(\vX)\) de nuevo en una estimación de
\(\vX\), es preciso revertir el efecto del módulo de incrustación
\(f_{\theta}^{\emb}\) utilizando el módulo de desincrustación
\(g_{\eta}^{\unemb}\). Así, volviendo a la forma funcional
del módulo de incrustación en \eqref{eq:definition_of_embedding_module}, es decir,
\begin{equation}
f_{\theta}^{\emb}(\vX) \doteq \mat{\vz_{\cls}^{1},
\vW^{\emb}f^{\patch}(\vX) + \vE^{\pos}}
\end{equation}
esto implica que nuestra operación inversa \(g_{\eta}^{\unemb}\) tiene
el siguiente aspecto:
\begin{align}
g_{\eta}^{\unemb}(\tilde{\vZ}) &\doteq
g_{\eta}^{\unemb}(\mat{\tilde{\vz}^{1}, \dots, \tilde{\vz}^{n}}) \\
&=
g^{\unpatch}(\vW^{\unemb}([\tilde{\vz}^{2}, \dots, \tilde{\vz}^{n}] -
\tilde{\vE}^{\pos})),
\end{align}
donde \(g^{\unpatch}\) realiza la operación inversa al proceso de
desenrollado y aplanamiento que lleva a cabo \(f^{\patch}\).\footnote{Una vez más,
la «codificación posicional inversa» \(\tilde{\vE}^{\pos}\) se aprende
para entradas grandes, y para entradas más pequeñas puede interpolarse. Incluso
es posible establecer directamente \(\tilde{\vE}^{\pos}\) igual a la
codificación posicional \(\vE^{\pos}\) y utilizar las mismas codificaciones posicionales interpoladas
para cada entrada tanto en el codificador como en el decodificador.}

Esta arquitectura es un auto-codificador white-box \((f_{\theta},
g_{\eta})\), donde (recordemos) \(f_{\theta} = f_{\theta}^{\backbone}
\circ f_{\theta}^{\emb}\) y \(g_{\eta} = g_{\eta}^{\unemb} \circ
g_{\eta}^{\backbone}\). En particular, podemos usarla para calcular una
estimación de una vista enmascarada \(\hat{\vX}_{\theta, \eta}(\vX_{m}) =
(g_{\eta} \circ f_{\eta})(\vX_{m})\), que debería aproximarse al propio
\(\vX\).

\subsection{Optimización}\label{sub:image_completion_optimization}

Al igual que en la \Cref{sub:image_classification_optimization}, se utiliza una configuración de optimización sencilla:
se toman muestras de imágenes y máscaras, se calcula la pérdida sobre
esas muestras y los gradientes de dicha pérdida y se actualiza los
parámetros utilizando un algoritmo de optimización genérico y los
gradientes mencionados anteriormente. Para cada paso temporal \(k\), se hace lo siguiente:
\begin{itemize}
\item Tomar una submuestra \(B\) de muestras diferentes \(\{\vX_{b}^{(k)}\}_{b =
    1}^{B} \subseteq \cI\).
\item Para cada muestra \(\vX_{b}^{(k)}\), calcular un recorte y una máscara aleatorios
    de tamaño diferente \(v_{b, m}^{(k)}\) y aplicarlos
    a \(\vX_{b}^{(k)}\) para obtener \(\vX_{b, m}^{(k)} \doteq v_{b,
    m}^{t}(\X_{b}^{(k)})\).
\item Calcular la auto-codificación estimada \(\hat{\vX}_{\theta,
        \eta}(\vX_{b, r}^{(k)}) \doteq (g_{\eta} \circ
    f_{\theta})(\vX_{b, r}^{(k)})\).
\item Formular la pérdida estocástica sustitutiva
    \begin{equation}
        \hat{\cL}_{\mathrm{MAE}}^{(k)}(\theta, \eta) \doteq
        \frac{1}{B}\sum_{b = 1}^{B}\norm{\hat{\vX}_{\theta,
        \eta}(\vX_{b, r}^{(k)}) - \vX_{b}^{(k)}}_{F}^{2}.
    \end{equation}
\item Calcular un paso de un algoritmo de optimización sobre \((\theta,
    \eta)\), lo que da lugar a la siguiente iteración:
    \begin{equation}
        (\theta^{(k + 1)}, \eta^{(k + 1)}) \doteq
        \textsc{OptUpdate}^{(k)}(\theta^{(k)}, \eta^{(k)};
        \nabla_{(\theta, \eta)}\hat{\cL}_{\mathrm{MAE}}^{(k)}).
    \end{equation}
\end{itemize}

\subsection{Evaluación} \label{sub:image_completion_optimization_1}

Esta es la primera red de auto-codificador que se analiza en este capítulo. Se utiliza
la misma vista recortada central \(v_{\cc}\) que en
\Cref{sub:contrastive_learning_evals,sub:image_classification_evals},
redimensionando la imagen final a un cuadrado con una longitud de lado de \(S_{\cc} =
S_{\mask}\) píxeles para que coincida con las formas de las imágenes de entrada observadas durante el entrenamiento.

Además de evaluar la pérdida de auto-codificación enmascarada en sí misma, también es
posible evaluar directamente las características \(\vZ_{\theta}(\vX_{\cc})\) de
la vista \(\vX_{\cc} \doteq v_{\cc}(\vX)\) de los datos \(\vX\)
. Para la evaluación de la fidelidad del mapa de atención, basta con obtener
\(\vZ_{\theta}(\vX_{\cc})\), pero para el sondeo lineal (linear probing)
se requiere extraer una característica resumida o agregada de
\(\vZ_{\theta}\). Para ello, se puede utilizar un mapa de extracción de características (sin parámetros)
que devuelva la característica correspondiente al token de clase, es decir,
\begin{equation}
f_{\theta}^{\ext}(\vZ) \doteq f_{\theta}^{\ext}([\vz^{1}, \dots,
\vz^{n}]) = \vz^{1},
\end{equation}
como (por ejemplo) en 
\Cref{sub:image_classification_objective,sub:image_classification_architecture}.
Con esto, se tiene una forma de obtener características agregadas
\(\vz_{\theta}(\vX_{\cc}) \doteq (f_{\theta}^{\ext} \circ
f_{\theta})(\vX_{\cc})\), momento en el que se pueden realizar sondeos lineales,
evaluaciones de segmentación, etc.

\subsection{Experimentos}\label{sub:image_completion_experiments}

Dado que CRATE-MAE se basa directamente en ViT-MAE, se comparan los parámetros óptimos
de ViT-MAE, tal y como se indican en \citep{he2022masked}, con los mismos
parámetros aplicados a CRATE-MAE, con el fin de realizar una comparación imparcial.

\paragraph{Arquitectura del modelo.} Durante el entrenamiento, el recorte enmascarado
\(v_{m}\) redimensiona toda la imagen de modo que el borde más corto tenga un tamaño de
\(256\) (es decir, \(S_{\rsz} = 256\)) antes de realizar un recorte aleatorio de
tamaño \(224 \times 224\) (es decir, \(S_{\mathrm{mask}} = 224\)), y
enmascarando \(p_{\mathrm{mask}} = \frac{3}{4}\) de los parches. Se toma
un tamaño de parche de \(16\) (es decir, \(P_{H} = P_{W} = 16\)). Se utilizan las variantes 'small'
y 'base' de la arquitectura ViT-MAE como incrustación y
backbone tanto para el codificador como para el decodificador, sustituyendo las componentes MHSA y
MLP por MSSA, ISTA y capas lineales, respectivamente. En el caso de MSSA,
se emplean el mismo número de cabezales y la misma dimensión de cabezales.
Sin embargo, el ViT-MAE original utiliza un codificador que emplea casi todas
las capas totales y un decodificador que solo utiliza unas pocas capas; nosotros
asignamos la mitad del número total de capas (que se mantiene igual de
ViT-MAE a CRATE-MAE) a nuestro codificador y decodificador, tal como sugiere el
marco conceptual y teórico del
\Cref{ch:consistent-representations}. Para
CRATE-MAE se establece \((\beta, \lambda) = (1, 0.1)\).

\paragraph{Conjuntos de datos y optimización.} Para el preentrenamiento, se utiliza el
conjunto de datos ImageNet-1K. Se emplea el optimizador AdamW para preentrenar tanto nuestra
replicación de ViT-MAE como CRATE-MAE. Se establece la tasa de aprendizaje
base en \(3 \times 10^{-5}\), la tasa de desvanecimiento de pesos en \(0.1\) y el tamaño del lote
en \(B = 4096\). Nuestro programa de tasas de aprendizaje aumenta la
tasa de aprendizaje en forma lineal hasta la tasa de aprendizaje base durante las primeras
\(40\) épocas, y la reduce a \(0\) siguiendo un programa coseno durante
las siguientes \(760\) épocas (entrenando todos los modelos durante \(800\) épocas
cada uno). Para el preentrenamiento, se aplica el régimen habitual de
aumentos de datos (volteos, desenfoques gaussianos, solarización, etc.) a los datos de imagen.

Para el sondeo lineal (linear probing), se utilizan varios conjuntos de datos de evaluación, tales como
CIFAR10, CIFAR100, Oxford-Flowers y Oxford-IIT-Pets. 
Se precalculan las características de todas las muestras del conjunto de datos
de destino y se aplica un solucionador rápido de regresión lineal, por ejemplo, de un
paquete estándar como Scikit-Learn.

\paragraph{Resultados experimentales.} La \Cref{tab:crate_mae_linear_probing}
demuestra que los modelos CRATE-MAE alcanzan, a grandes rasgos, un rendimiento equivalente
al de la popular arquitectura ViT-MAE con un número similar de parámetros,
y también que el rendimiento en el aprendizaje de características (medido por
el rendimiento en tareas de clasificación posteriores) aumenta con la escala.
Por otra parte, la \Cref{fig:crate_mae_semantic_heads} muestra que los
mapas de saliencia del codificador (y, por lo tanto, las características de alta resolución 
aprendidas por el codificador) realmente aíslan y resaltan las partes clave de
la imagen de entrada.

\begin{table}
\caption{\small\textbf{Precisión en la clasificación mediante sondeo lineal de
    CRATE-MAE y ViT-MAE} en diversos conjuntos de datos con diferentes tamaños de modelo
    cuando la red base se ha entrenado previamente para la auto-codificación enmascarada en
    ImageNet-1K. Con el mismo número de parámetros, CRATE-MAE alcanza
    un rendimiento más o menos similar, al tiempo que cuenta con un diseño arquitectónico más sencillo
y basado en principios.}
\centering
\resizebox{1\columnwidth}{!}
{
    \begin{tabular}{@{}lcc|cc@{}}
        \toprule
        \textbf{Modelo} & CRATE-MAE-S(mall) & CRATE-MAE-B(ase) &
        {\color{gray} ViT-MAE-S} & {\color{gray} ViT-MAE-B} \\
        \midrule
        \midrule
        \# parameters & 25.4M & 44.6M & 47.6M & {\color{gray}143.8M} \\
        \midrule
        CIFAR10 & 79.4 & 80.9 & {\color{gray} 79.9} & {\color{gray} 87.9} \\
        CIFAR100 & 56.6 & 60.1 & {\color{gray} 62.3} & {\color{gray} 68.0} \\
        Oxford Flowers-102 & 57.7 & 61.8 & {\color{gray} 66.8} &
        {\color{gray} 66.4} \\
        Oxford-IIIT-Pets & 40.6 & 46.2 & {\color{gray} 51.8} &
        {\color{gray} 80.1} \\
        \bottomrule
    \end{tabular}
}

\label{tab:crate_mae_linear_probing}
\end{table}

\begin{figure}
\centering
\includegraphics[width=\textwidth]{\toplevelprefix/chapters/chapter8/figs/crate_mae_semantic_heads.pdf}
\caption{\small\textbf{Mapas de saliencia de CRATE-MAE.} Cada par de
    imágenes consta de la imagen original (izquierda) y un mapa de saliencia
    seleccionado (derecha) correspondiente a una cabeza de atención de la última capa. Como
    es habitual en los modelos CRATE, pero inusual en los modelos generales de tipo
transformer, los mapas de saliencia se corresponden con los objetos de la imagen de entrada.}
\label{fig:crate_mae_semantic_heads}
\end{figure}

\section{Generación de imágenes mediante auto-codificación y muestreo}
\label{sec:auto-encode-img-gen}

En la sección anterior, hemos presentado la auto-codificación enmascarada (MAE) como un enfoque 
para aprender una auto-codificación de imágenes naturales $\x \sim p(\x)$:
\begin{equation}
\x\xrightarrow{\hspace{1mm} \mathcal{E} = f \hspace{1mm}} \z
\xrightarrow{\hspace{1mm} \mathcal{D} = g \hspace{1mm}} \hat{\x}.
\end{equation}
donde $\z = f(\x)$ es un determinado espacio de características. Este auto-codificador está diseñado 
para aprender la correlación, o coocurrencia, entre los parches que componen una imagen natural.

Sin embargo, el MAE no impone que la distribución de las características aprendidas $\z$ tenga 
estructuras o propiedades explícitas, lo cual puede ser muy importante si queremos utilizar las 
características así aprendidas para ciertas tareas posteriores, como el reconocimiento o la 
generación controlada.Además, la pérdida MMSE simple que se utiliza para completar la imagen no 
respeta necesariamente las posibles estructuras no lineales de la variedad de la imagen y, por 
lo general, genera un valor esperado---y, por lo tanto, una versión borrosa---para los parches 
de imagen que faltan, tal como hemos mostrado en la Figura \ref{fig:mae_autoencoding-small} del 
anterior Capítulo \ref{ch:conditional-inference}.

Por lo tanto, en esta sección presentaremos otros métodos para aprender representaciones de (la 
distribución de) imágenes naturales mediante auto-codificación, cuyo objetivo es, en gran medida, 
subsanar los dos problemas mencionados anteriormente. Más concretamente, queremos que la 
distribución de las características codificadas presente ciertas estructuras deseadas y que las 
imágenes decodificadas a partir de dichas características sean mucho más fieles a la distribución 
de las imágenes naturales; es decir, que tengan una calidad visual mucho mejor.


\subsection{Auto-codificación variacional de datos de imágenes}
% As we have introduced in Section \ref{sec:vae} of Chapter \ref{ch:consistent-representations}, variational auto-encoding is a method that aims to promote the distribution of learned features $\z$ and the conditional distributions between $\x$ and $\z$ to be more Gaussian like.\footnote{One may view such prior is to promote the learned the features to be diverse, in a similar spirit to maximizing the coding rate of all features.} 

% Here we introduce a popular implementation of a VAE for imagery data...




% \subsection{Variational Auto-Encoding of Imagery Data}
\label{sec:VAE-implement}


%refer back to Chapter 6 VAE here.
Como hemos presentado en la Sección~\ref{sec:vae} del Capítulo~\ref{ch:consistent-representations}, 
la auto-codificación variacional es un método que tiene como objetivo hacer que la distribución de 
las características aprendidas $\z$ y las distribuciones condicionales entre $\x$ y $\z$ se asemejen 
más a una distribución gaussiana.\footnote{Se puede considerar que dicha distribución a priori 
tiene como objetivo promover la diversidad de las características aprendidas, en un espíritu similar 
al de maximizar la tasa de codificación de todas las características.}

Aquí presentamos una implementación muy utilizada de una VAE para datos de imágenes, siguiendo el 
modelo de compresión perceptual empleado en Stable Diffusion~\citep{rombach2022high}. La idea clave 
es entrenar un auto-codificador que mapee imágenes de alta resolución en un espacio latente de baja 
dimensión, al tiempo que se impone una regularización KL sobre la distribución latente para fomentar 
que se acerque a una gaussiana estándar $\mathcal{N}(\boldsymbol{0}, \boldsymbol{I})$. Este espacio 
latente puede funcionar posteriormente como un espacio de trabajo eficaz para modelos generativos posteriores, como 
la difusión latente. Tanto el codificador como el decodificador se entrenan conjuntamente de extremo 
a extremo. A continuación, describimos los detalles.

\paragraph{El codificador ($\mathcal{E}$).}
El codificador $\mathcal{E}$ mapea una imagen RGB de entrada $\x \in \mathbb{R}^{H \times W \times 3}$ 
a un espacio latente de baja dimensión. Su arquitectura, adoptada de~\citep{esser2021taming}, es una 
red totalmente convolucional. Comienza con una convolución de entrada que proyecta la entrada RGB de 3 
canales en un mapa de características con un recuento de canales base de $128$. A continuación, se desarrollan cuatro 
etapas, cada una de las cuales consta de dos bloques ResNet~\citep{he2016deep}. La dimensión del canal 
se escala mediante multiplicadores $(1, 2, 4, 4)$ a lo largo de las etapas, lo que da como resultado 
mapas de características con $128, 256, 512, 512$ canales, respectivamente. Se realiza un submuestreo 
espacial por un factor de $2\times$ entre etapas consecutivas mediante convoluciones con paso, lo que 
da como resultado tres submuestreos y un factor de submuestreo global de $f = 2^3 = 8$. Tras las cuatro 
etapas, un bloque intermedio compuesto por dos bloques ResNet adicionales refina aún más las características. 
Por último, una convolución de salida proyecta las características en $2c$ canales, que se dividen en las 
dos magnitudes que parametrizan la distribución latente (descrita a continuación).

A diferencia de un auto-codificador estándar---y esto es fundamental---el codificador no genera un único 
vector latente determinista. En su lugar, parametriza una distribución gaussiana diagonal posterior 
$q(\z | \x) = \mathcal{N}(\boldsymbol{\mu}, \boldsymbol{\sigma}^2 \boldsymbol{I})$ generando dos magnitudes:
\begin{itemize}
    \item Una media $\boldsymbol{\mu} = \mathcal{E}_\mu(\x) \in \mathbb{R}^{h \times w \times c}$ y
    \item Una varianza logarítmica $\log \boldsymbol{\sigma}^2 = \mathcal{E}_\sigma(\x) \in \mathbb{R}^{h \times w \times c}$.
\end{itemize}
La representación latente se obtiene entonces mediante lo que se conoce como \emph{«truco de reparametrización»}~\citep{kingma2013auto}:
\begin{equation}
    \z = \boldsymbol{\mu} + \boldsymbol{\sigma} \hada \boldsymbol{\epsilon}, \quad \boldsymbol{\epsilon} \sim \mathcal{N}(\boldsymbol{0}, \boldsymbol{I}),
\end{equation}
donde $\hada$ denota la multiplicación elemento por elemento. Esta formulación permite la optimización 
basada en el gradiente mediante el paso de muestreo estocástico.

\paragraph{El Decodificador ($\mathcal{D}$).}
El decodificador $\mathcal{D}$ aprende el mapeo inverso del espacio latente al espacio 
de píxeles, generando la imagen reconstruida $\tilde{\x} = \mathcal{D}(\z) \in \mathbb{R}^{H \times W \times 3}$. 
Su arquitectura es aproximadamente simétrica a la del codificador: una convolución de entrada 
proyecta primero el latente de $c$ canales en el espacio de características, seguida de un bloque 
intermedio de dos bloques ResNet. A continuación, cuatro etapas---cada una con dos bloques ResNet---aumentan 
progresivamente la resolución espacial, con unas dimensiones de canales que reflejan las del codificador en 
orden inverso ($512, 512, 256, 128$). Una convolución de salida final mapea las características de vuelta a 
3 canales RGB. Tanto la arquitectura del codificador como la del decodificador siguen el diseño convolucional 
introducido en ~\citep{esser2021taming}.

\subsubsection{Datos}
Suponemos que se dispone de un conjunto de datos $\mathcal{D} = \{\x_i\}_{i=1}^N$ de muestras procedentes 
de la distribución original de los datos. Cada muestra es una imagen de la forma $\x \in \mathbb{R}^{H \times W \times 3}$, 
donde $H$ y $W$ son la altura y el ancho, respectivamente. En~\citep{rombach2022high}, el auto-codificador se 
entrena en el conjunto de datos OpenImages con imágenes recortadas a una resolución de $256 \times 256$.

Dado el factor de submuestreo $f = 8$ y el número de canales latentes $c = 4$, el codificador mapea cada imagen 
de entrada a un tensor latente $\z \in \mathbb{R}^{H/f \times W/f \times c}$. Para la resolución de entrenamiento 
predeterminada, esto da como resultado un tensor latente de tamaño $32 \times 32 \times 4$. En el momento de la 
inferencia, el auto-codificador se generaliza a otras resoluciones; por ejemplo, una entrada 
de $512 \times 512 \times 3$ se codifica en un latente de $64 \times 64 \times 4$.

\subsubsection{Objetivo y métricas de evaluación}
Para entrenar el VAE, nuestro objetivo es optimizar conjuntamente el codificador $\mathcal{E}$ y el 
decodificador $\mathcal{D}$ de tal manera que la reconstrucción 
$\tilde{\x} = \mathcal{D}(\mathcal{E}(\x))$ se acerquen lo más posible a los datos originales 
$\x$, al tiempo que regularizamos la distribución latente para que se aproxime a una distribución 
gaussiana a priori estándar. Una elección empíricamente eficaz de la pérdida de reconstrucción 
consiste en una combinación ponderada de las pérdidas L1, LPIPS~\citep{zhang2018unreasonable} y 
las pérdidas adversarias, tal como se hace en las GAN~\citep{goodfellow2014generative}. Además, 
se incluye un término de divergencia de KL para regularizar el espacio latente. Formalmente, el 
objetivo de entrenamiento puede expresarse como:

\begin{align}
\label{eq:vae_loss}
\min_{\theta_{\mathcal{E}}, \theta_{\mathcal{D}}} \; \mathcal{L}_{\mathrm{VAE}}(\theta_{\mathcal{E}}, \theta_{\mathcal{D}}) 
&= \mathbb{E}_{\x \sim \mathcal{D}} \Big[
    \|\tilde{\x} - \x\|_1 + \lambda_1 \, \mathrm{LPIPS}(\tilde{\x}, \x) \\
    &\qquad\quad + \lambda_2 \, \mathcal{L}_{\mathrm{GAN}}(\tilde{\x}, \x) + \; \lambda_3 \, D_{KL}\!\big(q(\z | \x) \,\big\|\, p(\z)\big)
\Big],
\end{align}
donde $\theta_{\mathcal{E}}$ y $\theta_{\mathcal{D}}$ son los parámetros del codificador y del 
decodificador, respectivamente, y $\lambda_1, \lambda_2, \lambda_3$ son hiperparámetros que 
equilibran las contribuciones de cada término de pérdida. Los tres primeros términos garantizan 
la consistencia a nivel de muestra de la misma manera que los objetivos de reconstrucción de la 
Sección~\ref{sec:RAE-implement}. El cuarto término---la divergencia KL entre la distribución a 
posteriori del codificador $q(\z | \x) = \mathcal{N}(\boldsymbol{\mu}, \boldsymbol{\sigma}^2 \boldsymbol{I})$ 
y la distribución a priori gaussiana estándar $p(\z) = \mathcal{N}(\boldsymbol{0}, \boldsymbol{I})$---regulariza 
el espacio latente. Para una distribución a posteriori gaussiana diagonal, esto admite una expresión en forma cerrada:
\begin{equation}
    D_{KL}\!\big(q(\z | \x) \,\big\|\, p(\z)\big) = \frac{1}{2} \sum_{j=1}^{d} \left( \mu_j^2 + \sigma_j^2 - \log \sigma_j^2 - 1 \right),
\end{equation}
donde $d = h \times w \times c$ es la dimensionalidad total del espacio latente.

El detalle de implementación de ~\citep{rombach2022high} es que el peso KL $\lambda_3$ se establece 
en un valor muy pequeño ($\approx 10^{-6}$). Esto se debe a que el objetivo principal del término KL 
aquí no es permitir el muestreo incondicional del espacio latente (como en una VAE tradicional), sino 
más bien evitar que el espacio latente desarrolle una varianza arbitrariamente alta, lo que degradaría 
la relación señal-ruido para los modelos posteriores que operan en este espacio.

Para la evaluación, medimos la calidad de la reconstrucción calculando la puntuación 
FID~\citep{heusel2017gans} en las imágenes reconstruidas del conjunto de validación ImageNet-1k, 
lo que denominamos rFID. Un valor rFID más bajo indica una mejor calidad de reconstrucción.

\subsection{Auto-codificación de representaciones de datos de imágenes}
\label{sec:RAE-implement}
% \yima{Conceptual connections and necessity... Practical
% implementations details... (Robin and Peter)}

% rae stuff
\definecolor{deemph}{gray}{0.6}
\newcommand{\gc}[1]{\textcolor{deemph}{#1}}
\newcommand{\tablestyle}[2]{\setlength{\tabcolsep}{#1}\renewcommand{\arraystretch}{#2}\centering\scriptsize}
% \newcommand{\round}[1]{\num[round-mode=places, round-precision=2]{#1}}

Como se discute en el \Cref{ch:lossy-compression}, el aprendizaje de
representaciones estructuradas de los datos es fundamental para muchas tareas posteriores. En
el aprendizaje de representaciones visuales, esto se suele lograr mediante
métodos de aprendizaje auto-supervisado como DINO
\cite{caron2021emerging} y SimDINO (analizados en la Sección
\ref{sec:contrastive_learning}). En el
\Cref{ch:consistent-representations}, postulamos
que una buena representación de los datos debe poder decodificarse de manera efectiva
hacia el espacio de datos original. Esto requiere un marco de auto-codificación
consistente que pueda reconstruir con precisión los datos originales
a partir de su representación aprendida. Por lo tanto, es lógico preguntarse: ¿es
consistente la representación aprendida por métodos como DINO y puede
decodificarse de manera efectiva en un marco de auto-codificación?

Para averiguarlo, podemos entrenar un \textit{auto-codificador de representación}
(RAE) que consta de un codificador de representación preentrenado (por ejemplo,
DINO) y un decodificador entrenable, siguiendo el trabajo:\\
\centerline{\url{https://rae-dit.github.io}.}\\
El codificador $f$ mapea una muestra de entrada
$\x$ a su representación aprendida $\z = f(\x)$, mientras que el decodificador
$g$ mapea la representación de vuelta al espacio de datos original
$\x'=g(\z)$. Tal y como se detalla en \Cref{ch:consistent-representations}, un objetivo sencillo
consiste en garantizar la consistencia a nivel de muestra entrenando al
decodificador para minimizar la pérdida de reconstrucción entre los datos originales
$\x$ y la reconstrucción $\x'$. A continuación, describimos los detalles.

\paragraph{El codificador.} Utilizamos un codificador de representación preentrenado
$f \colon \cX \to \cZ$ que asigna una muestra de datos de entrada $\x \in \cX$ a
su representación aprendida $\z = f(\x) \in \cZ$. El codificador se mantiene
congelado durante el entrenamiento. En este caso, asumimos que utilizamos un codificador basado en Vision Transformer (ViT), tal como
se presenta en \Cref{sub:contrastive_learning_architecture}.

\paragraph{El decodificador.} Para simplificar la arquitectura, utilizamos un
decodificador entrenable basado en ViT $g \colon \cZ \to \cX$ que mapea la
representación aprendida $\z$ de vuelta al espacio de datos original,
generando la reconstrucción $\x' = g(\z)$.

\subsubsection{Datos}
Suponemos que tenemos acceso a un conjunto de datos $\cD = \{\x_i\}_{i=1}^N$ de muestras
procedentes de la distribución de datos original $\cX$. Cada muestra es una imagen de
la forma $\x \in \R^{3 \times H \times W}$, donde $H$ y $W$ son la
altura y el ancho, respectivamente. El codificador $f$ se configura con un tamaño de parche
de $p_e$ y una dimensión de características de $d$. En consecuencia, para cada
imagen de entrada $\x$, obtendríamos $N = HW/p_e^2$ tokens de dimensión
$d$. El decodificador $g$, con un tamaño de parche de $p_d$, mapea entonces estos tokens visuales
para reconstruir $\x' \in \R^{3 \times H\frac{p_d}{p_e}
\times W\frac{p_d}{p_e}}$ en el espacio de píxeles original. Por defecto,
establecemos $p_d = p_e$ para garantizar que la imagen reconstruida tenga la misma
resolución que la imagen original.

\subsubsection{Objetivo y métricas de evaluación} Para entrenar RAE, nuestro objetivo
es entrenar un decodificador $g$ de tal manera que la reconstrucción $\x' =
g(f(\x))$ sean lo más parecidos posible a los datos originales $\x$. Esto se puede
lograr minimizando una pérdida de consistencia por muestra sobre el
conjunto de datos $\cD$. Una opción empíricamente eficaz para la función de pérdida
es una combinación ponderada de L1 y LPIPS \cite {zhang2018unreasonable}
\footnote{Se trata de una pérdida perceptual basada en el cálculo de las distancias
entre características en diferentes capas de una red preentrenada. Esta
captura empíricamente la similitud perceptiva humana mejor que
las pérdidas basadas en píxeles. Para más detalles, consulte el artículo original.} y
pérdidas adversarias como en las GAN \cite{goodfellow2014generative}.
Formalmente, el objetivo de entrenamiento puede expresarse como:

\begin{equation}
\min_{\theta} \cL_{\mathrm{RAE}}(\theta) = \Ex_{\x \sim \cD} \left[
    \norm{\x' - \x}_1 + \lambda_1 \mathrm{LPIPS}(\x', \x) + \lambda_2
\cL_{\mathrm{GAN}}(\x', \x) \right],
\end{equation}
donde $\theta$ son los parámetros del decodificador $g$, y $\lambda_1,
\lambda_2$ son hiperparámetros que equilibran las contribuciones de cada término de pérdida.

Para la evaluación, medimos la calidad de la reconstrucción calculando
la puntuación FID \cite{heusel2017gans} en las imágenes reconstruidas del
conjunto de validación ImageNet-1k, lo que denominamos rFID. Un valor rFID más bajo
indica una mejor calidad de reconstrucción.

\subsubsection{Resultados experimentales}
Podemos comparar RAE entrenados con diferentes codificadores preentrenados,
entre ellos DINOv2, MAE y SigLIP2. Como referencia, también consideramos un
auto-codificador variacional (VAE) básico entrenado directamente sobre los píxeles de la imagen
sin ningún codificador preentrenado. Una opción común es el SD-VAE.
Los decodificadores RAE se entrenan en el conjunto de entrenamiento ImageNet-1K y
se evalúan en el conjunto de validación ImageNet-1K. Los resultados se
resumen en \Cref{tab:rae_recon_results}. En la \Cref{tab:decoder_scaling_no_noise}, 
podemos observar que los RAE con
codificadores preentrenados superan a los SD-VAE en cuanto a calidad de reconstrucción.
Además, los decodificadores de mayor tamaño conducen a una mejor calidad de reconstrucción,
como se muestra en \Cref{tab:decoder_scaling_no_noise}, al tiempo que
siguen siendo más eficientes que los VAE. \Cref{tab:encoder_scaling_dinov2}
muestra que la calidad de reconstrucción es estable en diferentes
tamaños de los codificadores DINO. Por último, \Cref{tab:linear_probing}
demuestra que las RAE también conservan representaciones de alta calidad,
alcanzando una precisión de búsqueda lineal mucho mayor con respecto a las VAE.

\begin{table*}[t]
\centering
{\captionsetup[subfloat]{labelfont=scriptsize}
    \caption{
    \small\textbf{Resultados de la reconstrucción y el sondeo lineal de los RAE.} }
    \begin{tabular}{@{}cc@{}}
        \subfloat[
            \scriptsize \textbf{Opción codificador.} Todos los codificadores superan a SD-VAE.
            \label{tab:recon_deterministic}
        ]
        {
            \begin{minipage}[t]{0.45\linewidth}{
                    \vspace{0pt}
                    \begin{center}
                        \tablestyle{4pt}{1.1}
                        \begin{tabular}{lc}
                            \toprule
                            Modelo & rFID \\
                            \midrule
                            \rowcolor{gray!20}
                            DINOv2-B & 0.49 \\
                            \rowcolor{white}
                            SigLIP2-B & 0.53 \\
                            MAE-B & 0.16 \\
                            \midrule
                            SD-VAE & \gc{0.62} \\
                            \bottomrule
                        \end{tabular}
                    \end{center}
                }
            \end{minipage}
        }
        &
        \subfloat[
            \scriptsize \textbf{Los decodificadores de mayor tamaño} mejoran el rFID y, al mismo tiempo,
            siguen siendo mucho más eficientes que los VAE.
            \label{tab:decoder_scaling_no_noise}
        ]
        {
            \begin{minipage}[t]{0.45\linewidth}{
                    \vspace{0pt}
                    \begin{center}
                        \tablestyle{4pt}{1.1}
                        \begin{tabular}{lcc}
                            \toprule
                            Decodificador & rFID & GFLOPs \\
                            \midrule
                            ViT-B & 0.58 & 22.2 \\
                            ViT-L & 0.50 & 78.1 \\
                            \rowcolor{gray!20}
                            ViT-XL & 0.49 & 106.7 \\
                            \midrule
                            \rowcolor{white}
                            SD-VAE & \gc{0.62} & \gc{310.4} \\
                            \bottomrule
                        \end{tabular}
                \end{center}}
            \end{minipage}
        }
        \\[1em]
        \subfloat[
            \scriptsize \textbf{Escalado del codificador.}  El rFID se mantiene estable
            en todos los tamaños de RAE.
            \label{tab:encoder_scaling_dinov2}
        ]{
            \begin{minipage}[t]{0.45\linewidth}{
                    \vspace{0pt}
                    \begin{center}
                        \tablestyle{4pt}{1.1}
                        \begin{tabular}{lc}
                            \toprule
                            Codificdor & rFID \\
                            \midrule
                            DINOv2-S & 0.52 \\
                            \rowcolor{gray!20}
                            DINOv2-B & 0.49 \\
                            \rowcolor{white}
                            DINOv2-L & 0.52 \\
                            \bottomrule
                            \multicolumn{2}{c}{~} \\
                            \multicolumn{2}{c}{~} \\
                        \end{tabular}
                        \vspace{-0.1cm}
                    \end{center}
                }
            \end{minipage}
        }
        &
        \subfloat[
            \scriptsize \textbf{Calidad de la representación.} Las RAE tienen
            una precisión de exploración lineal mucho mayor que las VAE.
            \label{tab:linear_probing}
        ]
        {
            \begin{minipage}[t]{0.45\linewidth}{
                    \vspace{0pt}
                    \begin{center}
                        \tablestyle{4pt}{1.1}
                        \begin{tabular}{lc}
                            \toprule
                            Modelo & Top-1 Acc. \\
                            \midrule

                            \rowcolor{gray!20} DINOv2-B & 84.5 \\
                            \rowcolor{white}
                            SigLIP2-B & 79.1 \\
                            MAE-B & 68.0 \\
                            \midrule
                            SD-VAE & \gc{8.0} \\
                            \bottomrule
                        \end{tabular}
                    \end{center}
                }
            \end{minipage}
        }
    \end{tabular}
}
% \vspace{-4pt}
\label{tab:rae_recon_results}
\end{table*}
Para complementar las métricas cuantitativas de la \Cref{tab:rae_recon_results}, la
\Cref{fig:rae_recon_qualitative} ofrece una comparación visual de
las reconstrucciones. Las reconstrucciones de RAE conservan las texturas de grano fino,
los límites de los objetos y los detalles semánticos con la misma fidelidad que el modelo de referencia SD-VAE,
lo cual coincide con sus mejores resultados en rFID.

\begin{figure}[t]
    \centering
    \includegraphics[width=\textwidth]{chapters/chapter8/figs/recon_grid.pdf}
    \caption{\small\textbf{Ejemplos de reconstrucción RAE.} De izquierda a derecha: imagen de entrada, RAE (DINOv2-B), RAE (Siglp-B), RAE (MAE-B), SD-VAE. Amplíe la imagen para ver los detalles.}
    \label{fig:rae_recon_qualitative}
\end{figure}

% \yima{Besides the quantitative results, please have some
% illustrations to show qualitative improvements... }

\subsection{Muestreo a partir de representaciones aprendidas mediante eliminación de ruido}

Una vez que hemos aprendido una representación consistente y estructurada de la distribución de las 
imágenes naturales, por ejemplo, mediante VAE o RAE, normalmente nos interesa generar de manera 
eficiente muestras $\z \in \cZ$ a partir de la distribución de la representación aprendida, de modo 
que podamos regenerar las imágenes naturales asociadas a estas muestras mediante el decodificador 
aprendido $\hat{\x} = g(\z) \in \cX$. Este proceso de auto-codificación se ilustra en la mitad 
superior de la Figura \ref{fig:mental-picture} del Capítulo \ref{ch:intro}.

Para conocer y por lo tanto muestrear la distribución de las características $\z$ en el espacio 
latente $\cZ$, podemos recurrir al método de eliminación de ruido y difusión presentado en el 
Capítulo \ref{ch:denoising}. Concretamente, supongamos que $p(\z)$ es la distribución de las 
características aprendidas $\z$ en el espacio de características. Podemos aprender un proceso 
de eliminación de ruido que, partiendo de una distribución gaussiana estándar inicial $\vg \sim \cN(0, \vI)$, 
converja hacia esta distribución $p(\z)$. Siguiendo el método presentado en el Capítulo \ref{ch:denoising}, 
queremos aprender un filtro de ruido $\bar{\z}$ tal que, de manera iterativa, elimine el ruido de una 
muestra de la distribución inicial $\cN(0, \vI)$ hasta la distribución objetivo de las representaciones 
aprendidas $p(\z)$. Este proceso se ilustra en la mitad inferior de la Figura \ref{fig:mental-picture} 
del Capítulo \ref{ch:intro}.

Para entender el proceso de eliminación de ruido, tomemos el caso sencillo de la coincidencia de flujos 
presentado en el Capítulo \ref{ch:denoising}. Podemos definir el filtro de ruidp como un campo vectorial 
dependiente del tiempo $\bar{\z}_{\theta}(t, \z_t)$, parametrizado por $\theta$, que transforma una muestra de la
distribución inicial a la distribución objetivo a lo largo del tiempo $t \in [0, 1]$. 
Aquí, $\z_t$ es una interpolación entre $\z_0$ y $\z_1$ en el tiempo $t$:
\begin{equation}
\z_t = (1 - t)\z_0 + t\z_1,
\end{equation}
donde $\z_0 \sim p(\z)$ es una muestra de la distribución de la representación aprendida
y $\z_1 \sim \cN(0, \vI)$.
El objetivo de entrenamiento para el filtro de ruido puede expresarse entonces como:
\begin{equation}
\min_{\theta} \cL_{\mathrm{FM}}(\theta) = \Ex_{\z_0 \sim p(\z), \z_1
\sim \cN(0, I), t \sim \mathrm{Unif}(0, 1)} \left[
\norm{\bar{\z}_{\theta}(t,\z_t) - (\z_1 - \z_0)}_2^2 \right].
\end{equation}

Por lo general, el campo vectorial de eliminación de ruido $\bar{\z}_{\theta}(t, \z_t)$ 
se modela mediante una red profunda. Las opciones más comunes para la red profunda son 
una U-Net o un transformer de imágenes por difusión (DiT). A continuación ofrecemos 
más detalles sobre cómo se pueden implementar en la práctica.

% %more detail about the UNet, and refer back to PCA. Add the references.
% \subsubsection{Realizing Denoising with a U-Net.}
% A common architecture for implementing the parametric denoiser
% $\bar{\z}_{\theta}(t, \z_t)$ is the \textit{U-Net}, an encoder-decoder 
% network with skip connections. The encoder progressively downsamples the 
% noisy input to extract multi-scale features, while the decoder upsamples 
% and refines predictions. Skip connections concatenate encoder features to 
% corresponding decoder layers, preserving spatial detail. To capture 
% long-range dependencies, self-attention layers are often inserted at 
% lower resolutions. The convolutional structure provides strong spatial 
% inductive biases, which can be beneficial when training data is limited.

% The timestep $t$ is typically embedded via 
% sinusoidal functions (analogous to positional encodings in transformers) 
% and injected through \textit{adaptive group normalization}: the scale and 
% shift parameters of normalization layers are regressed from the timestep 
% embedding. For class-conditional generation, a class embedding is simply 
% added to the timestep embedding before injection. For text conditioning, 
% a pretrained text encoder (e.g., CLIP) produces a sequence of token 
% embeddings, which are incorporated via \textit{cross-attention}: at 
% selected layers, the image features serve as queries while the text 
% embeddings serve as keys and values. This allows the network to attend to 
% relevant parts of the text prompt when denoising each spatial location.

\subsection{Impementación de la eliminación de ruido con una red U-Net}
\label{sec:UNet-implement}

La arquitectura predominante para implementar la red de eliminación de ruido $\boldsymbol{\epsilon}_\theta$ 
en los modelos de difusión ha sido la \textit{U-Net}~\citep{ronneberger2015unet}, una red convolucional de 
codificador-decodificador con conexiones de salto desarrollada originalmente para la segmentación de 
imágenes biomédicas. Ho et al.~\citep{ho2020denoising} fueron los primeros en adoptar la red U-Net como 
estructura principal de los modelos probabilísticos de difusión con eliminación de ruido (DDPM), 
consolidándola como la arquitectura predeterminada en este campo. Posteriormente, Dhariwal 
y Nichol~\citep{dhariwal2021diffusion} llevaron a cabo una serie sistemática de ablaciones arquitectónicas 
que mejoraron sustancialmente el diseño de U-Net, dando lugar a lo que denominan 
el \emph{Modelo de Difusión Ablacionado} (ADM). Su modelo mejorado fue el primer modelo de difusión 
en superar a las GAN en cuanto a calidad de las muestras en ImageNet. Partiendo de estas bases, 
Rombach et al.~\citep{rombach2022high} trasladaron la U-Net del espacio de píxeles al espacio latente 
de una VAE preentrenada (Sección~\ref{sec:VAE-implement}) e introdujeron capas de atención cruzada para 
un condicionamiento multimodal flexible, dando lugar al \emph{Modelo de Difusión Latente} (LDM) que 
sustenta Stable Diffusion. En esta sección, repasamos la evolución de la estructura central de U-Net 
a lo largo de estas tres etapas, antes de analizar las limitaciones arquitectónicas que han motivado 
la reciente transición hacia alternativas basadas en transformers (Sección~\ref{sec:DiT-implement}).

\paragraph{Estructura de codificador-decodificador con conexiones de salto.}
U-Net sigue una estructura de codificador-decodificador. El codificador reduce progresivamente la 
resolución de la entrada a través de una serie de etapas de resolución, extrayendo características 
a múltiples escalas espaciales; luego, el decodificador vuelve a aumentar la resolución de estas 
características hasta la resolución original para producir la predicción de salida. La característica 
distintiva de la red U-Net son sus \emph{conexiones de salto}: en cada nivel de resolución, el mapa de 
características del codificador se concatena por canales con el mapa de características correspondiente 
del decodificador. Esto permite al decodificador combinar características generales y semánticamente 
ricas procedentes del cuello de botella con detalles espaciales de mayor precisión procedentes de capas 
anteriores del codificador, lo cual es esencial para generar predicciones espacialmente precisas.

En la red U-Net DDPM~\citep{ho2020denoising}, la arquitectura comienza con una convolución de entrada 
que proyecta la entrada con ruido en un espacio de características con un número básico de canales 
(por ejemplo, 128). A continuación, el codificador pasa por múltiples etapas de resolución---cuatro 
para entradas de $32 \times 32$, o seis para entradas de $256 \times 256$---donde cada etapa consta 
de dos bloques residuales convolucionales. Entre etapas consecutivas se realiza una reducción espacial 
mediante convoluciones con paso, resultando en tres reducciones y un factor de submuestreo 
global $f = 2^3 = 8$. Las capas de auto-atención se insertan en el nivel de resolución de 
$16 \times 16$ entre los bloques convolucionales, lo que 
proporciona un modelado de largo alcance limitado a baja resolución. En el cuello de botella (la 
resolución más baja), un bloque intermedio con bloques residuales adicionales y auto-atención refina 
aún más las características. El decodificador es un reflejo del codificador: bloques residuales en 
cada nivel de resolución seguidos de un remuestreo mediante convoluciones transpuestas, con 
conexiones de salto que concatenan las activaciones correspondientes del codificador. La 
arquitectura utiliza la normalización por grupos~\citep{wu2018group} en todo momento, 
sustituyendo la normalización de pesos en la que se basa el diseño.

Cuando U-Net opera en el espacio de píxeles, su entrada y su salida son, respectivamente, la 
imagen con ruido $\x_t \in \mathbb{R}^{H \times W \times 3}$ y el ruido predicho 
$\hat{\boldsymbol{\epsilon}} \in \mathbb{R}^{H \times W \times 3}$. Cuando se opera en 
el espacio latente de una VAE preentrenada (como en LDM), la entrada se convierte 
en $\z_t \in \mathbb{R}^{h \times w \times c}$ y la salida es 
$\hat{\boldsymbol{\epsilon}} \in \mathbb{R}^{h \times w \times c}$, con $h = H/f$, $w = W/f$ 
tal como se define en la Sección~\ref{sec:VAE-implement}.

\paragraph{Condicionamiento del paso de tiempo mediante la normalización adaptativa por grupos.}
La red de eliminación de ruido debe condicionarse al paso de tiempo de difusión $t$ para que 
pueda adaptar su comportamiento a los distintos niveles de ruido. Siguiendo el esquema de 
codificación posicional del Transformer~\citep{vaswani2017attention}, el paso de tiempo $t$ 
se asigna primero a una incrustación de frecuencia sinusoidal:
\begin{equation}
    \mathrm{Embed}(t)_{2i} = \sin\!\left(\frac{t}{10000^{2i/d_e}}\right), \quad
    \mathrm{Embed}(t)_{2i+1} = \cos\!\left(\frac{t}{10000^{2i/d_e}}\right),
\end{equation}
donde $d_e$ es la dimensión de la incrustación e $i$ indexa sus entradas. A continuación, esta 
incrustación se pasa a través de una pequeña red neuronal multicapas (MLP) (normalmente dos capas 
lineales con una función no lineal) para generar un vector de paso de 
tiempo $\boldsymbol{e}_t \in \mathbb{R}^{d_e}$.

En lugar de limitarse a concatenar $\boldsymbol{e}_t$ a la entrada, la red U-Net inyecta la 
información del paso de tiempo en cada bloque residual mediante la \emph{normalización 
adaptativa por grupos} (AdaGN)~\citep{dhariwal2021diffusion}. Después de la primera convolución 
dentro de cada bloque residual, se aplica una operación de normalización de grupo, seguida de 
una modulación afín cuyos parámetros de escala y desplazamiento se regresan a partir de la 
incrustación del paso de tiempo:
\begin{equation}
    \mathrm{AdaGN}(\boldsymbol{h}, \boldsymbol{e}_t) = \boldsymbol{y}_s \hada \mathrm{GroupNorm}(\boldsymbol{h}) + \boldsymbol{y}_b,
    \label{eq:adagn}
\end{equation}
donde $\boldsymbol{h}$ denota el mapa de características intermedio, 
$[\boldsymbol{y}_s, \boldsymbol{y}_b] = \mathrm{Linear}(\boldsymbol{e}_t)$ 
son los vectores de escala y desplazamiento obtenidos a partir de una proyección lineal 
de la incrustación del paso de tiempo y $\hada$ denota la multiplicación elemento por 
elemento. Este mecanismo es análogo a la normalización adaptativa de 
instancias~\citep{dumoulin2017learned} y al condicionamiento FiLM~\citep{perez2018film}, 
y permite que cada capa modifique sus activaciones en función del nivel de ruido actual. 
Los experimentos de ablación en~\citep{dhariwal2021diffusion} muestran que la eliminación 
de AdaGN reduce el FID en más de 2 puntos, lo que confirma su importancia.

Para la generación condicionada por clase, simplemente se añade una incrustación de clase 
$\mathrm{Embed}(c)$ (una tabla de consulta aprendida) a la incrustación del paso temporal 
antes de la inyección, de modo que $\boldsymbol{e}_t$ se reemplaza por 
$\boldsymbol{e}_t + \mathrm{Embed}(c)$ en la ecuación ~\eqref{eq:adagn}. Este mecanismo es 
ligero, pero se limita a etiquetas categóricas; las modalidades de condicionamiento más 
ricas, como el texto, requieren el mecanismo de atención cruzada introducido en 
la Sección~\ref{sec:unet-ldm}.

\subsubsection{Mejoras en la arquitectura: de DDPM a ADM}
\label{sec:unet-adm}

Dhariwal y Nichol~\citep{dhariwal2021diffusion} llevaron a cabo un estudio sistemático de 
ablación para determinar qué modificaciones arquitectónicas mejoran más la calidad de las 
muestras. Partiendo del modelo de referencia DDPM U-Net~\citep{ho2020denoising}, exploraron 
los siguientes cambios en ImageNet $128 \times 128$: (a) ampliar la auto-atención de la 
resolución $16 \times 16$ a tres resoluciones ($32 \times 32$, $16 \times 16$, $8 \times 8$); 
(b) aumentar el número de cabezales de atención, utilizando 64 canales por cabezal para 
ajustarse mejor a la convención de Transformer; (c) la adopción del bloque residual 
BigGAN~\citep{brock2019large} tanto para el upsampling como para el downsampling, lo que 
mejora el flujo de gradientes; y (d) la inclusión por defecto de la capa de condicionamiento 
AdaGN (descrita anteriormente) en cada bloque residual. Todos estos cambios mejoran el FID y 
sus efectos se potencian de manera positiva. El modelo resultante, al que denominan ADM, 
sirve como nuevo diseño predeterminado de U-Net para trabajos posteriores.

Además de las mejoras arquitectónicas, Dhariwal y Nichol introdujeron la \emph{orientación 
por clasificador} como mecanismo para equilibrar la diversidad de la muestra con la fidelidad. 
Se entrena un clasificador $p_\phi(c|\z_t, t)$ con imágenes ruidosas, y sus gradientes se utilizan 
para dirigir la trayectoria de muestreo hacia regiones de alta probabilidad de la clase deseada:
\begin{equation}
    \hat{\boldsymbol{\mu}}_\theta(\z_t, t) = \boldsymbol{\mu}_\theta(\z_t, t) + s \, \boldsymbol{\Sigma}_\theta(\z_t, t) \, \nabla_{\z_t} \log p_\phi(c|\z_t),
\end{equation}
donde $s > 1$ es una escala de guía que agudiza la distribución efectiva de clases. Con esta 
técnica, ADM alcanzó un FID récord del estado del arte de 4,59 en ImageNet $256 \times 256$ condicional a 
la clase, superando por primera vez el referente de GAN que se había mantenido durante mucho 
tiempo, BigGAN-deep (FID 6,95). El modelo guiado (ADM-G) mejoró aún más el FID hasta 3,94 
cuando se combinó con una pila de supermuestreo (ADM-U).

El principal inconveniente de la guía basada en clasificadores es que requiere entrenar un 
clasificador independiente sobre imágenes con ruido. Esta limitación se resolvió posteriormente 
mediante la \emph{guía sin clasificadores}~\citep{ho2022classifier}, que prescinde por completo 
del clasificador externo al aprender en forma conjunta la eliminación de ruido condicional e 
incondicional dentro de un único modelo, una técnica que describimos en detalle en 
la Sección~\ref{sec:DiT-implement}.

\subsubsection{La red U-Net en la difusión latente (Stable Diffusion)}
\label{sec:unet-ldm}

Mientras que ADM opera directamente en el espacio de píxeles, Rombach et al.~\citep{rombach2022high} 
observaron que gran parte de la carga computacional se dedica a modelar detalles de alta frecuencia 
que son irrelevantes desde el punto de vista perceptivo. Su idea clave es separar 
la \emph{compresión perceptual} del \emph{modelado generativo}: un auto-codificador preentrenado 
(Sección~\ref{sec:VAE-implement}) primero comprime la imagen en un espacio latente de baja dimensión 
y luego el modelo de difusión opera completamente dentro de este espacio latente. Para una imagen de 
$256 \times 256$ con un factor de submuestreo $f = 8$, U-Net procesa un tensor latente de 
$32 \times 32 \times 4$ en lugar de un tensor de píxeles de $256 \times 256 \times 3$, lo que supone 
una reducción de aproximadamente $48\times$ en los elementos espaciales. Esto produce una aceleración 
de al menos 2,7 veces tanto en el rendimiento del entrenamiento como en el del muestreo, al tiempo 
que mejora el FID.

Más allá del ahorro computacional, el marco de difusión latente introduce un \emph{mecanismo de 
condicionamiento de atención cruzada} de uso general que permite a la red U-Net condicionarse en 
función de modalidades arbitrarias, como indicaciones de texto, mapas semánticos o cajas 
delimitadoras. En resumen, un codificador específico del dominio $\tau_\theta$ mapea la entrada 
de condicionamiento $\boldsymbol{y}$ a una representación intermedia 
$\tau_\theta(\boldsymbol{y}) \in \mathbb{R}^{M \times d_\tau}$, que luego se inyecta en la U-Net 
en múltiples niveles de resolución a través de capas de atención cruzada: las características de 
la imagen sirven como consultas, mientras que las incrustaciones de condicionamiento sirven como 
claves y valores, lo que permite que cada ubicación espacial preste atención a las partes relevantes 
de la señal de condicionamiento. Dejamos para la Sección~\ref{subsec:basic_realization} la descripción 
matemática completa de este mecanismo de atención cruzada y del objetivo de entrenamiento condicional.

En Stable Diffusion~\citep{rombach2022high}, el codificador de texto $\tau_\theta$ se implementa como 
el modelo CLIP ViT-L/14 preentrenado~\citep{radford2021learning}, que mapea las indicaciones de texto 
a una secuencia de incrustación de tokens. El propio U-Net conserva las mejoras de ADM---bloques 
residuales de BigGAN, atención en tres niveles de resolución, AdaGN para la inyección de pasos 
temporales---y, además, intercala capas de atención cruzada en estos mismos niveles de resolución 
para el condicionamiento de texto. El modelo completo contiene aproximadamente 860 millones de 
parámetros en la red U-Net y 123 millones en el codificador de texto.

\subsubsection{Limitaciones de la arquitectura U-Net}
\label{sec:unet-limitations}

A pesar de su éxito, la red U-Net presenta varias limitaciones que han motivado el desarrollo 
de estructuras alternativas. En primer lugar, \textbf{escalar} una red U-Net resulta complicado: 
su estructura de codificador-decodificador implica numerosas decisiones de diseño interrelacionadas 
(multiplicadores de canales, ubicación de la atención, tipos de bloques de muestreo 
ascendente/descendente), y no está claro cómo asignar de manera óptima la capacidad de cálculo 
adicional. Por el contrario, el Transformer básico se escala a lo largo de dos ejes simples, la 
profundidad y la anchura, y presenta leyes de escalado bien caracterizadas según estudios 
empíricos~\citep{kaplan2020scaling}. En segundo lugar, las capas convolucionales de U-Net 
codifican un fuerte \textbf{sesgo inductivo espacial} (localidad, invariancia a la traslación) 
que resulta útil en entornos con pocos datos, pero que limita el modelado de largo alcance en 
conjuntos de datos de gran tamaño. La auto-atención resuelve parcialmente este problema, pero en 
los diseños estándar de U-Net solo se aplica en las etapas de baja resolución debido a su coste 
cuadrático, lo que hace que las características de alta resolución sean totalmente locales. En 
tercer lugar, U-Net es una arquitectura \textbf{específica para imágenes} que no se generaliza 
fácilmente a otras modalidades. El modelo Transformer, cuya eficacia ha quedado demostrada en 
entornos lingüísticos, de visión, de audio y multimodales~\citep{vaswani2017attention}, ofrece una 
estructura central unificada capaz de aprovechar los conocimientos sobre escalabilidad entre dominios.

Estas consideraciones dieron lugar al Diffusion Transformer (DiT)~\citep{peebles2023scalable}, que 
sustituye por completo a la red U-Net por un Transformer básico que opera sobre parches latentes. 
Tal como describimos en la Sección~\ref{sec:DiT-implement}, DiT iguala y supera a los mejores modelos 
basados en U-Net, al tiempo que presenta un comportamiento de escalabilidad más limpio (correlación 
Gflops-FID de $-0,93$) y una arquitectura más simple.

% % refer the transformer arch back to previous Chapters. Add the references.
% \subsubsection{Realizing Denoising with a DiT.}
% An alternative to the U-Net is the \textit{Diffusion Transformer (DiT)}, 
% which replaces convolutions entirely with transformer blocks. Following 
% the Vision Transformer, the input is divided into patches and linearly 
% embedded as token sequences. Unlike the U-Net, the DiT has no explicit 
% spatial inductive bias---all structure is learned from data. This 
% generality enables the same architecture to handle images, video, and 
% other modalities with minimal modification. Perhaps most importantly, DiT 
% follows favorable scaling laws: larger models and more compute yield 
% predictable improvements in sample quality.

% Timestep and class conditioning in DiT is 
% achieved through \textit{adaptive layer normalization} (adaLN-Zero): the 
% scale and shift parameters of each layer normalization are regressed from 
% the conditioning embedding (formed by summing timestep and class 
% embeddings), with these parameters initialized near zero for training 
% stability. For text conditioning, two approaches are common: (1) 
% cross-attention, as in the U-Net, where text token embeddings from a 
% frozen encoder are attended to by the image tokens; or (2) prepending the 
% text tokens directly to the image token sequence, allowing the 
% transformer's self-attention to jointly process both modalities.



% Once we have trained the denoiser, we can use it to generate new
% samples from the learned representation distribution via the denoising
% process. The generated feature samples can then be decoded back to
% the image space with the corresponding decoder $g$.

\subsection{Implementación de la eliminación de ruido con un transformer de difusión}
\label{sec:DiT-implement}

Una alternativa a U-Net para implementar la red de eliminación de ruido 
$\boldsymbol{\epsilon}_\theta$ es el \textit{Diffusion Transformer (DiT)}~\citep{peebles2023scalable}, 
que sustituye por completo las convoluciones por bloques de transformer. La conclusión clave es que 
el sesgo inductivo espacial de U-Net \emph{no es crucial} para el rendimiento de los modelos de 
difusión: un transformer estándar que opera sobre secuencias de parches latentes puede igualar e 
incluso superar a U-Net, al tiempo que hereda las propiedades de escalabilidad favorables de la 
clase de arquitectura de transformers~\citep{vaswani2017attention}. A diferencia de U-Net, DiT no 
tiene un sesgo inductivo espacial explícito: toda la estructura espacial se aprende a partir de los 
datos. Esta generalidad permite que la misma arquitectura pueda procesar imágenes, vídeos y otras 
modalidades con modificaciones mínimas.

DiT opera dentro del espacio latente de un VAE preentrenado (Sección~\ref{sec:VAE-implement}), lo que 
convierte al proceso completo de generación de imágenes en una arquitectura híbrida: un VAE 
convolucional para la compresión perceptual y un DDPM basado en transformers para el modelado 
generativo. Siguiendo el Vision Transformer (ViT)~\citep{dosovitskiy2020image}, el latente de 
entrada se divide en parches y se incrusta linealmente como una secuencia de tokens. A continuación 
describimos la arquitectura en detalle.

\paragraph{División en parches y la codificación posicional.}
La entrada a DiT es una variable latente con ruido $\z_t \in \mathbb{R}^{h \times w \times c}$, 
donde $h = H/f$ y $w = W/f$ son las dimensiones espaciales del espacio latente definido por el 
codificador VAE (por ejemplo, $32 \times 32 \times 4$ para imágenes de $256 \times 256$ con 
$f = 8$, $c = 4$). La primera capa, llamada \emph{patchify}, divide $\z_t$ en una cuadrícula de 
parches no superpuestos de tamaño $p \times p$ e incrusta linealmente cada parche en una 
representación de tokens de $d$ dimensiones. Esto produce una secuencia de
\begin{equation}
    T = \left(\frac{h}{p}\right) \times \left(\frac{w}{p}\right)
\end{equation}
tokens en $\mathbb{R}^d$. Se añaden incrustaciones posicionales sinusoidales estándar (la variante 
seno-coseno de ViT) a cada token para codificar la posición espacial. El tamaño del parche $p$ es 
un parámetro de diseño clave: reducir $p$ a la mitad cuadruplica $T$ y, por lo tanto, cuadruplica 
al menos el total de Gflops del transformer, mientras que el número de parámetros permanece 
prácticamente sin cambios. En~\citep{peebles2023scalable}, el espacio de diseño incluye 
$p \in \{2, 4, 8\}$; la configuración más fina $p = 2$ produce la mejor calidad de generación.

\paragraph{Bloque DiT con adaLN-Zero.}
Tras la «patchificación», la secuencia de tokens se procesa mediante una pila de $N$ bloques de 
transformer. Cada bloque sigue la estructura estándar del transformer~\citep{vaswani2017attention}: 
normalización de capa, auto-atención multicanal (MHSA), una segunda normalización de capa y una red 
de propagación hacia adelante punto a punto (MLP con activaciones GELU).

La cuestión central del diseño es cómo incorporar la información de condicionamiento---el intervalo 
de tiempo de difusión $t$ y la etiqueta de clase $c$---en cada bloque. Peebles y Xie~\citep{peebles2023scalable} 
comparan cuatro estrategias: (1) \emph{condicionamiento en contexto}, que añade $t$ y $c$ como tokens 
adicionales en la secuencia de entrada; (2) \emph{atención cruzada}, que introduce una capa adicional 
de atención cruzada multicanal que presta atención a las incrustaciones de condicionamiento; (3) 
\emph{normalización adaptativa de capas (adaLN)}, que realiza una regresión de los parámetros de escala 
y desplazamiento de cada normalización de capa a partir de la señal de condicionamiento; y (4) 
\emph{adaLN-Zero}, que amplía adaLN con un mecanismo de activación adicional inicializado en cero. 
De entre ellas, adaLN-Zero alcanza el FID más bajo y es la más eficiente en términos de computación, 
por lo que se adopta como opción predeterminada.

En el bloque adaLN-Zero, la información de condicionamiento se forma primero sumando las representaciones del paso temporal y la etiqueta de clase:
\begin{equation}
    \boldsymbol{e} = \mathrm{Embed}(t) + \mathrm{Embed}(c),
\end{equation}
donde $\mathrm{Embed}(t)$ se obtiene pasando una incrustación de frecuencia sinusoidal 
de $t$ a través de una red neuronal multicapas (MLP) de dos capas con activaciones SiLU 
y $\mathrm{Embed}(c)$ es una incrustación de clase aprendida. A partir de $\boldsymbol{e}$, 
una capa lineal realiza una regresión de seis conjuntos de parámetros dimensionales: 
$\boldsymbol{\gamma}_1, \boldsymbol{\beta}_1, \boldsymbol{\alpha}_1$ para el subbloque de 
auto-atención y $\boldsymbol{\gamma}_2, \boldsymbol{\beta}_2, \boldsymbol{\alpha}_2$ para 
el subbloque de alimentación directa. La normalización de la capa adaptativa se aplica de 
la siguiente manera:
\begin{equation}
    \mathrm{adaLN}(\boldsymbol{h}, \boldsymbol{\gamma}, \boldsymbol{\beta}) = \boldsymbol{\gamma} \hada \mathrm{LayerNorm}(\boldsymbol{h}) + \boldsymbol{\beta},
\end{equation}
donde $\hada$ denota la multiplicación elemento por elemento. Los parámetros $\boldsymbol{\alpha}_1$ 
y $\boldsymbol{\alpha}_2$ actúan como \emph{factores de control} por dimensión que se aplican 
inmediatamente antes de cada conexión residual:
\begin{equation}
\begin{aligned}
    \boldsymbol{h}' &= \boldsymbol{h} + \boldsymbol{\alpha}_1 \hada \mathrm{MHSA}\!\big(\mathrm{adaLN}(\boldsymbol{h}, \boldsymbol{\gamma}_1, \boldsymbol{\beta}_1)\big), \\
    \boldsymbol{h}'' &= \boldsymbol{h}' + \boldsymbol{\alpha}_2 \hada \mathrm{MLP}\!\big(\mathrm{adaLN}(\boldsymbol{h}', \boldsymbol{\gamma}_2, \boldsymbol{\beta}_2)\big).
\end{aligned}
\end{equation}
El «Zero» de adaLN-Zero hace referencia a la inicialización: todos los parámetros 
$\boldsymbol{\alpha}$ son inicializados al valor cero por la MLP que los regresa, de modo 
que cada bloque DiT actúa como la \emph{función identidad} al inicio del entrenamiento. 
Esta estrategia de inicialización a cero, inspirada en prácticas similares en 
ResNets~\citep{goyal2017accurate} y U-Nets de difusión, es importante para la estabilidad 
del entrenamiento y produce un FID significativamente más bajo que el adaLN básico.

Observamos que adaLN aplica la \emph{misma} función de condicionamiento a todos los tokens 
de manera uniforme. Para la generación de texto a imagen, donde la señal de condicionamiento 
es una secuencia de tokens de texto de longitud variable, son comunes dos enfoques alternativos: 
(1) la atención cruzada, donde los tokens de imagen prestan atención a las incrustaciones de 
tokens de texto de un codificador congelado; o (2) anteponer los tokens de texto directamente a 
la secuencia de tokens de imagen, lo que permite que la auto-atención del transformer procese 
conjuntamente ambas modalidades.

\paragraph{Proyección de la salida.}
Tras el último bloque DiT, el modelo debe decodificar la secuencia de representaciones de 
tokens para convertirla de nuevo en una predicción espacial. Se aplica una normalización de 
capa adaptativa final (condicionada por $\boldsymbol{e}$), seguida de una capa lineal que 
proyecta cada token en un tensor en forma $p \times p \times 2c$. El factor de $2c$ da 
cuenta de las dos magnitudes que predice el modelo: la estimación del ruido 
$\hat{\boldsymbol{\epsilon}} \in \mathbb{R}^{h \times w \times c}$ y la covarianza diagonal 
$\boldsymbol{\Sigma}_\theta \in \mathbb{R}^{h \times w \times c}$. A continuación, los tokens 
decodificados se reorganizan (se les elimina el parchado) para volver a su disposición 
espacial original de $\mathbb{R}^{h \times w \times c}$ para cada salida.

\subsubsection{Entrenamiento}
El DiT se entrena como un modelo de difusión latente condicional a la clase en el conjunto 
de datos ImageNet~\citep{deng2009imagenet} con resoluciones de imagen de $256 \times 256$ y 
$512 \times 512$. El proceso de difusión opera sobre las representaciones latentes generadas 
por el VAE preentrenado de la Sección~\ref{sec:VAE-implement}: para imágenes de 
$256 \times 256$, la entrada al DiT es $\z_t \in \mathbb{R}^{32 \times 32 \times 4}$; para 
imágenes de $512 \times 512$, es $\z_t \in \mathbb{R}^{64 \times 64 \times 4}$.

El objetivo de entrenamiento combina dos términos. La pérdida de predicción de ruido entrena 
$\boldsymbol{\epsilon}_\theta$ mediante el error cuadrático medio simple:
\begin{equation}
    \mathcal{L}_{\mathrm{simple}}(\theta) = \mathbb{E}_{\z_0, \boldsymbol{\epsilon}, t}\left[\|\boldsymbol{\epsilon}_\theta(\z_t, t) - \boldsymbol{\epsilon}\|_2^2\right],
\end{equation}
donde $\z_0$ es la variable latente sin ruido, $\boldsymbol{\epsilon} \sim \mathcal{N}(\boldsymbol{0}, 
\boldsymbol{I})$ es el ruido muestreado, y 
$\z_t = \sqrt{\bar{\alpha}_t}\,\z_0 + \sqrt {1 - \bar{\alpha}_t}\,\boldsymbol{\epsilon}$ es la latente 
con ruido en el paso de tiempo $t$. La covarianza aprendida $\boldsymbol{\Sigma}_\theta$ se entrena 
adicionalmente con el término completo del límite inferior variacional $D_{KL}$, siguiendo el enfoque 
de Nichol y Dhariwal~\citep{nichol2021improved}. Se utiliza un programa de varianza lineal con 
$t_{\max} = 1000$ pasos, que va desde $1 \times 10^{-4}$ hasta $2 \times 10^{-2}$.

Todos los modelos se entrenan con AdamW~\citep{loshchilov2019decoupled} utilizando una tasa de 
aprendizaje constante de $1 \times 10^{-4}$, sin decaimiento de pesos y un tamaño de lote de $256$. 
El aumento de datos consiste únicamente en volteos horizontales aleatorios. Se mantiene una media 
móvil exponencial (EMA) de los pesos del modelo con una decaimiento de $0.9999$; todos los resultados 
reportados utilizan el modelo EMA. Cabe destacar que, a diferencia del entrenamiento supervisado de ViT, 
DiT \emph{no} requiere un período de calentamiento de la tasa de aprendizaje ni \emph{ninguna} 
regularización: el entrenamiento es muy estable en todas las configuraciones del modelo, sin que 
se observen picos de pérdida.

\subsubsection{Guía sin clasificador}
Para generar un muestreo controlado por clases, DiT puede emplear fácilmente la \emph{orientación 
sin clasificador}~\citep{ho2022classifier}. Durante el entrenamiento, la etiqueta de clase $c$ se 
omite aleatoriamente con cierta probabilidad y se reemplaza por una incrustación nula aprendida 
$\varnothing$, de modo que el modelo aprende tanto las predicciones de ruido condicionales 
$\boldsymbol{\epsilon}_\theta(\z_t, c)$ como las incondicionales 
$\boldsymbol{\epsilon}_\theta(\z_t, \varnothing)$. En el momento de la inferencia, la 
predicción guiada de la perturbación se calcula como:
\begin{equation}
    \hat{\boldsymbol{\epsilon}}_\theta(\z_t, c) = \boldsymbol{\epsilon}_\theta(\z_t, \varnothing) + s \cdot \big(\boldsymbol{\epsilon}_\theta(\z_t, c) - \boldsymbol{\epsilon}_\theta(\z_t, \varnothing)\big),
    \label{eq:cfg}
\end{equation}
donde $s > 1$ es la escala de guía ($s = 1$ corresponde al muestreo estándar sin guía). 
La justificación se deriva de la regla de Bayes: dado que 
$\nabla_{\z_t} \log p(c | \z_t) \propto \nabla_{\z_t} \log p(\z_t | c) - \nabla_{\z_t} \log p(\z_t)$, 
la estimación guiada en la ecuación ~\eqref{eq:cfg} dirige efectivamente la trayectoria de 
muestreo hacia regiones de alta verosimilitud condicional $p(c | \z_t)$. Esta técnica tiene 
un efecto notable en la calidad de la generación: DiT-XL/2 alcanza un FID de 9,62 sin orientación, 
que mejora a 2,27 con una escala de orientación de s = 1,5, superando a todos los modelos de difusión 
anteriores en el banco de pruebas ImageNet de 256 × 256 con condiciones de clase.

\subsubsection{Generación}
Una vez entrenada la red de eliminación de ruido $\boldsymbol{\epsilon}_\theta$ (es decir, la DiT), 
se pueden generar nuevas imágenes ejecutando el proceso inverso aprendido. Partiendo de un ruido 
puro $\z_{t_{\max}} \sim \mathcal{N}(\boldsymbol{0}, \boldsymbol{I})$, el modelo elimina el ruido 
de manera iterativa mediante el muestreo de $\z_{t-1} \sim p_\theta(\z_{t-1} | \z_t)$ en cada paso, 
utilizando el ruido predicho $\hat{\boldsymbol{\epsilon}}_\theta$ y la covarianza $\boldsymbol{\Sigma}_\theta$ 
para parametrizar la transición inversa. Tras $t_{\max}$ pasos de eliminación de ruido, la variable 
latente limpia resultante $\z_0$ se decodifica de nuevo al espacio de la imagen mediante el decodificador 
VAE: $\tilde{\x} = \mathcal{D}(\z_0) \in \mathbb{R}^{H \times W \times 3}$.


\section{Representación y generación condicionada de imágenes}
\label{sec:cond_gen}
%\yima{From the above latent diffusion model for images, it is natural to transition to conditioned image generation. Text condition via CLIP or SimCLIP... and DPS implementation, classifier-free, or cross-attention etc. (Robin, Jingfeng, Daniel, and Peter, decide what to feature...)}

% --- draft by Daniel Yeh ---
En las secciones anteriores, establecimos los fundamentos
del modelado generativo: una arquitectura de «autoencoder» para codificar los datos $\x$ en una variable latente $\z$ (comprimida y más estructurada):
\begin{equation}
\x\xrightarrow{\hspace{1mm} \mathcal{E} \hspace{1mm}} \z
\xrightarrow{\hspace{1mm} \mathcal{D} \hspace{1mm}} \hat{\x}.
\end{equation}

A continuación, se aprende un proceso generativo para modelar la distribución de estas variables 
latentes $p(\z)$. En esta sección, nos enfocamos en el \textit{control} condicional en lugar de la 
\textit{generación incondicional}. Suponemos que la representación $\z$ ya se ha aprendido (mediante los métodos 
detallados en el Capítulo \ref{ch:conditional-inference}); nuestra tarea específica es guiar 
el proceso de muestreo utilizando una señal de ``control'' $\vc$ proporcionada por el usuario. Es 
decir, nuestro objetivo es obtener muestras de la distribución condicional $p(\z | \vc)$. En la 
práctica, $\vc$ podría representar atributos parciales de la imagen de muestra deseada: una imagen 
enmascarada, una máscara semántica, información de clase o una leyenda que describa el contenido deseado.

\subsection{Formulación de la tarea}
\label{subsec:cond_problem}

Formalmente, el objetivo de la generación condicionada es aprender una correspondencia que transforme 
una distribución de origen simple (normalmente ruido gaussiano) en una distribución condicional de 
destino. La característica generada $\hat{\z}$ se decodifica posteriormente en una imagen correspondiente 
$\hat{\x}$. Tal y como se ilustra a continuación, el sistema se define formalmente por sus entradas y salidas:
\begin{enumerate}
    \item \textbf{Entrada estocástica ($\epsilon$):} Una variable de ruido muestreada a partir de una 
          distribución normal estándar previa $\epsilon \sim \mathcal{N}(\bm{0}, \I)$. Esto proporciona 
          la «semilla» para la diversidad, lo que garantiza que podamos generar múltiples variaciones del mismo concepto.
    \item \textbf{Señal de control ($\vc$):} Una entrada estructurada que define el contenido semántico o espacial deseado. 
          Puede tratarse de un texto descriptivo (por ejemplo, «un auto rojo»), una máscara de segmentación semántica o un diseño.
    \item \textbf{Salida ($\hat{\x}$):} Una muestra de datos de alta fidelidad $\hat{\x} = \mathcal{D}(\hat{\z})$ que se ajusta a 
          la intención del usuario $\vc$ sin salirse de la variedad (manifold) natural de imágenes.
\end{enumerate}
Utilizamos un \textbf{codificador de condiciones} específico del dominio, denotado como $\tau_\phi$, 
para proyectar la señal de control $\vc$ en un espacio de incrustación útil. Por lo tanto, el proceso 
general de generación controlada puede describirse formalmente como una composición de un generador 
$G_\theta$ y el decodificador $\mathcal{D}$:
\begin{equation}
    \underbrace{\epsilon \sim \mathcal{N}(\bm{0}, \I)}_{\text{Source}}, \vc 
    \xrightarrow{\quad G_\theta(\epsilon, \tau_\phi(\vc)) \quad} 
    \hat{\z} \sim p(\z\mid \vc) 
    \xrightarrow{\quad \mathcal{D}(\hat{\z}) \quad} 
    \hat{\x} \approx \x.
\end{equation}

\paragraph{La diferencia modal.}
Aunque la tarea es conceptualmente sencilla---establecer una correspondencia entre 
$(\epsilon, \vc)$ y $\x$---, requiere tender un puente entre espacios topológicos disjuntos, 
especialmente en la generación de imágenes a partir de texto. Los datos visuales $\x$ se 
encuentran en un espacio de píxeles continuo y de alta dimensión $\mathbb{R}^{H \times W \times 3}$, 
mientras que la condición $\vc$ suele residir en un espacio simbólico discreto (el lenguaje). 
No existe una métrica euclidiana natural que permita medir directamente la distancia $\| \x - \vc \|$.
Por lo tanto, la componente crítica $\tau_\phi(\vc)$ mencionada anteriormente debe alinear las 
modalidades. Tal y como se explica en el Capítulo \ref{ch:conditional-inference}, los enfoques 
modernos se basan en codificadores contrastivos preentrenados (como CLIP) o en grandes modelos 
de lenguaje (como T5~\citep{raffel2020exploring}) para mapear tokens de lenguaje discretos en un 
espacio de incrustación semántico continuo. La red troncal generativa $G_\theta$ (es decir, 
Denoising U-Net) debe entonces aprender a mapear la distribución del ruido a la distribución 
de los datos, guiada por estas incrustaciones semánticas a través de un mecanismo de fusión 
específico, como la atención cruzada.

\paragraph{El problema inverso uno-a-muchos.}
El aprendizaje por condiciones es intrínsecamente mal planteado porque la señal de control suele 
ser \textit{subdeterminada}. Una indicación textual $\vc=$ «un perro» proporciona información 
sobre la categoría semántica, pero omite la postura del perro, la iluminación, el fondo y el 
patrón del pelaje. En consecuencia, la correspondencia inversa $\vc \to \x$ no es un único punto, 
sino una distribución multimodal compleja. Si utilizáramos un modelo de regresión determinista 
simple para minimizar el error euclidiano $\|\x - G_\theta(\vc)\|^2$, el modelo convergería hacia 
la media condicional $\mathbb{E}[\x|\vc]$---un promedio de todos los perros posibles---, lo que 
daría lugar a un resultado borroso y poco realista. Para generar resultados de alta fidelidad, 
podemos utilizar marcos generativos (como U-Net o DiT, mencionados en 
la Sección~\ref{sec:auto-encode-img-gen}) para modelar la distribución condicional completa 
$p(\z|\vc)$, lo que nos permite obtener muestras de instancias diversas y nítidas que satisfacen la condición. 

% To generate high-fidelity results, we
% must utilize the Flow Matching framework (Eq. 8.7 \yima{what is this?}) to model the full
% conditional distribution $p(\z|\vc)$, allowing us to sample diverse,
% sharp instances that satisfy the condition.

% \subsection{The Basic Realization: Text-to-Image Generation}
% \label{subsec:basic_realization}

% To realize the conditional framework defined in Section
% \ref{subsec:cond_problem}, we require a concrete architecture capable
% of bridging the modality gap between discrete text tokens and
% continuous pixel arrays. As illustrated in Figure
% \ref{fig:sd_pipeline}, modern state-of-the-art models (such as Stable
% Diffusion~\citep{rombach2022high}, DALL-E
% 2~\citep{ramesh2022hierarchical}, and
% Imagen~\citep{saharia2022photorealistic}) universally follow a
% pipeline consisting of three decoupled components: (1) a Text Encoder
% to interpret the prompt, (2) a Generation Model to synthesize a
% compressed ``latent representation'' of the image, and (3) a Decoder to map
% this representation back to a high-fidelity image in the pixel domain. We detail the
% implementation of each component below, using the Latent Diffusion
% Model (LDM)~\citep{rombach2022high} as the primary case study.

% \begin{figure}[h]
% \centering
% \includegraphics[width=\textwidth]{\toplevelprefix/chapters/chapter8/figs/sd_pipeline.pdf}
% \caption{\textbf{The three-stage inference pipeline of Stable Diffusion.}
%     (1) \textbf{Text Encoder} (e.g., CLIP) converts the input prompt
%     $\vc$ into semantic embeddings.
%     (2) \textbf{Generation Model} (a time-conditional U-Net)
%     iteratively denoises a random latent tensor $\z_T$ to produce a
%     clean intermediate representation $\hat{\z}$, conditioned on the
%     text embeddings via cross-attention.
%     (3) \textbf{Decoder} (from a VAE) maps the denoised latent
%     $\hat{\z}_0$ back to the pixel space to synthesize the final
% high-resolution image $\hat{\x}$.}
% \label{fig:sd_pipeline}
% \end{figure}

% \paragraph{The Text Encoder.}
% The first component converts the user's text prompt $\vc$ (e.g., ``A
% fluffy ginger cat on the couch'') into a machine-readable vector
% sequence. Its role is semantic understanding: because the generative
% model's knowledge is limited to its training data, a powerful
% pre-trained encoder is crucial for generalizing to novel concepts or
% rare words. In practice, Stable Diffusion typically utilizes a frozen
% CLIP text encoder~\citep{radford2021learning}, producing a sequence
% of semantically rich embeddings $\tau_\theta(\vc) \in \mathbb{R}^{d_\tau \times M}$ (where $d_\tau$ is typically 768 or 1024). \yima{Please refer to Section \ref{sec:CLIP}, and make sure notation and dimension match!} Recent
% works like Imagen have demonstrated that scaling up this encoder
% (e.g., using T5-XXL \yima{citation?}) yields significantly larger improvements in
% image fidelity than scaling the generative backbone itself.

% \paragraph{The Generation Model (Latent Representation).}
% The core generative process does not produce the final image
% directly. Instead, to avoid the prohibitive computational cost of
% modeling high-resolution pixels (e.g., $1024 \times 1024$), it
% generates a compressed ``intermediate product''—a latent representation $\z$.
% \begin{itemize}
% \item \textbf{Data Format:} The input image $\x \in \mathbb{R}^{H
%     \times W \times 3}$ is compressed by a factor of $f=8$ into a
%     latent tensor $\z_0 \in \mathbb{R}^{h \times w \times 4}$.\yima{How is it ``compressed''? Via a VAE? Do you follow some work here? citation?} For a
%     standard $512^2$ image, this results in a manageable $64 \times
%     64 \times 4$ latent space.
% \item \textbf{Conditioning Mechanism:} The generation model is
%     implemented as a Time-Conditional U-Net. It predicts the noise
%     $\epsilon$ added to the latent $\z_t$, conditioned on the text
%     embeddings via {\em Cross-Attention}. To enable this interaction,
%     the spatial latent features $\phi_i(\z_t)$ \yima{what ``spatial'' mean here? it was not introduced. what is $\phi_i(\cdot)$? how it is specified?} are flattened into a
%     sequence $\mathbb{R}^{N \times d}$ to act as the Query ($Q$),
%     while the text embeddings serve as Key ($K$) and Value ($V$):
%     \begin{equation}
%         Q = W_Q^{(i)} \cdot \phi_i(\z_t), \quad K = W_K^{(i)} \cdot
%         \tau_\theta(\vc), \quad V = W_V^{(i)} \cdot \tau_\theta(\vc).
%     \end{equation}
%     \yima{is $i$ here the index for layer?}
%     This mechanism allows the model to spatially align specific
%     visual regions (e.g., ``cat's eye'') with specific text tokens.
% \end{itemize}

% Training involves a noise prediction task. We sample a random
% timestep $t$, add Gaussian noise $\epsilon$ to the clean latent
% $\z_0$ to create a noisy latent $\z_t$, and train the U-Net backbone
% $\epsilon_\theta$ to predict the added noise. The model minimizes the
% standard reweighted diffusion objective:
% \begin{equation}
% \mathcal{L}_{\text{LDM}} = \mathbb{E}_{\z_0, \vc, \epsilon \sim
% \mathcal{N}(0,1), t} \left[ \| \epsilon - \epsilon_\theta(\z_t, t,
% \tau_\theta(\vc)) \|^2 \right].
% \end{equation}
% During inference (sampling), we start with pure Gaussian noise $\z_T
% \sim \mathcal{N}(\bm{0}, \I)$. In each step, the U-Net predicts the
% noise component, which is then subtracted from the current state to
% gradually refine the latent representation $\z_t \to \z_{t-1}$ until
% a clean latent $\hat{\z}_0$ is formed.



% \paragraph{The Generation Model (Latent Representation).}
% The core generative process does not produce the final image
% directly. Instead, to avoid the prohibitive computational cost of
% modeling high-resolution pixels (e.g., $1024 \times 1024$), it
% generates a compressed "intermediate product"—a latent representation $\z$.
% \begin{itemize}
% \item \textbf{Data Format:} The input image $\x \in \mathbb{R}^{H
%     \times W \times 3}$ is compressed by a factor of $f=8$ into a
%     latent tensor $\z_0 \in \mathbb{R}^{h \times w \times 4}$.
%     This compression is performed by a pre-trained Variational Autoencoder (VAE)~\citep{rombach2022high} that learns a perceptually equivalent but computationally efficient lower-dimensional space.
%     For a standard RGB input image of size $512 \times 512 \times 3$ (i.e., $H=W=512$), this results in a manageable $64 \times
%     64 \times 4$ latent space.

% \item \textbf{Conditioning Mechanism:} The generation model is
%     implemented as a Time-Conditional U-Net. It predicts the noise
%     $\epsilon$ added to the latent $\z_t$, conditioned on the text
%     embeddings via \textbf{Cross-Attention}.
%     To enable this interaction, we treat the intermediate visual features from the U-Net layer $i$ as the Query source, and the text embeddings as the Key/Value source.
%     Let $\phi_i(\z_t) \in \mathbb{R}^{d_\epsilon^{(i)} \times N}$ denote the flattened spatial features at layer $i$ (where $d_\epsilon^{(i)}$ varies by layer, e.g., 320, 640, or 1280 channels), and let $\tau_\theta(\vc) \in \mathbb{R}^{d_\tau \times M}$ denote the sequence of $M$ text embeddings (typically $d_\tau=768$).
%     Since these dimensions differ, we project both into a common subspace of dimension $d$ (head dimension) using learnable weights.

%     We first define the Query ($\boldsymbol{Q}$), Key ($\boldsymbol{K}$), and Value ($\boldsymbol{V}$) matrices using projection weights $\boldsymbol{U}^{(i)}_{Q} \in \mathbb{R}^{d \times d_\epsilon^{(i)}}$ and $\boldsymbol{U}^{(i)}_{K}, \boldsymbol{U}^{(i)}_{V} \in \mathbb{R}^{d \times d_\tau}$:
%     \begin{equation}
%         \boldsymbol{Q} = \boldsymbol{U}^{(i)}_{Q} \phi_i(\z_t), \quad
%         \boldsymbol{K} = \boldsymbol{U}^{(i)}_{K} \tau_\theta(\vc), \quad
%         \boldsymbol{V} = \boldsymbol{U}^{(i)}_{V} \tau_\theta(\vc).
%     \end{equation}
%     Adapting the column-wise formulation from [placeholder], the cross-attention operation is then computed as:
%     \begin{equation}
%         \text{CrossAttn}(\boldsymbol{Q}, \boldsymbol{K}, \boldsymbol{V}) \doteq \boldsymbol{V} \cdot \operatorname{softmax}\left( \frac{\boldsymbol{K}^\top \boldsymbol{Q}}{\sqrt{d}} \right),
%     \end{equation}
%     where the $\operatorname{softmax}$ is defined column-wise as
%     \begin{align}
%         \operatorname{softmax}(\boldsymbol{A})
%         &\doteq \begin{bmatrix} \operatorname{softmax}(\boldsymbol{a}_{1}) & \cdots & \operatorname{softmax}(\boldsymbol{a}_{N}) \end{bmatrix}, \\
%         \forall \boldsymbol{A}
%         &= \begin{bmatrix} \boldsymbol{a}_{1}, \dots, \boldsymbol{a}_{N} \end{bmatrix} \in \mathbb{R}^{M \times N}.
%     \end{align}
%     In this context, the matrix $\boldsymbol{A} = \frac{\boldsymbol{K}^\top \boldsymbol{Q}}{\sqrt{d}}$ represents the alignment scores, and the column-wise softmax normalizes the attention weights over the $M$ text tokens for each of the $N$ visual regions. This mechanism allows the model to spatially align specific visual regions (e.g., ``cat's eye'') with specific text tokens.
% \end{itemize}

% Training involves a noise prediction task. We uniformly sample a random
% timestep $t \in \{1, \dots, T\}$, add Gaussian noise $\epsilon$ to the clean latent
% $\z_0$ to create a noisy latent $\z_t$. To inform the model of the noise level, the scalar timestep $t$ is first mapped to a sinusoidal embedding (encoding low and high frequency components), followed by a Multi-Layer Perceptron (MLP) composed of `Linear` $\to$ `SiLU` $\to$ `Linear` layers. This produces a rich time embedding vector (e.g., 1280-dim) which is added to the feature maps in each ResNet block of the U-Net backbone $\epsilon_\theta$. The model then minimizes the standard reweighted diffusion objective:
% \begin{equation}
% \mathcal{L}_{\text{LDM}} = \mathbb{E}_{\z_0, \vc, \epsilon \sim
% \mathcal{N}(0,1), t} \left[ \| \epsilon - \epsilon_\theta(\z_t, t,
% \tau_\theta(\vc)) \|^2 \right].
% \end{equation}
% During inference (sampling), we start with pure Gaussian noise $\z_T
% \sim \mathcal{N}(\bm{0}, \I)$. In each step, the U-Net predicts the
% noise component, which is then subtracted from the current state to
% gradually refine the latent representation $\z_t \to \z_{t-1}$ until
% a clean latent $\hat{\z}_0$ is formed.



% \paragraph{The Decoder (Image Reconstruction).}
% The final component is a neural network trained to translate the sample latent vector $\hat{\z}$ back into a natural image
% $\hat{\x}$. Unlike the generation model which requires image-text
% pairs, this Decoder is part of a VAE (Variational Autoencoder) that
% can be trained separately on large-scale unlabeled image data.
% During the training phase, the VAE minimizes a reconstruction error
% (combining L1 loss and Perceptual Loss) to ensure the decoder
% $\mathcal{D}$ can accurately map compressed latents back to pixels
% ($\mathcal{D}(\mathcal{E}(\x)) \approx \x$). Once trained, the
% Encoder is discarded from the inference pipeline, and the Decoder is
% frozen. It simply takes the denoised latent $\hat{\z}_0$ produced by
% the Generation Model and upsamples it to the final pixel result $\hat{\x}$.

% \textit{Note on Visualization:} When visualizing the generation for every time-step using Stable Diffusion, users often observe the image evolving from a
% "blurry blob" to a sharp result. This distinctive "blur-to-sharp"
% transition occurs because, at every step of the sampling process, the
% system passes the current noisy latent $\z_t$ through the Decoder to
% visualize progress. Since the Decoder is trained on clean data, it
% interprets the noise in the latent space as high-frequency texture
% information, rendering the noisy latent as a blurry, low-detail image
% rather than static noise.



\subsection{La implementación básica: generación de imágenes a partir de texto}
\label{subsec:basic_realization}

Para implementar el marco condicional definido en \ref{subsec:cond_problem}, necesitamos 
una arquitectura concreta capaz de puentear la brecha modal entre los tokens de texto 
discretos y las matrices de píxeles continuas. Tal y como se ilustra en la 
Figura~\ref{fig:sd_pipeline}, los modelos modernos de vanguardia (como Stable 
Diffusion~\citep{rombach2022high}) descomponen esta compleja tarea en tres componentes 
independientes: (1) un módulo de \textbf{compresión perceptual} (VAE) para traducir entre 
píxeles y latentes, (2) una interfaz de \textbf{condicionamiento semántico} (codificador 
de texto CLIP) para interpretar las instrucciones de texto (prompts) y (3) una red troncal (backbone) de 
\textbf{generación de latentes} (U-Net) para sintetizar el contenido. A continuación 
detallamos la implementación de cada componente.

\begin{figure}[h]
\centering
\includegraphics[width=\textwidth]{\toplevelprefix/chapters/chapter8/figs/sd_pipeline.pdf}
\caption{\textbf{El proceso de inferencia en tres etapas de Stable Diffusion.} 
    (1) \textbf{Codificador de texto} (p. ej., CLIP) convierte la indicación de entrada 
    $\boldsymbol{c}$ en incrustaciones semánticas. 
    (2) \textbf{Modelo de generación} (una red U-Net condicional en el tiempo $\epsilon_\theta$) 
    elimina el ruido de manera iterativa de un tensor latente aleatorio $\boldsymbol{z}_T$ para 
    producir una representación intermedia limpia $\hat{\boldsymbol{z}}$, condicionada a las 
    incrustaciones de texto mediante atención cruzada. 
    (3) El \textbf{decodificador} (de una VAE) mapea la variable latente sin ruido 
    $\hat{\boldsymbol{z}}$ de vuelta al espacio de píxeles para sintetizar la imagen final de 
    alta resolución $\hat{\boldsymbol{x}}$. \textit{Nota: El codificador de la VAE solo se utiliza 
    durante el entrenamiento y se omite en este proceso de inferencia.}}
\label{fig:sd_pipeline}
\end{figure}


\paragraph{Componente 1: Compresión perceptual (el VAE).}
El primer pilar de la arquitectura es un auto-codificador variacional (VAE) entrenado para 
aprender a abstraer los detalles de alta frecuencia que, aunque son insignificantes perceptualmente,
resultan computacionalmente costosos de modelar. Tal 
y como se presentó en la Sección~\ref{sec:VAE-implement}, esta componente define el 
\textbf{espacio latente} en el que tiene lugar la generación.
\begin{itemize}
    \item \textbf{El codificador ($\mathcal{E}$):} El codificador mapea una imagen RGB de 
          entrada $\boldsymbol{x} \in \mathbb{R}^{H \times W \times 3}$ a una representación 
          latente de baja dimensión 
          $\boldsymbol{z} = \mathcal{E}(\boldsymbol{x}) \in \mathbb{R}^{h \times w \times c}$. 
          Siguiendo a \cite{rombach2022high}, utiliza un codificador regularizado por KL con un 
          factor de submuestreo de $f=H/h=W/w=8$. Para una entrada estándar de 
          $512 \times 512 \times 3$, esto produce un tensor latente compacto de dimensiones 
          $64 \times 64 \times 4$. En el contexto del proceso de difusión (Componente 3), esta 
          representación codificada sirve como el latente de referencia «limpio» $\boldsymbol{z}_0$.
    \item \textbf{El decodificador ($\mathcal{D}$):} El decodificador aprende la correspondencia 
          inversa para reconstruir la imagen a partir del espacio latente, lo cual se denota como 
          $\tilde{\boldsymbol{x}} = \mathcal{D}(\boldsymbol{z}) = \mathcal{D}(\mathcal{E}(\boldsymbol{x}))$. 
          Está entrenado para minimizar la diferencia perceptual entre 
          $\boldsymbol{x}$ y $\tilde{\boldsymbol{x}}$. Durante la fase de inferencia del modelo de 
          difusión, este decodificador se mantiene congelado y se utiliza para mapear los vectores latentes 
          generados (sin ruido) de vuelta al espacio de píxeles de alta resolución.
\end{itemize}
Cabe señalar que el VAE se entrena previamente con una combinación de pérdida perceptual y pérdida 
adversaria basada en parches para garantizar la calidad de las reconstrucciones. Una vez entrenado, 
sus pesos se congelan y actúa como una interfaz estática para el modelo de difusión.

\paragraph{Componente 2: Condicionamiento semántico (el codificador de texto).}
Para guiar el proceso de generación, necesitamos una representación semántica robusta de la 
solicitud del usuario $\vc$ (por ejemplo, «Un gato peludo de color jengibre»). En lugar de 
entrenar un codificador de texto desde cero, los modelos de difusión latente suelen aprovechar 
el codificador de texto preentrenado de CLIP~\citep{radford2021learning}. Este transformer 
preentrenado mapea los tokens del texto de entrada a una secuencia de incrustaciones semánticas 
$\tau_\theta(\vc) \in \mathbb{R}^{d_\tau \times M}$, donde $M$ es la longitud de la secuencia 
(por ejemplo, 77 tokens) y $d_\tau$ es la dimensión de la incrustación (normalmente 768 o 1024). 
Estas incrustaciones sirven como fuentes semánticas de «Clave» (Key) y «Valor» (Value) para la estructura 
central de generación.

\paragraph{Componente 3: Generación latente (la U-Net condicional al tiempo).}
El motor de generación principal se implementa como una red U-Net condicional en el tiempo 
$\epsilon_\theta$. A diferencia de los modelos en el espacio de píxeles, esta red opera 
íntegramente dentro del espacio latente comprimido. Su arquitectura se compone de una serie 
de bloques ResNet para la extracción de características y capas de atención cruzada para el 
condicionamiento.
\begin{itemize}
    \item \textbf{Flujos de entrada:} El modelo procesa tres señales distintas:
    \begin{enumerate}
        \item \textbf{Flujo visual:} La variable latente ruidosa $\z_t \in \mathbb{R}^{h \times w \times 4}$, 
              que representa la imagen en el paso de difusión actual.
        \item \textbf{Intervalo de tiempo ($t$):} Para informar al modelo del nivel de ruido actual, el paso 
              de tiempo escalar $t \in \{1, \dots, T\}$ se codifica mediante una incrustación sinusoidal (que codifica 
              componentes de baja y alta frecuencia), seguida de una red neuronal multicapa (MLP) compuesta por 
              capas `Linear` $\to$ `SiLU` $\to$ `Linear`. Esto produce un vector de incrustación temporal global 
              (por ejemplo, de 1280 dimensiones) que se inyecta en cada bloque ResNet de la U-Net.
        \item \textbf{Flujo contextual:} Las incrustaciones de texto $\tau_\theta(\vc)$ de la componente 2.
    \end{enumerate}
    \item \textbf{Mecanismo de condicionamiento (atención cruzada):} La interacción entre el contenido visual y la 
          indicación de texto se gestiona mediante capas de atención cruzada intercaladas a lo largo de la red U-Net. 
          Sea $\phi_i(\z_t) \in \mathbb{R}^{d_\epsilon^{(i)} \times N}$ las características espaciales aplanadas en 
          la capa $i$ de U-Net (donde $N = h_i \times w_i$ es el número de tokens visuales y $d_\epsilon^{(i)}$ es 
          la profundidad del canal, p. ej., 320, 640 o 1280) y sea $\tau_\theta(\vc) \in \mathbb{R}^{d_\tau \times M}$ 
          las incrustaciones de texto.
          Dado que la dimensión visual $d_\epsilon^{(i)}$ y la dimensión textual $d_\tau$ son diferentes, proyectamos 
          ambas en un subespacio común de dimensión $d$ (la dimensión de la cabeza de atención) utilizando pesos 
          aprendibles. En primer lugar, definimos las matrices de consulta ($\boldsymbol{Q}$), 
          clave ($\boldsymbol{K}$) y valor ($\boldsymbol{V}$):
          \begin{equation}
              \boldsymbol{Q} = \boldsymbol{U}^{(i)}_{Q} \phi_i(\z_t), \quad
              \boldsymbol{K} = \boldsymbol{U}^{(i)}_{K} \tau_\theta(\vc), \quad
              \boldsymbol{V} = \boldsymbol{U}^{(i)}_{V} \tau_\theta(\vc),
          \end{equation}
          donde $\boldsymbol{U}^{(i)}_{Q} \in \mathbb{R}^{d \times d_\epsilon^{(i)}}$ y 
          $\boldsymbol{U}^{(i)}_{K}, \boldsymbol{U}^{(i)}_{V} \in \mathbb{R}^{d \times d_\tau}$. 
          Adaptando la formulación por columnas de \eqref{sub:image_classification_architecture}, 
          la operación de atención cruzada se calcula como:
          \begin{equation}
              \text{CrossAttn}(\boldsymbol{Q}, \boldsymbol{K}, \boldsymbol{V}) \doteq \boldsymbol{V} \cdot \operatorname{softmax}\left( \frac{\boldsymbol{K}^\top \boldsymbol{Q}}{\sqrt{d}} \right),
          \end{equation}
          donde el operador $\operatorname{softmax}$ se define por columnas como
          \begin{align}
              \operatorname{softmax}(\boldsymbol{A})
              &\doteq \begin{bmatrix} \operatorname{softmax}(\boldsymbol{a}_{1}) & \cdots & \operatorname{softmax}(\boldsymbol{a}_{N}) \end{bmatrix}, \\
              \forall \boldsymbol{A}
              &= \begin{bmatrix} \boldsymbol{a}_{1}, \dots, \boldsymbol{a}_{N} \end{bmatrix} \in \mathbb{R}^{M \times N}.
          \end{align}
          En este contexto, la matriz $\boldsymbol{A} = \frac{\boldsymbol{K}^\top \boldsymbol{Q}}{\sqrt{d}}$ 
          representa las puntuaciones de alineación y la función softmax aplicada columna a columna normaliza los pesos de 
          atención sobre los $M$ tokens de texto para cada una de las $N$ regiones visuales. Este mecanismo 
          permite al modelo alinear espacialmente regiones visuales específicas (por ejemplo, «ojo de gato») 
          con tokens de texto específicos.
\end{itemize}

\paragraph{Objetivo de entrenamiento.}
El proceso de entrenamiento unifica estas componentes. Tomamos una imagen aleatoria $\x$, la 
comprimimos utilizando el codificador congelado $\z_0 = \mathcal{E}(\x)$ y le añadimos ruido 
gaussiano $\epsilon$ para generar $\z_t$. A continuación, se entrena la red U-Net $\epsilon_\theta$ 
para predecir y corregir este ruido, condicionado por el texto condicionador $\vc$ y el paso de tiempo de difusión $t$:
\begin{equation}
\mathcal{L}_{\text{LDM}} = \mathbb{E}_{\mathcal{E}(\x), \vc, \epsilon \sim \mathcal{N}(0,1), t} \left[ \| \epsilon - \epsilon_\theta(\z_t, t, \tau_\theta(\vc)) \|^2 \right].
\end{equation}
Durante la inferencia, se descarta el codificador. El proceso comienza a partir de ruido puro 
$\z_T \sim \mathcal{N}(\bm{0}, \I)$, lo depura iterativamente utilizando la red U-Net y, finalmente, 
pasa el $\hat{\z}_0$ resultante a través del decodificador $\mathcal{D}$ para sintetizar la imagen.

\paragraph{Métricas de evaluación.}
Dado que la generación condicional carece de una única «verdad de referencia» (ground truth) correcta, nos basamos en 
métricas estadísticas para evaluar el rendimiento. Denotamos la distribución de las imágenes reales 
como $p_r$ y la distribución de las imágenes generadas como $p_g$.

\begin{itemize}
    \item \textbf{FID (Fréchet Inception Distance):} Esta métrica cuantifica el realismo y la 
          diversidad de las muestras generadas midiendo la distancia estadística entre las distribuciones $p_r$ y $p_g$ en el 
          espacio de características de una red Inception-v3 preentrenada. Aproximamos estas 
          distribuciones de características como gaussianas multidimensionales 
          $\mathcal{N}(\bm{\mu}_r, \bm{\Sigma}_r)$ y $\mathcal{N}(\bm{\mu}_g, \bm{\Sigma}_g)$. 
          El valor o la métrica FID se define entonces como la distancia de Fréchet entre ellas:
          \begin{equation}
              \text{FID}(p_r, p_g) = \| \bm{\mu}_r - \bm{\mu}_g \|^2 + \text{Tr}\left( \bm{\Sigma}_r + \bm{\Sigma}_g - 2(\bm{\Sigma}_r \bm{\Sigma}_g)^{1/2} \right).
          \end{equation}
          Un valor más bajo del FID indica que la distribución generada es estadísticamente más cercana 
          a la distribución real, lo que implica una mayor calidad 
          y diversidad de la imagen.
    \item \textbf{Puntuación CLIP:} Esta métrica mide la alineación semántica entre una imagen 
          generada $\hat{\x}$ y su prompt condicionante $\vc$. Sea $E_{\text{img}}$ y $E_{\text{txt}}$ 
          los codificadores de imagen y texto de un modelo CLIP preentrenado, tal como se explica en el 
          Capítulo \ref{ch:conditional-inference}. La puntuación se calcula como la similitud coseno 
          entre sus incrustaciones:
          \begin{equation}
              \text{CLIP-Score}(\hat{\x}, \vc) = \max \left( 0, \frac{E_{\text{img}}(\hat{\x}) \cdot E_{\text{txt}}(\vc)}{\|E_{\text{img}}(\hat{\x})\| \|E_{\text{txt}}(\vc)\|} \right) \cdot w,
          \end{equation}
          donde $w$ es un factor de escala (normalmente 2,5). Una puntuación más alta indica que el 
          contenido visual refleja fielmente la descripción textual.
\end{itemize}

\paragraph{Arquitecturas avanzadas de condicionamiento por texto (SDXL y Flux).}
Para hacer frente a las limitaciones en cuanto a resolución y adherencia a las instrucciones textuales (prompts), se han 
desarrollado arquitecturas avanzadas que van más allá del LDM básico. SDXL~\citep{podell2023sdxl} 
introduce una arquitectura de doble codificador para capturar matices textuales más ricos, calculando 
representaciones tanto de CLIP ViT-L como de OpenCLIP ViT-bigG y concatenándolas a lo largo del 
eje de canales. En cuanto al manejo de datos, SDXL utiliza el \textit{«Aspect Ratio Bucketing»}, 
que agrupa en lotes las imágenes de entrenamiento con relaciones de aspecto similares para evitar 
que el modelo aprenda a recortar los sujetos de manera arbitraria. Además, implementa el 
\textit{«Micro-Conditioning»}, en el que se incorporan metadatos como la resolución de la 
imagen y las coordenadas de recorte y estos se concatenan con la condición de texto $\vc$, lo 
que permite al modelo distinguir explícitamente entre sujetos centrados y recortados.

Por su parte, Flux~\citep{flux2024} supone una transformación paradigmática, transitando de una arquitectura 
U-Net a un transformer de difusión multimodal (MMDiT). Procesa tokens de texto y visuales en 
un flujo unificado mediante bloques de atención unificada, gracias a componentes modernos de 
transformers como las incrustaciones posicionales rotativas (Rotary Positional Embeddings, RoPE) y RMSNorm. Matemáticamente, 
Flux sustituye la difusión estándar por un \textbf{Rectified Flow Matching}. A diferencia de la 
difusión estándar, que predice el ruido $\epsilon$, Flow Matching simplifica el objetivo de 
entrenamiento para predecir una trayectoria lineal de velocidad $v_\theta$ entre el ruido y los datos. 
Definimos el proceso en un tiempo continuo $t \in [0, 1]$, donde $t=0$ representa la distribución 
del ruido ($\z_0$) y $t=1$ representa la distribución de los datos ($\z_1$).\footnote{Tenga en 
cuenta que esta convención de tiempo ($0 \to \text{Ruido}, 1 \to \text{Datos}$) es la inversa de 
la notación de difusión estándar utilizada en secciones anteriores.} El objetivo es minimizar:
\begin{equation}
    \mathcal{L}_{\text{RFM}} = \mathbb{E}_{t, \z_1, \z_0} \left[ \| v_\theta(\z_t, t, \vc) - (\z_1 - \z_0) \|^2 \right],
\end{equation}
donde $\z_t = t \z_1 + (1-t) \z_0$ es el estado latente interpolado. El modelo $v_\theta$ aprende 
a predecir el campo de velocidades $(\z_1 - \z_0)$ que transporta el ruido $\z_0$ 
directamente a los datos $\z_1$ a lo largo de una trayectoria en línea recta. Esta linealidad mejora 
significativamente la eficiencia de muestreo y la adherencia a las instrucciones en comparación con las 
trayectorias estocásticas curvas de los modelos de difusión tradicionales.


\subsection{Mecanismos avanzados de condicionamiento}
\label{subsec:advanced_conditioning}

\begin{figure}[h]
    \centering
    \includegraphics[width=\textwidth]{\toplevelprefix/chapters/chapter8/figs/cnet_pipeline.pdf}
    \caption{\textbf{Flujo de la arquitectura ControlNet.} Una copia entrenable del codificador y los bloques 
             intermedios de la red U-Net de Stable Diffusion (SD), condicionada a una imagen de referencia a 
             través del codificador de referencia $\mathcal{E}_{\text{hint}}$, inyecta mapas de características 
             en el decodificador U-Net de SD congelado mediante convoluciones cero para permitir el control espacial. (Ejemplo de 
             imagen de \url{https://github.com/huggingface/diffusers/blob/main/docs/source/en/using-diffusers/controlnet.md}).}
    \label{fig:controlnet_pipeline}
\end{figure}

Si bien el condicionamiento de texto proporciona control semántico, presenta dificultades a la 
hora de manejar restricciones espaciales detalladas (por ejemplo, poses específicas o mapas de 
bordes). Para abordar este problema, ControlNet~\citep{zhang2023adding} introduce una arquitectura 
neuronal que añade control espacial a los grandes modelos de difusión de texto a imagen 
preentrenados, tal y como se muestra en la Figura~\ref{fig:controlnet_pipeline}.

\textbf{Arquitectura y conectividad.}
ControlNet funciona clonando los parámetros entrenables de un modelo de difusión preentrenado. 
Específicamente, crea una copia entrenable de los 12 bloques codificadores y 1 bloque intermedio 
de la red U-Net de Stable Diffusion, mientras mantiene los parámetros originales $\Theta$ 
completamente bloqueados (ejecutado en modo sin gradiente con \texttt{torch.no\_grad()}), conservando la calidad 
de generación aprendida a partir de miles de millones de imágenes, mientras  permite al modelo 
aprender tareas condicionales sobre conjuntos de datos pequeños.

Para adaptarse a la resolución espacial de los mapas de características latentes 
$\z \in \mathbb{R}^{h \times w \times c}$ (normalmente $64 \times 64$, tal como se define en 
la Sección~\ref{subsec:basic_realization}), la condición espacial $\vc_s$ (por ejemplo, un mapa 
de bordes de Canny de $512 \times 512$) se procesa primero mediante un \textbf{codificador de pistas} 
ligero $\mathcal{E}_{\text{hint}}(\cdot)$. A diferencia del VAE, este codificador es más profundo 
pero más estrecho, y consta de \textbf{ocho} capas convolucionales con núcleos de $3 \times 3$ y 
activaciones SiLU. Utiliza tres capas con paso 2 para reducir la resolución, aumentando gradualmente 
los canales ($16 \to 32 \to 96 \to 256$) antes de proyectar finalmente a la dimensión del modelo 
(320 canales). Es crucial señalar que la capa final de este codificador se inicializa en ceros.

La innovación clave que permite la fusión de estas dos ramas es la 
\textbf{convolución cero}: una capa convolucional de $1 \times 1$ $\mathcal{Z}(\cdot; \theta_z)$ 
inicializada con los pesos $W$ y el sesgo $B$ fijados estrictamente en cero. La conectividad sigue 
un patrón específico:
\begin{enumerate}
    \item \textbf{Inyección de entradas:} La pista codificada $\mathcal{E}_{\text{hint}}(\vc_s)$ se añade \textbf{únicamente} al primer bloque de entrada de la copia de ControlNet.
    \item \textbf{Inyección de salida:} La salida de cada bloque entrenable de la copia de ControlNet se pasa a través de una convolución cero y, a continuación, se suma a la \textbf{conexión de salto} correspondiente del decodificador Stable Diffusion congelado.
\end{enumerate}

Formalmente, consideremos el $i$-ésimo bloque de la red U-Net. Sea $\boldsymbol{h}_i$ el mapa 
de características de entrada a esta capa. Abstraemos el funcionamiento de este bloque congelado 
como una función $\mathcal{F}_i(\cdot; \Theta_i)$. En consecuencia, sea 
$\mathcal{F}_{\text{trainable}}^{(i)}(\cdot; \Theta_{c,i})$ su copia entrenable en la rama 
ControlNet. La salida $\boldsymbol{y}_i$ se calcula como:
\begin{equation}
    \boldsymbol{y}_i = \mathcal{F}_i(\boldsymbol{h}_i; \Theta_i) + \mathcal{Z}\left( \mathcal{F}_{\text{trainable}}^{(i)}(\boldsymbol{h}_i + \mathcal{E}_{\text{hint}}(\vc_s); \Theta_{c,i}); \theta_z \right).
    \label{eq:controlnet}
\end{equation}


Dado que la convolución por ceros se inicializa en ceros, el término de control desaparece al inicio 
del entrenamiento ($\mathcal{Z}(\cdot) = 0$), lo que garantiza que el modelo se comporte 
exactamente igual que el Stable Diffusion original. Sin embargo, los gradientes de los 
pesos $W$ son distintos de cero:
\begin{equation}
    \frac{\partial \mathcal{Z}(I; \{W, B\})}{\partial W_{i,j}} = I_{p,i} \neq 0,
\end{equation}
donde $I$ es el mapa de características de entrada. Esto garantiza que, aunque la influencia 
inicial sea cero, el aprendizaje comience de inmediato en la primera pasada hacia atrás.

\textbf{Dinámicas de entrenamiento y experimentos.}
ControlNet se entrena para minimizar el objetivo de difusión estándar, pero con la condición 
espacial adicional $\vc_s$ incorporada en el predictor de ruido. El objetivo de aprendizaje específico es:
\begin{equation}
    \mathcal{L} = \mathbb{E}_{\z_0, t, \vc_t, \vc_s, \epsilon \sim \mathcal{N}(0,1)} \left[ \| \epsilon - \epsilon_\theta(\z_t, t, \vc_t, \vc_s) \|^2 \right],
\end{equation}
donde $\vc_t$ representa la condición del prompt de texto y $\z_t$ es la variable latente con 
ruido. Un fenómeno singular que se observa es la «convergencia repentina»: el modelo ignora 
inicialmente la condición espacial y, de repente, se ajusta bruscamente hacia la alineación, a medida que el 
error desciende drásticamente. Para obligar al modelo a basarse en señales estructurales en 
lugar de solo en el texto (que podría contradecir el mapa de bordes), una estrategia clave 
consiste en el \textbf{CFG-Dropout}: sustituir aleatoriamente el 50\% de las indicaciones de 
texto $\vc_t$ por cadenas vacías durante el entrenamiento. Esto obliga a la red a extraer el 
significado semántico directamente de la condición espacial $\vc_s$. Según los resultados 
empíricos, ControlNet es muy eficiente, ya que solo requiere aproximadamente un 23\% más de 
memoria de GPU en comparación con el ajuste fino estándar y sin embargo aprende de manera 
efectiva un control robusto a partir de conjuntos de datos de tan solo 1 000 muestras.


\paragraph{Personalización basada en sujetos.}
Aunque los modelos fundacionales cuentan con un amplio prior semántico, carecen de 
la capacidad de sintetizar instancias de sujetos específicos (por ejemplo, el gato concreto de un 
usuario) de manera coherente en diferentes contextos. Para solucionar esto, DreamBooth 
\citep{ruiz2023dreambooth} introduce un protocolo de ajuste fino de «personalización». El principal 
desafío consiste en introducir un nuevo objeto en el dominio de salida del modelo sin que se 
produzca un «olvido catastrófico»---en el que el modelo pierde su comprensión de la clase general 
(por ejemplo, olvida cómo es un gato en general)---o sufra una «deriva lingüística» (sobrescribiendo 
la palabra «gato» con las características específicas de ese gato en concreto).


\begin{figure}[h]
    \centering
    \includegraphics[width=\textwidth]{\toplevelprefix/chapters/chapter8/figs/db_pipeline.pdf}
    \caption{\textbf{El proceso de entrenamiento e inferencia de DreamBooth.} Durante el 
             \textit{ajuste fino basado en el sujeto}, el modelo de texto a imagen se actualiza 
             utilizando una pérdida combinada: una pérdida de reconstrucción para aprender el sujeto 
             específico (p. ej., «[V] gato») y una pérdida de preservación de priores para mantener 
             el conocimiento de la clase genérica (p. ej., «un gato»). En la \textit{inferencia}, el 
             modelo personalizado puede generar el sujeto específico en contextos nuevos (p. ej., 
             «acostado sobre una mesa de madera») utilizando el identificador único.}
    \label{fig:dreambooth_pipeline}
\end{figure}


\textbf{Estrategia de identificadores y datos.}
A diferencia del entrenamiento estándar, que requiere millones de pares, DreamBooth funciona con un 
conjunto de datos de «pocas muestras» proporcionado por el usuario, que suele contener solo entre 3 y 5 
imágenes del sujeto. Para vincular a este sujeto con el espacio de texto del modelo, se requiere un 
identificador único $[V]$. Los autores proponen una estrategia de selección de tokens poco frecuentes: 
en lugar de utilizar caracteres aleatorios (que podrían tokenizarse en letras comunes con priores 
fuertes), el método busca en el vocabulario del tokenizador secuencias de tokens poco frecuentes (por 
ejemplo, en el rango del tokenizador T5-XXL $\{5000, \dots, 10000\}$) que tengan asociaciones semánticas 
débiles. Este identificador siempre se empareja con un sustantivo de clase general (p. ej., «un gato $[V]$») 
en el prompt. Esta «anclaje» es fundamental: aprovecha la sólida información previa del modelo sobre 
la clase (saber cómo representar un «gato» en diversas posturas) al tiempo que vincula las características 
de identidad específicas al identificador $[V]$.

\textbf{Entrenamiento con preservación de priores.}
Para evitar que el modelo se adapte excesivamente a las imágenes de pocos ejemplos y 
pierda diversidad, DreamBooth emplea una pérdida de preservación de priores autógena 
específica para cada clase. Tal y como se ilustra en la Figura \ref{fig:dreambooth_pipeline}, 
el proceso de entrenamiento minimiza un objetivo combinado:
\begin{equation}
    \mathcal{L}_{\text{DB}} = \mathbb{E}_{\z_t, \vc, \epsilon, t} \left[ \underbrace{\| \epsilon - \epsilon_\theta(\z_t, t, \vc_{\text{subj}}) \|^2}_{\text{Reconstruction Loss}} + \lambda \underbrace{\| \epsilon - \epsilon_\theta(\z'_t, t, \vc_{\text{prior}}) \|^2}_{\text{Prior Preservation Loss}} \right].
    \label{eq:dreambooth}
\end{equation}
El primer término ajusta con precisión los pesos del modelo (normalmente toda la red U-Net y el 
codificador de texto) para reconstruir las imágenes del sujeto condicionadas por la indicación 
única $\vc_{\text{subj}}$ (por ejemplo, «un $[V]$ gato»). El segundo término actúa como 
regularizador: supervisa el modelo con sus propias muestras generadas de la clase genérica 
$\vc_{\text{prior}}$ (por ejemplo, «un gato»). Esto obliga al modelo a conservar su 
interpretación original de la distribución de clases, lo que garantiza que el identificador 
$[V]$ aprenda los matices específicos del sujeto, mientras que el sustantivo de clase «gato» 
sigue siendo genérico.


\textbf{Personalización multiconcepto.}
Aunque DreamBooth es eficaz, el ajuste fino de todo el modelo requiere un gran esfuerzo computacional 
y consume mucho espacio de almacenamiento. Custom Diffusion \cite{kumari2023multi} mejora la eficiencia 
al actualizar únicamente las matrices de proyección de claves ($W_K$) y valores ($W_V$) en las capas 
de atención cruzada, con lo que demuestra que estos pesos específicos son suficientes para codificar 
nuevos conceptos.


\begin{figure}[h]
    \centering
    \includegraphics[width=\textwidth]{\toplevelprefix/chapters/chapter8/figs/gen4gen_teaser.pdf}
    \caption{\textbf{El proceso Gen4Gen para la personalización multiconcepto.} Para superar las 
             limitaciones de los conjuntos de datos existentes en escenarios con múltiples sujetos, 
             Gen4Gen \cite{yeh2024gen4gen} automatiza la creación de un conjunto de datos de 
             referencia de alta calidad (MyCanvas). El entrenamiento con estos datos mejora 
             significativamente la capacidad del modelo para generar escenas complejas con múltiples 
             sujetos específicos (por ejemplo, «un $[V1]$ gato y un $[V2]$ perro»).}
    \label{fig:gen4gen_pipeline}
\end{figure}

Sin embargo, sigue siendo difícil sintetizar con precisión múltiples conceptos personalizados 
(por ejemplo, «un gato $[V1]$ y un perro $[V2]$») en una sola imagen debido a la escasez de datos 
de entrenamiento de alta calidad, con alineación entre texto e imagen, para composiciones complejas 
durante la etapa de preentrenamiento de Stable Diffusion. Para abordar esto, 
Gen4Gen \cite{yeh2024gen4gen} propone una solución centrada en los datos ilustrada en la 
Figura \ref{fig:gen4gen_pipeline}. En lugar de modificar la arquitectura del modelo, introduce un 
flujo de datos generativo que aprovecha los modelos base para construir «MyCanvas», un conjunto de 
datos de referencia que incluye escenas complejas con múltiples objetos, acompañadas de descripciones 
detalladas y densas. Al realizar ajustes finos sobre estos datos sintetizados, en los que las 
identidades de los objetos y las relaciones espaciales se han definido en forma explícita, los 
modelos pueden mejorar significativamente su capacidad para distinguir y conservar múltiples sujetos 
personalizados al mismo tiempo (por ejemplo, un gato y un perro concretos que interactúan en una 
escena compleja) sin que se produzcan mezclas de identidades u omisiones.

Más recientemente, el campo se ha orientado hacia arquitecturas que no requieren ajuste (tuning-free), como 
PhotoMaker \cite{li2024photomaker} e InstantID \cite{wang2024instantid}. En lugar de la optimización 
iterativa, estos métodos emplean una estrategia de «apilamiento»: un codificador de identidad externo 
(a menudo basado en estructuras de reconocimiento facial) extrae una incrustación de identidad de alta 
fidelidad de la imagen de referencia, que luego se inyecta directamente en el proceso de difusión 
mediante atención cruzada desacoplada o apilamiento de prompts. Esto permite la personalización sin 
entrenamiento previo de sujetos arbitrarios en una sola pasada hacia adelante sin actualizaciones de gradientes.


\subsection{Ampliación a la generación de vídeos}
\label{subsec:video_gen}

Los principios de la generación condicional en 2D se extienden en forma natural a la generación de 
vídeo al tratar el vídeo como una secuencia de cuadros correlacionados temporalmente 
$\boldsymbol{x} \in \mathbb{R}^{T \times H \times W \times 3}$. Sin embargo, modelar directamente 
este espacio de alta dimensión resulta computacionalmente prohibitivo. Por lo tanto, los enfoques 
modernos de vanguardia suelen operar dentro del espacio latente comprimido definido por la VAE 
(Sec.~\ref{subsec:basic_realization}), donde el video se representa como un tensor 
$\boldsymbol{z} \in \mathbb{R}^{T \times h \times w \times c}$.

\paragraph{Expansión espacio-temporal.}
Para aprovechar los sólidos conocimientos previos visuales establecidos por los generadores de 
imágenes 2D, arquitecturas como VideoLDM~\citep{blattmann2023align} y 
AnimateDiff~\citep{guo2023animatediff} introducen el concepto de «expandir» las capas 2D 
en capas pseudo-3D. Esto se logra reestructurando el tensor de entrada 
$\boldsymbol{z}$ para alternar entre el procesamiento espacial y el temporal:
\begin{itemize}
    \item \textbf{Capas espaciales (congeladas):} La estructura principal 2D (ResNet y atención 
          espacial) procesa los cuadros de manera independiente remodelando la entrada a 
          $(B \cdot T, h, w, c)$. De este modo se conservan las capacidades de generación de 
          texturas y objetos preentrenadas.
    \item \textbf{Capas temporales (entrenables):} Para modelar el movimiento, se reestructura 
          el tensor con el fin de aislar el eje temporal para cada ubicación espacial. 
          Consideramos los puntos de la grilla espacial ($B \cdot h \cdot w$) como la dimensión 
          del lote y el tiempo ($T$) como la longitud de la secuencia. A continuación, se aplica 
          una capa de \textbf{auto-atención temporal} para intercambiar información entre cuadros.
\end{itemize}

Matemáticamente, sea $\boldsymbol{u} \in \mathbb{R}^{c \times T}$ el vector de características en 
una ubicación espacial específica de un píxel a lo largo del tiempo (análogo a las características 
espaciales $\phi_i(\z_t)$ de la Sección~\ref{subsec:basic_realization}). Definimos las matrices 
temporales de consulta ($\boldsymbol{Q}$), clave ($\boldsymbol{K}$) y valor ($\boldsymbol{V}$) 
utilizando pesos de proyección aprendibles $\boldsymbol{U}^{(temp)}$:
\begin{equation}
    \boldsymbol{Q} = \boldsymbol{U}^{(temp)}_Q \boldsymbol{u}, \quad \boldsymbol{K} = \boldsymbol{U}^{(temp)}_K \boldsymbol{u}, \quad \boldsymbol{V} = \boldsymbol{U}^{(temp)}_V \boldsymbol{u},
\end{equation}
donde $\boldsymbol{U}^{(temp)}_{Q}, \boldsymbol{U}^{(temp)}_{K}, \boldsymbol{U}^{(temp)}_{V} \in \mathbb{R}^{c \times c}$. 
Siguiendo la formulación por columnas establecida en la Sección~\ref{subsec:basic_realization}, la auto-atención temporal se calcula como:
\begin{equation}
    \text{TempAttn}(\boldsymbol{Q}, \boldsymbol{K}, \boldsymbol{V}) \doteq \boldsymbol{V} \cdot \operatorname{softmax}\left( \frac{\boldsymbol{K}^\top \boldsymbol{Q}}{\sqrt{c}} \right) + \boldsymbol{u}.
\end{equation}
Aquí, el término $\boldsymbol{K}^\top \boldsymbol{Q} \in \mathbb{R}^{T \times T}$ representa 
el mapa de atención temporal, que relaciona cada cuadro con todos los demás. Es fundamental 
que la capa de proyección de salida de este bloque de atención se inicialice típicamente en 
cero. Esto garantiza que, al inicio del entrenamiento, la contribución de las capas temporales 
desaparezca, lo que permite que el modelo se comporte exactamente como el generador 2D 
preentrenado y aprenda gradualmente la dinámica temporal.

\paragraph{Modelos de mundos interactivos.}
Los avances recientes han ido más allá de la síntesis pasiva de texto a video para adentrarse 
en la simulación de mundos interactivos, lo que permite considerar al generador de video como 
un modelo predictivo de entornos físicos. Formalmente, el objetivo es modelar la probabilidad 
de transición latente condicionada a la acción de un agente $\boldsymbol{a}_t$:
\begin{equation}
    p(\boldsymbol{z}_{t+1} \mid \boldsymbol{z}_{1:t}, \boldsymbol{a}_t),
\end{equation}
donde $\boldsymbol{a}_t$ representa la señal de control (por ejemplo, la entrada del teclado o 
la orden al motor) en el instante $t$. En marcos como Genie~\citep{bruce2024genie}, la acción 
$\boldsymbol{a}_t$ se discretiza en un token y se inyecta en la red utilizando los mismos 
mecanismos de condicionamiento que se emplean para el texto. Por ejemplo, $\boldsymbol{a}_t$ 
puede asignarse a un vector de incrustación y añadirse a la incrustación del paso de tiempo o 
utilizarse como un par clave/valor en capas de atención cruzada. Esto permite que el modelo 
generativo funcione como un motor de física guiado por datos, prediciendo las consecuencias 
visuales de las acciones (por ejemplo, el salto de un personaje) mediante la minimización de 
los objetivos de reconstrucción o de coincidencia de flujos. Implementaciones a gran escala 
como Sora~\citep{openai2024sora} sugieren además que el escalado de esta predicción auto-regresiva 
o basada en la difusión da lugar a una consistencia 3D emergente y a la permanencia de los objetos.


\section{Generación de objetos 3D condicionada por imágenes y texto}
\label{sec:Michelangelo}\label{sec:object-generation}
%\yima{Zibo and Sam, please move the image-text conditioned 3D shape
% generation material (Michelangelo) to this section. Try to be
% self-contained and consistent with the rest of the book for
% notation and terminologies, especially with the next Section for
% generation-based 3D reconstruction. }

Hasta ahora, en las secciones anteriores, hemos mostrado principalmente cómo aprender
representaciones de la distribución de imágenes {\em bidimensionales} (2D),
así como muchas aplicaciones o tareas basadas en imágenes que
dichas representaciones facilitan: como el reconocimiento de objetos a partir de imágenes,
la segmentación o el completamiento de imágenes, y la
generación de imágenes.\footnote{Basado en técnicas presentadas anteriormente: el
transformer visual (ViT)~\cite{NIPS2017_3f5ee243}, el aprendizaje
contrastivo de lenguaje e imagen (CLIP)~\cite{pmlr-v139-radford21a} y los modelos de
difusión latente (VAE-SD)~\cite{rombach2022high}.}

Sin embargo, nuestro mundo es tridimensional (3D).\footnote{o
cuatridimensional (4D) si consideramos el tiempo como una dimensión adicional, ya que
el mundo es dinámico.} Aunque nuestros ojos solo perciben proyecciones bidimensionales
del mundo---pares estereoscópicos a través de nuestros dos ojos---, somos
capaces de desarrollar una memoria visual que representa un modelo tridimensional completo
de nuestro entorno y de los objetos que lo componen.

En esta sección se demuestra que, utilizando los métodos presentados en este
libro, es posible aprender la distribución de las formas tridimensionales de
objetos naturales y, a continuación, generar nuevas formas de objetos tridimensionales a partir de dicha
distribución, basándose en una imagen bidimensional de entrada o en una descripción
textual del objeto. En concreto, el trabajo que aquí se presenta se
basa en el proyecto Michelangelo~\cite{NEURIPS2023_ea1a7f7b}, que
generaliza en gran medida las técnicas que hemos introducido para
a datos de formas 3D
\footnote{En términos generales, las mallas poligonales y los datos de video pueden
denominarse «datos 3D»; sin embargo, en esta sección, el término
«datos 3D» se refiere a mallas poligonales o sus representaciones equivalentes,
explícitas o implícitas.}.

En primer lugar, ofrecemos una breve introducción a los formatos de datos más comunes que
representan formas 3D y los conjuntos de datos correspondientes. A continuación, siguiendo el
proyecto de código abierto denominado \textbf{Michelangelo}  \cite{NEURIPS2023_ea1a7b}:\\
\centerline{\url{https://github.com/NeuralCarver/Michelangelo},}\\
presentamos un marco computacional, los objetivos y las
arquitecturas de red asociadas que intentan resolver la tarea de la generación de objetos 3D
condicionada por imágenes y texto. Por último, presentamos
los resultados experimentales y el análisis para verificar la solución propuesta, así
como una reflexión sobre posibles mejoras.

% Introduce basic 3D representations and common used 3D dataset.
\subsection{Representaciones de formas en 3D, conjuntos de datos y tokenización}

\paragraph{Representaciones de formas en 3D}
\label{sec:3D representation}
A diferencia de las imágenes 3D, las formas y los datos 3D no se han
presentado formalmente hasta ahora en este libro. Por lo tanto, en este
apartado se ofrece una breve introducción a los datos 3D.
Se recomienda a los lectores que deseen obtener más información que consulten
las fuentes a través de las referencias citadas.

Los datos 3D a los que nos referimos aquí se refieren principalmente a formas físicas reales y continuas
en 3D que, por lo general, se digitalizan en formas discretas y computables
en un determinado espacio de representación. Encontrar buenas representaciones
de formas 3D constituye uno de los problemas centrales en los campos de la
visión computacional y los gráficos computacionales. Los cuatro tipos de representaciones de formas 3D más comúnmente utilizados,
y en gran medida equivalentes matemáticamente,
son: nubes de puntos, vóxeles, mallas poligonales y
funciones implícitas,\footnote{como la función de señal de distancia
(DSF).} tal y como se muestra en
Figura~\ref{fig:3d_representations}~\footnote{Las imágenes de las tres primeras
columnas proceden de \cite{shi2022deep} y la imagen de la última columna
procede de
https://medium.com/@sunil7545/implicit-vs-parametric-3d-shape-representation-9d4c01c8c60c}.

\begin{figure}[h]
\centering
\includegraphics[width=\textwidth]{\toplevelprefix/chapters/chapter8/figs/3drep.pdf}
\caption{\textbf{Ilustraciones de diferentes representaciones 3D 
    del «Stanford Bunny».} (a) la nube de puntos; (b) el vóxel; (c) la malla poligonal;
    (d) la función implícita.}
\label{fig:3d_representations}
\end{figure}

\textbf{Las nubes de puntos} caracterizan con precisión la forma tridimensional de un
objeto mediante puntos muestreados en forma discreta en la superficie del objeto.
Registran directamente las coordenadas espaciales tridimensionales $(x,y,z)$ de los puntos de
la superficie (y tal vez alguna información fotométrica adicional en
esos puntos), tal como se ilustra en
la Figura~\ref{fig:3d_representations}(a) Nube de puntos.
Una nube de puntos se define como un conjunto de coordenadas discretas en el espacio euclidiano 3D:
\begin{equation}
\label{eq:pointcloud}
P = \{ \vp_i \in \mathbb{R}^d\ |\ i=1,2,\dots,N \}
\end{equation}
donde $d \geq 3$ denota la dimensionalidad de las características, lo que puede incluir
(concatenaciones de) coordenadas $(x_i,y_i,z_i)$, colores
$(r_i,g_i,b_i)$ y normales de superficie $(n_{xi},n_{yi},n_{zi})$.

\textbf{Vóxeles}
representan unidades volumétricas en el espacio 3D, lo que constituye una
extensión natural del concepto de píxeles 2D a tres dimensiones. Dicha
representación se obtiene discretizando un espacio 3D en una malla estructurada.
Los modelos típicos de vóxeles dividen el espacio 3D en celdas cúbicas uniformes,
donde cada vóxel está indexado espacialmente y almacena atributos específicos de la región,
tales como ocupación, densidad, material o color, como se muestra
en la Fig.~\ref{fig:3d_representations}(b) Vóxel.

\textbf{Mallas poligonales}
se consideran, posiblemente, la representación más utilizada para las formas tridimensionales
en los ámbitos de los gráficos por computadora y la simulación física. Aproximan
las superficies continuas mediante vértices, aristas y caras poligonales
discretas, siendo las mallas triangulares la variante predominante, tal como
se muestra en la Fig.~\ref{fig:3d_representations}(c) Malla poligonal.
Geométricamente, las mallas poligonales consisten en coordenadas de vértices en $\mathbb{R}^3$,
complementadas por su conectividad topológica. Formalmente, una malla $M$ es una tupla
$(V,F)$, donde el conjunto de vértices $V=\{\vv_i\in\mathbb{R}^3\}_{i=1}^n$
denota las coordenadas de los vértices y el conjunto de caras $F=\{f_j\}_{j=1}^m$
define la topología.

\textbf{Funciones implícitas}
representar superficies mediante conjuntos de niveles de una función escalar. Este enfoque
ofrece un marco unificado para la descripción matemática de
superficies complejas (Fig.~\ref{fig:3d_representations}(d) Función implícita).
La formulación implícita estándar es:
\begin{equation}
\label{eq:implicit_function}
F: \mathbb{R}^3 \rightarrow \mathbb{R}, \quad F(x,y,z) = d,
\end{equation}
donde $(x,y,z)$ denotan las coordenadas de un punto de consulta, y $d$
denota la distancia entre el punto de consulta y una superficie de
interés.\footnote{Obsérvese que hemos utilizado el mismo concepto para
describir el soporte de baja dimensión de una distribución en un
espacio de alta dimensión, en la \Cref{sec:manifold-distance}.} Por lo tanto, la
superficie de interés se define implícitamente como:
\begin{equation}
\label{eq:implicit_function_surface}
\mathcal{S} = \{ (x,y,z) \in \mathbb{R}^3 \mid F(x,y,z)=0. \}
\end{equation}
La conversión de una representación funcional implícita de una forma 3D
a una malla poligonal suele realizarse mediante métodos de extracción de isosuperficies
como el de los cubos marchantes~\cite{10.1145/280811.281026}.

Desde el punto de vista matemático, las cuatro formas de representación son, en gran medida,
equivalentes en cuanto a su capacidad de representación.
No obstante, cada formato de datos puede presentar diferentes
ventajas o limitaciones computacionales. Por ello, a menudo se
prefiere una representación u otra para distintas tareas. En esta sección, utilizaremos la
representación tradicional de nubes de puntos y mostraremos cómo aprender la
distribución de las formas 3D a partir de este tipo de datos. En la siguiente
sección, mostraremos cómo utilizar una representación (generalizada) de vóxeles
para aprender la distribución de las formas y las posturas 3D.

\paragraph{Conjuntos de datos 3D.}
Existen numerosos conjuntos de datos de acceso público para cada una de las
representaciones mencionadas anteriormente. Estos varían en cuanto a
cantidad y calidad. En esta sección, utilizaremos
ShapeNet~\cite{chang2015shapenet} para estudiar una representación conjunta de forma, imagen y texto,
así como su rendimiento en la generación. ShapeNet ofrece
aproximadamente 55 000 mallas generadas en 55 categorías.
Cada malla cuenta con una etiqueta de categoría y los textos correspondientes, tales como
categorías detalladas o breves descripciones proporcionadas por el creador.
Para preparar los datos de tripletas (forma 3D, imagen, texto), los textos proporcionados
se amplían de dos maneras.
En primer lugar, la etiqueta de forma y la descripción correspondiente se combinan en
el formato «un modelo 3D de (etiqueta de forma), al estilo de
(descripción)» o «un modelo 3D de (etiqueta de forma), (descripción)».
Luego, inspirándonos en ULIP~\cite{10203465}, se utilizan múltiples plantillas que contienen
65 frases predefinidas para proporcionar más información textual
durante el entrenamiento.
En cuanto a los datos de imagen, cada malla se renderiza en cuatro posiciones de cámara,
lo que aumenta y mejora la diversidad del renderizado mediante la red
ControlNet condicionada por profundidad~\cite{10377881}.

\paragraph{Tokenización de datos de formas 3D.}
\label{sec:3d_tokenization}
Como se ha comentado anteriormente, los datos 3D
se presentan en diversas formas, entre ellas nubes de puntos, vóxeles, mallas
poligonales y representaciones neuronales implícitas. Fundamentalmente, las
nubes de puntos, los vóxeles y las mallas sirven como aproximaciones discretas de
variedades 3D continuas. En los gráficos por computadora y la ingeniería,
las mallas poligonales son particularmente omnipresentes porque conservan explícitamente
la topología de la superficie, ofreciendo un equilibrio eficiente entre
la fidelidad visual y la estructura geométrica.

Sin embargo, la aplicación directa de las modernas arquitecturas de aprendizaje profundo
---en concreto, los Transformers---a estas representaciones sin procesar
plantea importantes retos. A diferencia de las imágenes 2D, que poseen una
estructura de cuadrícula regular que se divide fácilmente en parches, los datos 3D suelen ser
esparsos, no estructurados e irregulares. Para aprovechar el poder del
modelado de secuencias y los mecanismos de auto-atención en 3D, debemos salvar
la brecha entre los datos geométricos sin procesar y las secuencias discretas y ordenadas
que requieren los Transformers. Este proceso se conoce como
\textit{tokenización 3D}.

\begin{itemize}
\item \textbf{Tokenización de nubes de puntos.} 
Las nubes de puntos sin procesar suelen tokenizarse en grupos locales para capturar
el contexto geométrico. Este proceso suele comenzar con el muestreo del punto más lejano (FPS)~\cite{eldar1997farthest},
un algoritmo iterativo que selecciona un subconjunto de puntos $\{x_1, \dots, x_m\}$
del conjunto original $\mathcal{P}$, de tal manera que cada punto seleccionado maximice
la distancia al conjunto de puntos seleccionados previamente:
$x_{i} = \arg\max_{x \in \mathcal{P}} \min_{j=1}^{i-1} \|x - x_j\|$.
Tras el FPS, los puntos se agrupan mediante k-Nearest Neighbors (k-NN) para formar clústeres locales.
Cada clúster sirve como un token para gestionar la invariancia de permutación de los puntos
dentro de un clúster, por lo que se aplica una red ligera como PointNet~\cite{qi2017pointnet}.
Matemáticamente, para un clúster de puntos $\{p_1, \dots, p_k\}$,
PointNet aproxima una función simétrica $f$ aplicando un
perceptrón multicapa (MLP) compartido $h$ seguido de un operador de agregación simétrico (normalmente max-pooling):
\begin{equation}
    \boldsymbol{t} = \gamma \left( \max_{i=1}^k \{ h(p_i) \} \right)
\end{equation}
donde $\boldsymbol{t}$ es el token de características resultante y $\gamma$ es una red neuronal multicapas (MLP) de posprocesamiento.
Esto convierte una nube de puntos no estructurada en una secuencia de vectores de características.


% {\color{blue} Update: We have added the implement details of Farthest Point
%     Sampling, and use math formula to describe the process of PointNet, and the output token format. The following is the original version, and is replaced by the above paragraph.}
    
% Point clouds are typically
%     tokenized using grouping techniques such as Farthest Point
%     Sampling (FPS) \yima{what is this?} followed by k-Nearest Neighbors (k-NN). Each
%     cluster of points is treated as a token, often processed by a
%     permutation-invariant networkwork (such as PointNet \yima{reference? what does it do exactly, in mathematical terms?} ) to extract a
%     feature vector.
    
\item \textbf{Tokenización de vóxeles.} 
De manera análoga a la estrategia de incrustación de parches utilizada en los Vision Transformers 
(ViT) para imágenes 2D (véase el Capítulo \ref{sub:contrastive_learning_architecture}), las mallas 
de vóxeles pueden dividirse en submallas 3D que no se superponen. Dada una cuadrícula de vóxeles 
$V \in \mathbb{R}^{D \times H \times W \times C}$, extraemos parches de tamaño $P \times P \times P$. 
Cada parche se aplana en un vector $v_p \in \mathbb{R}^{P^3 \cdot C}$ y se mapea mediante una 
proyección lineal aprendible a la dimensión de los tokens $d_{model}$. El resultado es una secuencia 
de tokens $Z \in \mathbb{R}^{L \times d_{model}}$ donde $L = \frac{DHW}{P^3}$. Aunque sencillo, este 
método hereda la complejidad cúbica de la resolución de vóxeles ($O(N^3)$), lo que a menudo crea 
cuellos de botella en la memoria para formas de alta resolución.


% Similar to 2D image patches, voxel
%     grids can be partitioned into 3D sub-grids or ``cublets." While
%     straightforward, this method often suffers from the cubic
%     complexity of voxel resolution (the $O(N^3)$ memory bottleneck). 
    
%     Extending the 2D patch embedding strategy to 3D, 
% we partition the input voxel grid $V \in \mathbb{R}^{D \times H \times W \times C}$ into 
% non-overlapping 3D patches or ``cubelets'' of size $p \times p \times p$. 
% These cubelets are flattened into 1D vectors and mapped via a learnable linear projection 
% to the token dimension $d_{model}$. 
% The resulting output is a sequence of tokens with length $L = \frac{DHW}{p^3}$. 
% However, this approach suffers from the cubic complexity of voxel resolution, 
% as the sequence length $L$ scales with $O(N^3)$, causing significant memory bottlenecks.
% {\color{blue} Update: We have added the implement details of the output tokens and the reference section about the patch embedding strategy in Vision Transformers for 2D images.}
    
%     \yima{Not clear at all, what are the output tokens? how are they generated in what format? If this process is similar to tokenizing 2D image, please refer to an earlier section that describes tokenizing 2D image, if that is what you mean here.}


\item \textbf{Tokenización de mallas.} 
La tokenización de mallas supone un reto debido a su estructura similar a un grafo, 
que incluye tanto vértices como caras.  A diferencia de las imágenes, las mallas carecen 
de un orden canónico. Para permitir el modelado secuencial, los enfoques habituales implican
una serialización explícita, en la que los elementos geométricos y topológicos sin procesar 
se cuantifican directamente y se ordenan en secuencias.

% Current methods typically serialize the mesh elements into a sequence of discrete tokens to enable autoregressive modeling (where the model predicts the next vertex/face token based on previous ones to generate valid topology).

% Prominent approaches like MeshGPT~\cite{siddiqui2024meshgpt} serialize the mesh into two interleaved sequences: (1) a **Vertex Sequence**, where vertex coordinates are quantized and ordered (often via z-order curve sorting), and (2) a **Face Sequence**, where faces are represented as indices pointing to the vertex list. Alternatively, graph-based methods may treat local mesh patches (geodesic neighborhoods) as tokens, encoding them via Graph Convolutional Networks (GCNs) into feature vectors.

Tomando como ejemplo PolyGen~\cite{nash2020polygen}, la malla se representa como un conjunto de
vértices $V = \{v_1, \dots, v_n\}$ (donde cada $v_i \in \mathbb{R}^3$) y caras
$F = \{f_1, \dots, f_m\}$ (donde cada $f_j$ es una tupla de índices de vértices).
El modelo aprende la distribución conjunta $P(V, F)$ al dividirla en dos etapas secuenciales:
un modelo de vértices que predice coordenadas cuantificadas y un modelo de caras que predice índices de vértices
\begin{equation}
    P(V, F) = P(V) \cdot P(F \mid V) = \prod_{i=1}^{n} P(v_i \mid v_{<i}) \cdot \prod_{j=1}^{m} P(f_j \mid f_{<j}, V).
\end{equation}
Fundamentalmente, esta factorización determina dos formatos de tokens distintos correspondientes a cada etapa.
Para el modelo de vértices ($P(V)$), los tokens son valores de coordenadas cuantificados;
las coordenadas continuas se discretizan (por ejemplo, en enteros de 8 bits),
creando una secuencia de tokens a partir de un vocabulario fijo de tamaño 256.
Por el contrario, para el modelo de caras ($P(F \mid V)$), los tokens son índices de vértices (punteros).
Aunque es eficaz, operar con estas primitivas sin procesar suele dar lugar a secuencias 
de longitud extremadamente grande, lo que limita la escalabilidad.

% {\color{blue} Update: We take the PolyGen as a representative work and add the implement details of the mesh tokenization process. The following is the original version, and is replaced by the above paragraph.}
% Tokenizing meshes presents unique
%     challenges due to their irregular connectivity and graph-like
%     structure. Common approaches involve traversing the surface to
%     serialize faces or vertices into sequences, or decomposing the
%     mesh into discrete primitives. These methods are particularly
%     prevalent in autoregressive mesh generation \yima{what is this, how it is related to the tokenization we are discussing here}, where the objective
%     is to generate not just the 3D geometry, but also a valid, clean topology.\yima{This paragraph does NOT describe exactly how to tokenize meshes at all, not even mention any existing method! }

\end{itemize}

Si bien estas estrategias de tokenización explícita discretizan eficazmente
las variedades continuas, operan directamente en el espacio de observación de alta dimensión,
lo que a menudo da lugar a secuencias excesivamente largas
y a redundancia computacional. Para hacer frente a estas limitaciones de escalabilidad
y lograr mayores índices de compresión, el campo se ha inclinado hacia
las representaciones discretas basadas en el aprendizaje. En lugar de tokenizar directamente las primitivas geométricas sin procesar,
los siguientes métodos proyectan primero los datos 3D
en un espacio latente compacto y de baja dimensión donde se agrega la información semántica,
lo que permite separar de manera efectiva la reconstrucción geométrica del modelado estructural.

\textbf{Cuantificación de vectores latentes estándar.} Este enfoque aprende un
libro de códigos discreto para representar la forma mediante la cuantización de vectores (VQ). 
Un método representativo es  AutoSDF~\cite {mittal2022autosdf}. En este marco, una entrada de alta dimensión
(por ejemplo, una malla de vóxeles $X$ o un campo de funciones implícito $F_\theta$)
es comprimida por un codificador $E$ en una
malla espacial de menor dimensión de vectores latentes $Z \in \mathbb{R}^{h \times w \times d \times c}$. 
A continuación, cada vector $z_i \in Z$ se sustituye por su vecino más cercano de un libro de códigos 
entrenable $\mathcal{C}=\{c_k\}_{k=1}^K$:
\begin{equation}
q(z_i) = \arg\min_{c_k \in \mathcal{C}} | z_i - c_k |_2^2.
\end{equation}
Esto reduce la forma a una secuencia de índices discretos del libro de códigos que conservan
la estructura espacial al tiempo que reducen drásticamente la dimensionalidad. Estos índices discretos
sirven como tokens para el procesamiento posterior del Transformer. Sin embargo, este enfoque basado 
en una cuadrícula regular suele tratar de manera uniforme tanto el espacio ocupado como el vacío, lo 
que da lugar a redundancia computacional.  This reduces the shape to a sequence of discrete codebook 
indices that preserve spatial structure while drastically reducing dimensionality. These discrete indices 
serve as the tokens for subsequent Transformer processing. However, this regular grid-based approach 
typically treats occupied and empty space uniformly, leading to computational redundancy.

% {\color{blue} Update: We add more details of latent vector quantization process and wew specify that the 3D shape is converted to a sequence of tokens.}
% \yima{Then each 3D shape is converted to a single vector, instead of a sequence of tokens?}

\textbf{De la optimización estructural a la alineación semántica.}
Para abordar la redundancia de las mallas regulares, enfoques como
3DILG~\cite{zhang20223dilg} proponen mallas latentes irregulares.
En lugar de una malla de vóxeles densa en la que los tokens se asignan de manera ineficiente
a espacios vacíos, 3DILG define la forma como un conjunto
de códigos latentes que flotan en posiciones arbitrarias y adaptativas
en un espacio continuo, normalmente concentrados en la superficie de la forma.
En concreto, dada una forma de entrada representada como una nube de puntos, el método primero 
la codifica en un conjunto de pares de características y posiciones 
$\mathcal{P} = \{ (x_i, f_i)\}_{i=1}^M$, donde $x_i \in \mathbb{R}^3$ son puntos clave 
seleccionados mediante muestreo del punto más lejano (FPS) y $f_i \in \mathbb{R}^C$ son sus 
vectores de características continuas asociados, extraídos por un codificador basado en 
PointNet~\cite{qi2017pointnet}. La conversión de este conjunto no estructurado en una 
secuencia de tokens implica tres pasos clave:
\begin{itemize}
\item \textbf{Cuantificación de coordenadas.} En primer lugar, para un conjunto de $M$ puntos clave
    seleccionados mediante el método de muestreo del punto más lejano (véase «Tokenización de nubes de puntos»),
    las coordenadas espaciales continuas $x_i \in
    \mathbb{R}^3$ se discretizan en enteros de 8 bits:
    \begin{equation}
        x_i \rightarrow (x_{i,1}, x_{i,2}, x_{i,3}), \quad
        \text{where } x_{i,k} \in {0, 1, \dots, 255}.
    \end{equation}
    Este paso convierte de manera efectiva las posiciones continuas en
    unidades discretas y tokenizables.

\item \textbf{Cuantificación de características.} Al mismo tiempo, el vector de características locales 
asociado a cada punto se comprime en un índice de libro de códigos discreto $z_i$ mediante cuantización vectorial (VQ):
\begin{equation}
    z_i \in \{0, 1, \dots, D-1\},
\end{equation}
donde $D$ representa el tamaño del libro de códigos.

\item \textbf{Construcción de secuencias.} Para facilitar la generación auto-regresiva con Transformers, 
es necesario serializar el conjunto no estructurado de tuplas. 3DILG emplea una ordenación lexicográfica 
basada en las coordenadas cuantificadas. La forma final se representa como una secuencia estructurada 
$\mathcal{S}$ que combina explícitamente la información de las coordenadas posicionales y las 
características geométricas locales:
\begin{equation}
    \mathcal{S} = \big\{ (x_{0,1}, x_{0,2}, x_{0,3}, z_0), \dots, (x_{M-1,1}, x_{M-1,2}, x_{M-1,3}, z_{M-1}) \big\}.
\end{equation}
\end{itemize}
Esta formulación híbrida reduce drásticamente la longitud de las secuencias en comparación con las mallas densas, 
al tiempo que conserva los detalles de la superficie con gran precisión.

Aunque 3DILG optimiza la estructura espacial de los tokens,
sigue siendo, en gran medida, un método de compresión geométrica.
A continuación, ofrecemos una exposición detallada
de Michelangelo, un trabajo que amplía este paradigma
al abordar la brecha semántica. Codifica formas tridimensionales
en una secuencia de tokens latentes que no solo
son geométricamente compactos, sino que también están alineados explícitamente
con el espacio de características de los modelos de visión-lenguaje bidimensionales
(p. ej., CLIP)~\cite{pmlr-v139-radford21a}), lo que
facilita la generación condicional de alta calidad.

% Introduce task and challenges
\subsection{Tarea}
Nuestro objetivo aquí es aprender un modelo generativo de formas 3D condicional
capaz de generar formas 3D a partir de la distribución aprendida, sujeto
a una condición determinada, como una imagen o un texto, tal y como
se ilustra en la Figura \ref{fig:michelangelo_network}.

Siguiendo el modelo de difusión latente~\cite{rombach2022high}, para aprovechar al máximo
las capacidades del modelo generativo, primero se entrena una VAE
para comprimir muestras de datos 3D en un espacio latente.
En general, una buena representación latente debe poseer dos características clave:
debe ser altamente compresible y, al mismo tiempo, permitir una reconstrucción fiel
de las muestras de datos 3D, además de facilitar el entrenamiento del modelo generativo.
Sin embargo, en comparación con los datos de imágenes 2D, los datos 3D siguen siendo
escasos~\footnote{Esta situación persiste incluso tras el lanzamiento de
Objaverse-XL~\cite{10.5555/3666122.3667676}.}.
En consecuencia, aprender una representación latente efectiva para datos 3D
supone un gran desafío. La Figura \ref{fig:michelangelo_network} muestra
el diagrama general del marco propuesto por el proyecto Michelangelo.
A continuación explicaremos los detalles de implementación de cada componente.
\begin{figure}[t]
\centering
\includegraphics[trim=0cm 0cm 0cm 0cm,
clip=true,width=.99\linewidth]{\toplevelprefix/chapters/chapter8/figs/michelangelo_network.pdf}
\caption{
    \textbf{Proceso de alineación previa a la generación}. El enfoque
    incluye dos modelos: el codificador automático variacional
    con alineación de forma, imagen y texto (SITA-VAE) y el modelo de
    difusión latente de formas alineadas (ASLDM). El SITA-VAE consta de cuatro módulos: un codificador visual,
    un codificador de texto, un codificador de formas 3D y un decodificador de formas 3D.
    Los codificadores codifican los pares de entrada en un espacio alineado, y
    el decodificador de formas 3D reconstruye formas 3D a partir de las incrustaciones del
    espacio alineado. El ASLDM mapea el estado de la imagen o el texto al
    espacio latente de formas alineadas para obtener una muestra de una incrustación de forma 3D
    de alta calidad, que posteriormente el decodificador de formas 3D
    reconstruye en formas 3D de alta fidelidad.
}
\label{fig:michelangelo_network}
\end{figure}

% Introduce network and loss func
\subsection{Arquitectura y objetivo}

\paragraph{Auto-codificador variacional alineado con forma, imagen y texto.}

Inspirado en CLIP~\cite{pmlr-v139-radford21a}, el enfoque adoptado por
el proyecto Michelangelo inyecta información semántica en el espacio latente
mediante una pérdida contrastiva para mitigar los problemas observados al
entrenar VAE exclusivamente a partir de datos 3D, resultando en el Autoencoder 
Variacional Alineado de Forma, Imagen y Texto (SITA-VAE).

El SITA-VAE consta de cuatro componentes: un codificador visual preentrenado y fijo CLIP 
$\mathcal{E}_{\mathrm{i}}$ y un codificador de texto CLIP
$\mathcal{E}_{\mathrm{t}}$, un codificador de formas 3D entrenable
$\mathcal{E}_{\mathrm{s}}$ y un decodificador de campo neuronal $\mathcal{D}_{\mathrm{s}}$.
El codificador visual y el codificador de texto de CLIP toman como entrada imágenes bidimensionales $\vI \in \R^{H \times
W \times 3}$ y textos tokenizados $\vT \in \R^{d_{{\mathrm{t}}} \times
L_{\mathrm{t}}}$, y generan tokens de imagen
$\vE_{\mathrm{i}} \in \R^{d \times (1 + L_{\mathrm{i}})}$ y tokens de texto
$\vE_{\mathrm{t}} \in \R^{d \times L_{\mathrm{t}}}$, donde
$(1+L_{\mathrm{i}})$ y $L_{\mathrm{t}}$ son la longitud de la secuencia de
los tokens de imagen $\vE_{\mathrm{i}}$ y de texto $\vE_{\mathrm{t}}$.
% Advantage
El enfoque aprovecha el codificador visual y
el codificador de texto preentrenados de CLIP. Estos dos codificadores se han entrenado con pares imagen-texto a gran escala
y son lo suficientemente robustos como para capturar un espacio visión-lenguaje bien alineado,
lo que enriquece la semántica de la representación de la forma 3D
tras la alineación multimodal mediante aprendizaje contrastivo.

% 3D shape encoder
\textbf{El «codificador de formas 3D»} tiene como objetivo extraer representaciones de características
poderosas para caracterizar cada forma 3D de manera eficaz. Para lograrlo,
primero se muestrean nubes de puntos $\vP \in \R^{(3 + C) \times N}$
de la superficie de las formas 3D, donde $N$ representa el número de
puntos, y $C$ denota características adicionales de los puntos, como la normal o
el color.

Para capturar mejor los detalles geométricos de alta frecuencia,
utilizamos la codificación posicional de Fourier para mapear las
coordenadas espaciales de baja dimensión
$\boldsymbol{v} \in \mathbb{R}^3$
de cada punto a un espacio de frecuencias de mayor dimensión.
Esta codificación $\gamma(\boldsymbol{v})$ se define como una
secuencia de funciones sinusoidales con frecuencias que
aumentan exponencialmente:
$$\gamma(\boldsymbol{v}) = [\sin(2^0\pi \boldsymbol{v}), \cos(2^0\pi \boldsymbol{v}), \dots, \sin(2^{L-1}\pi \boldsymbol{v}), \cos(2^{L-1}\pi \boldsymbol{v})],$$
donde $L$ es el número de bandas de frecuencia.
Para construir la representación de entrada final, descomponemos explícitamente
la nube de puntos de entrada $\vP \in \mathbb{R}^{(3+C)
\times N}$ en coordenadas espaciales $\boldsymbol{v}$ y
características auxiliares $\boldsymbol{f} \in \mathbb{R}^{C \times
N}$. Las incrustaciones posicionales de alta dimensión
$\gamma(\boldsymbol{v})$ se concatenan entonces con las
características conservadas $\boldsymbol{f}$ y se proyectan a través de una
capa lineal aprendible. Formalmente, la entrada del codificador de formas 3D
$\boldsymbol{X} \in \mathbb{R}^{d \times N}$ se obtiene como:
$$\boldsymbol{X} = \text{Linear}([\gamma(\boldsymbol{v}); \boldsymbol{f}])$$
donde $[\cdot ; \cdot]$ denota la operación de concatenación a lo largo de la 
dimensión de la característica.

% Next, a linear layer is used to project the concatenation of
% the Fourier positional encoded point clouds $\vP$ to the 3D shape
% encoder input $\vX \in \R^{d \times N}$. 
% \yima{What does ``Fourier positional encoded'' mean? How is $\X$ related to $\vP$, exactly?} {\color{blue} Update: We have added the formula for Fourier positional encoding and provided a mathematical expression describing the process of transforming from $\vP$ to $\X$.}

Para procesar esta secuencia de características densa $\vX$ de manera eficiente,
% \yima{Which data, $\vP$ or $\X$?} {\color{blue} Update: We have changed the expression to avoid the ambiguity.} 
el codificador adopta una
arquitectura basada en Perceiver~\cite{pmlr-v139-jaegle21a}. A diferencia de
los Transformers estándar que escalan cuadráticamente con la longitud de la entrada,
el Perceiver emplea un mecanismo de cuello de botella latente para desacoplar
la complejidad del procesamiento del tamaño de la entrada. En lugar de tokenizar los
puntos de entrada directamente para la auto-atención, el modelo utiliza un conjunto fijo
de consultas latentes pequeñas y aprendibles $\vQ \in \R^{(1+L_{\mathrm{s}}) \times
d}$, donde $1 + L_{\mathrm{s}}$ es el número de tokens de consulta.
Desde el punto de vista mecánico, una capa de atención cruzada proyecta el espacio de entrada de alta dimensión
en este espacio latente compacto. Formalmente, siguiendo la formulación del Capítulo ~\ref{subsec:basic_realization}, sean $\boldsymbol{U}^{(temp)}_{Q}, \boldsymbol{U}^{(temp)}_{K}, \boldsymbol{U}^{(temp)}_{V}$ las matrices de proyección aprendibles. La representación latente inicial
$\vZ_0$ se calcula como:
$$\vZ_0 = \text{Atención cruzada}
(\vQ, \vX) = (\boldsymbol{U}^{(temp)}_{V}\vX) \cdot \text{softmax}\left(\frac{(\boldsymbol{U}^{(temp)}_{K}\vX)^\top(\boldsymbol{U}^{(temp)}_{Q}\vQ)}
{\sqrt{d_k}}\right).$$
Aquí, la nube de puntos de entrada $\vX$ sirve
como claves y valores, mientras que las consultas entrenables $\vQ$ actúan como
consultas. Esta operación comprime la entrada de longitud variable $N$ en una
secuencia latente de longitud fija $\vZ_0 \in \R^{(1+L_{\mathrm{s}}) \times
d}$.

En el contexto de Michelangelo, estas consultas están estructuradas para
capturar la semántica jerárquica: consisten en un token global $\vQ_g \in \R^{1 \times d}$ con semántica de alto nivel y
$L_{\mathrm{s}}$ tokens locales $\vQ_l \in \R^{L_{\mathrm{s}} \times d}$
que contienen información de estructura geométrica de bajo nivel. A continuación, se utilizan varios
bloques de auto-atención para mejorar iterativamente la representación de características
y obtener las incrustaciones de forma finales, $\vE_{\mathrm{s}}
\in \R^{(1+L_{\mathrm{s}}) \times d}$.

% Drawing inspiration from
% previous transformer-based architectures for point cloud
% understanding~\cite{pmlr-v139-jaegle21a}, the 3D shape encoder is
% built on a perceiver-based transformer. Specifically, a
% cross-attention layer is used to inject the 3D shape information from
% the input $\vX$ into a series of learnable query tokens $\vQ \in
% \R^{d \times (1+L_{\mathrm{s}})}$, where $1 + L_{\mathrm{s}}$ is the
% length of query tokens $\vQ$, consisting of one global head token
% $\vQ_g \in \R^{d}$ with high-level semantics and $L_{\mathrm{s}}$
% local tokens $\vQ_l \in \R^{d \times L_{\mathrm{s}}}$ containing
% low-level geometric structure information. Then, several
% self-attention blocks are used to iteratively improve the feature
% representation and obtain the final shape embeddings,
% $\vE_{\mathrm{s}} \in \R^{d \times (1+L_{\mathrm{s}})}$.

% \yima{Probably more detail or illustration for the above description about how to tokenize 3D shape data, say point clouds, through a perceiver. How is it trained, separately or jointly, with what architecture and data etc.? How is it related to the ``3D tokenization'' methods introduced earlier?} {\color{blue} Update: We have supplemented the description of how the input is tokenized through the perceiver-based architecture, and the entire process is trained jointly. }

% Learn aligned latent
\textbf{La alineación entre formas 3D, imágenes y textos} desempeña un papel crucial
en SITA-VAE y en los modelos generativos condicionales. Dado que los datos 3D
son un orden de magnitud más pequeños que los datos de imágenes y texto, para aprender
una representación mejor alineada entre formas 3D, imágenes y textos,
el codificador de formas 3D está restringido a estar cerca de un espacio visión-lenguaje prealineado,
el cual está preentrenado en pares imagen-texto a gran escala
con representaciones ricas de imagen y texto al aprovechar la
estrategia de aprendizaje contrastivo. Consideremos un triplete de entrada compuesto por formas 3D
$\vX$, imágenes $\vI$ y textos tokenizados $\vT$. Los codificadores de tripletes
generan la incrustación de forma correspondiente $\ve_{\mathrm{s}}$, la
$\ve_{\mathrm{i}}$ y la incrustación de texto $\ve_{\mathrm{t}}$ mediante la
proyección de los tokens de forma extraídos $\vE_{\mathrm{s}}$, los tokens de imagen
$\vE_{\mathrm{i}}$ y los tokens de texto $\vE_{\mathrm{t}}$ como tres
vectores de la misma dimensión, lo cual se expresa como:
$\ve_{\mathrm{s}} = \mathcal{F}_{\mathrm{s}}(\vE_{\mathrm{s}}),
\ve_{\mathrm{i}} = \mathcal{F}_{\mathrm{i}}(\vE_{\mathrm{i}})$, and
$\ve_{\mathrm{t}} = \mathcal{F}_{\mathrm{t}}(\vE_{\mathrm{t}})$,
donde $\mathcal{F}_{\mathrm{s}}$ es un proyector de incrustación de formas
aprendible, mientras que el proyector de incrustación de imágenes $\mathcal{F}_{\mathrm{i}}$ y
el proyector de incrustación de texto $\mathcal{F}_{\mathrm{t}}$ están preentrenados
y se mantienen fijos durante el entrenamiento y la inferencia. La pérdida contrastiva es:

\begin{align}
    &\mathcal{L}_{(\mathrm{s},\mathrm{i})} = - \frac{1}{2} \sum\limits_{(j,k)}
    \left(
        \log \frac{\exp (\ve_{\mathrm{s}}^j \ve_{\mathrm{i}}^k) }{
        \sum\limits_l \exp (\ve_{\mathrm{s}}^j \ve_{\mathrm{i}}^l) } +
        \log \frac{\exp (\ve_{\mathrm{s}}^j \ve_{\mathrm{i}}^k) }{
        \sum\limits_l \exp (\ve_{\mathrm{s}}^l \ve_{\mathrm{i}}^k) }
    \right), \\
    &\mathcal{L}_{(\mathrm{s},\mathrm{t})}  = - \frac{1}{2} \sum\limits_{(j,k)}
    \left(
        \log \frac{\exp (\ve_{\mathrm{s}}^j \ve_{\mathrm{t}}^k) }{
        \sum\limits_l \exp (\ve_{\mathrm{s}}^j \ve_{\mathrm{t}}^l) } +
        \log \frac{\exp (\ve_{\mathrm{s}}^j \ve_{\mathrm{t}}^k) }{
        \sum\limits_l \exp (\ve_{\mathrm{s}}^l \ve_{\mathrm{t}}^k) }
    \right),
\end{align}
donde $(j,k)$ indica el par positivo en los lotes de entrenamiento y
dado que se utilizan los codificadores preentrenados de CLIP, el modelo está exento
de la restricción $\mathcal{L}_{(\mathrm{i}, \mathrm{t})}$.

\textbf{El decodificador de formas 3D}, $\mathcal{D}_{\mathrm{s}}$, toma como entradas las
incrustaciones de formas $\vE_{\mathrm{s}}$ para reconstruir el campo neuronal 3D
con alta calidad. Se utiliza la pérdida de divergencia KL
$\mathcal{L}_{KL}$ para facilitar el proceso generativo y
mantener el espacio latente como una distribución continua.
% pre-fc, post-fc
Además, se utiliza una capa de proyección para comprimir la variable latente de
la dimensión $d$ a dimensiones inferiores $d_0$ y obtener así una representación compacta.
A continuación, se emplea otra capa de proyección para transformar la variable latente
muestreada de la dimensión $d_0$ de vuelta a la dimensión alta $d$ con el fin de
reconstruir los campos neuronales correspondientes de las formas 3D.
% decode neural field
Al igual que el codificador, el modelo decodificador también se basa en un transformer con el
mecanismo de atención cruzada. Dado un punto 3D de consulta $\vx \in \R^{3}$ en el campo
y sus correspondientes incrustaciones latentes de forma $\vE_{\mathrm{s}}$, el decodificador
calcula la atención cruzada de forma iterativa para predecir la ocupación del punto de consulta
$\mathcal{O} (\vx)$. La pérdida de entrenamiento se expresa como:

\begin{equation}
\mathcal{L}_r = \mathbb{E}_{\vx \in \R^3} [\text{BCE}(
\mathcal{D}(\vx \mid \vE_{\mathrm{s}} ), \mathcal{O}(\vx) )],
\end{equation}
donde BCE es la pérdida de entropía cruzada binaria, y la pérdida total para
el entrenamiento del Auto-codificador Variacional Alineado de Forma, Imagen y Texto (SITA) se expresa como:

\begin{equation}
\mathcal{L}_{SITA} = \lambda_c (\mathcal{L}_{(\mathrm{s},\mathrm{i})}
+ \mathcal{L}_{(\mathrm{s},\mathrm{t})}) + \mathcal{L}_r +
\lambda_{KL} \mathcal{L}_{KL}.
\label{eq:sita_vae}
\end{equation}

\paragraph{Modelo de difusión latente de formas alineadas.}
Tras entrenar el SITA-VAE, se obtiene un espacio de alineación entre formas 3D,
imágenes y textos, así como un codificador y
decodificador de formas 3D que comprime la forma 3D en representaciones latentes de formas de baja dimensión
y reconstruye dichas representaciones latentes en un campo neuronal
con alta calidad. Aprovechando el éxito del modelo de difusión latente
(LDM)~\cite{rombach2022high} en la generación de texto a imagen, que
logra un equilibrio entre la carga computacional y la calidad de la generación,
se propone un modelo de difusión latente de formas en el espacio alineado
para aprender un mapeo probabilístico más preciso de imágenes 2D o textos
a incrustaciones latentes de formas 3D. Al aprovechar el espacio de alineación y
el modelo de difusión latente de formas, es posible generar
formas 3D de alta calidad que se ajustan mejor a las entradas condicionales visuales o textuales.

El modelo de difusión latente de formas alineadas (ASLDM) se basa en un
transformer similar a UNet~\cite{10203178}, con el objetivo de ajustar una distribución
de las incrustaciones latentes de formas, acompañado de un auto-codificador para
codificar muestras de datos en el espacio latente y reconstruir las
muestras de datos a partir del latente muestreado. Al aprender en el espacio latente,
el modelo de difusión latente es computacionalmente eficiente, y
aprovechar una representación tan compacta permite al modelo ajustar la
distribución objetivo más rápidamente. En concreto, el modelo
$\epsilon_{\theta}$ se centra en generar incrustaciones latentes de forma
$\vE_{\mathrm{s}}$ condicionadas a $\vC$, que está representada por el
codificador de imágenes o texto CLIP. Siguiendo el LDM~\cite{rombach2022high}, el
objetivo es
\begin{equation}
\mathcal{L} = \mathbb{E}_{\vE_{\mathrm{s}}, \boldsymbol{\epsilon} \sim
\mathcal{N}(\Zero,\vI),t} \left[ \Vert \boldsymbol{\epsilon}
- \epsilon_\theta(\vE_{\mathrm{s}}^{(t)}, \vC, t) \Vert^2_2 \right],
\end{equation}
donde $t$ se muestrea uniformemente de $\{1, ..., T\}$ y
$\vE_{\mathrm{s}}^{(t)}$ es una versión borrosa de
$\vE_{\mathrm{s}}^{(0)}$. Durante la inferencia, tras muestrear ruido gaussiano,
el modelo elimina gradualmente el ruido de la señal hasta alcanzar
$\vE_{\mathrm{s}}^{(0)}$. Siguiendo la guía sin clasificador
(CFG)~\cite{ho2021classifierfree}, el modelo de difusión latente condicional
se entrena con guía sin clasificador. En la fase de entrenamiento,
la condición $\vC$ se convierte aleatoriamente en un conjunto vacío
$\emptyset$ con una probabilidad fija del 10\%. A continuación, el muestreo se
realiza mediante la combinación lineal de muestras condicionales e incondicionales:
\begin{equation}
\epsilon_\theta ( \vE_{\mathrm{s}}^{(t)}, \vC, t ) = \epsilon_\theta(
\vE_{\mathrm{s}}^{(t)}, \emptyset, t ) + \lambda
( \epsilon_\theta( \vE_{\mathrm{s}}^{(t)}, \vC, t ) -\epsilon_\theta(
\vE_{\mathrm{s}}^{(t)}, \emptyset, t ) ),
\label{eq:cfg-michelangelo}
\end{equation}
donde $\lambda$ es la escala de la guía para equilibrar la
fidelidad y la diversidad del muestreo.

\subsection{Experimentos y análisis}

\textbf{Métricas.}
La intersección sobre la unión (IoU) se utiliza para reflejar la precisión de
las reconstrucciones. Se proponen dos nuevas métricas para evaluar los métodos de generación de formas 3D.
La primera es una puntuación de forma-imagen (SI-S), que
utiliza un codificador de formas 3D y un codificador visual para extraer las correspondientes
incrustaciones de forma e imagen y calcular la similitud coseno de estas
dos modalidades. Otra es la puntuación de forma-texto (ST-S), que calcula
la similitud entre la forma 3D generada y el texto condicional
introducido en el espacio alineado de incrustaciones de formas e incrustaciones de texto.
Ambas métricas evalúan la similitud entre los resultados y sus
condiciones correspondientes. Además, tanto el modelo preentrenado
ULIP~\cite{10203465} como SITA se utilizan para calcular SI-S y ST-S, en
términos de SI-S (ULIP), ST-S (ULIP), SI-S (SITA) y ST-S (SITA),
respectivamente. Además, se siguen las métricas P-IS y P-FID presentadas en
Point-E~\cite{nichol2022point}, utilizando un modelo preentrenado
PointNet++~\cite{NIPS2017_d8bf84be} para calcular la puntuación de Inception
análoga a la nube de puntos~\cite{NIPS2016_8a3363ab} y
el FID~\cite{NIPS2017_8a1d6947} para evaluar la diversidad y la calidad de
las formas 3D generadas.

\textbf{Líneas de base.}
Comparamos el enfoque de Michelangelo con otros métodos, entre ellos
OccNet~\cite{8953655}, ConvOcc~\cite{10.1007/978-3-030-58580-8_31},
IF-Net~\cite{9156292}, 3DILG~\cite{NEURIPS2022_894ca1c4} y
Learnable Query Version of 3DS2V~\cite{10.1145/3592442},  en
tareas de reconstrucción para validar la capacidad del modelo para recuperar
un campo neuronal a partir de incrustaciones de formas en el conjunto de datos ShapeNet
~\cite{chang2015shapenet}. Para la etapa de generación condicional,
se han elegido como modelos de referencia dos potentes métodos recientes de generación 3D: 3DILG y 3DS2V.
En primer lugar, sus módulos de representación de formas se ajustan
con precisión utilizando un conjunto de datos mixto compuesto por ShapeNet y 3D Cartoon
Monster. A continuación, los modelos generativos condicionales de texto e imagen de
3DILG y 3DS2V se vuelven a entrenar siguiendo los mismos protocolos.

\textbf{Comparación numérica.}
Los resultados numéricos se recogen en
la Tabla~\ref{tab:michelangelo_recon} y la Tabla~\ref{tab:michelangelo_gen}.
La Tabla~\ref{tab:michelangelo_recon} muestra que Michelangelo tiene el
mejor rendimiento de reconstrucción en 55 categorías generales, superando
al resto. Los resultados de las categorías seleccionadas demuestran además que
el modelo puede reconstruir fielmente formas 3D en cada una de las 55 categorías.
La Tabla~\ref{tab:michelangelo_gen} presenta los resultados numéricos para la
generación condicional de formas 3D.
Michelangelo obtiene los mejores resultados en todos los indicadores SI-S y ST-S,
lo que indica que es capaz de mapear la información de la imagen o
del texto a la información de forma 3D correspondiente para generar
resultados de alta fidelidad. Además, el P-FID demuestra que el
modelo puede producir tokens de formas de alta calidad para generar
formas 3D realistas, y el P-IS indica la diversidad de las muestras.
En concreto, las cuatro columnas de la izquierda muestran que Michelangelo supera
los valores de referencia en la generación condicionada por imágenes, lo que demuestra que
es capaz de mapear mejor la información visual a formas 3D.
Las cuatro columnas de la derecha validan la calidad generativa de
la generación condicionada por texto. Dado que el lenguaje natural, en comparación con las imágenes 2D,
suele proporcionar información limitada y abstracta, entrenar un
modelo para mapear la información textual a formas 3D supone un reto. Sin embargo,
gracias al entrenamiento en el espacio latente alineado, Michelangelo
mejora significativamente la generación condicionada por texto, como se muestra en
las columnas de la derecha de la Tabla~\ref{tab:michelangelo_gen}, lo que refleja
que el modelo mapea eficazmente la información del lenguaje natural a formas 3D
y genera resultados diversos y de alta calidad.

\begin{table}[th]
\centering
\resizebox{1\columnwidth}{!}
{
    \begin{tabular}{rcccccccc}
        \hline
        & General & Seleccionadas & Mesa & Silla & Avión & Coche & Rifle
        & Lámpara   \\ \hline
        OccNet~\cite{8953655}
        & 0.825 & 0.81 & 0.823 & 0.803 & 0.835 & 0.911 & 0.755 & 0.735  \\
        ConvOccNet~\cite{10.1007/978-3-030-58580-8_31}
        & 0.888 & 0.873 & 0.847 & 0.856 & 0.881 & 0.921 & 0.871 & 0.859  \\
        IF-Net~\cite{9156292}
        & 0.934 & 0.924 & 0.901 & 0.927 & 0.937 & 0.952 & 0.914 & 0.914 \\
        3DILG~\cite{NEURIPS2022_894ca1c4}
        & 0.950 & 0.948 & 0.963 & 0.95 & 0.952 & 0.961 & 0.938 & 0.926 \\
        3DS2V(LQ)~\cite{10.1145/3592442}
        & 0.955 & 0.955 & \textbf{0.965} & 0.957 & 0.962 & 0.966 &
        0.947 & 0.931 \\
        Ours
        & \textbf{0.966} & \textbf{0.964} & \textbf{0.965} &
        \textbf{0.966} & \textbf{0.966} & \textbf{0.969} &
        \textbf{0.967} & \textbf{0.95} \\ \hline
    \end{tabular}
}
\caption{\textbf{Resultados numéricos para la comparación de reconstrucciones en
    IoU ($\uparrow$, cuanto mayor sea el valor, mejor).} Los resultados muestran que
    Michelangelo obtiene los mejores resultados en 55 categorías. Los resultados de algunas
    categorías seleccionadas demuestran además que el modelo puede reconstruir cada
categoría con fidelidad.}
\label{tab:michelangelo_recon}
\end{table}

\begin{table}[th]
\resizebox{1\columnwidth}{!}{%
    \centering
    \begin{tabular}{rcccccccc}
        \hline
        & \multicolumn{4}{c}{Image-Conditioned} &
        \multicolumn{4}{c}{Text-Conditioned} \\
        \cmidrule(lr){2-5}\cmidrule(lr){6-9}
        & SI-S (ULIP)$\uparrow$ & SI-S (SITA)$\uparrow$   &
        P-FID$\downarrow$ & P-IS$\uparrow$ & ST-S (ULIP)$\uparrow$ &
        ST-S (SITA)$\uparrow$  & P-FID$\downarrow$ & P-IS$\uparrow$ \\ \hline
        3DILG
        & 9.134  & 11.703 & 4.592 & 12.247 & 10.293 & 6.878  & 10.283 & 12.921\\
        3DS2V
        & 13.289 & 15.156  & 2.921 & 12.92 & 12.934 & 9.833  & 5.704 & 13.149\\
        Ours
        & \textbf{13.818} & \textbf{15.206}  & \textbf{1.586} &
        \textbf{13.233} & \textbf{16.647} & \textbf{13.128}  &
        \textbf{2.075} & \textbf{13.558}\\ \hline
    \end{tabular}
}
\caption{\textbf{Resultados numéricos para la generación condicional
    comparación}. Los resultados muestran que Michelangelo alcanza el mejor
    rendimiento generativo. Los índices SI-S y ST-S indican que el modelo
    genera resultados de alta fidelidad al asignar eficazmente la información de la condición
    a sus formas 3D correspondientes. Además, el P-FID indica que
    el modelo genera las formas 3D más realistas, y el P-IS indica
    que las muestras generadas son diversas. $\uparrow$ significa que un valor
más alto es mejor y $\downarrow$ de lo contrario.}
\label{tab:michelangelo_gen}
\end{table}

\textbf{Comparación visual}.
Las comparaciones visuales de las generaciones de formas 3D condicionadas por imágenes o texto
se ilustran en
la Figura~\ref{fig:michelangelo_img_cond_gen} y
la Figura~\ref{fig:michelangelo_text_cond_gen}.
La Figura~\ref{fig:michelangelo_img_cond_gen} muestra que
3DILG~\cite{NEURIPS2022_894ca1c4} presta más atención a la forma global
en el proceso de generación auto-regresiva, donde sus resultados
carecen de representaciones de los detalles de las formas 3D. Si bien
3DS2V~\cite{10.1145/3592442} genera formas 3D con mayor detalle, los
resultados presentan superficies discontinuas o con ruido. Además, ambos métodos
tienen dificultades para generar una forma completa cuando las condiciones dadas
corresponden a un objeto complejo, una máquina de precisión o un monstruo poco común.
La Figura~\ref{fig:michelangelo_text_cond_gen} muestra la comparación visual
de la generación condicionada por texto. En las filas de la mitad superior,
se muestran los resultados obtenidos con conceptos simples y abstractos, mientras que en
las filas de la mitad inferior, se muestran los resultados obtenidos con textos detallados, como
descripciones de partes determinadas de la forma objetivo.
De manera similar a la observación anterior, 3DILG~\cite{NEURIPS2022_894ca1c4}
genera superficies de forma excesivamente suaves con menos detalles, y
3DS2V~\cite{10.1145/3592442} produce menos detalles en superficies de objetos discontinuas.
Por el contrario, Michelangelo genera formas correctas
que se ajustan a los conceptos o descripciones detalladas proporcionados, con
detalles delicados sobre superficies lisas.

\begin{figure}[th]
\centering
\includegraphics[width=.98\linewidth]{\toplevelprefix/chapters/chapter8/figs/michelangelo_img_cond_gen.pdf}
\caption{
    \textbf{Resultados visuales de la comparación de la generación condicionada por imágenes}. 
    La figura muestra que 3DILG~\cite{NEURIPS2022_894ca1c4} genera superficies excesivamente lisas
    y carece de detalles de formas, mientras que 3DS2V~\cite{10.1145/3592442}
    genera pocos detalles con superficies ruidosas y discontinuas de
    las formas. A diferencia de las líneas de base, Michelangelo produce superficies suaves
    y representa los detalles de las formas. Por favor, amplíe la imagen para ver más detalles visuales.
}
\label{fig:michelangelo_img_cond_gen}
\end{figure}

\begin{figure}[th]
\centering
\includegraphics[width=.98\linewidth]{\toplevelprefix/chapters/chapter8/figs/michelangelo_text_cond_gen.pdf}
\caption{
    \textbf{Resultados visuales de la comparación de generación condicionada por texto}.
    En las dos primeras filas, los modelos se prueban con textos abstractos, y
    el resultado muestra que solo Michelangelo genera una forma 3D que
    se ajusta al texto de referencia, con una superficie lisa y detalles precisos.
    Las dos últimas filas muestran los resultados obtenidos con textos que contienen descripciones detalladas,
    lo que demuestra aún más que Michelangelo es capaz de
    captar tanto la información condicional global como la información local
    para generar formas 3D de alta fidelidad. Las palabras clave aparecen
resaltadas en rojo; amplíe la imagen para ver más detalles.}
\label{fig:michelangelo_text_cond_gen}
\end{figure}

\textbf{La eficacia del entrenamiento del modelo generativo en el espacio alineado}. 
Se realiza una comparación visual para evaluar la eficacia
del entrenamiento del modelo generativo en el espacio alineado, tal y como se ilustra
en la Figura~\ref{fig:michelangelo_ablate_align}. Las muestras superiores se
generan a partir del modelo generativo entrenado en el espacio alineado,
mientras que las muestras inferiores se generan a partir del modelo generativo
entrenado sin el espacio alineado. Los resultados demuestran que las
muestras superiores se ajustan al texto dado, mientras que las inferiores
no, lo que indica que entrenar el modelo generativo en el
espacio alineado da lugar a muestras de alta fidelidad.

\begin{figure}[h]
\centering
\includegraphics[width=.98\linewidth]{\toplevelprefix/chapters/chapter8/figs/michelangelo_1aligned_space_for_generation.pdf}
\caption{
    \textbf{Estudio de ablación sobre la eficacia del entrenamiento del modelo generativo
    en el espacio alineado}. Esta figura ilustra comparaciones visuales
    para los estudios de ablación sobre la eficacia del entrenamiento
    del modelo generativo en el espacio alineado. En comparación con las
    muestras inferiores basadas en los textos condicionales, las muestras superiores
    se acercan más a las condiciones desde el punto de vista semántico, lo que indica la
    eficacia del entrenamiento del modelo generativo en el espacio alineado.
}
\label{fig:michelangelo_ablate_align}
\end{figure}

\textbf{La eficacia de los modelos de visión y lenguaje}. Además del
conocido modelo de visión y lenguaje (VLM)
CLIP~\cite{pmlr-v139-radford21a}, se introduce otro modelo de visión-lenguaje (VLM)
SLIP~\cite{10.1007/978-3-031-19809-0_30} para entrenar
el SITA-VAE con el fin de realizar una comparación exhaustiva. En primer lugar, se evalúa el impacto del
modelo de visión-lenguaje en la capacidad de reconstrucción de SITA-VAE,
y los resultados se muestran en la
Figura~\ref{fig:michelangelo_ablate_vlm}. Los resultados muestran que el
modelo creado con CLIP alcanza el mejor rendimiento. A continuación, se
evalúa el impacto del modelo de visión-lenguaje en la capacidad de
alinear el espacio multimodal. Se seleccionan tareas de clasificación
estándar y de aprendizaje sin datos previos para reflejar el impacto
del modelo de visión-lenguaje. Tenga en cuenta
que la clasificación se lleva a cabo mediante una operación de coincidencia de características,
en la que se proporcionan múltiples formas 3D y frases al SITA-VAE; este
devuelve la similitud entre las formas 3D y cada frase como
resultados de clasificación, lo que indica que cuanto más alineado esté
el espacio multimodal, mayor será la precisión de la clasificación. Los
resultados muestran que el modelo compuesto con CLIP alcanza el mejor rendimiento.

\textbf{El impacto de las incrustaciones de consultas aprendibles}. Se lleva a cabo un estudio de ablación
de las incrustaciones de consultas aprendibles utilizando los mismos
experimentos que los anteriores y los resultados muestran que el uso de 512
incrustaciones de consultas aprendibles ofrece el mejor rendimiento en las tareas de reconstrucción y
clasificación.

\begin{figure}
\centering
\includegraphics[width=1.\linewidth]{\toplevelprefix/chapters/chapter8/figs/michelangelo_2aligned_space_for_generation.pdf}
\caption{
    \textbf{Estudio de ablación sobre la eficacia de los modelos de visión-lenguaje
    y el impacto de las incrustaciones de consultas aprendibles}. Esta figura
    muestra el estudio de ablación sobre la eficacia del
    modelo de visión-lenguaje y el impacto de las incrustaciones de consultas
    aprendibles. Según la tabla, el modelo con CLIP y 512
    incrustaciones de consulta aprendibles logra el mejor rendimiento de reconstrucción y
    clasificación, lo que indica su capacidad para recuperar formas 3D
    y alinear el espacio multimodal.
}
\label{fig:michelangelo_ablate_vlm}
\end{figure}

\paragraph{Análisis del vecino más cercano.}
Para verificar si el modelo aprende a generar resultados basándose en las
condiciones dadas o si memoriza la partición de entrenamiento del conjunto de datos, se
recorre el conjunto de entrenamiento para recuperar los vecinos más cercanos de las
muestras generadas, tal y como se ilustra en la
Figura~\ref{fig:michelangelo_nearestneighbor}. En concreto, se muestran los tres
vecinos más cercanos de cada muestra a efectos comparativos.
Según la visualización, las formas 3D generadas difieren de
cada uno de los vecinos más cercanos recuperados, lo que demuestra que el modelo
crea nuevas formas 3D en lugar de sufrir un sobreajuste a los conjuntos de entrenamiento.

\begin{figure}[th]
\centering
\includegraphics[width=1.\linewidth]{\toplevelprefix/chapters/chapter8/figs/michelangelo_nearestneighbor.pdf}
\caption{
    \textbf{Análisis del vecino más cercano}. Recorremos todo el
    conjunto de entrenamiento para encontrar los tres vecinos más cercanos de las formas 3D
    generadas, y los resultados revelan que nuestro modelo es capaz de producir formas 3D
    nuevas a partir de imágenes dadas, en lugar de memorizar formas específicas.
}
\label{fig:michelangelo_nearestneighbor}
\end{figure}

\paragraph{Discusión.}
Aunque el proyecto Michelangelo ha demostrado que podemos lograr
resultados excelentes en la generación de objetos 3D, aún presenta algunas
limitaciones. En primer lugar, requiere una gran cantidad de muestras de
formas 3D de referencia para el entrenamiento y el aprendizaje de la distribución, mientras que
la obtención de esos datos de formas 3D es mucho más costosa y los conjuntos de datos
disponibles suelen ser, por lo general, varios órdenes de magnitud más pequeños que los de las
imágenes 2D. Como hemos mencionado en la \Cref{sec:3D-visual-model} al
final del capítulo anterior, nuestro cerebro es capaz de aprender un modelo visual completo en 3D (o
4D) partiendo únicamente de imágenes en 2D, ya que este tipo de datos son mucho
más naturales y económicos de obtener. Por lo tanto, en última instancia, para ampliar verdaderamente
el modelo 3D generativo, debemos desarrollar métodos más avanzados
que nos permitan aprender la representación de la forma (dentro de un espacio 3D
alineado entre forma, imagen y texto) a partir únicamente de imágenes 2D (multivista), por ejemplo,
mediante renderización diferenciable.
Además, dado que representamos cada forma 3D como un campo de ocupación,
es necesario convertir la malla 3D en una malla cerrada, lo que
inevitablemente degradará la calidad original de la malla 3D.

\section{Reconstrucción de formas y poses en 3D basada 
en la generación}\label{sec:CUPID}\label{sec:shape-and-pose}

% refactor
\newcommand{\pose}{\boldsymbol{\theta}}
\newcommand{\object}{\mathcal{O}}
\newcommand{\Image}{\boldsymbol{I}}
\newcommand{\image}{\vI}
% \newcommand{\mask}{\boldsymbol{M}^{\textrm{occ}}}
\newcommand{\render}{\boldsymbol{\mathcal{P}}}
\newcommand{\latent}{\boldsymbol{z}}
\newcommand{\timestep}{t}
\newcommand{\latentt}{\boldsymbol{z}_{\timestep}}
\newcommand{\latentzero}{\boldsymbol{z}_{0}}
\newcommand{\noise}{\boldsymbol{\epsilon}}
\newcommand{\velocity}{\boldsymbol{v}}
\newcommand{\parameter}{\phi}
\newcommand{\encoder}{\varphi}
\newcommand{\point}{\boldsymbol{x}}
\newcommand{\pixcoord}{\boldsymbol{u}}
\newcommand{\feat}{\boldsymbol{f}}
\newcommand{\sflow}{\boldsymbol{\mathcal{G}_{S}}}
\newcommand{\lflow}{\boldsymbol{\mathcal{G}_{L}}}
\newcommand{\transpose}{\operatorname{T}}
\newcommand{\Grid}{\mG}
\newcommand{\sgrid}{\mS}
\newcommand{\numpts}{L}
% \newcommand{\res}{N}
\newcommand{\sres}{D}
\newcommand{\vaeencoder}{\psi}
\newcommand{\trellis}{{\textsc{TRELLIS}}\xspace}
% \newcommand{\dino}{\textsc{DINO}\xspace}
\newcommand{\interp}{\operatorname{BilinearInterp}}
\newcommand{\slatenc}{\operatorname{SlatEncoder}}
\newcommand{\twotothree}{$2$D $\to$ $3$D\xspace}
\newcommand{\threetotwo}{$2$D $\to$ $3$D\xspace}
\newcommand{\low}{\mathrm{l}}
\newcommand{\high}{\mathrm{h}}



% \paragraph{Generation-based 3D Reconstruction}
La razón fundamental por la que desarrollamos el sentido de la vista es para aprender a percibir 
nuestro entorno tridimensional (3D). Nuestro cerebro ha dedicado una parte importante de sí mismo, 
{\em la corteza visual}, a realizar esta tarea. Esto nos permite reconstruir de manera precisa y 
eficiente las formas de los objetos 3D y sus relaciones espaciales
\footnote{es decir, las posiciones relativas con respecto a nuestro punto de vista o entre los 
propios objetos} a partir de las imágenes 2D que perciben nuestros ojos. Existe una larga y rica 
tradición dedicada a comprender matemáticamente cómo, en principio, podemos reconstruir 
completamente la geometría tridimensional a partir de múltiples imágenes 2D. 
Los lectores interesados pueden consultar el libro de texto \cite{MaY2003} para obtener 
una explicación detallada de este importante tema.

\begin{figure}[th]
\centering
\includegraphics[width=0.8\linewidth]{chapters/chapter8/figs/cupid_definition.pdf}
\caption{\textbf{Tarea: inferir la forma 3D de un objeto con una
    postura canónica respecto a la postura de la cámara a partir de una sola imagen.}
    Izquierda: Imagen de entrada. Centro: El resultado recupera un modelo 3D canónico
    del objeto (con propiedades intrínsecas como la forma y
    la textura) junto con la postura exacta de la cámara (propiedad extrínseca)
    que genera la imagen de entrada.
    Derecha: La nueva representación del modelo 3D a partir de la pose estimada
reproduce fielmente la imagen de entrada.}
\label{fig:cupid_teaser}
\end{figure}

No obstante, tal y como hemos mencionado en la Sección \ref{sec:3D-visual-model} del 
capítulo anterior \ref{ch:conditional-inference}, idealmentela reconstrucción 
3D debería llevarse a cabo de manera bayesiana, al igual que lo hace nuestro 
cerebro. Nuestra percepción del entorno tridimensional parece ser tan natural, principalmente 
porque reconstruimos una escena tridimensional y reconocemos objetos tridimensionales 
basándonos en una rica «memoria visual». Más concretamente, gracias a nuestros años de 
experiencia acumulada, hemos aprendido a comprender la distribución de nuestro entorno 
3D y de los objetos que contiene. Dicha distribución es bastante simple y estructurada, 
por lo que podemos reconstruir de forma correcta y eficiente la geometría 3D de una escena 
y de los objetos que la componen, incluso si solo hemos visto una pequeña parte de ella, 
por ejemplo, desde una única perspectiva, tal y como ilustra el ejemplo de la 
Figura \ref{fig:cupid_teaser}. Desde un punto de vista conceptual, esta capacidad de 
«completar» la información que falta es muy similar a la forma en que las imágenes 
naturales pueden completarse a partir de una pequeña fracción de parches, tal como 
hemos visto en la Sección \ref{sec:image_completion}.

Como distribución, nuestra memoria visual tridimensional no solo es de baja dimensión, 
sino que también está muy estructurada. Por ejemplo, parece que nuestro cerebro almacena 
la semántica (la identidad de los objetos) de una escena en la corteza temporal inferior 
(corteza IT) y la información espacial tridimensional de la escena en el hipocampo. Por 
ejemplo, parece que nuestro cerebro almacena la semántica (la identidad de los objetos) 
de una escena en la corteza temporal inferior (corteza IT) y la información espacial 
tridimensional de la escena en el hipocampo. Nuestro cerebro no solo reconstruye la 
geometría tridimensional de la escena,\footnote{Por ejemplo, la geometría 3D que se 
representa en formas diferentes, pero en gran medida equivalentes: nubes de puntos, 
mallas, funciones de distancia con signo, campos de radiancia neuronal o salpicaduras 
gaussianas.} sino que también desglosa su contenido en objetos individuales y analiza 
sus relaciones espaciales con el observador y entre los propios objetos. Esta 
representación estructurada de «qué está dónde» nos permite realizar inferencias 
espaciales dentro de la escena desde una perspectiva {\em egocéntrica, objetocéntrica 
o alocéntrica}. Estas capacidades son necesarias y cruciales para tareas como orientarse 
en una escena, interactuar con los objetos que contiene y manipularlos.

Por lo tanto, si el objetivo de aprender un modelo de la escena 3D es facilitar la 
interacción física con un entorno y los objetos que contiene,\footnote{lo cual es el 
objetivo de la denominada «IA incorporada» para agentes inteligentes o robots}, debemos 
aprender una representación del entorno 3D que comparta las mismas propiedades y 
características que nuestra memoria visual mencionada anteriormente.


\subsection{Tarea y objetivo}
\label{sub:cupid_objective}

\paragraph{Forma y posturas de objetos 3D a partir de una sola vista 2D}
En esta sección, damos un primer paso para mostrar cómo se puede
lograr esto mediante la ilustración de un problema complejo y aparentemente mal planteado:
{\em inferir la forma y la pose 3D completas, incluyendo
tanto la pose centrada en el observador como aquella centrada en el objeto, a partir de una sola
imagen 2D.} Este problema, por sí mismo, se encuentra en el corazón de la visión
computacional, con el objetivo de revertir los procesos físicos de formación de imágenes. Una imagen
$\image$ es el resultado de un proceso de renderización,
\[
\image = \render(\object, \pose),
\]
donde $\object$ representa el objeto 3D en un sistema de coordenadas canónico,
independiente de la vista, y $\pose$ denota el punto de vista de la cámara. Si solo se dispone
de $\image$, la tarea de separar las propiedades intrínsecas del objeto
(por ejemplo, la forma o la textura) de la posición extrínseca de la cámara es
fundamentalmente ambigua.

Los enfoques anteriores suelen simplificar en exceso esta ambigüedad. Los modelos generativos 
suelen dar prioridad al objeto canónico $\object$ en detrimento de la pose de la cámara $\pose$, 
lo que da lugar a reconstrucciones que no se alinean fielmente con la vista de entrada. Por el 
contrario, los métodos tradicionales centrados en la vista fijan la pose, lo que limita su 
capacidad para generar o «rellenar» las regiones ocluidas e integrar a la perfección el 
objeto reconstruido en la imagen original.

A continuación presentamos un enfoque conocido como \textbf{Cupid}~\citep{huang2025cupid}:
\\
\centerline{\url{https://cupid3d.github.io},}\\ 
un marco de código abierto que modela conjuntamente la distribución del objeto 3D y 
de la pose de la cámara (Figura~\ref{fig:cupid_teaser}). Al generar simultáneamente el 
modelo 3D y realizar un razonamiento sobre el punto de vista, Cupid garantiza que el resultado 
se pueda \textit{pegar} con precisión en la imagen original. Este enfoque replantea la 
reconstrucción a partir de una sola vista como un problema de muestreo condicional, logrando 
la fidelidad geométrica mediante la combinación de datos observacionales con priores generativos 
aprendidos, lo que permite una reconstrucción eficiente tanto a nivel de objetos como de 
escenas (Figura~\ref{fig:cupid_results}).

\begin{figure}
    \centering
    \includegraphics[width=0.95\textwidth]{chapters/chapter8/figs/cupid/figures/results.pdf}
    \caption{\textbf{Resultados de la \textit{reconstrucción 3D generativa} a partir de una sola imagen de prueba.} A partir de una imagen de entrada (arriba a la izquierda), Cupid estima la orientación de la cámara (abajo a la izquierda) y reconstruye un modelo 3D (abajo a la derecha), volviendo a renderizar la imagen de entrada (arriba a la derecha). Es robusto ante cambios de escala, ubicación e iluminación, al tiempo que conserva los detalles más sutiles, y admite la reconstrucción de escenas con alineación de componentes (fila inferior). Todos los resultados se obtienen en segundos mediante el muestreo feed-forward del modelo entrenado. }
    \label{fig:cupid_results}
    \vspace{-5mm}
\end{figure}

\paragraph{Formulación de la tarea y objetivo.}
Formulamos la reconstrucción generativa en 3D como la tarea de estimar
la distribución a posteriori conjunta
\[
p(\object, \pose \mid \image)
\]
bajo la restricción de que la observación $\image$ sea una representación
del objeto $\object$ desde la pose $\pose$. Para modelar esta distribución
condicional, empleamos un modelo generativo basado en flujos. Concretamente,
primero mapeamos el objeto 3D y la pose de la cámara a una representación
latente volumétrica unificada
\(\latent = \encoder(\object, \pose)\).
A continuación, utilizamos un modelo Rectified Flow~\citep{lipmanflow} para aprender la generación condicional
de esta variable latente. El modelo define una interpolación lineal
entre una muestra de la distribución de datos,
$\latentzero$, y una muestra de una distribución de ruido, $\noise$, a lo largo de
un paso de tiempo normalizado $\timestep \in [0,1]$:
\begin{equation}
\latentt = (1 - \timestep) \latentzero + \timestep \noise.
\end{equation}
El proceso generativo consiste en invertir esta trayectoria mediante el aprendizaje
de un campo de velocidad dependiente del tiempo
\[
\velocity(\latentt, \image, \timestep) = \nabla_{\timestep} \latentt
\]
que transporta las muestras ruidosas hacia la variedad de datos, condicionado
por la imagen de entrada $\image$. Parametrizamos este campo de velocidad con
una red neuronal $\velocity_{\parameter}$ y la entrenamos utilizando el
objetivo de coincidencia de flujo condicional (CFM)~\cite{lipmanflow}:
\begin{equation}
\mathcal{L}_{\mathrm{CFM}}(\parameter) \doteq \mathbb{E}_{\timestep,
\latentzero, \noise} \left\| \velocity_{\parameter}(\latentt, \image,
\timestep) - (\noise - \latentzero) \right\|_{2}^{2}.
\end{equation}
Al tomar muestras de este modelo entrenado, podemos generar un código latente
$\latent$ que luego puede decodificarse para obtener el objeto 3D final y la pose de la cámara.

\subsection{Representación y tokenización de formas 3D}

Uno de los principales retos en la generación de formas 3D es encontrar una tokenización que 
sea a la vez compacta y expresiva, capaz de permitir una reconstrucción de alta fidelidad. Si 
bien \textbf{Michelangelo}~\citep{zhao2023michelangelo} utiliza una representación vectorial 
latente aprendida---que codifica las formas en una secuencia de 1024 tokens---, estos vectores 
carecen de coordenadas espaciales 3D explícitas. En este proyecto, buscamos una representación 
latente que conserve la información posicional 3D explícita, lo que nos permite alinear los 
tokens de forma directamente con los píxeles de la imagen. Al establecer esta correspondencia 
espacial, podemos aprovechar la información local de los píxeles para mejorar los detalles de 
la generación. En consecuencia, \textbf{Cupid} adopta un enfoque de representación latente y 
tokenización basado en vóxeles 3D, tal como se describe en la Sección \ref{sec:3d_tokenization}. 
Aquí representamos una forma mediante vóxeles dispersos acompañados de características, que 
suelen definirse mediante el color RGB o características DINO más avanzadas~\citep{DINOv2} 
procedentes de modelos base 2D. A continuación, entrenamos un decodificador para que mapee 
dicha representación de vuelta a la forma, parametrizada por la SDF y a la apariencia, 
parametrizada por los splats gaussianos~\citep{kerbl20233d}.

En concreto, dado un conjunto de características de vóxeles dispersas 
$\boldsymbol{F} = \{\boldsymbol{f}_i\}_{i=1}^N$ ubicadas en las coordenadas 
$\boldsymbol{x}_i$, empleamos dos decodificadores distintos, 
$\mathcal{D}_{geo}$ y $\mathcal{D}_{app}$, para separar la geometría de la 
apariencia. El proceso de decodificación se formula como:
\begin{equation}
s_i = \mathcal{D}_{geo}(\boldsymbol{f}_i), \quad \mathcal{G}i = \mathcal{D}_{app}(\boldsymbol{f}_i),
\end{equation}
donde $s_i \in \mathbb{R}$ representa el valor del campo de distancia con signo (SDF) utilizado para la reconstrucción de la superficie y $\mathcal{G}_i = \{\boldsymbol{c}_i, \alpha_i, \boldsymbol{s}_i, \boldsymbol{q}_i\}$ denota el conjunto de parámetros gaussianos 3D---que comprenden color, opacidad, escala y rotación---requeridos para el renderizado volumétrico. A continuación, utilizaremos las características de vóxeles dispersos como representación latente para la coincidencia de flujos latentes~\citep{rombach2022high}.

\subsection{Arquitectura e implementación}
\label{sub:3d_reconstruction_architecture}
Dado que utilizamos vóxeles 3D para la representación de objetos, aplicar la difusión de 
manera simplista da lugar a un consumo excesivo de memoria, ya que una gran parte de la 
malla está compuesta por espacio vacío innecesario, tal y como se describe en la 
Sección~\ref{sec:3d_tokenization}. Para mitigar esto, la arquitectura \textbf{Cupid} está 
diseñada como un modelo de flujo en cascada de dos etapas que refina progresivamente la 
reconstrucción 3D~\citep{xiang2025structured}. La primera etapa genera una forma geométrica 
aproximada y estima la orientación de la cámara, mientras que la segunda etapa genera una 
geometría y una apariencia de alta fidelidad. Una innovación clave de este marco radica en 
la representación explícita de la orientación de la cámara y su integración en el proceso 
generativo. El flujo de trabajo general se muestra en la \Cref{fig:cupid_architecture}.

\paragraph{Representación desacoplada de objetos y poses.}

Como se explica en la Sección~\ref{sec:3D representation}, el objeto 3D $\object$ se 
tensoriza en una representación basada en vóxeles dispersos:
\begin{equation}
    \object \triangleq \{\point_i, \feat_i\}_{i=1}^{\numpts},
\end{equation}
donde $\point_i \in \mathbb{R}^3$ denota la coordenada de un vóxel activo, y $\feat_i$ es 
un vector de características que codifica la geometría y la apariencia locales. Estas 
características se obtienen mediante la compresión de una malla 3D densa de características 
DINO~\citep{oquab2023dinov2} de múltiples vistas a través de una VAE 3D. Posteriormente, el 
decodificador VAE se entrena para reconstruir $\object$ en un campo de distancia con signo 
(SDF) para la extracción de mallas. Esta representación latente en 3D se ilustra en la parte 
superior del bloque azul central de la Figura~\ref{fig:cupid_architecture}.

La pose de la cámara $\pose$, representada por su matriz de proyección
$\boldsymbol{P} = \boldsymbol{K}[\boldsymbol{R}|\boldsymbol{t}]$, se reparametriza
de una manera novedosa para facilitar la generación conjunta. En lugar de un
vector compacto, representamos la pose como un campo denso de correspondencias de 3D a 2D,
lo que denominamos una \textit{distribución intracubo de píxeles}:
\[
\pose \triangleq \{\point_i, \pixcoord_i\}_{i=1}^{\numpts},
\]
donde $\pixcoord_i = (u_i, v_i) \in [0,1]^2$ son las coordenadas normalizadas en 2D
del píxel correspondientes al centro del vóxel 3D $\point_i$.
Estas correspondencias se obtienen mediante una proyección en perspectiva:
\[
\pixcoord_i = \pi(\boldsymbol{P}, \point_i).
\]
Este conjunto de correspondencias se puede visualizar como un cubo UV en 3D, en el que
los valores $(u,v)$ actúan como colores que dependen de la perspectiva en una cuadrícula tridimensional. Esto
se ilustra en la parte inferior del bloque azul situado en el centro de
la Figura \ref{fig:cupid_architecture}.

Dado un conjunto generado de correspondencias, la matriz global de la cámara
$\boldsymbol{P}^{*}$ se recupera resolviendo un problema de «Perspectiva de n puntos»
(PnP) mediante un algoritmo de mínimos cuadrados~\citep{abdel2015direct}:
\begin{equation}
\boldsymbol{P}^{*} = \argmin_{\boldsymbol{P}} \sum_{i=1}^{\numpts}
\big\|\pi(\boldsymbol{P},\point_i) - \pixcoord_i\big\|^2.
\label{eq:pnp}
\end{equation}
Esta formulación transforma de manera elegante el problema conjunto de la generación de la pose de un objeto
en uno de generación de una forma/apariencia 3D centrada en el objeto
con una pose centrada del observador/cámara asociada al
objeto (a través de las coordenadas UV). Esto se ilustra en la
esquina inferior izquierda del bloque rojo situado a la derecha de la Figura
\ref{fig:cupid_architecture}.


\begin{remark}[Modelado 3D estructurado]
Obsérvese que las representaciones de tipo vóxel elegidas aquí para la forma y
la pose son diferentes de las representaciones de tipo nube de puntos adoptadas
en la tarea de generación de formas condicionadas estudiada en la
\Cref{sec:Michelangelo} anterior. En esa tarea, el objetivo es generar un
objeto que solo tenga que estar fuertemente correlacionado con la imagen
(o el texto) dada. No requiere una alineación precisa de la pose entre el
objeto (en su pose canónica) y la imagen dada (punto de vista). Sin embargo, para
la tarea que nos ocupa, queremos «reconstruir» un
objeto cuya forma, apariencia y postura coincidan fielmente con la
imagen dada. Por lo tanto, es necesario que la representación interna
modele explícitamente la forma canónica del objeto y
su pose relativa de manera simultánea. Tenga en cuenta que tal «desacoplamiento» de la
forma/apariencia (canónica) de un objeto y la pose desde la que se
observa imita de cerca las funciones antes mencionadas de la corteza IT
y el hipocampo en nuestro cerebro, respectivamente. Como veremos más adelante,
una representación desacoplada de este tipo nos permitirá obtener un modelo 3D
estructurado y coherente de una escena que puede contener múltiples objetos
desde múltiples perspectivas. Sostenemos que un modelo 3D así estructurado es
necesario para que podamos llevar a cabo razonamientos centrados en la vista o en el objeto
,\footnote{Por ejemplo, responder preguntas tales como «¿qué
    objeto está delante, detrás, a la izquierda o a la derecha de, o encima de
qué objeto?»} generar y interactuar con la escena.
\end{remark}

\begin{figure*}[t]
\centering
\includegraphics[width=\textwidth]{chapters/chapter8/figs/cupid_method.pdf}
\caption{\textbf{El proceso de reconstrucción generativa en dos etapas de Cupid}. 
    A partir de una imagen de entrada, la primera etapa ($\sflow$) genera
    un cubo de ocupación aproximado y un cubo UV, que codifica las correspondencias de 3D a 2D
    . Un solucionador PnP recupera la pose de la cámara
    $\boldsymbol{P}^{*}$ a partir de estas correspondencias. La segunda etapa
    ($\lflow$) se basa en esta pose recuperada. Inyecta
    características alineadas con los píxeles (extraídas de DINOv2 y de mapas de características
    de bajo nivel) en el proceso generativo para sintetizar una geometría y una apariencia de alta fidelidad
para el objeto 3D final.}
\label{fig:cupid_architecture}
\end{figure*}


\subsection{Modelado de flujo en cascada}
\label{sec:flow_model}

Utilizamos un modelo de flujo en cascada de dos etapas para muestrear conjuntamente un objeto 
3D y su pose de cámara correspondiente, denotada como 
$\latent = \{ (\point_i, \feat_i), (\point_i,\pixcoord_i) \}_{i=1}^{\numpts}$. En la primera 
etapa, el modelo generador de ocupación y pose ($\sflow$) produce dos resultados clave: 
(i) un cubo de ocupación que identifica los vóxeles activos y (ii) un cubo UV que codifica 
la pose de la cámara desde la perspectiva del objeto. En la segunda etapa, a partir de estas 
predicciones, el modelo de geometría y apariencia alineado con la pose ($\lflow$) sintetiza 
las características DINO~\citep{oquab2023dinov2} $\feat_i$ para cada vóxel activo, con el 
fin de obtener la variable latente final $\latent$.

\paragraph{Ocupación y generación de poses.}
Dada una imagen de condicionamiento \(\image\), esta etapa genera un cubo de ocupación y un 
cubo UV, ambos con una resolución de $r$. El cubo de ocupación, 
\(\boldsymbol{G}_o \in \{0,1\}^{r \times r \times r \times 1}\), contiene valores binarios 
que distinguen los vóxeles activos de los inactivos. El cubo UV, 
\(\boldsymbol{G}_{uv} \in [0,1]^{r \times r \times r \times 2}\), almacena las coordenadas 
de píxeles normalizadas $(u,v)$ para cada vóxel; cabe destacar que los vóxeles que comparten 
coordenadas de píxel idénticas residen sobre el mismo haz (ray) de la cámara.\footnote{No 
    modelamos explícitamente la oclusión, ya que estudios recientes sugieren que los 
    transformers modelan eficazmente el transporte de la luz de forma 
implícita~\citep{jin2024lvsm, zeng2025renderformer}.}

Para mejorar la eficiencia computacional, entrenamos una VAE 3D para comprimir el cubo UV en 
una malla de características de baja resolución 
\(\boldsymbol{S}_{uv} \in \mathbb{R}^{s\times s\times s \times C}\), logrando una recuperación 
de la pose con pérdidas mínimas (RRE/RTE medio $< 0,5^\circ$). Para afinar a TRELLIS, 
concatenamos la malla de características original 
\(\boldsymbol{S}_o\) con \(\boldsymbol{S}_{uv}\) e introducimos capas de proyección lineal 
tanto en la entrada como en la salida de la red de flujo $\sflow$. Una vez generados los cubos 
de ocupación y UV, extraemos el conjunto 
$\{ (\point_i, \pixcoord_i(\pose)) \}_{i=1}^{\numpts}$ recopilando vóxeles activos y resolviendo 
la pose de la cámara mediante la ecuación~\ref{eq:pnp}.

% {\color{blue} Updated: We have included 3D tokenization in section \ref{sec:3d_tokenization}.}
% \yima{How these voxel-type 3D data are tokenized. Discuss
% relationship or difference from the Perceiver for tokenizing point
% clouds in the Michelangelo project. Also, describe basic deep
% architectures used for the two flow-matching based generative models learned. }


\paragraph{Generación de geometría y apariencia alineadas a la pose.}
En la segunda etapa, generamos características latentes 3D detalladas 
$\{ \feat_i \}_{i=1}^{\numpts}$ únicamente en las ubicaciones de los vóxeles activos. 
Los experimentos indican que el modelo estándar 
$\lflow$ de TRELLIS~\citep{xiang2025structured}, que se basa en información de imagen con 
atención global, está sujeto a deriva de color y a la pérdida de detalles finos.
Mitigamos esto utilizando la pose de la cámara derivada en la primera etapa para inyectar 
información de píxeles con atención local en cada vóxel.

En concreto, calculamos las características del vóxel $i$ utilizando la pose estimada de la siguiente manera:
\begin{equation}
\begin{aligned}
\feat^\dino_{i} &= \interp(\pixcoord_i, \dino(\image)) \in \mathbb{R}^{1024}, \\
\{\feat^{\high}_{i}\}_{i=1}^{\numpts} &= \slatenc\big(\{\point_i, \feat^\dino_i\}_{i=1}^{\numpts}\big), \quad \feat^{\high}_{i} \in \mathbb{R}^{8},
\end{aligned}
\label{eq:latent}
\end{equation}
donde $\pixcoord_i$ representa la proyección del centro del vóxel 3D $i$-ésimo sobre el 
plano de la imagen, $\interp$ representa la interpolación bilineal y $\slatenc$ es el codificador 
3D VAE. Aunque DINO~\citep{oquab2023dinov2} capta la semántica de alto nivel, frecuentemente carece 
de las pistas de bajo nivel necesarias para una reconstrucción precisa. Para compensar esto, 
extraemos características complementarias de bajo nivel $\feat^{\low}$ de $\image$ utilizando 
un núcleo convolucional ligero, y las muestreamos en $\pixcoord_i$ mediante $\interp$. 
Finalmente, en cada paso temporal $t$, fusionamos la característica de vóxel ruidosa
$\feat_i^{t}$ con las características alineadas con los píxeles a través de una capa lineal 
antes de introducirlas en el transformer de flujo:
\begin{equation}
l_t = \mathrm{Linear}\big([\feat_i^{t}\,\oplus\, \feat^{\high}_{i} \,\oplus\, \feat^{\low}_{i}]\big).
\end{equation}
Esta estrategia de fusión basada en la alineación de poses mejora significativamente tanto la precisión 
geométrica en 3D como la fidelidad de la apariencia en comparación con la imagen de entrada.

\begin{remark}[Procesos generativos desacoplados]
Obsérvese que el marco de dos etapas anterior desacopla la distribución
de los objetos 3D en varias componentes interrelacionadas ---forma, pose y
apariencia--- y aprende cada uno de forma independiente. Aunque esto implica que debemos
aprender dos modelos generativos de eliminación de ruido al mismo tiempo, las representaciones aprendidas
resultantes están mucho mejor estructuradas y permiten un razonamiento espacial y una generación mucho más
precisos y flexibles.
\end{remark}

\subsection{Experimentos}
\label{sec:experiments}
\label{sec:Evaluation}

Para validar la arquitectura, evaluamos el sistema en tres aspectos fundamentales: predicción 
de la geometría monocular, coherencia entre la vista de entrada y la vista de salida, así como 
reconstrucción completa de una sola imagen a 3D. El proceso de evaluación comparativa se centró en 
sistemas basados en el aprendizaje, excluyendo los métodos de optimización por escena para garantizar 
una comparación imparcial de las capacidades de inferencia. Eliminamos el condicionamiento por pose.

\paragraph{El panorama comparativo.}
Hemos realizado una comparación con tres paradigmas distintos que actualmente dominan el campo:
\begin{enumerate}[label=(\roman*)]
    \item \textit{Regresión de mapas de puntos}, como VGGT~\citep{wang2025vggt} y MoGe~\citep{wang2025moge}, 
          que predicen simultáneamente puntos 3D por píxel y las poses de la cámara, pero que a menudo carecen 
          de priores robustos para la geometría oculta;
    \item \textit{Reconstrucción centrada en la vista}, incluyendo LRM~\citep{hong2023lrm} y LaRa~\citep{LaRa}, 
          que generan modelos 3D directamente en el espacio de la vista para evitar la estimación de la pose 
          durante la inferencia; y
    \item \textit{Generación desacoplada}, ejemplificada por OnePoseGen~\citep{geng2025one}, que separa la 
          generación 3D canónica de la estimación de la pose, basándose en una alineación post-hoc.
\end{enumerate}

\newcommand{\cmark}{\checkmark} % Checkmark
\newcommand{\xmark}{\times} % Cross

\begin{table*}[t]
\renewcommand{\arraystretch}{1.0}
\centering
\caption{
\textbf{Monocular geometry accuracy.} Cupid outperforms all 3D reconstruction and generation baselines and matches point-map regression methods that predict only partial geometry. Note that VGGT uses a ground-truth object mask, which may overestimate accuracy.
}
\label{tab:geometry}
\resizebox{0.99\textwidth}{!}{
\begin{tabular}{l c c c c c c | c c c c c}
\toprule
\multirow{2}{*}{\textbf{Method}} & \multirow{2}{*}{\textbf{3D}} &
\multicolumn{5}{c|}{\textbf{Toys4k}} &
\multicolumn{5}{c}{\textbf{GSO}} \\
& &
\textbf{mIOU} & \textbf{CD} & \textbf{CD} & \textbf{F-score} & \textbf{F-score} &
\textbf{mIOU} & \textbf{CD} & \textbf{CD} & \textbf{F-score} & \textbf{F-score} \\
& &
(avg)$\uparrow$ & (avg)$\downarrow$ & (med)$\downarrow$ & (0.01)$\uparrow$ & (0.05)$\uparrow$ &
(avg)$\uparrow$ & (avg)$\downarrow$ & (med)$\downarrow$ & (0.01)$\uparrow$ & (0.05)$\uparrow$ \\
\midrule
\textcolor{gray}{VGGT} & $\xmark$ &
\textcolor{gray}{--} & \textcolor{gray}{1.144} & \textcolor{gray}{0.498} & \textcolor{gray}{61.85} & \textcolor{gray}{95.90} &
\textcolor{gray}{--} & \textcolor{gray}{1.396} & \textcolor{gray}{0.388} & \textcolor{gray}{65.98} & \textcolor{gray}{95.95} \\
\textcolor{gray}{MoGe} & $\xmark$ &
\textcolor{gray}{92.80} & \textcolor{gray}{1.284} & \textcolor{gray}{0.581} & \textcolor{gray}{58.54} & \textcolor{gray}{95.31} &
\textcolor{gray}{96.18} & \textcolor{gray}{1.743} & \textcolor{gray}{0.575} & \textcolor{gray}{58.99} & \textcolor{gray}{94.68} \\ 
OnePoseGen & $\cmark$ &
9.34 & 153.2 & 59.92 & 6.11 & 24.10 &
12.16 & 116.2 & 60.56 & 7.28 & 25.77 \\
LaRa & $\cmark$ &
68.11 & 32.15 & 16.59 & 18.57 & 57.67 &
70.63 & 34.23 & 19.36 & 13.48 & 49.95 \\
OpenLRM & $\cmark$ &
86.26 & 2.726 & 1.291 & 40.42 & 90.60 &
91.35 & 3.741 & 1.858 & 34.14 & 87.20 \\
Ours & $\cmark$ &
\textbf{92.43} & \textbf{2.534} & \textbf{0.236} & \textbf{69.82} & \textbf{97.76} &
\textbf{95.27} & \textbf{1.823} & \textbf{0.434} & \textbf{61.01} & \textbf{95.59} \\
\bottomrule
\end{tabular}
}
\end{table*}

\begin{table}[t]
\vspace{6pt}
\centering
\caption{\textbf{Comparison on full 3D quality. }We report CLIP image scores of novel views following \citep{gao2024cat3d}.}
\label{tab:clip}

\renewcommand{\arraystretch}{1.0}
\resizebox{1.0\linewidth}{!}{
\begin{tabular}{l|ccccc}
\toprule
 & OnePoseGen & LaRa & OpenLRM & TRELLIS & Cupid \\
\midrule
\textbf{ViT-B/16} & 0.7933 & 0.8334 & 0.8939 & 0.9465 & \textbf{0.9501} \\
\textbf{ViT-L/14} & 0.7193 & 0.7682 & 0.8410 & 0.9210 & \textbf{0.9291} \\
\bottomrule
\end{tabular}
}
\end{table}




\begin{figure}[ht]
    \centering
    \includegraphics[width=1.0\linewidth]{chapters/chapter8/figs/cupid/figures/nvs.pdf}
    \caption{
\textbf{
Comparación cualitativa en cuanto a la consistencia con la vista de entrada.
} Renderizamos la vista de entrada utilizando la pose de cámara generada.
En el caso de los métodos centrados en la vista (LRM, LaRa), utilizamos los parámetros 
intrínsecos de referencia para el renderizado, ya que estos métodos no modelan dichos parámetros.
Nuestro método produce la geometría y la apariencia de mayor fidelidad; LRM genera detalles 
erróneos, LaRa resulta excesivamente borrosa debido a inconsistencias en la difusión 2D y el 
método de generación 3D OnePoseGen a menudo no logra registrar la pose de manera confiable.
}
\label{fig:view_consistency}
\end{figure}
\renewcommand{\arraystretch}{1.0}
\begin{table}[t]
\centering
\caption{
\textbf{Input-view consistency.} Cupid achieves superior input view consistency, producing accurate appearance alignment.
}
\label{tab:appearance}
\centering
\resizebox{\linewidth}{!}{
\begin{tabular}{lc|ccc|ccc}
\toprule
\textbf{Dataset} & &
\multicolumn{3}{c|}{\textbf{Toys4K}} &
\multicolumn{3}{c}{\textbf{GSO}} \\
\textbf{Method} & \textbf{Pose} & \textbf{PSNR}$\uparrow$ & \textbf{SSIM}$\uparrow$ & \textbf{LPIPS}$\downarrow$ & \textbf{PSNR}$\uparrow$ & \textbf{SSIM}$\uparrow$ & \textbf{LPIPS}$\downarrow$ \\
\midrule
LaRa & $\xmark$ & 22.00 & 93.42 & 0.0884 & 19.81 & 91.61 & 0.1119 \\
OpenLRM & $\xmark$ & 26.41 & 80.17 & 0.1156 & 25.79 & 78.80 & 0.1268 \\
OnePoseGen & $\cmark$ & 17.43 & 89.37 & 0.1174 & 14.87 & 86.46 & 0.1386 \\
Cupid & $\cmark$ & \textbf{30.05} & \textbf{96.81} & \textbf{0.0251} & \textbf{28.68} & \textbf{95.49} & \textbf{0.0354} \\
\bottomrule
\end{tabular}
}
\end{table}

\paragraph{Precisión geométrica y simplificación de la secuencia de proceso.}
En las tareas de geometría monocular, el modelo propuesto demostró una clara superioridad 
frente a los enfoques centrados en la vista (Tabla~\ref{tab:geometry}). En el conjunto de 
datos GSO~\citep{downs2022google}, redujo la distancia chamfer (CD) promedio en un 50\% 
en comparación con LRM~\citep{hong2023lrm}. Esta mejora se atribuye a la doble capacidad del 
modelo: modela con precisión las intrínsecas de los objetos mediante prior 
generativo, al tiempo que predice con exactitud los parámetros extrínsecos de la cámara.

Lo más importante es que el modelo alcanzó una precisión geométrica comparable a la de regresores 
especializados en mapas de puntos, como MoGe~\citep{wang2025moge}. Este hallazgo tiene importantes 
implicaciones para el diseño de flujos de trabajo: sugiere que los flujos de trabajo complejos y 
de múltiples etapas---que tradicionalmente utilizan la estimación geométrica parcial como paso 
previo~\citep{yao2025castcomponentaligned3dscene}---pueden simplificarse. Nuestro enfoque alcanza 
una precisión comparable en un solo paso, lo que hace innecesarias esas representaciones intermedias.

\paragraph{Fidelidad y consistencia visuales.}
Evaluamos la consistencia entre la entrada y la vista para medir hasta qué grado la reconstrucción 
3D se alinea con la imagen de origen (Tabla~\ref{tab:appearance}). Los resultados pusieron de 
manifiesto debilidades específicas en los modelos de referencia: la calidad de la textura de LaRa 
se degradó debido a la generación inconsistente de nuevas vistas, mientras que OnePoseGen adolecía de 
una grave desalineación de texturas y de cambios de color durante el registro.

Por el contrario, nuestro método mantuvo una alta precisión en la pose sin sacrificar los detalles 
de textura (Figura~\ref{fig:view_consistency}). Reconstruyó con éxito objetos de alta fidelidad 
tanto a partir de conjuntos de datos sintéticos (Toys4K) como de imágenes de entornos naturales complejos que 
planteaban un gran desafío (Figura~\ref{fig:cupid_results}), superando de manera consistente a 
las líneas de base en métricas de alineación semántica como la similitud 
CLIP~\citep{radford2021learning} (Tabla~\ref{tab:clip}).


\paragraph{Estudios de ablación}
\label{sec:ablation}

\begin{figure}
    \centering
    \includegraphics[width=1\linewidth]{chapters/chapter8/figs/cupid/figures/ablation.pdf}
    \caption{
\textbf{Comparación cualitativa de diversos métodos de condicionamiento con alineación de poses.} Nuestro método (e) ofrece la mejor calidad visual en cuanto a fidelidad del color y nivel de detalle.
}\label{fig:abl_appearance}
\end{figure}
% \begin{table}[t]
% \centering
% \renewcommand{\arraystretch}{1.0}
% \begin{minipage}[t]{\linewidth}
% \vspace{0pt} % Ensures baseline alignment
% \captionof{table}{Comparison on CLIP image scores of novel views.}
% \label{tab:clip}
% \centering
% \resizebox{0.55\linewidth}{!}{
% \begin{tabular}{l|cc}
% \toprule
% \textbf{Method} & \textbf{ViT-B/16} & \textbf{ViT-L/14} \\
% \midrule
% OnePoseGen   & 0.7933 & 0.7193 \\
% LaRa         & 0.8334 & 0.7682 \\
% OpenLRM      & 0.8939 & 0.8410 \\
% TRELLIS      & 0.9465 & 0.9210 \\
% Ours         & \textbf{0.9501} & \textbf{0.9291} \\
% \bottomrule
% \end{tabular}
% }
% \end{minipage}

% \hfill

% \renewcommand{\arraystretch}{1.065}

% \begin{minipage}[t]{\linewidth}
% \vspace{0pt} % Ensures baseline alignment
% \captionof{table}{Ablation studies of pose-aligned conditioning.}
% \label{tab:abl_appearance}
% \centering
% \resizebox{\linewidth}{!}{
% \begin{tabular}{l|ccc|ccc}
% \toprule
% \multirow{2}{*}{\textbf{Method}} & \multicolumn{3}{c|}{\textbf{GT Geo \& Pose}} & \multicolumn{3}{c}{\textbf{Sampled Geo \& Pose}} \\
%  & \textbf{PSNR} & \textbf{SSIM} & \textbf{LPIPS} & \textbf{PSNR} & \textbf{SSIM} & \textbf{LPIPS} \\
% \midrule
% (a) Baseline (TRELLIS)   & 31.84 & 97.50 & 0.0219 & 27.47 & 95.64 & 0.0327 \\
% (b) Position Embedding   & 32.07 & 97.58 & 0.0211 & 27.56 & 95.67 & 0.0323 \\
% (c) Latent (w/o Occ.)    & 32.37 & 97.72 & 0.0201 & 27.85 & 95.87 & 0.0309 \\
% (d) Latent (Occ.)        & 32.39 & 97.77 & 0.0199 & 27.74 & 95.80 & 0.0313 \\
% (e) Latent (Visual Feat.)& \textbf{34.86} & \textbf{98.24} & \textbf{0.0168} & \textbf{30.05} & \textbf{96.81} & \textbf{0.0251} \\
% \bottomrule
% \end{tabular}
% }
% \end{minipage}
% \end{table}


\begin{table}[t]
\centering
\caption{\textbf{Ablation studies of pose-aligned conditioning.}}
\label{tab:abl_appearance}
\renewcommand{\arraystretch}{1.065}
\resizebox{1.0\linewidth}{!}{
\begin{tabular}{l|ccc|ccc}
\toprule
\multirow{2}{*}{\textbf{Method}} & \multicolumn{3}{c|}{\textbf{GT Geo \& Pose}} & \multicolumn{3}{c}{\textbf{Sampled Geo \& Pose}} \\
 & \textbf{PSNR} & \textbf{SSIM} & \textbf{LPIPS} & \textbf{PSNR} & \textbf{SSIM} & \textbf{LPIPS} \\
\midrule
(a) Baseline (w/o PAC)   & 31.84 & 97.50 & 0.0219 & 27.47 & 95.64 & 0.0327 \\
(b) Position Embedding   & 32.07 & 97.58 & 0.0211 & 27.56 & 95.67 & 0.0323 \\
(c) Latent (w/o Occ.)    & 32.37 & 97.72 & 0.0201 & 27.85 & 95.87 & 0.0309 \\
(d) Latent (Occ.)        & 32.39 & 97.77 & 0.0199 & 27.74 & 95.80 & 0.0313 \\
(e) Latent (Visual Feat.)& \textbf{34.86} & \textbf{98.24} & \textbf{0.0168} & \textbf{30.05} & \textbf{96.81} & \textbf{0.0251} \\
\bottomrule
\end{tabular}
}
\end{table}

Los estudios de ablación realizados en la segunda fase de refinamiento confirmaron el papel 
fundamental del condicionamiento alineado con la pose (PAC). Probamos variantes que iban desde 
desde líneas base estándar de latente hasta otras que incorporaban volúmenes de características DINOv2.

Como se muestra en la Tabla~\ref{tab:abl_appearance}, la configuración más eficaz resultó ser 
un enfoque híbrido: la combinación de características de DINOv2 con indicios visuales de bajo 
nivel extraídos mediante capas convolucionales. Esta combinación proporcionó la mayor calidad 
visual, uniendo de manera eficaz la comprensión semántica con los detalles a nivel de píxel 
(Figura~\ref{fig:abl_appearance}). Además, este mecanismo de condicionamiento demostró una 
solidez notable, ya que mantuvo un alto rendimiento incluso cuando la geometría de entrada 
presentaba perturbaciones estocásticas desde la primera etapa de generación.


\paragraph{Aplicación I: Reconstrucción de escenas compuestas.}

\begin{figure*}[t]
\centering
\includegraphics[width=\textwidth]{chapters/chapter8/figs/cupid_scene.pdf}
\caption{\small\textbf{Reconstrucción de escenas alineada por componentes.} Cada
    objeto se reconstruye de forma independiente con \textbf{Cupid}; las correspondencias explícitas
entre 3D y 2D permiten su colocación precisa en un marco compartido.}
\label{fig:cupid_scene_reconstruction}
\end{figure*}


\label{sub:cupid_application}
Para extender Cupid a la reconstrucción a nivel de escena, utilizamos modelos de fundación (foundation models) para descomponer 
la escena en $K$ objetos con máscaras $\{ \boldsymbol{M}^{(k)} \}_{k=1}^K$ y estimar un mapa de 
puntos métricos global $\mathcal{P}$~\citep{wang2025moge}.

En la primera etapa, reconstruimos cada objeto de forma independiente para recuperar su geometría 
canónica, adaptando el generador para manejar las oclusiones mutuas mediante un ajuste fino con 
máscaras aleatorias~\citep{wu2025amodal3r}. Para el objeto $k$-ésimo, esto da como resultado un 
conjunto de puntos 3D canónicos $\boldsymbol{x}_i^{(k)}$ correspondientes a los píxeles 
$\boldsymbol{u}_i \in \boldsymbol{M}^{(k)}$. Para resolver la ambigüedad de escala inherente a 
la generación independiente, formulamos la alineación como un problema de correspondencia entre 
el espacio canónico generado y el espacio de la cámara estimado:
\begin{equation}
    \boldsymbol{p}_i = \mathcal{P}(\boldsymbol{u}_i), \quad \forall \boldsymbol{u}_i \in \boldsymbol{M}^{(k)},
\end{equation}
donde $\boldsymbol{p}_i$ representa las coordenadas métricas del objetivo en el sistema de referencia 
de la cámara, derivadas del modelo de base de profundidad.

En la segunda etapa, calculamos una transformación de similitud por objeto 
$\mathcal{T}_{k} = \{s_{k}, \boldsymbol{R}_{k}, \boldsymbol{t}_{k}\}$ para situar todos las componentes 
en un sistema de coordenadas unificado. Determinamos estos parámetros utilizando el método de 
Umeyama~\citep{umeyama2002least} minimizando el error de alineación sobre los píxeles visibles de cada objeto:
\begin{equation}
    \mathcal{T}_{k}^* = \argmin_{s_k, \boldsymbol{R}_k, \boldsymbol{t}_k} \sum_{i} \left\| \boldsymbol{p}_i - \left(s_k \boldsymbol{R}_k \boldsymbol{x}_i^{(k)} + \boldsymbol{t}_k\right) \right\|^2.
\end{equation}
La aplicación de estas transformaciones óptimas da como resultado una composición de escena 
métricamente consistente, en la que todas las componentes generados están correctamente posicionados 
y escalados entre sí. Consulte la Figura~\ref{fig:cupid_scene} para ver más ejemplos de composición de escenas.

\begin{figure}
    \centering
    \includegraphics[width=0.95\linewidth]{chapters/chapter8/figs/cupid/figures/scene_result.pdf}
    %
    \caption{\textbf{Ejemplos adicionales de reconstrucción de escenas alineadas por componentes.} 
             En cada uno de los ejemplos mostrados, los paneles muestran: (arriba a la izquierda) 
             la imagen de entrada, (arriba a la derecha o abajo a la izquierda) el resultado final 
             renderizado y (abajo) las componentes individuales reconstruidas, codificadas por 
             colores para mayor claridad.
    }
    \label{fig:cupid_scene}
    \vspace{-5mm}
\end{figure}


\paragraph{Aplicación II: Reconstrucción consistente a partir de múltiples vistas.}

\begin{figure}
    \centering
    \includegraphics[width=\textwidth]{chapters/chapter8/figs/cupid/figures/mv.pdf}
    \caption{
    \textbf{Condicionamiento multivista.} Nuestro modelo de articulaciones desacopladas admite de 
             forma natural el condicionamiento multivista. Al disponer de múltiples vistas de 
             entrada, fusionamos el objeto latente compartido---independiente de la 
             vista---a lo largo de las rutas de flujo (al igual que 
             MultiDiffusion~\citep{bar2023multidiffusion}), lo que permite el refinamiento de los 
             objetos y las cámaras en todas las vistas. Arriba: entradas; en el centro: objeto 3D 
             reconstruido y poses de la cámara; abajo: imágenes renderizadas y geometría.}
    \label{fig:mv_images}
\end{figure}

\label{sub:cupid_application_mv}
% To implement the multi-view capability described in Section~\ref{sec:flow_model}, we generalize the sampling process to accommodate $K$ input views $\{ \image^{(k)} \}_{k=1}^K$. 
Gracias a nuestra formulación bayesiana, nuestro modelo se generaliza a escenarios con múltiples 
vistas incluso cuando se entrena únicamente con imágenes individuales. Para implementar la 
capacidad de múltiples vistas descrita en la Sección~\ref{sec:flow_model}, extendemos el proceso 
de muestreo para acomodar $K$ vistas de entrada ${ \image^{(k)} }_{k=1}^K$.

En la primera etapa, nuestro objetivo es recuperar una única geometría canónica y, al mismo 
tiempo, estimar $K$ poses distintas de la cámara. Inicializamos una latente de geometría 
compartida $\boldsymbol{S}_{o}$ y $K$ latentes de pose independientes 
$\{ \boldsymbol{S}_{uv}^{(k)} \}_{k=1}^K$. En cada paso de eliminación de ruido $t$, 
predecimos los campos de flujo tanto para la geometría como para la pose en paralelo para 
cada vista. Sea la salida de la red de flujo para la vista $k$-ésima un vector concatenado 
de actualizaciones de geometría y pose:
\begin{equation}
    [d_{o, t}^{(k)}, d_{uv, t}^{(k)}] = \sflow(\text{concat}(\boldsymbol{S}_{o}^t, \boldsymbol{S}_{uv}^{(k), t}), t, \image^{(k)}).
\end{equation}
Para garantizar la coherencia geométrica en todas las observaciones, agregamos las 
actualizaciones geométricas al tiempo que mantenemos las trayectorias de pose específicas de 
cada vista. La regla de actualización en el paso $t$ queda así:
\begin{align}
    \boldsymbol{S}_{o}^{t-1} &\leftarrow \text{Step}\left(\boldsymbol{S}_{o}^{t}, \frac{1}{K} \sum_{k=1}^K d_{o, t}^{(k)}\right), \\
    \boldsymbol{S}_{uv}^{(k), t-1} &\leftarrow \text{Step}\left(\boldsymbol{S}_{uv}^{(k), t}, d_{uv, t}^{(k)}\right),
\end{align}
donde $\text{Step}(\cdot)$ denota la regla de actualización del muestreador de flujo. Las mallas 
de características resultantes se decodificarán en una malla de ocupación gruesa compartida 
$\boldsymbol{G}_o$ y en las poses $\{\pose^{k}\}_{k=1}^{K}$.

En la segunda etapa, con la geometría fijada, aplicamos una estrategia de difusión 
múltiple~\citep{bar2023multidiffusion} a la latente de apariencia $\boldsymbol{F}_{o}$. 
Mantenemos una única latente compartida y la actualizamos utilizando las predicciones de 
flujo promediadas de todas las $K$ vistas, condicionadas a la geometría gruesa 
$\boldsymbol{G}_o$ y a las poses $\{\pose^{k}\}_{k=1}^{K}$ estimadas en la primera etapa:
\begin{equation}
    d_{\text{feat}, t}^{(k)} = \lflow(\boldsymbol{F}_{o}^t, \boldsymbol{G}_o, \pose^{k}, t, \image^{(k)}),
\end{equation}
\begin{equation}
    \boldsymbol{F}_{o}^{t-1} \leftarrow \text{Step}\left(\boldsymbol{F}_{o}^{t}, \frac{1}{K} \sum_{k=1}^K d_{\text{feat}, t}^{(k)}\right).
\end{equation}
Esta operación de promediado garantiza que la representación 3D generada cumpla simultáneamente 
con las restricciones visuales impuestas por todas las vistas de entrada, logrando una consistencia 
similar a la del bundle adjustment SfM que se muestra en la Figura~\ref{fig:mv_images}.


\subsection{Conclusión}
En esta sección se ha presentado un marco unificado para la generación y reconstrucción 3D, 
que operacionaliza la distinción fundamental entre identidad del objeto («qué») y 
contexto espacial («dónde»). Al modelar conjuntamente el objeto 3D canónico $\object$ y la 
orientación de la cámara $\pose$, nuestro enfoque desacopla eficazmente la forma 
intrínseca del punto de vista extrínseco.

Aprovechando los priores de postura y el condicionamiento alineado, el modelo realiza una 
inferencia similar a la bayesiana para reconstruir la geometría 3D completa a partir de entradas 
limitadas de una sola vista. Esto garantiza la coherencia geométrica y la robustez frente a 
las variaciones del punto de vista. En última instancia, esta representación estructurada 
mejora las capacidades de la IA incorporada, dotando a los agentes de la comprensión espacial 
necesaria para desplazarse e interactuar de manera eficaz en entornos 3D.

\section{Generación condicional de movimientos del esqueleto humano}
\label{sec:egoallo}
% \yima{EgoAllo (Brent, Annie and Chen), move the stuff from Chapter
% 6 to here and elaborate on the methods and implementation more.}
% \subsection{Body Pose Generation Conditioned on Head and
% Hands}\label{sub:ego-allo}

%\input{chapters/chapter8/applications_egoallo}
% This type of
% conditional estimation or generation problem arises rather naturally in many
% practical applications.
% A typical problem of this kind is how to estimate and generate body
% pose and hand gesture conditioned on a given head pose and
% egocentric images, as illustrated in Figure \ref{fig:pose-teaser}.
% This is often the problem we need to solve when one is wearing a
% head-mounted device such as the Vision Pro from Apple or the
% Project Aria from Meta. The pose of the whole body and the gesture
% of the hands need to be inferred so that we can use the information
% to control virtual objects that the person interacts with.
% \begin{figure*}[t]
%   \centering
%   \includegraphics[width=\linewidth]{\toplevelprefix/chapters/chapter7/figs/pose-teaser.pdf}
%     \caption{
% A system that estimates human body height, pose, and hand
% parameters (middle), conditioned on egocentric SLAM poses and
% images (left). Outputs capture the wearer's actions in the
% allocentric reference frame of the scene, which we visualize here
% with 3D reconstructions (right).
%   }
%   \label{fig:pose-teaser}

% \end{figure*}

% Notice that in this case, one only has the head pose provided by
% the device and a very limited field of view for part of one's hands
% and upper limbs. The pose of the rest of the body needs to be
% ``inferred'' or ``completed'' based on such partial information.
% The only way one can estimate the body pose over time is by
% learning the joint distribution of the head and body pose sequences
% in advance and then sampling this prior distribution conditioned on
% the real-time partial inputs. Figure \ref{fig:pose-method} outlines
% a system called EgoAllo \cite{yi2024egoallo} to solve this problem
% based on a learned conditional diffusion-denoising model.
% \begin{figure*}[t]
%   \centering
%   \includegraphics[width=\linewidth]{\toplevelprefix/chapters/chapter7/figs/pose-method.pdf}
% \caption{
%     \textbf{Overview of technical components of EgoAllo \cite{yi2024egoallo}.}
%     A diffusion model is pretrained that can generate body pose
% sequence based on local body parameters (middle).
%     An invariant parameterization $g(\cdot)$ of SLAM poses (left)
% is used to condition the diffusion model. These can be placed into
% the global coordinate frame via global alignment to input poses.
%     When available, egocentric video is used for hand detection
% (left) via HaMeR~\cite{pavlakos2023reconstructing}, which can be
% incorporated into samples via guidance by the generated gesture.
%   }
%   \label{fig:pose-method}
% \end{figure*}

% Figure \ref{fig:comparison} compares some ground truth motion
% sequences with sampled results generated by the EgoAllo. Although
% the figure shows one result for each input head pose sequence,
% different runs can generate different body pose sequences that are
% consistent with the given head pose, all drawing from the
% distribution of natural full-body motion sequences.
% \begin{figure}[t]
% \centering
% \begin{subfigure}{0.45\textwidth}
%   \centering
%   \includegraphics[width=\linewidth]{\toplevelprefix/chapters/chapter7/figs/qual0_gt.png}
%   \caption{\centering Ground-truth}
% \end{subfigure}
% \hfill
% \begin{subfigure}{0.45\textwidth}
%   \centering
%   \includegraphics[width=\linewidth]{\toplevelprefix/chapters/chapter7/figs/qual0_ours.png}
%   \caption{EgoAllo}
% \end{subfigure}

% \begin{subfigure}{0.45\textwidth}
%   \centering
%   \includegraphics[width=\linewidth]{\toplevelprefix/chapters/chapter7/figs/qual1_gt.png}
%   \caption{\centering Ground-truth}
% \end{subfigure}
% \hfill
% \begin{subfigure}{0.45\textwidth}
%   \centering
%   \includegraphics[width=\linewidth]{\toplevelprefix/chapters/chapter7/figs/qual1_ours.png}
%   \caption{EgoAllo}
% \end{subfigure}

% \caption{
% \textbf{Egocentric human motion estimation for a running (top) and
% squatting (bottom) sequence.}
% The ground-truth motion is compared with one output from EgoAllo
% that is consistent with the given head pose sequence.
% }
% \label{fig:comparison}
% \end{figure}

% Strictly speaking, the solution proposed in EgoAllo \cite{yi2024egoallo} does
% not enforce measurement matching using the techniques introduced
% above. Instead
% it heuristically enforces the condition by utilizing  the cross-attention
% mechanism in a transformer architecture. As we will describe with
% more precision
% in the paired data setting in \Cref{sub:text-cond}, there is reason
% to believe that the
% cross-attention mechanism is in a way approximately realizing the conditional
% sampling of the denoising a posteriori. We believe the more principled
% techniques introduced here, if properly implemented, can lead to
% better methods
% that further improve the body pose and hand gesture estimation.

%EgoAllo: Egocentric Human Motion Estimation
% Section for Chapter 7 - Learning Representations for Real-World Data
%
% Following the structure of the Cupid (3D Object Generation) section

% \section{Egocentric Human Motion Estimation}
% One must know that the primary reason why our brain senses the
% external 3D world to develop a visual memory is for us to negotiate
% with the physical environment with our body and interact with
% objects in it with our hand. Hence, our body movement and hand
% motion are arguable another most important data that we would like
% to learn about. In fact, our brain dedicates a significant portion
% of it, the motion cortex, to record and coordinate our own body
% motion. In this section, we show how to use the methods introduced
% in this book to learn the distribution of human motions in advance
% and how it can be used to facilitate the task of sensing and
% estimating our own motion in real time.

Hasta ahora, en las secciones anteriores, hemos mostrado cómo aprender las
distribuciones de los datos visuales, incluidas las imágenes en 2D y las formas 
de los objetos en 3D.  Cabe destacar que el objetivo principal de percibir el 
mundo tridimensional y desarrollar un modelo del mismo es la navegación, la interacción y 
la manipulación de objetos en él. A partir de la información que percibe sobre el 
entorno, nuestro cerebro elabora planes y lleva a cabo acciones en función de las 
tareas que debe realizar. El nivel de destreza que se puede alcanzar mediante la 
«coordinación ojo-mano» es, sin duda alguna, la forma más elevada de inteligencia 
física que existe en la naturaleza. Por lo tanto, calcular y controlar con precisión 
los movimientos de nuestro cuerpo y nuestras manos es fundamental para interactuar con 
los objetos del entorno. Nuestro cerebro dedica importantes recursos a este fin: la 
{\em corteza motora}, situada en el lóbulo frontal del cerebro, registra y coordina
los movimientos de nuestro cuerpo, mientras que las señales propioceptivas procedentes del
sistema nervioso de nuestro cuerpo proporcionan información continua sobre la posición de 
nuestras extremidades.  Por lo tanto, si queremos que los robots, especialmente los de 
aspecto humanoide, imiten las acciones humanas, deben aprender (la distribución de) nuestros
movimientos corporales.

En esta sección, mostramos cómo el marco de inferencia condicional del
Capítulo~\ref{ch:conditional-inference} puede aplicarse para aprender la
distribución del movimiento humano y, posteriormente, aprovecharse para la tarea de
estimación de la postura corporal. Para ello, consideramos una configuración sencilla pero natural:
una persona lleva un dispositivo montado en la cabeza, como un casco de RA o unas
gafas inteligentes, que captura vídeo frontal y  rastrea su pose tridimensional
mediante Estructura a partir del Movimiento (SFM) o Localización y
Mapeo Simultáneos (SLAM). Estos métodos estiman la posición del dispositivo mediante
el seguimiento y la triangulación de características visuales de la escena para reconstruir la estructura y la pose en 3D.
\footnote{Los lectores interesados pueden consultar el
libro de texto \cite{MaY2003} para obtener una introducción a esta línea de investigación.}
Nuestro objetivo es estimar la pose de todo el cuerpo, la altura y la configuración de las manos
del usuario a partir de estas observaciones egocéntricas, aunque el
cuerpo rara vez sea visible en el video capturado. La observación clave es
que \emph{el movimiento de la cabeza codifica una gran cantidad de información sobre el movimiento del cuerpo}.
Cuando caminamos, nuestra cabeza se balancea; cuando estiramos el brazo, se inclina; cuando nos sentamos,
describe un arco característico. Estas correlaciones, obtenidas a partir de datos de captura de movimiento,
forman una distribución a priori condicional $p(\text{cuerpo} \mid
\text{cabeza})$. El método que presentamos en esta sección se basa en
un proyecto de código abierto llamado \textbf{EgoAllo} (por
«estimación de egocéntrico a alocéntrico»):\\
\centerline{\url{https://egoallo.github.io}},\\
que incorpora mediciones adicionales como la detección de las manos y
el contacto con el suelo mediante muestreo guiado. La Figura~\ref{fig:egoallo_teaser}
ilustra la configuración general y la tarea.

\begin{figure}[t]
    \centering
    \includegraphics[width=\textwidth]{\toplevelprefix/chapters/chapter8/figs/egoallo/teaser_oct9.pdf}
    \caption{\textbf{Estimación egocéntrica del movimiento humano.} A partir de las
        posturas de la cabeza obtenidas mediante SLAM y de imágenes egocéntricas captadas por un dispositivo montado en la cabeza
        (izquierda), EgoAllo estima la postura completa del cuerpo, la altura y los
        parámetros de las manos en el sistema de referencia mundial (alocéntrico) (derecha).
        Los cuerpos estimados están anclados en la escena, con los pies
    en contacto con el suelo y las manos en posiciones físicamente plausibles.}
    \label{fig:egoallo_teaser}
\end{figure}

% ==============================================================================
\subsection{Representación de cuerpos humanos}
\label{sub:egoallo_body_models}
% ==============================================================================

Antes de formular el problema de estimación, necesitamos una
representación matemática del cuerpo humano. El esqueleto de nuestro cuerpo constituye un
árbol cinemático: huesos rígidos unidos por articulaciones, con la pelvis como
raíz y cadenas que se extienden hasta la cabeza, las manos y los pies. Cada articulación
tiene grados de libertad rotacionales: el hombro gira en tres
ejes, mientras que el codo gira principalmente en uno. Para una introducción más formal y detallada
a la representación cinemática de sistemas de cuerpos rígidos articulados, como
el cuerpo humano, los lectores interesados pueden consultar el clásico libro de texto
\cite{MurrayR1994}.

\paragraph{El modelo corporal SMPL.}
SMPL (modelo lineal multi-persona con piel) \citep{loper2015smpl}
ofrece una representación compacta y diferenciable del cuerpo humano,
tanto en lo que respecta a su postura (cinemática) como a su forma. Divide 
la configuración del cuerpo en dos componentes:
\begin{itemize}
    \item \textbf{Parámetros de postura} $\theta$: Rotaciones articulares
        que especifican la articulación del cuerpo. SMPL utiliza 23 articulaciones (más
        la raíz), cada una parametrizada por una rotación 3D. En conjunto,
        estas determinan si una persona está de pie, sentada,
        alcanzando algo, caminando, etc.
    \item \textbf{Parámetros de forma} $\beta$: Un vector de baja dimensión,
        normalmente compuesto por 10 coeficientes, que controla las proporciones
        corporales, como la altura, el peso, la longitud de las extremidades y el ancho
        del tórax. Estos se derivan del análisis de componentes principales de
        un amplio corpus de escaneos corporales en 3D.
\end{itemize}

Dada una pose $\theta$ y una forma $\beta$, SMPL genera una malla 3D mediante
el método de blend skinning aprendido. También proporciona las posiciones de las articulaciones a través de
la cinemática directa: partiendo de la posición de la raíz (pelvis) y
aplicando la rotación de cada articulación en secuencia a lo largo de la cadena cinemática,
podemos calcular dónde termina cada articulación en el espacio 3D. Por ejemplo, para
encontrar la posición de la muñeca, combinamos las rotaciones de la pelvis,
la columna vertebral, el hombro y el codo, acumulando las traslaciones definidas por
las longitudes de los huesos.

\paragraph{Agregar manos: SMPL-H y MANO.}
El modelo SMPL original representa las manos como simples puntos finales. Para
aplicaciones que requieren posturas detalladas de las manos, como agarrar objetos, hacer gestos
o escribir, esto resulta insuficiente. SMPL-H \citep{romero2017embodied}
amplía SMPL con manos articuladas basadas en el modelo de manos MANO,
añadiendo 15 articulaciones por mano (tres por dedo), lo que suma un total de 52 articulaciones.

Las rotaciones de la mano $\phi$ siguen la misma parametrización que las articulaciones
del cuerpo. El estado completo incluye las rotaciones del cuerpo $\theta$, las rotaciones
de la mano $\phi$ y los parámetros de forma $\beta$.

\paragraph{Por qué son importantes los modelos corporales.}
¿Por qué utilizar SMPL-H en lugar de, por ejemplo, realizar una regresión directa de las posiciones de las articulaciones
o de los vértices de la malla? Los modelos paramétricos ofrecen ventajas clave. Los resultados son
siempre cuerpos humanos válidos, sin configuraciones anatómicamente imposibles.
El modelo es diferenciable, lo que permite la optimización basada en gradientes.
La parametrización es compacta, ya que solo requiere unos pocos
cientos de números en comparación con una malla completa. Y obtenemos una estructura semántica,
razonando sobre la «muñeca izquierda» en lugar de índices de vértices.

% ==============================================================================
\subsection{Tarea y configuración}
\label{sub:egoallo_objective}
% ==============================================================================

\paragraph{Entradas.}
Las entradas consisten en dos flujos procedentes de un dispositivo montado en la cabeza:
\begin{itemize}
    \item \textbf{Posturas de la cabeza} $\boldsymbol{h}_{1:T} \in \mathsf{SE}(3)^T$\footnote{Aquí, 
            $\mathsf{SE}(3)$ representa el grupo euclidiano especial en $\mathbb{R}^3$, es decir, el 
        grupo de transformaciones que consisten en rotaciones y traslaciones.}:
        Posición y orientación del dispositivo en el marco de refrencia mundial
        en cada paso de tiempo, a partir de SLAM visual-inercial. Aquí $\mathsf{SE}(3)$
        denota el grupo de transformaciones rígidas en 3D, que comprende
        rotaciones y traslaciones. \footnote{Los lectores interesados pueden
            consultar \cite{MaY2003} para obtener una introducción formal al movimiento de
        cuerpos rígidos, así como a su representación y propiedades.} Estas
        posturas tienen escala métrica, lo que significa que conocemos las distancias reales en
        metros, en lugar de solo la geometría relativa, y están
        alineadas con la gravedad, por lo que sabemos cuál es la dirección hacia arriba.
    \item \textbf{Imágenes egocéntricas} $I_{1:T}$: Vídeo orientado hacia delante.
        El cuerpo del usuario no suele verse, pero las manos aparecen con frecuencia
        al alcanzar o manipular objetos.
\end{itemize}

A partir de imágenes, un estimador de manos estándar como HaMeR
\citep{pavlakos2023hamer} genera mediciones por cuadro cuando las manos
son visibles:
\begin{equation}
    \boldsymbol{y}^{\text{hand},t} = \big(\hat{\phi}^{\text{MANO},t},\,
    \hat{\boldsymbol{q}}^{3D,t},\, \hat{\boldsymbol{q}}^{2D,t}\big),
\end{equation}
que incluye rotaciones estimadas de la mano, puntos clave 3D en el encuadre de la cámara
y puntos clave 2D.

\paragraph{Salidas.}
EgoAllo estima una secuencia de movimientos a lo largo de $T$ pasos temporales:
\begin{equation}
    \boldsymbol{x}_0 = \{\boldsymbol{x}_0^{1}, \dots, \boldsymbol{x}_0^{T}\},
    \quad\text{where}\quad
    \boldsymbol{x}_0^{t} = \big(\theta^t,\, \phi^t,\, \beta,\, s\big).
\end{equation}
Aquí, $\theta^t$ son las rotaciones de las articulaciones del cuerpo, $\phi^t$ las rotaciones de las manos,
$\beta$ los parámetros de forma compartidos a lo largo del tiempo, y $s$ un parámetro de altura
que también se comparte. El subíndice $0$ denota el movimiento limpio,
siguiendo la notación de difusión, donde $\boldsymbol{x}_n$ es la versión con ruido
en el paso $n$.

EgoAllo calcula la altura además de la postura, lo cual es esencial para
la fijación al suelo. Sin ella, el cuerpo flota en una posición vertical
arbitraria; con ella, los pies tocan el suelo y la cabeza se alinea con la cámara.

\paragraph{Formulación como posterior.}
Following the Bayesian framework from
Chapter~\ref{ch:conditional-inference}, estimation is formulated as
sampling from a posterior. Let $\boldsymbol{y}$ denote all observations
and $\boldsymbol{c} = g(\boldsymbol{h}_{1:T})$ be a processed representation of
head poses (to be explained soon in
Section~\ref{sub:egoallo_representation}). The posterior factorizes
as:\footnote{For this factorization to be valid, we essentially
assume that given $\boldsymbol{x}_0$, $\boldsymbol{y}^{\text{hand}},
\boldsymbol{y}^{\text{phys}}$ is conditionally independent of the head
pose $\boldsymbol{c}$.}

Siguiendo el marco bayesiano del
Capítulo~\ref{ch:conditional-inference}, la estimación se formula como
un muestreo desed una distribución a posteriori. Sea $\boldsymbol{y}$ el conjunto de todas las observaciones
y $\boldsymbol{c} = g(\boldsymbol{h}_{1:T})$ una representación procesada de
las poses de la cabeza (que se explicará más adelante en
la Sección~\ref{sub:egoallo_representation}). El posterior se descompone
en factores:\footnote{Para que esta descomposición sea válida, suponemos esencialmente
    que, dado $\boldsymbol{x}_0$ y $\boldsymbol{y}^{\text{hand}},
    \boldsymbol{y}^{\text{phys}}$ es condicionalmente independiente de la postura de la cabeza
$\boldsymbol{c}$.}
\begin{equation}
    p(\boldsymbol{x}_0 \mid \boldsymbol{y}) \;\propto\; p_\theta(\boldsymbol{x}_0 \mid
    \boldsymbol{c}) \cdot p(\boldsymbol{y}^{\text{hand}},
    \boldsymbol{y}^{\text{phys}} \mid \boldsymbol{x}_0),
\end{equation}
donde:
\begin{itemize}
    \item $p_\theta(\boldsymbol{x}_0 \mid \boldsymbol{c})$ es una distribución a priori condicional
        aprendida, concretamente un modelo de eliminación de ruido por difusión que
        genera movimientos corporales plausibles a partir del movimiento de la cabeza.
    \item $p(\boldsymbol{y}^{\text{mano}}, \boldsymbol{y}^{\text{físico}} \mid
        \boldsymbol{x}_0)$ es la probabilidad de las observaciones de la mano y
        las restricciones físicas, como el contacto con el suelo y la ausencia de
        deslizamiento de los pies, dado el movimiento del cuerpo.
\end{itemize}

La probabilidad se expresa mediante una energía de guía $\mathcal{L}_{\text{guide}}$:
\begin{equation}
    p(\boldsymbol{y}^{\text{hand}}, \boldsymbol{y}^{\text{phys}} \mid
    \boldsymbol{x}_0) \propto \exp\big(-\mathcal{L}_{\text{guide}}(\boldsymbol{x}_0)\big),
\end{equation}
por lo que el posterior queda así:
\begin{equation}
    \label{eq:egoallo_posterior}
    p(\boldsymbol{x}_0 \mid \boldsymbol{y}) \;\propto\; p_\theta(\boldsymbol{x}_0 \mid
    \boldsymbol{c}) \exp\big(-\mathcal{L}_{\text{guide}}(\boldsymbol{x}_0)\big).
\end{equation}
Esta es la forma descrita en el Capítulo~\ref{ch:conditional-inference}: una
distribución a priori aprendida combinada con restricciones de medición a través de una
verosimilitud basada en la energía.

% ==============================================================================
\subsection{Diseño de la representación condicional}
\label{sub:egoallo_representation}
% ==============================================================================

La distribución a priori condicional $p_\theta(\boldsymbol{x}_0 \mid \boldsymbol{c})$ genera
el movimiento del cuerpo dada una representación $\boldsymbol{c}$ del movimiento de la cabeza. ¿Cómo
debería ser esta representación? Como se ha destacado a lo largo de este
libro, la elección de la representación afecta de manera fundamental al aprendizaje y a la generalización.

Un enfoque ingenuo se basaría directamente en las poses brutas de la cabeza
$\boldsymbol{h}_{1:T} \in \mathsf{SE}(3)^T$, es decir, la posición y la orientación absolutas
en cada paso temporal. Esto no funciona bien. Para entender por qué, pensemos en qué
propiedades debería cumplir la representación.

\paragraph{Invarianza espacial.}
La correspondencia entre el movimiento de la cabeza y el movimiento del cuerpo no debería depender de
el lugar concreto en el que se produzca el movimiento. Caminar en la cocina debería
generar las mismas estimaciones de pose corporal que caminar en el salón.
Lo único que importa es el patrón de movimiento local, no la posición absoluta.

El condicionamiento en poses absolutas incumple este principio. Como
ilustra la Figura~\ref{fig:egoallo_conditioning}, un movimiento corporal
idéntico corresponde a diferentes trayectorias absolutas de la cabeza dependiendo
de la ubicación en el mundo. Un modelo condicionado a las poses absolutas debe aprender
por separado que «la cabeza en $(3,2,1)$ moviéndose hacia adelante» y «la cabeza en
$(10,5,2)$ moviéndose hacia adelante» corresponden ambas a caminar, una
carga innecesaria que perjudica la capacidad de generalización.

\paragraph{Invarianza temporal.}
El mapeo tampoco debe depender del momento en que se produzca el movimiento dentro de una ventana de observación.
El procesamiento de los cuadros $0$--$100$ debe comportarse de la
misma manera que el procesamiento de los cuadros $50$--$150$ si contienen el mismo
movimiento subyacente. Esto descarta las magnitudes acumulativas, como el desplazamiento total
desde el inicio de la secuencia, así como los índices de tiempo absolutos.

Se podría intentar canonizar al primer intervalo de tiempo colocando el
primer cuadro en el origen. Pero esto viola la invariancia temporal:
diferentes segmentos temporales del mismo movimiento, canonizados a sus
respectivos primeros cuadros, producen trayectorias diferentes. El modelo
percibe entradas diferentes para el mismo movimiento corporal subyacente.

\paragraph{La solución: canonización por paso de tiempo.}
Para lograr la invariancia tanto espacial como temporal, EgoAllo define un
sistema de coordenadas canónicas en \emph{cada} paso de tiempo, en lugar de solo
al inicio. La representación $\boldsymbol{c} = \{\boldsymbol{c}^{1}, \dots,
\boldsymbol{c}^{T}\}$ relaciona:
\begin{itemize}
    \item Movimiento relativo de la cabeza $\Delta\boldsymbol{T}_{\text{head}}^{t}$
        entre cuadros consecutivos, que refleja cómo se movió la cabeza desde
        $t-1$ hasta $t$.
    \item Pose canónica por paso de tiempo
        $\boldsymbol{T}_{\text{canonical}}^{t}$: la orientación de la cabeza
        relativa a la gravedad y su altura sobre el suelo, eliminando así 
        la posición en el plano del suelo y el giro horizontal (yaw).
\end{itemize}

Estas magnitudes son independientes del lugar del mundo en el que se produzca el movimiento
y del intervalo temporal que se analice. Dos secuencias con
un movimiento local idéntico producen tokens de condicionamiento idénticos
$\boldsymbol{c}$, invariantes respecto a la ubicación del movimiento en el 
mundo y al intervalo temporal seleccionado para su extracción.

\begin{figure}[t]
    \centering
    \includegraphics[width=2in]{\toplevelprefix/chapters/chapter8/figs/egoallo/conditioning_oct9.pdf}
    \caption{\textbf{Acondicionamiento del movimiento de la cabeza invariante.} La
        representación del condicionamiento debe ser invariante tanto respecto a las transformaciones espaciales
        (el mismo movimiento en diferentes ubicaciones del mundo) como a los
        desplazamientos temporales (el mismo movimiento en diferentes momentos absolutos).
    La canonicización por paso de tiempo cumple ambas propiedades.}
    \label{fig:egoallo_conditioning}
\end{figure}

% ==============================================================================
\subsection{La distribución a priori condicional}
\label{sub:egoallo_architecture}
% ==============================================================================

La distribución a priori condicional $p_\theta(\boldsymbol{x}_0 \mid \boldsymbol{c})$ es un
modelo de difusión que genera secuencias de movimiento corporal a partir del movimiento de la cabeza
como condición. La arquitectura sigue el diseño basado en transformers
descrito en la Sección~\ref{sec:image_classification}, adaptado
para datos de movimiento secuenciales.

\paragraph{Arquitectura.}
El modelo toma tres entradas:
\begin{itemize}
    \item La secuencia de movimiento ruidosa $\boldsymbol{x}_n$ en el paso de difusión
        $n$, tokenizada por intervalos de tiempo.
    \item La incrustación del paso de tiempo de difusión, que indica el nivel de ruido.
    \item Los tokens de condicionamiento de la cabeza invariantes $\boldsymbol{c}$, uno por paso temporal.
\end{itemize}

Los tokens de cabeza se procesan primero mediante un codificador compuesto por seis
bloques de transformer con auto-atención para generar representaciones
contextualizadas. Estas alimentan un decodificador, compuesto por otros seis
bloques de transformer, que elimina el ruido de la secuencia de movimiento. El condicionamiento se introduce
mediante atención cruzada: en cada capa del decodificador, los tokens de movimiento forman
consultas que prestan atención a los tokens de la cabeza como claves y valores, lo que permite
al modelo seleccionar dinámicamente qué aspectos del movimiento de la cabeza son
relevantes para predecir cada articulación corporal en cada paso temporal. El modelo
genera una estimación directa del movimiento limpio $\boldsymbol{x}_0$, en lugar
de predecir el ruido, parametrizado como rotaciones articulares en un
sistema de coordenadas relativo a la cabeza.

\paragraph{Datos de entrenamiento.}
El modelo se entrena con AMASS \citep{mahmood2019amass}, un conjunto de datos de captura de movimiento a gran escala
agregando múltiples fuentes en un formato unificado
SMPL-H. AMASS contiene diversos movimientos humanos---caminar,
correr, bailar, sentarse, estirarse, hacer ejercicio---pero no incluye grabaciones de dispositivos
montados en la cabeza. Las trayectorias de la cabeza se simulan extrayendo la
posición y la orientación de la articulación de la cabeza a partir de los datos de captura de movimiento, como si un
dispositivo estuviera fijado rígidamente a la cabeza.

\paragraph{Objetivo del entrenamiento.}
El entrenamiento utiliza la pérdida de coincidencia de puntuaciones con eliminación de ruido del \Cref{ch:denoising}:
\begin{equation}
    \mathcal{L}_{\text{diff}} = \mathbb{E}_{\boldsymbol{x}_0, \boldsymbol{n}} \left[ w_n
    \left\| D_\theta(\boldsymbol{x}_n, \boldsymbol{n}, \boldsymbol{c}) - \boldsymbol{x}_0 \right\|^2 \right],
\end{equation}
donde $D_\theta$ es la red de eliminación de ruido, $w_n$ son pesos que dependen del ruido
y $\boldsymbol{x}_n$ se obtiene añadiendo ruido al movimiento limpio
$\boldsymbol{x}_0$ según el proceso de difusión hacia adelante. El modelo
aprende a predecir el movimiento limpio a partir del movimiento ruidoso, condicionado por las trayectorias
de la cabeza.

% ==============================================================================
\subsection{Incorporación de medidas mediante guías}
\label{sub:egoallo_guidance}
% ==============================================================================

La distribución a priori condicional refleja la relación entre el movimiento de la cabeza
y el movimiento del cuerpo, pero no utiliza observaciones de imágenes ni restricciones físicas.
Para incorporarlas, EgoAllo define una energía de guía
$\mathcal{L}_{\text{guide}}$ que mide la coherencia entre una
hipótesis de movimiento y las observaciones.

\paragraph{La energía de guía.}
La energía de guía total combina varias componentes:
\begin{equation}
    \mathcal{L}_{\text{guide}}(\boldsymbol{x}_0) =
    \lambda_{3D}\mathcal{L}_{3D} + \lambda_{2D}\mathcal{L}_{2D} +
    \lambda_{\text{skate}}\mathcal{L}_{\text{skate}} +
    \lambda_{\text{prior}}\mathcal{L}_{\text{prior}},
\end{equation}
donde:
\begin{itemize}
    \item $\mathcal{L}_{3D}$: Alineación entre los puntos clave estimados de la mano
        y las detecciones 3D de la mano obtenidas con HaMeR, en el sistema de coordenadas global.
    \item $\mathcal{L}_{2D}$: Error de reproyección. Los puntos clave de la mano estimados,
        al proyectarse sobre la imagen, deben coincidir con las detecciones en 2D.
    \item $\mathcal{L}_{\text{skate}}$: Penalización de deslizamiento de pies (foot skating penalty). Los pies
        no deben deslizarse cuando están en contacto con el suelo.
    \item $\mathcal{L}_{\text{prior}}$: Regularización suave hacia
        la predicción del prior de difusión, evitando que la guía
        se aleje demasiado de los movimientos plausibles.
\end{itemize}

Siguiendo la interpretación basada en la energía del
Capítulo~\ref{ch:conditional-inference}, esta energía de guía
corresponde a una log-verosimilitud negativa: $p(\boldsymbol{y} \mid
\boldsymbol{x}_0) \propto \exp(-\mathcal{L}_{\text{guide}}(\boldsymbol{x}_0))$. Una
energía baja implica consistencia con las observaciones; una energía alta indica inconsistencia.

\paragraph{Por qué es importante la reproyección.}
Los estimadores de manos de cuadro único, como HaMeR, generan poses de manos en 3D a partir de
imágenes monoculares, pero sufren de ambigüedad de escala y profundidad. Una mano pequeña
cerca de la cámara aparece idéntica a una mano grande que está lejos. Por lo tanto, las
coordenadas en 3D, por tanto, no son confiables en términos absolutos, incluso cuando
la pose local de la mano, como la articulación de los dedos, es precisa.

La pérdida de reproyección 2D resuelve este problema mediante la supervisión en el espacio de la imagen,
donde se observó realmente la mano. Las posiciones 3D estimadas de la mano
se vuelven a proyectar sobre la imagen y se comparan contra las detecciones de
puntos clave en 2D. Esto elude la ambigüedad de profundidad manteniendo al mismo tiempo
una señal útil sobre la localización de la mano.

% ==============================================================================
\subsection{Algoritmo de muestreo condicional}
\label{sub:egoallo_sampling}
% ==============================================================================

El marco de coincidencia de mediciones del
Capítulo~\ref{ch:conditional-inference} nos indica cómo realizar un muestreo a partir
de la distribución a posteriori: hay que modificar el filtro de ruido para incluir una corrección de gradiente
hacia la coherencia de las medidas. Recordemos que el filtro de ruido condicional
se descompone de la siguiente manera:
\begin{equation}
    \mathbb{E}[\boldsymbol{x}_0 \mid \boldsymbol{x}_n, \boldsymbol{y}]
    =
    \mathbb{E}[\boldsymbol{x}_0 \mid \boldsymbol{x}_n]
    +
    \sigma_n^2 \nabla_{\boldsymbol{x}_n} \log p(\boldsymbol{y} \mid \boldsymbol{x}_n),
\end{equation}
donde el primer término corresponde al filtro de ruido anterior y el segundo, a la
corrección de ajuste de medidas.

En nuestro contexto, el filtro de ruido previo está condicionado por el movimiento de la cabeza:
$D_n^{\text{head}}(\cdot; \boldsymbol{c})$. La corrección de medidas utiliza 
la energía como guía. El filtro de ruido condicional completo queda así:
\begin{equation}
    D_n^{\text{guided}}(\boldsymbol{z})
    =
    D_n^{\text{head}}(\boldsymbol{z};\boldsymbol{c})
    -
    \sigma_n^2
    \nabla_{\boldsymbol{z}}
    \mathcal{L}_{\text{guide}}(\boldsymbol{x}_0(\boldsymbol{z})),
\end{equation}
utilizando la relación $\nabla \log p(\boldsymbol{y} \mid \boldsymbol{x}) =
-\nabla \mathcal{L}_{\text{guide}}(\boldsymbol{x})$.

\paragraph{Optimización alterna.}
En la práctica, en lugar de calcular los gradientes a partir de toda la
energía de guía en cada paso, EgoAllo alterna entre dos operaciones:

\begin{enumerate}
    \item \textbf{Paso previo de eliminación de ruido.} Aplicar una actualización al estilo DDIM
        utilizando el eliminador de ruido condicionado por la cabeza
        $D_n^{\text{head}}(\cdot;\boldsymbol{c})$ para obtener una propuesta
        $\tilde{\boldsymbol{x}}_{n-1}$. DDIM (Modelos implícitos de difusión con eliminación de ruido) 
        proporciona un procedimiento de muestreo determinista que mapea
        muestras ruidosas hacia otras más limpias; a diferencia de los muestreadores estocásticos,
        cada paso de eliminación de ruido es una función determinista del
        estado actual, lo que permite un muestreo más rápido con menos iteraciones.

    \item \textbf{Paso de ajuste de medidas.} Refinar la propuesta resolviendo
        un problema de optimización proximal:
        \begin{equation}
            \boldsymbol{x}_{n-1}
            =
            \operatorname*{arg\,min}_{\boldsymbol{x}}
            \left\{
                \frac{1}{2}\|\boldsymbol{x}-\tilde{\boldsymbol{x}}_{n-1}\|^2
                +
                \gamma_n\,\mathcal{L}_{\text{guide}}(\boldsymbol{x})
            \right\},
        \end{equation}
        donde $\gamma_n$ es un parámetro de tamaño del paso (step-size). Esto localiza un movimiento
        cercano a la predicción de la distribución a priori, pero más coherente con las observaciones.
\end{enumerate}

La optimización proximal se resuelve mediante Levenberg-Marquardt, un
algoritmo clásico para problemas de mínimos cuadrados no lineales. Este
interpola entre el descenso de gradiente y la optimización de Gauss-Newton,
adaptando el tamaño de su paso en dependencia de cuán bien la linealización
aproxima al objetivo real. Esto resulta muy adecuado para la estimación de parámetros SMPL-H,
donde el objetivo implica una cinemática directa diferenciable. Bastan unas pocas iteraciones 
en cada paso de eliminación de ruido.

\begin{figure}[t]
    \centering
    \includegraphics[width=\linewidth]{\toplevelprefix/chapters/chapter8/figs/egoallo/egoallo-overview.png}
    \caption{\textbf{Resumen de las componentes del marco EgoAllo}. 
        El modelo de difusión se limita a los parámetros corporales
        locales. Se utiliza una parametrización invariante
        g(·) de las poses SLAM (de la cabeza) para
        condicionar un modelo de difusión. Estas pueden situarse en el
        sistema de coordenadas global
        mediante una alineación global con las poses de entrada. Cuando está disponible, se utiliza un video egocéntrico
        para la detección de manos, por ejemplo, mediante HaMeR
        \citep{pavlakos2023hamer}, que puede incorporarse a las muestras
    mediante la guía.}
    \label{fig:egoallo_framework}
\end{figure}

\paragraph{Alineación global.}
El modelo de difusión opera en un sistema de coordenadas relativo a la cabeza,
prediciendo la postura del cuerpo en relación con la cabeza, en lugar de hacerlo en coordenadas absolutas
del mundo. Tras el último paso de eliminación de ruido, EgoAllo ancla el
cuerpo en el espacio componiéndolo con las posturas de la cabeza observadas por SLAM:
\begin{equation}
    \boldsymbol{T}_{\text{world},\text{root}}^t
    =
    \boldsymbol{T}_{\text{world},\text{head}}^t
    \cdot
    \boldsymbol{T}_{\text{head},\text{root}}^{(\theta^t, \beta)},
\end{equation}
donde $\boldsymbol{T}_{\text{head},\text{root}}$ se calcula mediante cinemática directa
a partir de la pose y la forma estimadas. Esto garantiza que el cuerpo
esté correctamente anclado: los pies en el suelo y la cabeza alineada con la cámara.

\paragraph{Manejo de secuencias largas.}
El modelo de difusión se entrena en secuencias de longitud fija (32-128 cuadros). 
Para grabaciones más largas, EgoAllo utiliza el enfoque MultiDiffusion
\citep{bar2023multidiffusion}: la secuencia se divide en
ventanas superpuestas y se elimina el ruido en paralelo. Tras cada paso de eliminación de ruido,
los resultados de las regiones superpuestas se mezclan por promedio
antes de proceder al siguiente paso. Esta combinación paso a paso garantiza
que las ventanas mantengan su coherencia mutuamente durante el
proceso de difusión, lo que produce un resultado temporalmente coherente sin
saltos o discontinuidades en los bordes de las ventanas.

% ==============================================================================
\subsection{Conclusiones extraídas de los experimentos}
\label{sub:egoallo_experiments}
% ==============================================================================

EgoAllo se evalúa en tres conjuntos de datos: AMASS
\citep{mahmood2019amass} (captura de movimiento con trayectorias sintéticas de la cabeza
), RICH \citep{huang2022rich} (captura de movimiento con
interacciones complejas entre personas y el entorno), y Aria Digital Twin
\citep{pan2023aria} (grabaciones reales de las gafas del Proyecto Aria con
poses corporales de referencia). Para la estimación de manos, se utiliza también EgoExo4D
\citep{grauman2023egoexo4d}.

\paragraph{La representación es crucial.}
La Tabla~\ref{tab:egoallo_cond_ablation} muestra el efecto de diferentes
representaciones de condicionamiento en la precisión de la estimación corporal, comparando
cinco parametrizaciones con diferentes propiedades de invariancia:
\begin{itemize}
    \item \textbf{Invariante por paso de tiempo}: Representación
        descrita en la sección~\ref{sub:egoallo_representation}, con
        invarianza tanto espacial como temporal.
    \item \textbf{Absoluto + relativo}: Posturas absolutas enriquecidas con
        el movimiento relativo en el marco local.
    \item \textbf{Deltas absolutas + globales}: Posturas absolutas con
        desplazamientos (deltas) de posición en coordenadas de mundo.
    \item \textbf{Canonización de secuencias}: alineación al primer
        cuadro, proporcionando invariancia espacial, pero no temporal.
    \item \textbf{Absoluto}: Poses brutas $\mathsf{SE}(3)$ sin propiedades de invariancia.
\end{itemize}

\begin{table}[t]
    \centering
    \caption{\textbf{Efecto de la representación de condicionamiento en AMASS.}
        Las diferentes parametrizaciones del condicionamiento del movimiento de la cabeza dan lugar a
        una precisión de estimación sustancialmente diferente. La representación
        invariante por paso de tiempo, que cumple tanto con la invarianza espacial como
        con la temporal, es la que ofrece mejores resultados. MPJPE es el error medio de posición por articulación
    en milímetros.}
    \label{tab:egoallo_cond_ablation}
    \small
    \begin{tabular}{lcccc}
        \toprule
        & \multicolumn{2}{c}{\textbf{Seq Length 32}} &
        \multicolumn{2}{c}{\textbf{Seq Length 128}} \\
        \cmidrule(lr){2-3} \cmidrule(lr){4-5}
        \textbf{Conditioning} & MPJPE & Effect & MPJPE & Effect \\
        \midrule
        Per-timestep invariant & \textbf{129.8} & --- & \textbf{119.7} & --- \\
        Absolute + relative & 133.0 & +2.4\% & 124.5 & +4.0\% \\
        Absolute + global deltas & 136.2 & +4.9\% & 127.4 & +6.4\% \\
        Sequence canonicalization & 153.1 & +17.9\% & 134.0 & +11.9\% \\
        Absolute & 159.9 & +23.2\% & 148.3 & +23.9\% \\
        \bottomrule
    \end{tabular}
\end{table}

Los resultados muestran que la representación de condicionamiento tiene un impacto
considerable en la precisión, con diferencias de hasta un 24\% en el MPJPE entre la
mejor y la peor elección. Esto refuerza una idea central de este libro:
la selección de la representación afecta fundamentalmente al aprendizaje y a la generalización.

\paragraph{La estimación de la altura permite la conexión al suelo.}
La Tabla~\ref{tab:egoallo_body_results} compara la estimación completa
(incluida la altura) con una variante sin estimación de la altura.
La estimación de la altura mejora el MPJPE en un 6-7\% y es clave para
el anclaje: sin ella, los cuerpos estimados tienden a «flotan» por encima del
suelo, ya que el sistema no puede determinar la posición vertical
correcta. La métrica de anclaje al suelo (GND) mide si los pies estimados
tocan el suelo en algún momento, y la estimación de la altura mejora esto sustancialmente.

\begin{table}[t]
    \centering
    \caption{\textbf{Resultados de la estimación corporal en distintos conjuntos de datos.} MPJPE
        y PA-MPJPE en milímetros; GND es la proporción de secuencias
        en las que los pies tocan el suelo. La estimación de la altura 
    facilita el anclaje.}
    \label{tab:egoallo_body_results}
    \small
    \begin{tabular}{llccc}
        \toprule
        \textbf{Dataset} & \textbf{Variant} &
        \textbf{MPJPE}$\downarrow$ & \textbf{PA-MPJPE}$\downarrow$ &
        \textbf{GND}$\uparrow$ \\
        \midrule
        AMASS & Full (with height) & \textbf{119.7} & \textbf{101.1}
        & \textbf{1.00} \\
        & Without height & 128.1 & 110.3 & 0.98 \\
        \midrule
        RICH  & Full (with height) & \textbf{176.2} & \textbf{160.1}
        & \textbf{0.96} \\
        & Without height & 185.7 & 169.9 & 0.82 \\
        \midrule
        ADT   & Full (with height) & \textbf{155.1} & \textbf{129.3} & 0.94 \\
        & Without height & 163.7 & 140.0 & \textbf{0.96} \\
        \bottomrule
    \end{tabular}
\end{table}

\paragraph{El contexto del cuerpo mejora la estimación de manos.}
Un hallazgo destacable es la relación simbiótica entre la estimación corporal y la manual.
Los estimadores monoculares de mano de un solo cuadro, como HaMeR, producen
posturas locales precisas de la mano, incluida la articulación de los dedos, pero no pueden
determinar la posición 3D absoluta de la mano debido a la ambigüedad de escala y profundidad. 
Al estimar primero el cuerpo completo, ganamos restricciones cinemáticas:
la mano está unida a la muñeca, que se conecta con el
antebrazo, el codo, el brazo y el torso. Esta cadena provee el contexto para
la localización global de la mano. La Figura \ref{fig:egoallo_hands} ilustra algunos
resultados cualitativos.

La Tabla ~\ref{tab:egoallo_hands} cuantifica este efecto en el conjunto de datos EgoExo4D. Las estimaciones sin procesar
de HaMeR presentan un error de 237,9 mm en el coordenadas de mundo. Al añadir el contexto corporal,
esto se reduce a 131,5 mm, lo que representa una mejora del 45\%, incluso utilizando solo
observaciones monoculares de las manos. Cuando se dispone de observaciones 3D adicionales de la muñeca
procedentes de cámaras estéreo, los errores disminuyen aún más, hasta alcanzar los 60,1 mm.

\begin{table}[t]
    \centering
    \caption{\textbf{Estimación de manos en EgoExo4D.} MPJPE mide
        la precisión en el marco global; PA-MPJPE mide la precisión de la pose local. El contexto corporal
        mejora notablemente la localización en el marco global al tiempo que
    mantiene la calidad de la pose local.}
    \label{tab:egoallo_hands}
    \small
    \begin{tabular}{lcc}
        \toprule
        \textbf{Método} & \textbf{MPJPE}$\downarrow$ &
        \textbf{PA-MPJPE}$\downarrow$ \\
        \midrule
        HaMeR (monocular, sin cuerpo) & 237.9 & \textbf{13.0} \\
        \midrule
        Con cuerpo (reproyección 2D) & 131.5 & 14.7 \\
        Con cuerpo (3D wrist from stereo) & \textbf{60.1} & 14.4 \\
        \bottomrule
    \end{tabular}
\end{table}

La relación es bidireccional: las observaciones de la mano también mejoran
la estimación corporal. Cuando se incorpora la orientación de la mano durante el muestreo,
el MPJPE corporal en AMASS mejora de 119,7 mm a 78,8 mm, lo que supone una reducción del 34\%.
Las restricciones de medición adicionales refinan la estimación a posteriori
y producen estimaciones más precisas a lo largo de toda la cadena cinemática.

\begin{figure}[t]
    \centering
    \includegraphics[width=\textwidth]{\toplevelprefix/chapters/chapter8/figs/egoallo/hands_oct1_v4.png}
    \caption{\textbf{El contexto corporal mejora la estimación de la mano.} Azul:
        estimaciones monoculares de la mano obtenidas con HaMeR, que presentan una pose local
        precisa pero una posición incorrecta en el marco mundial debido a la ambigüedad de escala/profundidad.
        Morado: estimaciones con contexto corporal, correctamente
    ancladas en el marco mundial.}
    \label{fig:egoallo_hands}
\end{figure}

\paragraph{Resultados cualitativos.}
La Figura ~\ref{fig:egoallo_qualitative} muestra las posturas corporales estimadas a partir de
grabaciones egocéntricas del mundo real realizadas con las gafas del Proyecto Aria. Los cuerpos estimados
están correctamente situados en el suelo: los pies tocan el suelo de manera adecuada, la altura del cuerpo
se coincide con el entorno y las posturas son coherentes con
el movimiento observado de la cabeza en diversas actividades, como caminar,
sentarse y manipular objetos.

\begin{figure}[t]
    \centering
    \includegraphics[width=\textwidth]{\toplevelprefix/chapters/chapter8/figs/egoallo/qual_grounded_outputs_v3.png}
    \caption{\textbf{Resultados cualitativos.} Posturas corporales estimadas a partir de
        grabaciones egocéntricas del mundo real, visualizadas en reconstrucciones de escenas en 3D.
        El método genera posturas físicamente plausibles
    con una base adecuada en diversas actividades.}
    \label{fig:egoallo_qualitative}
\end{figure}

% ==============================================================================
\subsection{Discusión}
\label{sub:egoallo_discussion}
% ==============================================================================

Esta aplicación demuestra cómo el marco de inferencia condicional
del Capítulo~\ref{ch:conditional-inference} se extiende más allá de las imágenes para abarcar
resultados estructurados y articulados, como el movimiento humano.

\paragraph{El diseño de la representación es fundamental.}
Un tema recurrente a lo largo de este libro es que las decisiones de representación
influyen fundamentalmente en el aprendizaje y la generalización. Este principio
resulta especialmente evidente en este caso. Al analizar qué invariancias debe
cumplir la representación de condicionamiento, derivamos una
parametrización que supera sustancialmente a las alternativas. La invarianza espacial
garantiza la independencia respecto a la ubicación en el mundo; la invarianza temporal
garantiza la independencia al desplazamiento temporal. La lección: antes de
entrenar un modelo, es crucial considerar cuidadosamente qué estructura
debe codificar la representación.

\paragraph{El patrón «Prior más guía» generaliza.}
El procedimiento de muestreo aplica directamente el marco de
correlación de mediciones: un prior aprendido---en este caso, el
modelo de difusión condicionado por la cabeza---se combina con
restricciones de medición, como observaciones de la mano y 
restricciones físicas, mediante un muestreo guiado. Este mismo patrón de
«prior más guía» se repite en diversos ámbitos, desde la
generación de imágenes hasta la estimación de movimiento. Las
componentes son modulares: podemos intercambiar diferentes priors o
términos de guía sin alterar el marco general.

\paragraph{La estimación holística supera a los flujos de proceso modulares.}
La relación simbiótica entre la estimación del cuerpo y la de la mano sugiere
que la inferencia conjunta sobre toda la cadena cinemática supera a
la estimación independiente de las componentes. El movimiento de la cabeza informa la postura del cuerpo;
la postura del cuerpo proporciona contexto para la localización de la mano; y las observaciones de la mano
refinan ambas estimaciones (cuerpo y mano). Este beneficio mutuo surge
en forma natural de la formulación a posteriori, donde todas las variables están
acopladas a través de la función de verosimilitud y la distribución a priori.

\paragraph{Limitaciones.}
Persisten varias limitaciones. En primer lugar, el método requiere un SLAM preciso
con alineación de escala métrica y gravedad; los errores en la estimación de la postura de la cabeza
se propagan directamente a las estimaciones del cuerpo. En segundo lugar,
la generalización está limitada por la diversidad de los datos de entrenamiento, por lo que las posturas o actividades inusuales,
no representadas en AMASS, pueden estimarse mal.
En tercer lugar, las manos solo se pueden refinar si son visibles en la vista egocéntrica;
las manos ocultas dependen exclusivamente del prior corporal. Por último, el método
asume suelos planos; las escaleras o los terrenos irregulares requerirían
un modelado adicional.

\paragraph{Problemas abiertos.}
Entre las líneas de trabajo prometedoras se encuentra la incorporación de la geometría de la escena para lograr una mejor
anclaje físico, de modo que las manos tengan contacto con las superficies y los cuerpos
eviten los obstáculos. Otras direcciones incluyen la habilitación de la inferencia en tiempo real
para aplicaciones de realidad aumentada o virtual en vivo y la extensión a escenarios con múltiples personas
en los que el usuario interactúe con otros. El marco es
lo suficientemente flexible como para acoger estas extensiones mediante términos de guía adicionales
o priores modificados.

\section{Representación y generación del lenguaje natural}\label{sec:clm_text}
Hasta ahora, hemos mostrado cómo aprender representaciones estructuradas profundas
para las distribuciones de datos asociadas con la percepción visual y
los movimientos corporales. En gran medida, las representaciones estructuradas así aprendidas
pueden considerarse como recuerdos almacenados en la corteza visual
y la corteza motora del cerebro. En conjunto, estos recuerdos pueden denominarse de manera general
un «modelo del mundo», o simplemente un «modelo mundial»
—que nos permite tanto percibir el mundo como realizar acciones. Obsérvese que
tanto los seres humanos como los animales tienen la capacidad de aprender y desarrollar
recuerdos del mundo y de su propio cuerpo, que registran lo que se puede
predecir a partir de lo percibido, así como las habilidades motoras cruciales para
su supervivencia en el entorno.

Además de ese modelo del mundo que abarca el entorno físico y el movimiento corporal,
el ser humano ha desarrollado los idiomas (hablados y escritos) para poder
comunicar a los demás los conocimientos adquiridos sobre el mundo;
las áreas de Broca y de Wernicke de nuestro cerebro son conocidas por
su función en el procesamiento del habla y la comprensión del lenguaje, respectivamente. En
gran medida, los idiomas registran una pequeña, pero importante, fracción de
nuestros modelos del mundo que vale la pena compartir entre nosotros. Por lo tanto,
los idiomas naturales, al ser un tipo de datos muy especial, deben tener estructuras muy
especiales que codifican una información muy rica. Si las máquinas
fueran capaces de aprender las estructuras de los idiomas naturales, podrían
comunicarse mejor con los seres humanos.

Por lo tanto, en esta sección mostramos cómo aplicar los métodos presentados
en este libro a la tarea del \textit{modelado del lenguaje}. O, más
concretamente, cómo entrenar modelos de lenguaje a gran escala (LLM) para que aprendan la
distribución de los lenguajes naturales y, a continuación, los generen. La tarea
y su configuración son muy similares a las utilizadas para entrenar GPT-2 y
muchos otros modelos de lenguaje.

\subsection{Datos} \label{sub:clm_text_data}

Los datos que utilizaremos para investigar el rendimiento de CRATE en
tareas lingüísticas serán OpenWebText (OWT)~\cite{Gokaslan2019OpenWeb},
una reproducción de código abierto del conjunto de datos WebText, aún sin publicar, utilizado por
OpenAI para entrenar a GPT2. Cada muestra de OWT es un documento web, normalmente
procedente de páginas web, blogs, artículos o debates en línea
de alta calidad, que está escrito en lenguaje natural bien formado. El
conjunto de datos OpenWebText contiene alrededor de 8,01 millones de documentos de
diferentes longitudes, lo que suma un total de unos 41,70 GB de texto. Para la evaluación, utilizaremos
varios conjuntos de datos, tales como
WikiText~\cite{merity2016pointer}\footnote{Para WikiText2 y
WikiText103~\cite{merity2016pointer}, las divisiones de prueba son las mismas,
por lo que las fusionamos en un único conjunto de datos denominado WikiText.},
LAMBADA~\cite{paperno2016lambadadatasetwordprediction}\footnote{Para
obtener la precisión en el conjunto de datos LAMBADA, utilizamos la decodificación codiciosa.},
y PTB~\cite{marcus-etal-1993-building}. PTB y OWT son, en general,
más sencillos en comparación con otros conjuntos de datos. PTB se centra en textos periodísticos más sencillos,
ideales para el modelado del lenguaje tradicional, mientras que OWT
es variado e informal, abarca diversos temas, pero presenta una menor
complejidad en la estructura del lenguaje o en las dependencias de largo alcance.
WikiText, con su estructura formal y su contenido específico de un ámbito,
requiere una comprensión más compleja que OWT, pero sigue siendo
manejable. LAMBADA es el más complejo, ya que implica
dependencias de largo alcance, lo que exige que el modelo comprenda
información contextual más amplia para completar las oraciones con precisión.

En un plano más formal, nuestros datos \(\vX\) serán texto, es decir, cadenas de
caracteres; definimos \(\cT\) como el conjunto de todas las cadenas de caracteres.

\subsection{Tarea y objetivo} \label{sub:clm_text_objective}

Para el preentrenamiento de modelos de lenguaje causales, la idea es que queremos
\textit{entrenar el modelo para que genere texto similar al humano}. La forma más
popular de hacerlo, con diferencia, es utilizar un proceso de entrenamiento
en dos etapas:\footnote{El entrenamiento de los modelos de lenguaje modernos cuenta con varios
pasos adicionales que requieren diferentes distribuciones de datos
y enfoques algorítmicos. Sin embargo, entrenar un modelo para que simplemente imite
la escritura humana solo requiere estos pocos pasos que se presentan aquí.}
\begin{itemize}
\item \textit{En primer lugar}, queremos \textit{aprender} una forma de codificar de manera óptima
    los documentos como una secuencia de cadenas básicas («bloques de construcción»)
    denominadas \textit{tokens}. A esto se le llama
    \textit{tokenización}, y vamos a crear un \textit{tokenizador}.
\item \textit{En segundo lugar}, queremos \textit{aprender} una forma de
    \textit{predecir la distribución de un token dados todos los tokens anteriores}. 
    A esto se le denominada \textit{predicción del siguiente token}, y para ello
    construimos un \textit{modelo de lenguaje}.
\end{itemize}
En realidad, este procedimiento se remonta a Markov, quien fue el primero en darse cuenta de que
el lenguaje natural podía modelarse mediante la estructura de la cadena de Markov 
que lleva su nombre \citep{markov2006example}, siempre que se aplicara una
tokenización adecuada; y posteriormente a Shannon, quien propuso aplicar esta misma
configuración de modelado del lenguaje con un tokenizador a nivel de caracteres (es decir, cada
carácter es un token) y los denominados «\(n\)-gramas» (es decir, una
tabla de consulta explícita, calculada a partir de los datos de entrenamiento, que indica la distribución de
un token dado los \(n\) tokens anteriores) en lugar del modelo de lenguaje
\citep{Shannon-1948}.\footnote{Un estudio reciente
\citep{liu2024infini} sobre la ampliación de modelos de \(n\)-gramas ha demostrado que
son capaces de modelar el texto con bastante precisión para valores grandes de \(n\), pero,
por supuesto, la memoria necesaria para almacenar dicha tabla de consulta es del orden de
\(V^{n}\) y, por lo tanto, completamente inmanejable.}

\subsubsection{Entrenamiento de un tokenizador}

Crear un tokenizador equivale a construir un vocabulario \(\cV\), que
es un conjunto de tokens y tiene un tamaño predefinido \(V\). Existen
varios métodos para hacerlo. Un algoritmo muy utilizado es el conocido como Byte
Pair Encoding (BPE), que puede describirse como:
\begin{itemize}
\item Empieza con una lista de todos los caracteres únicos de tus datos de entrenamiento
    y sus frecuencias. Asegúrate de que haya menos de
    \(V\) caracteres de este tipo y añade cada carácter como una cadena independiente
    («token») al vocabulario junto con su frecuencia.
\item Hasta que haya \(V\) tokens en el vocabulario:
    \begin{itemize}
        \item Crear un token tomando los dos tokens existentes más frecuentes
            y fusionándolos.
        \item Calcular la frecuencia de este token en el conjunto de datos.
        \item Añadirlo al vocabulario (junto con su frecuencia).
    \end{itemize}
\item En este punto, la información sobre la frecuencia ya no es necesaria
    y puede descartarse.
\end{itemize}
El proceso general del BPE se muestra en la \Cref{fig:BPE}. Cabe señalar que este
procedimiento es una modificación de un procedimiento clásico de compresión basado en la teoría de la información
para \textit{aprender una codificación sin pérdidas} de
un flujo de bytes (como el texto) y, como tal, puede interpretarse como
la búsqueda de una compresión sin pérdidas óptima de los datos. Obsérvese que esto
es posible porque (a diferencia de las imágenes), los datos aquí son fundamentalmente
discretos y libres de ruido.
\begin{figure}
\centering
\includegraphics[width=0.45\textwidth]{\toplevelprefix/chapters/chapter8/figs/BPE1.png}\hspace{0.3in}
\includegraphics[width=0.45\textwidth]{\toplevelprefix/chapters/chapter8/figs/BPE2.png}
\caption{\small {\bf El proceso de tokenización de datos de texto mediante BPE.}
    (Crédito de la imagen:
    \url{https://huggingface.co/learn/nlp-course/chapter6/5}). (Izquierda)
    Comenzamos analizando el corpus de texto dado y construyendo un
    vocabulario inicial compuesto por caracteres individuales (o bytes en el
    caso del BPE a nivel de byte). A continuación, calculamos las frecuencias de los
    pares de caracteres adyacentes en el corpus. Esto implica analizar todo el texto
    y contar cuántas veces aparece cada secuencia de dos caracteres (bigram).
    (Derecha) Tras calcular las frecuencias de los pares de caracteres adyacentes,
    identificamos el par más frecuente en el corpus. A continuación, este par se
    fusiona en una nueva unidad de subpalabra, que se añade al vocabulario como un
    único token. Este proceso se repite en forma iterativa hasta que
se alcanza el tamaño de vocabulario predefinido. }
\label{fig:BPE}
\end{figure}
Una vez creado dicho vocabulario, un tokenizador puede dividir un
documento en tokens (es decir, «tokenizarlo»). BPE utiliza un procedimiento similar
al del entrenamiento para tokenizar los datos:
\begin{itemize}
\item Divide el documento en una larga lista de tokens de un solo carácter.
    Es decir, si el documento es «Hello», la lista inicial
    será «H», «e», «l», «l», «o».
\item Siempre que dos tokens adyacentes puedan concatenarse y su
    concatenación resulte en otro token, lo hacemos, es decir, reemplazamos este
    par de tokens por el token fusionado. Concretamente, si «He» es un token
    del vocabulario, «H», «e», «l», «l», «o» se convertiría en «He», «l», «l», «o».
\item Repite el proceso anterior hasta que ya no se puedan realizar más fusiones. En
    ese momento, el documento queda dividido en la lista
    (secuencia) final de tokens.
\end{itemize}

Hay muchas consideraciones prácticas y relacionadas con la eficiencia que hay que
tener en cuenta durante la tokenización. El algoritmo anterior, tal y como se presenta,
está \textit{muy lejos de ser óptimo} si se implementa de manera simplista, por
ejemplo. No tratamos este tema en gran detalle; hay muchos
recursos en línea para aprender más, como
\href{https://huggingface.co/learn/nlp-course/en/chapter6/5}{los tutoriales de HuggingFace}.

Por ejemplo, cada token tiene un \textit{índice} correspondiente, que no es
más que su índice en el vocabulario (que, al fin y al cabo, es simplemente una lista de
longitud \(V\)). Por lo tanto, la salida de la mayoría de los tokenizadores es una lista de
índices, es decir, un elemento de \([V]^{*}\). Hay que tener en cuenta que estos
corresponden a subcadenas del documento original, como se ha mostrado anteriormente.

Una vez que se ha entrenado un tokenizador, cualquier
modelo de lenguaje puede utilizarlo como una caja negra. Por ejemplo, muchos modelos tienen el mismo
tokenizador (basado en OpenAI) que se basa en la biblioteca \texttt{tiktoken}. En
el resto de esta sección, utilizaremos dicho tokenizador fijo y precompilado
para todo, y así identificaremos cada documento de texto \(\vX
\in \cT\) con su versión tokenizada en \([V]^{*}\). Por lo tanto,
podemos considerar el espacio textual \(\cT\) como \textit{equivalente} al
espacio de secuencias de tokens \([V]^{*}\) (sin perder nada esencial).

\subsubsection{Entrenamiento de un modelo de lenguaje}

Una vez que tenemos cada documento como una secuencia de tokens \(\vX \in [V]^{N}
\subseteq [V]^{*} = \cT\), queremos realizar la predicción del siguiente token.
Es decir, dado un \textit{contexto} \(\vX_{:n} \in [V]^{n}\) (es decir, los
primeros \(n\) tokens \(\vx_{1}, \dots, \vx_{n} \in [V]\) del
documento)\footnote{Obsérvese la incongruencia con la notación de Python: aquí
la notación \textit{incluye} el índice \(n\).}, queremos predecir el
token \(\vx_{n + 1} \in [V]\) en la posición \(n + 1\). Para ello,
calculamos la característica agregada de \(\vX_{:n}\) mediante
\(\vz_{\theta}(\vX_{:n}) \doteq (f_{\theta}^{\ext} \circ
f_{\theta})(\vX_{:n}) \in \R^{d}\), y utilizamos un cabezal de clasificación
\(h_{\theta} \colon \R^{d} \to \Delta_{V}\) (implementado como una
capa lineal, una MLP o algo un poco más complejo) para proyectar
esta característica en el simplex de probabilidad de dimensión \(V\)
\(\Delta_{V}\). Esta proyección \(\vp_{\theta}(\vX_{:n}) \doteq
h_{\theta}(\vz_{\theta}(\vX_{:n}))\) sirve como una distribución de probabilidad estimada
del siguiente token. Entonces, utilizando la notación
\(\vone(\vx_{n + 1}) \in \Delta_{V}\) para indicar que es \(1\) en la componente \(\vx_{n +
1}\) y \(0\) en el resto, la pérdida del modelo de lenguaje causal es
\begin{equation}\label{eq:clm_loss}
\min_{\theta}\bc{\cL_{\mathrm{CLM}}(\theta) \doteq
    \Ex_{\vX}\rs{\frac{1}{N - 1}\sum_{n = 1}^{N - 1}\CE(\vone(\vx_{n +
1}), \vp_{\theta}(\vX_{:n}))}}
\end{equation}
Fíjate en lo similar que es esto a una pérdida de clasificación (por ejemplo, para imágenes);
se utiliza la entropía cruzada y se intenta alinear un vector de probabilidades
predichas con la verdad de referencia (ground truth). La principal diferencia entre estas dos
pérdidas es que, en este caso, calculamos la pérdida sobre una secuencia completa,
en la que cada predicción está correlacionada con las demás (a diferencia del
caso de clasificación de variables independientes e idénticamente distribuidas).

La optimización de esta pérdida se denomina habitualmente «preentrenamiento» en la
comunidad de modelos de lenguaje (en contraposición al «postentrenamiento» y, más
recientemente, al «entrenamiento intermedio», que son metodologías para modificar un
predictor del siguiente token con el fin de realizar tareas útiles).

\textit{Nota al margen:} ¿Por qué el primer término de \eqref{eq:clm_loss}
predice \(\vone(\vx_{2})\), y no hay ningún término que mida la
pérdida para predecir el primer token? Es porque si quisiéramos predecir
el primer token, tendríamos la \textit{secuencia vacía} como
contexto y, por lo tanto, haríamos esta predicción del primer token utilizando un
mecanismo cualitativamente diferente al que se aplica a los
demás tokens. Así que, en realidad, este modelo no está entrenado para predecir el
\textit{primer} token de ningún documento. La razón por la que esto funciona es
un detalle de implementación del tokenizador: a menudo, tras
crear el tokenizador, insertamos un token \textit{especial} en su
vocabulario, denominado token de \textit{inicio de cadena (o documento)}
y etiquetado como \texttt{<|bos|>}.\footnote{Por lo general, existen
varios tokens especiales para distintos fines. El texto existente
que contiene los tokens especiales se procesa de manera especial.} A continuación, mientras
procesamos cada documento, añadimos el token \texttt{<|bos|>} al
principio de la secuencia de tokens del documento, aumentando la longitud de
la secuencia tokenizada en \(1\). Por lo tanto, el objetivo de modelado de lenguaje causal anterior
tiene un término que implica intentar predecir el
primer token del documento dado únicamente el token \texttt{<|bos|>} como
contexto, por lo que es una pérdida conceptualmente correcta.

\subsection{Arquitectura: CRATE causal}

En cuanto a la arquitectura, utilizamos un transformer estándar al estilo GPT-2,
sustituyendo las capas del transformer por capas CRATE.\footnote{En
contravención directa de las convenciones de este libro y de las de
muchas otras comunidades, la comunidad de PLN denomina a estos transformers al estilo GPT-2
(que abarcan casi todos los LLM actuales) transformers «solo decodificadores». Los 
transformers «solo codificador» tienen una arquitectura diferente,
mientras que los transformers «codificador-decodificador» concatenan un
transformer «solo codificador» con un transformer «solo decodificador».
¡Y esto a pesar de que los transformers «solo decodificador»
\textit{también} calculan una codificación de los datos!} Para completar la información,
especificamos aquí la arquitectura.

\paragraph{Encaje.} Primero incrustamos la secuencia de tokens \(\vX \in
[V]^{N}\) en el espacio euclídeo. Esto se suele hacer asociando cada
índice en \([V]\) con un vector en \(\R^{d}\) utilizando una matriz
\textit{masiva}\footnote{Por «masiva» nos referimos a que dicha
estructura suele constituir una gran fracción del tamaño total.} 
\(\vE \in \R^{V \times d}\), y formando directamente la
secuencia \([\vE_{\vx_{1}}, \dots, \vE_{\vx_{N}}] \in \R^{d \times
N}\). El mapa de incrustación completo \(f_{\theta}^{\emb}\) también aplica una
codificación posicional \(\vE^{\pos} \in \R^{d  \times N_{\max}}\), donde
\(N_{\max}\) es el número máximo de tokens que es posible
procesar,\footnote{Los métodos modernos de codificación posicional han resuelto
este problema y permiten (en teoría) una extrapolación infinita,
pero dichos métodos son más complejos de desarrollar, y por
razones de simplicidad aquí solo introducimos la codificación posicional
aditiva absoluta.} lo que da como resultado el mapa de incrustación
\begin{equation}
f_{\theta}^{\emb}(\vX) \doteq [\vE_{\vx_{1}}, \dots, \vE_{\vx_{N}}] +
\vE_{:N}^{\pos}
\end{equation}
Los parámetros \(\vE\) y \(\vE^{\pos}\) se pueden entrenar directamente.
Dado que \(\vE\) es tan grande (y la actualización del gradiente es muy escasa
con respecto a él, ya que solo se utiliza una pequeña fracción del vocabulario en
cada muestra), se utiliza software especializado para garantizar que las actualizaciones de memoria
no sean demasiado onerosas. Obsérvese también que no utilizamos un token de clase
como en las otras secciones; más adelante se profundizará en este tema.

\paragraph{Red troncal.} Procesamos las incrustaciones utilizando una red troncal similar a CRATE
que utiliza enmascaramiento causal. Para entender el enmascaramiento causal,
consideremos la pérdida de modelado causal del lenguaje \(\cL_{\mathrm{CLM}}\)
definida en \eqref{eq:clm_loss}. La implementación más simple nos obligaría
a calcular la pasada hacia adelante \(N\) veces para poder
realizar una sola retropropagación. Obviamente, esto resulta extremadamente ineficiente, ya que
\(N\) suele ser del orden de miles. Para poder escalar el entrenamiento con
esta pérdida de manera eficiente, imponemos una restricción \textit{causal}, es decir,
\begin{equation}\label{eq:causal_backbone_def}
\vZ_{\theta}(\vX_{:n}) = \vZ_{\theta}(\vX)_{:n}
\end{equation}
es decir, las \(n\) columnas de las características de los tokens
\(\vZ_{\theta}(\vX_{:n}) \in \R^{d \times n}\) deben ser las mismas que
las primeras \(n\) columnas de las características de los tokens \(\vZ_{\theta}(\vX) \in
\R^{d \times N}\), independientemente de los valores positivos de \(n\) y
\(N\) tales que \(N \geq n\). En efecto, esto significa que podemos aplicar la
estructura principal \textit{una vez} a toda la secuencia y calcular
\(\vZ_{\theta}(\vX)\), y luego aplicar \(f_{\theta}^{\ext}\) a cada
subconjunto creciente \(\vZ_{\theta}(\vX_{:n}) = \vZ_{\theta}(\vX)_{:n}\)
a medida que \(n\) crece hasta la longitud de la secuencia \(N\). Entonces podemos usar todos
ellos para calcular la pérdida.

Ahora que queremos una arquitectura causal para la red troncal, ¿cómo
podemos conseguirla? Dado que la MLP y las normalizaciones de capa dentro de cada
capa del transformer afectan a cada token en forma individual, lo único que
importa para la causalidad es el bloque de atención (o \(\MSSA\) en el
caso de CRATE). Para que \(\MSSA\) sea causal, definimos el
bloque \(\mathrm{CMSSA}\) como
\begin{align}
&\operatorname{CMSSA}_{\theta}^{\ell}(\vZ) \doteq
\vU_{\out}^{\ell}\mat{\operatorname{CSA}([\vU^{1,
        \ell}]^{\top}\vZ, [\vU^{1, \ell}]^{\top}\vZ, [\vU^{1,
    \ell}]^{\top}\vZ) \\ \vdots \\ \operatorname{CSA}([\vU^{K,
        \ell}]^{\top}\vZ, [\vU^{K, \ell}]^{\top}\vZ, [\vU^{1,
\ell}]^{\top}\vZ)} + \vb_{\out}^{\ell}\vone_{N}^{\top} \\
&\text{where} \quad  \operatorname{CausalSA}(\vQ, \vK, \vV) \doteq
\vV\softmax\rp{\frac{\operatorname{CausalMask}(\vK^{\top}\vQ)}{\sqrt{p}}}, \\
&\text{where} \quad  \operatorname{CausalMask}(\vM)_{ij} =
\casework{M_{ij}, & \text{if}\ i \geq j, \\ -\infty, & \text{if}\ i < j}.
\end{align}
En este contexto, los especialistas afirman que la máscara causal \textit{permite que los tokens futuros
\(i\) presten atención a los tokens pasados \(j\), pero no al revés}. Para
entender por qué, escribamos la expresión correspondiente a la \(t\)-ésima columna de
\(\operatorname{CSA}(\vQ, \vK, \vV)\):
\begin{equation}
\operatorname{CSA}(\vQ, \vK, \vV)_{t} = \sum_{i =
1}^{t}\vV_{i}\softmax\rp{[\vK_{:t}]^{\top}\vQ_{t}}_{i}
\end{equation}
(donde el subíndice que no es dos puntos denota la columna). Esta
expresión para el token \(t\)-ésimo no utiliza información sobre ningún token
más allá del índice \(t\). Por lo tanto, \(\operatorname{CSA}\), y por lo tanto
\(\operatorname{CMSSA}\), y por lo tanto toda la estructura causal de CRATE
es causal en términos de la definición de
\eqref{eq:causal_backbone_def} y desbloqueamos las considerables
ganancias de eficiencia que se nos prometieron.

\paragraph{Extractor de características.} Utilizamos un paso de posprocesamiento
\(f_{\theta}^{\ext}\) que extrae el \textit{vector de características del
token válido más reciente} con el fin de predecir el siguiente token. En teoría,
esto significa que cada token \(\vZ_{\theta}(\vX)_{n}\) debería contener
información detallada sobre todos los tokens que preceden al índice \(n\) o se encuentran en él,
es decir, \(\vx_{1}, \dots, \vx_{n}\), ya que toda esta información debería
estar disponible para predecir el siguiente token en el índice \(n + 1\). En
la práctica, solo se necesitan unos pocos de estos tokens para cada
tarea de predicción. De todos modos, la ecuación para \(f_{\theta}^{\ext}\) es
\begin{equation}
f_{\theta}^{\ext}(\vZ_{\theta}(\vX_{:n})) \doteq (\vZ_{\theta}(\vX))_{n}
\end{equation}
donde (una vez más) el subíndice que no es dos puntos corresponde a la columna. En este caso, tal y como
prometimos, simplemente extraemos directamente el vector de características del último
token de la secuencia.

\paragraph{Cabezal específico para la tarea.} Para nuestro cabezal de clasificación
\(h_{\theta}\), la arquitectura GPT-2 utiliza una capa lineal simple y
una función softmax para obtener los vectores de probabilidad deseados:
\begin{equation}
h_{\theta}(\vz) \doteq \softmax(\vW^{\out}\vz + \vb^{\out}),
\end{equation}
donde \(\vW^{\out} \in \R^{V \times d}, \vb^{\out} \in \R^{V}\). Algunas
otras arquitecturas más modernas utilizan MLP pequeñas y
normalizaciones de capas, pero la idea es prácticamente la misma. Tenga en cuenta que estas
capas lineales también consumen mucha memoria (porque \(V\) es muy
grande) y constituyen un cuello de botella en el entrenamiento; se han realizado importantes
esfuerzos para intentar sortearlo.

Todas estas decisiones arquitectónicas hacen que el entrenamiento causal sea
extremadamente eficiente en comparación con el entrenamiento no causal:
\begin{itemize}
\item Solo necesitamos \textit{una pasada hacia adelante} a través de la red principal para
    calcular la pérdida de toda la secuencia.
\item La extracción de características es, básicamente, \textit{gratuita}.
\item Todos los tokens pueden enviarse a la cola específica de la tarea
    \textit{en paralelo}.
\end{itemize}

\subsection{Estrategia de optimización}

Entrenamos nuestro modelo de lenguaje mediante optimización estocástica de extremo a extremo.
Un desafío residual es que, en la práctica, los documentos distintos dentro de un lote poseen 
longitudes variables (en función del número de tokens requeridos para cada secuencia). Sin embargo, 
al momento de la redacción de este libro, los principales marcos de aprendizaje profundo admiten, por 
lo general, únicamente tensores «rectangulares», los cuales no se acomodan a este comportamiento. 
Para mitigar este problema, insertamos un token de relleno especial |pad|> para todos los ejemplos 
más cortos del lote, permitiendo así procesar todo el lote mediante tensores rectangulares.
En cada paso temporal \(k\), hacemos lo siguiente:
\begin{itemize}
\item Seleccionar una submuestra de \(B\) documentos tokenizados diferentes
    \(\{\vX_{b}^{(k)}\}_{b = 1}^{B} \subseteq \cT = [V]^{*}\), cada uno
    con una longitud de \(N_{b}^{(k)}\).
\item Calcular \(N_{\max}^{(k)} \doteq \max_{b \in [B]}N_{b}^{(k)}\)
    y rellenar cada \(\vX_{b}^{(k)}\) hasta una longitud de \(N_{\max}^{(k)}\) utilizando
    un token de relleno especial.
\item Calcular las características \(\vZ_{\theta}(\vX_{b}^{(k)})\).
\item Calcular las distribuciones predichas \(\vp_{\theta}(\vX_{b,
        :n}^{(k)}) \doteq (h_{\theta} \circ
    f_{\theta}^{\ext})(\vZ_{\theta}(\vX_{b}^{(k)})_{:n})\).
\item Construir la pérdida estocástica sustitutiva
    \begin{equation}
        \hat{\cL}_{\mathrm{CLM}}^{(k)}(\theta) \doteq
        \frac{1}{B(N_{\max}^{(k)} - 1)}\sum_{b = 1}^{B}\sum_{n =
        1}^{N_{\max}^{(k)}-1}\CE(\vone(\vx_{b, n + 1}^{(k)}),
    \vp_{\theta}(\vX_{b, :n}^{(k)}))).
\end{equation}
\item Ejecutar un paso de un algoritmo de optimización sobre \(\theta\),
obteniendo la siguiente iteración:
\begin{equation}
    \theta^{(k + 1)} \doteq \textsc{OptUpdate}^{(k)}(\theta^{(k)};
    \nabla_{\theta}\hat{\cL}_{\mathrm{CLM}}^{(k)}).
\end{equation}
\end{itemize}

\subsection{Metodología de evaluación} \label{sub:clm_text_evals}

Hay varias formas de evaluar un modelo de lenguaje Transformer entrenado.
\begin{itemize}
\item En un conjunto de datos holdout (de prueba) compuesto por texto arbitrario, podemos evaluar
\(\cL_{\mathrm{CLM}}\) sobre él; las pérdidas más bajas son preferibles, ya que refjan
que el modelo ha aprendido bien esa distribución.
\item En un conjunto de datos de preguntas de opción múltiple, para cada pregunta podemos
utilizarla como contexto y comprobar la probabilidad estimada de que
se genere la respuesta correcta.
\item También podemos probar las capacidades de \textit{generación de texto}.
Específicamente, podemos muestrear repetidamente desde la distribución de probabilidad
del modelo sobre el siguiente token, condicionada al contexto. Cada vez que
tomamos una muestra, generamos un nuevo token, que imprimimos y añadimos al
contexto. Esto nos permite muestrear el LLM y evaluar las
muestras generadas según nuestros criterios.\footnote{Tener que volver a ejecutar el
    modelo con cada token cada vez puede resultar prohibitivamente costoso.
    El almacenamiento inteligente de diferentes características internas del modelo de lenguaje
    (como la denominada \textit{\(K\)-\(V\) cache}), junto con la
causalidad de la arquitectura, puede reducir drásticamente el costo del muestreo.}
\end{itemize}
En esta sección, realizamos exclusivamente este al primer tipo de evaluación.

\subsection{Configuración experimental y resultados}

Dado que nuestra arquitectura CRATE causal se construye directamente en GPT-2,
comparamos los parámetros óptimos para GPT-2 especificados en el repositorio NanoGPT
\citep{nanogpt} con la aplicación de los mismos parámetros a CRATE
para realizar una comparación imparcial.

\paragraph{Arquitectura del modelo.} Utilizamos el tokenizador de GPT-2, que tiene
un tamaño de vocabulario de \(V = 50257\), incluyendo un token especial para
\texttt{<|pad|>}.\footnote{El token \texttt{<|bos|>} no se incluye
en esta configuración, aunque es muy común en los modelos de lenguaje modernos.}
La longitud del contexto es \(N_{\max} = 1024\). El modelo base sigue
la arquitectura GPT2-Base \citep{radford2019language} con las
modificaciones adecuadas para tener capas CRATE causales, y comparamos
con GPT2-Small y GPT2-Base.

\paragraph{Conjuntos de datos y optimización.} Para entrenar el modelo causal CRATE,
seguimos las implementaciones del repositorio NanoGPT \citep{nanogpt}.
Específicamente, utilizamos un tamaño de lote de 384 y entrenamos durante 600 000 pasos
con el optimizador Adam~\citep{kingma2014adam}. Para el optimizador Adam,
utilizamos $(\beta_1, \beta_2)=(0.9, 0.95)$ y una decaimiento de peso
de $0.1$. Para el programa de tasa de aprendizaje, aplicamos un calentamiento lineal
y una disminución cosinusoidal, con un valor máximo de $\eta=6\times 10^{-4}$ en la
iteración 2000, y un valor mínimo de $6\times 10^{-5}$. Las
pérdidas de entrenamiento y validación a lo largo de las iteraciones se muestran en
\Cref{fig:crate-text-evals}. La pérdida de entrenamiento/validación converge
alrededor de 3,37 tras el entrenamiento con un tamaño de lote de 384 y 600 000
iteraciones. En comparación, la implementación abierta de GPT-2 está
preentrenada en OpenWebText con un tamaño de lote de 512 y 600 000
pasos, y converge a una pérdida de validación de 2,85 \citep{nanogpt}.

\begin{figure}
\centering
\includegraphics[width=0.5\textwidth]{\toplevelprefix/chapters/chapter8/figs/gpt-loss.pdf}
\caption{\bf La curva de pérdida de CRATE-GPT-Base entrenado con el
conjunto de datos OpenWebText.}
\label{fig:crate-text-evals}
\end{figure}

\paragraph{Resultados experimentales.}

La \Cref{tab:gpt-eval} demuestra que los modelos CRATE alcanzan un rendimiento razonable
en la pérdida de modelado causal del lenguaje en una variedad de
conjuntos de datos, en comparación con los modelos GPT-2 con un número similar de parámetros y
arquitecturas similares.

\begin{table}
\def\arraystretch{1.1}
\small
\caption{\small Pérdida de entropía cruzada sin entrenamiento previo (zero-shot) del modelo CRATE-GPT2-Base
y de los modelos GPT2-Small y GPT2-Base, evaluada en la parte de prueba de
los conjuntos de datos ($\downarrow$: cuanto más bajo, mejor).
}
\centering
\resizebox{1\columnwidth}{!}
{
\begin{tabular}{ccccccc}
    \hline
    & \#parameters & \textbf{OWT} & \textbf{LAMBADA} &
    \textbf{WikiText} & \textbf{PTB} & \textbf{Avg} \\
    \hline
    GPT2-Base  & {124M} & 2.85$\downarrow$ & 4.12$\downarrow$ &
    3.89$\downarrow$ & 4.63$\downarrow$ & 3.87$\downarrow$ \\
    {GPT2-Small } &  {64M} & {3.04} & {4.49} & {4.31} & {5.15} & {4.25} \\
    Causal-CRATE-Base & {60M} & 3.37 & 4.91 & 4.61 & 5.53 & 4.61 \\
    \hline
\end{tabular}
}
\label{tab:gpt-eval}
\end{table}

\section{Escalado y mejora de los transformers white-box}\label{sec:scalable}\label{sec:improvements}

En esta última sección, discutiremos diversos métodos mediante los cuales se pueden 
escalar diversas componentes de los modelos de tipo CRATE o optimizarlos para tareas 
específicas, manteniendo su naturaleza completamente interpretable white-box.
Estos avances combinan conocimientos tanto conceptuales como empíricos,
y pueden considerarse como estudios de caso sobre cómo utilizar
la idea de white-box para mejorar los modelos de aprendizaje profundo en la práctica.
Las tareas que utilizaremos para evaluar los métodos serán la clasificación de imágenes
y la predicción del siguiente token; los datos serán ImageNet
y OpenWebText, respectivamente; el procedimiento de optimización será el
mismo de retropropagación, y lo único que cambia es la arquitectura.

\subsection{Aumento del ancho de la red:
CRATE-\texorpdfstring{\(\alpha\)}{alpha}}\label{sub:crate_alpha_experiments}

\begin{figure}[t]
\centering
\includegraphics[width=\textwidth]{\toplevelprefix/chapters/chapter8/figs/crate_alpha_backbone.pdf}
\caption{\small\textbf{Una capa de la estructura CRATE-\(\alpha\).}
La diferencia con respecto a CRATE es que el bloque \(\ISTA_{\theta}^{\ell}\)
se sustituye por el bloque \(\operatorname{ODL}_{\theta}^{\ell}\),
que realiza varios pasos \(\ISTA\) con un diccionario sobrecompleto.}
\label{fig:crate_alpha_backbone}
\end{figure}

Una decisión de diseño impuesta por el marco CRATE es el
\textit{ancho} de la función de activación ( ono linealidad) en la red. En un transformer
convencional, el ancho suele fijarse en \(4\), \(8\) o
\(\frac{11}{3}\) veces la dimensión de las características. Sin embargo, CRATE exige
que el ancho sea exactamente igual a la dimensión de las características; esto implica que los
diccionarios \(\vD^{\ell}\) sean cuadrados, lo cual podría resultar en un desempeño reducido.
La razón fundamental por la cual el marco CRATE
restringe nuestra elección es la siguiente:
\begin{itemize}
\item El bloque ISTA realiza \textit{un solo} paso de optimización para
el aprendizaje de diccionarios.
\item Generalmente, un solo paso de cualquier algoritmo de optimización iterativo no logra
optimizar eficazmente la función objetivo. Entonces, ¿cómo es posible que funcione?
\item Los algoritmos de optimización suelen converger muy rápidamente si y
solo si cuentan con una buena inicialización, o un \textit{warm start}. El
bloque ISTA posee un \textit{warm start}: trata las características de entrada como una
inicialización para los códigos dispersos resultantes.
\item Esto garantiza que las características de entrada y los códigos dispersos tengan la
misma dimensión. Es decir, ISTA aprende un diccionario de dispersión completo
(véase el \Cref{ch:classic-models}).
\end{itemize}
Por consiguiente, si deseamos emplear un diccionario ancho, necesitamos que ISTA realice
el aprendizaje de diccionarios \textit{sobrecompletos}. Esto implica que no podemos mantener
el mismo «warm start», ya que los códigos dispersos resultantes tienen una dimensión mayor que
las características de entrada; por tanto, requiriremos más iteraciones para converger a un código disperso.
En consecuencia, el paso de las características \(\vZ_{\theta}^{\ell + 1/2}\) hacia los códigos dispersos
\(\vZ_{\theta}^{\ell + 1}\) ya no se define como
\begin{equation}
\vZ_{\theta}^{\ell + 1} = \ISTA_{\theta}^{\ell}(\vZ_{\theta}^{\ell +
1/2} \mid \vZ_{\theta}^{\ell + 1/2})
\end{equation}
donde la función \(\ISTA_{\theta}^{\ell}\) utilizando una convención de notación 
(un poco abusiva) introducida en secciones anteriores, definida como:
\begin{equation}
\ISTA_{\theta}^{\ell}(\vZ \mid \vY) \doteq \ReLU(\vZ - \beta
(\vD^{\ell})^{\top}(\vD^{\ell}\vZ - \vY) + \beta \lambda
\vone_{s}\vone_{n}^{\top})
\end{equation}
sino más bien la siguiente iteración:
\begin{equation}
\vZ_{\theta}^{\ell + 1} = \vA_{\theta}^{\ell, T}; \qquad
\vA_{\theta}^{\ell, t + 1} =
\ISTA_{\theta}^{\ell}(\vA_{\theta}^{\ell, t} \mid \vZ_{\theta}^{\ell
+ 1/2}) \quad \forall 0 \leq t < T; \qquad \vA_{\theta}^{\ell, 0} =
\vzero_{s \times n},
\end{equation}
es decir, aplicando el gradiente proximal sobre el objetivo LASSO durante \(T \geq
1\) pasos en la pasada hacia adelante en cada capa, inicializado en
\(\vzero_{s \times n}\). En esta circunstancia, el diccionario puede ser
tan amplio como sea necesario, es decir, \(\vD^{\ell} \in \R^{s \times d}\) donde \(s
\geq d\) (por lo general, en la práctica se toma \(s = 4d\)).

\begin{table}[b]
\caption{\small \textbf{Detección de objetos y segmentación detallada
mediante MaskCut en COCO val2017~\citep{lin2014microsoft}}.
En este caso, todos los modelos se entrenan con un tamaño de parche de \(8\) en lugar de \(16\).
En comparación con modelos existentes como CRATE y ViT, la
familia de modelos CRATE-\(\alpha\) presenta un rendimiento notablemente mejorado,
así como una mayor escalabilidad.}
\centering
\begin{tabular}{@{}lcccccccc@{}}
\toprule
&  & \multicolumn{3}{c}{Detection} &  \multicolumn{3}{c}{Segmentation} \\
Model & AP$_{50} \uparrow $ & AP$_{75} \uparrow $ & AP $\uparrow$ &
AP$_{50} \uparrow$ & AP$_{75} \uparrow $ & AP $\uparrow$ \\
\midrule
\midrule
CRATE-\(\alpha\)-B/8 & 3.5 & 1.1 & 1.5 & 2.2 & 1.0 & 1.1 \\
CRATE-\(\alpha\)-L/8 & \textbf{4.0} & \textbf{1.7} & \textbf{2.0} &
\textbf{2.7} & \textbf{1.1} & \textbf{1.4} \\
\midrule
\color{gray}CRATE-B/8 & \color{gray}2.9 & \color{gray}1.0 &
\color{gray}1.3 & \color{gray}2.2 & \color{gray}0.7 & \color{gray}1.0 \\
\color{gray}ViT-B/8 & \color{gray}0.8 & \color{gray}0.2 &
\color{gray}0.4 & \color{gray}0.7 & \color{gray}0.5 & \color{gray}0.4 \\
\bottomrule
\end{tabular}
\label{tab:crate_alpha_detection_segmentation}
\end{table}

Sin embargo, esto presenta un problema práctico (empírico). Al utilizar la configuración anterior,
si \(\vZ^{\ell + 1/2} \in \R^{d \times n}\), entonces
\(\vZ^{\ell + 1} \in \R^{s \times n}\), lo que podría resultar en una dimensión de características
arbitrariamente grande. En la práctica, queremos que la dimensión de las características
en cada capa sea la misma. Así, esto plantea una tricotomía práctica
en el diseño de la red; \textit{no es posible} satisfacer simultáneamente
\textit{todos} los siguientes requisitos deseables:
\begin{enumerate}
\item La dimensión de las características es constante en todas las capas.
\item El diccionario es ancho, es decir, sobrecompleto.
\item La salida de la no linealidad corresponde a los códigos dispersos de 
      la entrada respecto al diccionario.
\end{enumerate}
En la práctica, renunciar a (1) resulta menos manejable por razones de eficiencia.
Renunciar a (2) conduce al marco CRATE habitual. Renunciar a (3) conduce
a una versión amplia de CRATE, es decir, CRATE-\(\alpha\), que tiene la
siguiente no linealidad para pasar de \(\vZ^{\ell + 1/2}\) a \(\vZ^{\ell + 1}\):
\begin{equation}
\vZ_{\theta}^{\ell + 1} = \vD^{\ell}\vA_{\theta}^{\ell, T}; \qquad
\vA_{\theta}^{\ell, t + 1} =
\ISTA_{\theta}^{\ell}(\vA_{\theta}^{\ell, t} \mid \vZ_{\theta}^{\ell
+ 1/2}); \qquad \vA_{\theta}^{\ell, 0} = \vzero,
\end{equation}
es decir, se toman los códigos dispersos resultantes del descenso de gradiente proximal 
y se multiplican por el diccionario para obtener una versión de la entrada sin ruido.
Por consiguiente, la no linealidad de CRATE-\(\alpha\) calcula una
versión sin ruido de la entrada que es apta para la codificación dispersa, en lugar
de los códigos dispersos reales en si mismos. El mapeo desde \(\vZ_{\theta}^{\ell
+ 1/2}\) a \(\vZ_{\theta}^{\ell + 1}\) se denomina aquí
bloque de aprendizaje de diccionarios sobrecompletos (ODL - Overcomplete Dictionary Learning) y se denota
\(\operatorname{ODL}_{\theta}^{\ell}\), es decir,
\begin{equation}
\vZ_{\theta}^{\ell + 1}(\vX) \doteq
\operatorname{ODL}_{\theta}^{\ell}(\vZ_{\theta}^{\ell + 1/2}(\vX)).
\end{equation}

\begin{figure}[t]
\centering
\includegraphics[width=\textwidth]{\toplevelprefix/chapters/chapter8/figs/crate_alpha_semantic_heads.pdf}
\caption{\small\textbf{Mapas de saliencia de CRATE-\(\alpha\) con un tamaño de
parche de \(8\).} Cada fila corresponde a una imagen diferente y cada columna
corresponde a una cabeza de atención diferente en la última capa. Se
observa que los mapas de saliencia se corresponden claramente con los objetos de
la imagen de entrada.}
\label{fig:crate_alpha_saliency_maps}
\end{figure}

\begin{table}
\caption{\small\textbf{Pérdida de validación en el modelado del lenguaje.} Aquí
todos los modelos se entrenan previamente con la mayor parte de OpenWebText, y la
pérdida de entropía cruzada de validación se mide en un subconjunto de validación de OpenWebText.
CRATE-\(\alpha\) muestra una mejora significativa con respecto al diseño de CRATE,
aunque sigue existiendo una brecha con respecto a los transformers tradicionales como GPT-2.}
\centering
\begin{tabular}{@{}lcccc@{}}
\toprule
Modelo & GPT-2-B(ase) & CRATE-B & CRATE-\(\alpha\)-S(mall) &
CRATE-\(\alpha\)-B \\
\midrule
\midrule
\# parámetros & 124M & 60M & 57M & 120M \\
OWT val.~loss & 2.85 & 3.37 & 3.28 & 3.14 \\
\bottomrule
\end{tabular}
\label{tab:crate_alpha_lm}
\end{table}

La capa CRATE-\(\alpha\) se muestra en
\Cref{fig:crate_alpha_backbone}. En la práctica, esta modificación de
CRATE rinde muy bien a escalas superiores. Por ejemplo, cuando
preentrenamos modelos CRATE-\(\alpha\) en ImageNet-21K, las tareas no supervisadas
como la segmentación (véase \Cref{fig:crate_alpha_saliency_maps} y la
\Cref{tab:crate_alpha_detection_segmentation}) suelen tener
un rendimiento significativamente mejorado en comparación con CRATE. Tendencia similar se observa 
en el entrenamiento de modelos de lenguaje utilizando la auto-atención causal
(véase \Cref{tab:crate_alpha_lm}). En general, se trata de un camino prometedor para
escalar el rendimiento e igualar a los modelos de caja negra como
los transformers.\footnote{Cabe señalar que los resultados experimentales de esta
sección utilizan una arquitectura de modelo ligeramente diferente, lo que aporta mejoras
empíricas muy leves. Los cambios son: (1) una conexión residual adicional
en el bloque ODL, (2) modificar \(\ISTA\) para que utilice dos
diccionarios independientes en lugar de \(\vD^{\ell}\) y \((\vD^{\ell})^{\top}\).}

\subsection{Transformadores de complejidad temporal lineal}\label{sub:tost_experiments}

En la práctica, los modelos de aprendizaje profundo sufren de cuellos de botella en cuanto a la complejidad espacial
y temporal, lo que se traduce en tamaños de problemas más allá de los cuales
no pueden escalar con recursos fijos. Uno de estos cuellos de botella, especialmente
relevante cuando se trata de datos en los que cada muestra es en sí misma
de alta dimensión y rica en información (como largas secuencias de texto o vídeos),
es la \textit{complejidad temporal} del procesamiento de secuencias largas de datos.
Con el fin de reducir la complejidad temporal del procesamiento de datos mediante
transformers, en \Cref{sub:tost} propusimos un operador de \textit{auto-atención basada en estadísticas de tokens}
\(\TSSA_{\theta}^{\ell}\). Ahora
construimos un \textit{transformer basado en estadísticas de tokens}, denominado ToST:\\
\centerline{\url{https://robinwu218.github.io/ToST/},} en torno a él, que podemos utilizar para tareas de contexto largo. En particular, podemos
utilizar la siguiente capa (representada en \Cref{fig:tost_backbone}) como
sustituto directo de una capa de backbone en CRATE.
\begin{align}
\vZ_{\theta}^{\ell + 1/2}(\vX)
&= \vZ_{\theta}^{\ell}(\vX) + \TSSA_{\theta}^{\ell}(\LN_{\theta}^{1,
\ell}(\vZ_{\theta}^{\ell}(\vX))) \\
\vZ_{\theta}^{\ell + 1}(\vX)
&= \vZ_{\theta}^{\ell + 1/2}(\vX) +
\MLP_{\theta}^{\ell}(\LN_{\theta}^{2, \ell}(\vZ_{\theta}^{\ell + 1/2}(\vX)))
\end{align}
donde el bloque \(\TSSA\) se define tal como se indica en \Cref{sub:tost}. Obsérvese
que esto es exactamente igual que la arquitectura del transformer visual
discutida en \Cref{sub:contrastive_learning_architecture}, excepto
que \(\TSSA\) reemplaza al bloque convencional de auto-atención multicanal
\(\MHSA\). En cualquier caso, la complejidad computacional de la
propagación hacia adelante de esta capa es lineal en todas las variables del problema ---
longitud de la secuencia, dimensión de las características, número de cabezales y dimensión de los cabezales.

\begin{figure}[!htbp]
\centering
\includegraphics[width=\textwidth]{\toplevelprefix/chapters/chapter8/figs/tost_backbone.pdf}
\caption{\small\textbf{Una capa de la estructura ToST}. Las representaciones de los tokens
pasan por normas de capa, el operador de auto-atención de estadísticas de tokens
(TSSA) y una red neuronal multicapas (MLP), con el fin de formar la
salida de la capa.}
\label{fig:tost_backbone}
\end{figure}

\begin{table}[t]
\centering
\resizebox{\textwidth}{!}{
\begin{tabular}{@{}lccc|cc|cc@{}}
    \toprule
    C.Datos & ToST-T(iny) & ToST-S(mall) & ToST-M(edium) &  {
    \color{gray} XCiT-S} &  { \color{gray}XCiT-M} &
    {\color{gray}ViT-S} & {\color{gray}ViT-B(ase)}\\
    \toprule
    \# parámetros & 5.8M & 22.6M & 68.1M &  { \color{gray}24.9M }& {
    \color{gray}80.2M } & { \color{gray} 22.1M } & { \color{gray}86.6 M } \\
    \midrule
    ImageNet & 67.3 & 77.9 & 80.3 &  { \color{gray} 80.5} & {
    \color{gray} 81.5 } & { \color{gray} 79.8} & { \color{gray} 81.8} \\
    ImageNet ReaL & 72.2  & 84.1 & 85.6 &  { \color{gray} 85.6 } & {
    \color{gray} 85.9} & { \color{gray} 85.6 } & { \color{gray} 86.7 }\\
    \midrule
    CIFAR10 & 95.5 & 96.5 & 97.5 &  { \color{gray} 98.1} & {
    \color{gray} 98.3} & { \color{gray} 98.6} & { \color{gray} 98.8}\\
    CIFAR100 & 78.3 & 82.7 & 84.5 &  { \color{gray} 86.1} & {
    \color{gray} 87.6} & { \color{gray} 88.8} & { \color{gray} 89.3}\\
    Oxford Flowers-102 & 88.6 & 92.8 & 94.2 &  { \color{gray} 93.9} &
    { \color{gray} 94.0} & { \color{gray} 94.0}& { \color{gray} 95.7}\\
    Oxford-IIIT-Pets & 85.6 & 91.1 & 92.8 &  { \color{gray} 92.9} & {
    \color{gray} 94.0} & { \color{gray} 92.8} & { \color{gray} 94.1} \\
    \bottomrule
\end{tabular}%
}
\caption{\small \textbf{Precisión de clasificación mediante sondeo lineal de
ToST} en diversos conjuntos de datos con diferentes tamaños de modelo cuando la
estructura principal se ha entrenado previamente para la clasificación de ImageNet-1K. Observamos
que, en comparación con XCiT (una arquitectura de tipo transformer diseñada empíricamente
y especializada en el procesamiento eficiente de secuencias largas)
y el ViT (Vision Transformer), ToST mantiene un rendimiento relativamente similar, incluso
a pesar de contar con ventajas como un tiempo de ejecución más rápido y un diseño transparente o de white-box.}
\label{tab:tost_linear_probing}
\end{table}

\begin{table}[!htbp]
\centering
\begin{tabular}{@{}lcccccc@{}}
\toprule
Modelo & \# parám. & OWT & Lambada & Wikitext & PTB & Avg $\downarrow$
\\ \midrule
GPT-2-Base & 124M & 2.84 & 4.32 & 4.13 & 5.75 & 4.26 \\
ToST-Base & 110M & 3.20 & 4.98 & 4.77 & 6.39 & 4.84 \\
ToST-Medium & 304M & 2.88 & 4.45 & 4.30 & 5.64 & 4.32 \\
ToST-Large & 655M & 2.72 & 4.32 & 3.99 & 5.03 & 4.02 \\ \bottomrule
\end{tabular}%
\caption{\small\textbf{Pérdida de validación del modelado del lenguaje} calculada sobre
(conjuntos de validación de) diversos conjuntos de datos de lenguaje natural, tras
el preentrenamiento del modelo en ese conjunto de datos. Observamos que ToST se adapta
bien, de modo que ToST-Large supera al modelo de referencia GPT-2-Base en el modelado causal
del lenguaje, al tiempo que ofrece una eficiencia superior en contextos largos.} 
\label{tab:tost_lm}
\end{table}

\begin{table}
\centering

\begin{tabular}{@{}lccccccc@{}}
\toprule
Modelo        & ListOps  & Texto    & Consulta & Image    &
Pathfinder & Avg      \\
\midrule
\midrule
Reformer              & \textbf{37.27} & 56.10             & 53.40
& 38.07             & 68.50               & 50.56             \\
BigBird               & 36.05             & 64.02             & 59.29
& 40.83             & 74.87                & 54.17             \\
LinFormer         & 16.13             & \underline{65.90} & 53.09
& 42.34             & \underline{75.30}               & 50.46             \\
Performer             & 18.01             & 65.40             & 53.82
& 42.77             & \textbf{77.05}                & 51.18             \\
Transformer           & 37.11             & 65.21             &
\underline{79.14}              & \underline{42.94}             &
71.83              & \underline{59.24}            \\
ToST & \underline{37.25}    & \textbf{66.75}    & \textbf{79.46}
& \textbf{46.62}    &    69.41      &     \textbf{59.90}\\

\bottomrule
\end{tabular}%
\caption{\small \textbf{Comparación del rendimiento en Long-Range Arena (LRA)
    entre ToST(-B) y las principales variantes de Transformer optimizadas para
contextos largos.} Long-Range Arena es una familia de pruebas de rendimiento que evalúa
la capacidad de modelado de secuencias largas de algoritmos y
arquitecturas, mediante el uso de un conjunto de datos y un mecanismo de evaluación fijos. ToST
se ubica en la parte superior del ranking, superando a todas las variantes conocidas
del transformer, incluyendo XCiT y el transformer regular (ViT)
(véase \Cref{tab:tost_linear_probing}). Además, ToST presenta
la inferencia con menor complejidad temporal y espacial. (En esta tabla, la
    mejor puntuación para un benchmark concreto aparece en negrita, y la segunda mejor
puntuación, subrayada.)}
\label{tab:tost_lra_results}
\end{table}

Además, la arquitectura propuesta, denominada ToST (por «Token Statistics Transformer», 
«Transformador de Estadísticas de Tokens»), ofrece un buen rendimiento en tareas de visión (véase, por ejemplo,
\Cref{tab:tost_linear_probing}) y en tareas lingüísticas (véase, por ejemplo,
\Cref{tab:tost_lm}). Esto es especialmente cierto en el caso de tareas con secuencias largas
(véase \Cref{tab:tost_lra_results}), donde ofrece un mayor
rendimiento y es mucho más eficiente que los transformers convencionales y
todas las demás arquitecturas similares a los transformers.

\subsection{Transformadores de solo atención} \label{sub:aot_experiments}

Otro cuello de botella que hay que eliminar de los modelos de aprendizaje profundo, concretamente
de las arquitecturas tipo transformer, es el cuello de botella de memoria que se deriva
de las multiplicaciones matriciales masivas en las MLP, donde la dimensión interna
es signifcativamente mayor que la dimensión de las características \(d\). Por lo tanto, resulta
interesante y relevante plantearse la siguiente pregunta: ¿necesitamos \textit{realmente}
la MLP dentro de un transformer, y qué nivel de rendimiento
se puede alcanzar sin ella? Para analizar esta cuestión, utilizamos la
arquitectura «Attention-Only Transformer» (AoT, transformer de solo atención) (véase \Cref{sub:aot}),
representada en \Cref{fig:aot_backbone}. Concretamente, cada capa tiene simplemente la forma
\begin{equation}
\vZ_{\theta}^{\ell + 1}(\vX) = \vZ_{\theta}^{\ell}(\vX) +
\MSSA_{\theta}^{\ell}(\LN_{\theta}^{\ell}(\vZ_{\theta}^{\ell}(\vX))).
\end{equation}
En nuestra implementación, también probamos a utilizar la auto-atención
multicabezal (MHSA) en lugar de la MSSA. Resulta que esta
arquitectura \textit{también} es viable, aunque la profundidad de la red
debe ser considerablemente mayor para lograr un rendimiento equivalente
al de la arquitectura CRATE o Transformer habitual. %

\begin{figure}[!htbp]
\centering
\includegraphics[width=0.7\textwidth]{\toplevelprefix/chapters/chapter8/figs/aot_backbone.pdf}
\caption{\small\textbf{Una capa de la red troncal AoT.} Las representaciones de los tokens
simplemente pasan por una normalización de capa y por el operador de auto-atención
multicabezal (subespacial) para formar la salida de la capa. Obsérvese
que no hay no linealidad a nivel de token, como en las redes MLP, ISTA u ODL.}
\label{fig:aot_backbone}
\end{figure}

Llevamos a cabo experimentos utilizando la arquitectura AoT propuesta y
demostramos su potencial. Preentrenamos los modelos AoT-MSSA y AoT-MHSA
de diferentes tamaños, junto con GPT-2, en
OpenWebText~\citep{Gokaslan2019OpenWeb}. Graficamos la pérdida de entrenamiento
y la pérdida de validación en función del número de iteraciones de entrenamiento en
\Cref{fig:loss}(a) y (b), respectivamente. Se observa que los modelos basados en AoT de tamaño mediano
y grande alcanzan pérdidas de entrenamiento y validación
comparables a las del modelo base GPT-2. Además,
en comparación con el modelo base GPT-2, el modelo AoT-MHSA es idéntico
al modelo base GPT-2, salvo por la ausencia de capas MLP en la
arquitectura. Como se muestra en \Cref{fig:loss}, la incorporación de capas MLP
puede acelerar el proceso de entrenamiento. Utilizando los modelos preentrenados
anteriores, calculamos la pérdida de validación de entropía cruzada sin entrenamiento
en diferentes conjuntos de datos en \Cref{tab:zeroshot}. Se observa que los
modelos AoT con un número de parámetros medio y grande pueden alcanzar
un rendimiento comparable al del modelo base GPT-2. Además, hemos constatado
que añadir capas de MLP a los modelos AoT no mejora el rendimiento
sin entrenamiento previo. Estos resultados ponen de relieve el potencial de los modelos basados únicamente en atención
para lograr resultados competitivos sin perder interpretabilidad.

\begin{table*}[!htbp]
\caption{Resultados «zero-shot» en varios conjuntos de datos y tareas de referencia lingüística:
Evaluación de diferentes tamaños de AoT con los operadores MSSA y MHSA
y comparación con el modelo GPT-2.}\vskip 0.1in
\centering
\resizebox{1\columnwidth}{!}
{
\begin{tabular}{l|cccccc}
    \toprule
    Modelos  & {\bf LAMBADA} & {\bf PTB} & {\bf WikiText} & {\bf
    LAMBADA} & {\bf CBT CN} & {\bf CBT NE} \\
    \# de parámetros  & (val loss) $\downarrow$ &  (val loss)
    $\downarrow$ &(val loss) $\downarrow$ & (acc) $\uparrow$ &(acc)
    $\uparrow$ &(acc) $\uparrow$ \\
    \midrule
    AoT-MSSA Base (102M) & 4.70 & 6.03 & 4.65 & 0.25 & 0.80 & 0.74\\
    AoT-MSSA Medium (182M) & 4.47 & 5.08 & 4.22 & 0.29 & 0.84 & 0.77 \\
    AoT-MHSA Base (122M) & 4.42 & 5.52 & 4.19 & 0.38 & 0.86 & 0.82\\
    GPT-2 Base (124M) & 4.32 & 5.75 & 4.13 &  0.40 &  0.87 &  0.84 \\
    \bottomrule
\end{tabular}
}
\label{tab:zeroshot}
\end{table*}
\vspace{-0.05in}

\begin{figure*}[t]
\begin{center}
\includegraphics[width=0.4\linewidth]{\toplevelprefix/chapters/chapter8/figs/training_loss.png}
\hspace{0.4in}
\includegraphics[width =
0.4\linewidth]{\toplevelprefix/chapters/chapter8/figs/val_loss.png}
\vspace{-0.15in}
\caption{\centering \textbf{Evaluación de modelos en tareas lingüísticas.}
    Representamos gráficamente la pérdida de entrenamiento (izquierda) y la pérdida de validación (derecha) de los modelos AoT
y GPT-2 preentrenados en OpenWebText.}  \label{fig:loss}
\end{center}
\vspace{-0.15in}
\end{figure*}

\section{Resumen y Notas}

Todo el trabajo de este capítulo se basa en la arquitectura Transformer,
introducida por \citet{vaswani2017attention}.
La arquitectura Transformer se describe formalmente en
\Cref{sec:contrastive_learning}. Una de las principales innovaciones empíricas de
los últimos años, impulsada por la prevalencia y el rendimiento de la
arquitectura Transformer, consiste en formular un problema de aprendizaje dado como
un problema de secuencia a secuencia y aplicar la arquitectura Transformer.
Esto ha permitido que la arquitectura Transformer sea
esencialmente omnipresente en (casi) todas las aplicaciones de aprendizaje profundo. Por
ello, las mejoras directas en el Transformer pueden extenderse hasta convertirse en
soluciones a muchos problemas y tener un impacto considerable; del mismo modo,
podemos aplicar nuestra comprensión interna (whte-box) de las arquitecturas similares al Transformer
a muchos problemas actuales. El material tratado en este
capítulo es solo una parte del trabajo que ya se ha realizado;
otros trabajos incluyen el completamiento enmascarado para datos de texto (p. ej., BERT)
\citep{devlin2019bert,yu2024white}, la interpretabilidad (mecánica) de
los modelos de lenguaje y visión \citep{bai2024improving}, y la
codificación con corrección de errores \citep{zheng2025white}. Queda mucho más por hacer.

También hay mucha más teoría específica sobre la práctica del
escalado de redes neuronales, que tiene una enorme viabilidad práctica, y
al menos lo mencionamos aquí. Esta línea de investigación se popularizó gracias a
la línea de trabajo de los «programas tensoriales» \citep{yang2022tensor}. La
regla básica es que queremos que las actualizaciones iniciales de los gradiente en un
transformer sean de tamaño constante y, al analizar con
detalle las ecuaciones de retropropagación (\Cref{app:optimization}), podemos
determinar la escala de las tasas de inicialización y de aprendizaje (elegidas
por capa) que se requieren para lograrlo. En la práctica, estas
recomendaciones aumentan considerablemente la estabilidad y la convergencia del
entrenamiento a gran escala; además, indican una forma de encontrar los
``hiperparámetros óptimos''\footnote{La palabra ``óptimos'' se utiliza entre comillas porque
este trabajo se limita a emplear ciertos criterios sobre el tamaño de los pesos,
el tamaño de las características y el tamaño del gradiente en la inicialización para determinar
la ``optimalidad'', en contraposición, por ejemplo, a la pérdida en el conjunto de prueba al alcanzar la convergencia.}
para el entrenamiento a gran escala utilizando únicamente
entrenamiento a pequeña escala. Los trabajos posteriores a este intentan adaptarse a la geometría de las características
\citep{bernstein2024oldoptimizernewnorm}, lo cual podría
basarse en el trabajo de este libro sobre el aprendizaje de representaciones.
Otros trabajos posteriores incorporan esta información sobre los pesos en el
propio optimizador para obtener estos beneficios de escalabilidad en forma automática,
lo que da lugar a optimizadores como Muon \citep{jordan6muon}, que
recientemente se han utilizado para entrenar modelos de billones de parámetros de manera extremadamente estable
\citep{moonshot2025kimi}. En general, los dos enfoques de la teoría del
aprendizaje profundo son ortogonales o complementarios.

\section{Ejercicios y extensiones}

\begin{exercise}
Lea el artículo de DINO \cite{caron2021emerging}.
\end{exercise}

\begin{exercise}
DINO v2 \cite{oquab2023dinov2} utiliza todas las componentes de DINO v1, pero además,
durante la fase de aumento de datos, enmascara aleatoriamente los parches dentro
de cada vista. Este tipo de aumento debería garantizar que las características
de las imágenes con información local similar sean también similares. Formule un
problema de optimización que favorezca esto en el codificador e impleméntelo.
\end{exercise}

\begin{exercise}
Este ejercicio analiza la aplicación de algoritmos de optimización estocástica
para minimizar las pérdidas que implican expectativas.
\begin{enumerate}[(a)]
\item Proponga una alternativa al término que involucra \(R_{\eps}\) en
    \eqref{eq:simdino_loss_teacherstudent_empirical} para
    aproximar el término de regularización de la covarianza en
    \eqref{eq:simdino_loss_teacherstudent}. Evalúe la complejidad temporal
    necesaria para calcular el término propuesto y su
    gradiente. Incluya un análisis del cálculo en un solo
    nodo de cómputo frente a múltiples nodos.
\item Evalúe la complejidad temporal necesaria para calcular el término existente
    en \eqref{eq:simdino_loss_teacherstudent_empirical} y su gradiente.
\end{enumerate}
\end{exercise}

\begin{exercise}
Demuestre que \eqref{eq:linear_probing} y
\eqref{eq:linear_probing_empirical} son problemas de optimización convexos.
\end{exercise}

\begin{exercise}
\phantom{}
\begin{enumerate}[(a)]
\item Implemente los modelos CRATE y CRATE-\(\alpha\).
\item Compare su rendimiento y eficiencia en el conjunto de datos CIFAR-10.
\item Compare su interpretabilidad de dos maneras:
Compare their interpretability in two ways:
    \begin{itemize}
        \item La dispersión \(\norm{\vZ}_{0}\) de la representación \(\vZ\)
        \item Los mapas de atención \(\va_{\theta}^{k, \ell}\)
    \end{itemize}
\end{enumerate}
\end{exercise}

\end{document}
